\documentclass[11pt]{article}
% Notation contract: H is the scalar graph-hull gap; T is the gap of the
% specified exact factorwise relaxation. rho_G is graph density; rho_B is
% physical box aspect ratio. Spatial order r means degree 2r.
\usepackage[T1]{fontenc}
\usepackage[utf8]{inputenc}
\usepackage{lmodern}
\usepackage{amsmath,amssymb,amsthm,mathtools}
\usepackage[margin=1in]{geometry}
\usepackage{microtype}
\usepackage{booktabs,array}
\usepackage[numbers,sort&compress]{natbib}
\usepackage[colorlinks=true,linkcolor=blue,citecolor=blue,urlcolor=blue]{hyperref}
\newtheorem{theorem}{Theorem}[section]
\newtheorem{lemma}[theorem]{Lemma}
\newtheorem{proposition}[theorem]{Proposition}
\newtheorem{corollary}[theorem]{Corollary}
\theoremstyle{definition}
\newtheorem{definition}[theorem]{Definition}
\newtheorem{example}[theorem]{Example}
\theoremstyle{remark}
\newtheorem{remark}[theorem]{Remark}
\DeclareMathOperator{\vex}{vex}
\DeclareMathOperator{\cav}{cav}
\DeclareMathOperator{\conv}{conv}
\DeclareMathOperator{\supp}{supp}
\DeclareMathOperator{\rank}{rank}
\DeclareMathOperator{\diag}{diag}
\DeclareMathOperator{\tw}{tw}
\DeclareMathOperator{\osc}{osc}
\DeclareMathOperator{\OPT}{OPT}
\newcommand{\R}{\mathbb R}
\newcommand{\N}{\mathbb N}
\newcommand{\E}{\mathbb E}
\newcommand{\Prob}{\mathbb P}
\newcommand{\ind}{\mathbf 1}
\newcommand{\hgap}{H}
\newcommand{\tgap}{T}
\newcommand{\cube}{[0,1]^n}
\newcommand{\pmcube}{\{-1,1\}^n}
\newcommand{\laws}{\mathcal P}
\newcommand{\dens}{\rho_G}
\newcommand{\aspect}{\rho_B}
\newcommand{\norm}[1]{\lVert #1\rVert}
\newcommand{\abs}[1]{\lvert #1\rvert}
\newcommand{\ip}[2]{\langle #1,#2\rangle}
\numberwithin{equation}{section}
\allowdisplaybreaks[1]

\title{Convex relaxation gaps and spatial certificates in nonlinear optimization}
\author{}
\date{}
\begin{document}
\maketitle
\begin{abstract}
Separate convexification of nonlinear terms can lose compatibility across
shared variables. We formulate this loss through probability laws on box
vertices and compare it with the resources needed for spatial optimality
certificates. For signed bilinear functions, the worst coefficient gap on
each interaction graph has the order of the square root of its largest
induced edge density, consistent with earlier Schur-multiplier theory.
For positive multilinear functions, we prove sharp leading growth
$\log d/\log\log d$ in degree, exact and asymptotically sharp structural
bounds, and leading growth equal to the physical box aspect ratio.
Finite cubic certificates, common-law rounding, and exact scalar-envelope
complexity delimit these comparisons. Separately, fractional-cardinality
and signed cubic XOR families require exponentially many certified regions
under specified spatial moment oracles. The latter bounds survive exact
full-box quadratic moment information and bounded monomial reformulations.
For the cardinality family, coordinatewise graph lifts admit a
degree-loss-free transfer. Matching
finite exponential upper certificates clarify the oracle scope. Supporting
geometric and complexity comparisons distinguish coordinate dependence,
global conic lift size, integer precision, and primitive bound propagation
from pointwise envelope gaps and spatial region counts.
\end{abstract}
\setcounter{tocdepth}{1}
\tableofcontents
\clearpage
\section{Introduction}\label{sec:introduction}
A convex relaxation can be exact for every individual term and still lose
information about their shared variables. Replacing each term by its graph
hull permits a separate convex representation for that term. The graph
hull of the sum requires a single representation that works simultaneously.
This paper studies the resulting width comparison and a complementary
question: how many regions must a specified spatial bounding rule certify
before it proves a desired optimization tolerance?

For a scalar function $f$ on a box $B$, write
$H_B(f;x)=\cav_B f(x)-\vex_B f(x)$. For a specified factorization
$f=\sum_e f_e$, write
$T_B(f;x)=\sum_e[\cav_B f_e(x)-\vex_B f_e(x)]$.
Always $T_B\ge H_B\ge0$, but only $T_B$ depends on the factorization.
Ratios are formed where $H_B>0$. A ratio of these pointwise widths is not
an approximation ratio for a constrained problem, and supplies no spatial
tree lower bound on its own. A spatial theorem must also specify its node
oracle, allowed branching coordinates, and how discarded regions are
certified.

McCormick's recursive convexification is the classical starting point
\citep{McCormick1976}. The vertex-law description of multilinear envelopes
and earlier gap comparisons appear in \citet{Luedtke2012}; their
positive-coefficient question motivates the universal laws studied here.
For signed bilinear functions, \citet{Boland2017} relate the gap to cut
ranges and establish square-root-dimensional growth. We give an elementary
per-graph density proof, while an explicit transfer from
\citet{DavidsonDonsig2007} identifies its established asymptotic content.
Classical symmetric envelopes, matching theory and balanced-matrix results
supply further ingredients, with the exact hypotheses recorded where used.

The positive multilinear results exploit a common upper-attaining law.
The remaining task is a simultaneous lower-envelope rounding guarantee
that preserves every mean, independently of the objective coefficients.
The universal degree and dimension suprema have sharp leading growth
$\log d/\log\log d$ and $\log n/\log\log n$, respectively. Explicit finite
bounds and cubic witnesses accompany these asymptotics. Incidence
orientations, marginal floors, feedback variables, and frequency two give
more precise laws; incidence treewidth two has exact constant two.
Positive physical boxes require separate coefficient arguments: their
aspect-ratio supremum lies between $\max\{2,\rho_B\}$ and $\rho_B+2$,
with a finite-dimensional improvement. A single product already has
NP-complete exact midpoint-envelope evaluation on arbitrarily narrow
unequal rational boxes, although its factorwise ratio is one. Thus
relaxation strength and exact evaluation complexity are distinct.

The spatial results use different examples and proof objects. A
fractional-cardinality quadratic supplies an elementary chord obstruction,
then explicit SDP--RLT and higher-order pseudomoments extend it. These
proofs build on established fractional-knapsack moment lower bounds
\citep{Grigoriev2001,Potechin2019}, and are compared with earlier spatial
lower bounds \citep{Jarre2018}. They extend to arbitrary coordinate domains,
fixed relative tolerance and coordinatewise graph lifts. A known coupled
quadratic cut closes that first family at the root. Signed cubic XOR
objectives, using the classical signed-character construction of
\citet{Schoenebeck2008}, give the stronger obstruction: exponentially many
regions remain necessary even with exact quadratic moment information for
the full coordinate box. The transfer survives bounded signed-monomial
lifts, with explicit degree and parity-rank accounting. These signed
objectives on $[-1,1]^n$ are separate from the positive-gap examples on
nonnegative boxes.

Table~\ref{tab:roadmap} gives a route through the proofs. Readers concerned
primarily with spatial methods can read the gap and oracle definitions in
Section~\ref{sec:foundations}, then start with
Section~\ref{sec:cardinality-spatial}; the spatial proofs do not require
the positive-coupling chapters. Appendices retain distinct predecessor
arguments, finite certificates, and supporting comparisons. The exact
check programs and their replay instructions accompany the manuscript;
finite checks supplement the general proofs.
\begin{table}[htbp]
\centering\small
\begin{tabular}{@{}p{.24\textwidth}p{.47\textwidth}p{.21\textwidth}@{}}
\toprule
Question & Main conclusion & Route \\
\midrule
Simultaneous envelopes & Signed graph-density order; sharp positive
 degree growth and finite cubic bounds & Sections~\ref{sec:foundations}--\ref{sec:equal-means} \\
Structure and physical boxes & Incidence and marginal laws; exact width
 two; aspect-ratio bounds and exact evaluation complexity &
 Sections~\ref{sec:incidence}--\ref{sec:single-product} \\
Spatial certificates & Cardinality and XOR cover lower bounds;
 graph-lift transfer and finite upper certificates &
 Sections~\ref{sec:cardinality-spatial}--\ref{sec:xor-upper-affine} \\
Other resources & Coordinate and scale dependence, global lift size,
 integer precision and primitive propagation &
 Section~\ref{sec:comparisons}, Appendices~\ref{app:point-packing}--\ref{app:integer-precision} \\
\bottomrule
\end{tabular}
\caption{Reader's roadmap. Each comparison retains its own domain,
representation and computational resource.}\label{tab:roadmap}
\end{table}

\section{Envelope gaps as probability problems}\label{sec:foundations}
\subsection{The scalar graph and the specified termwise relaxation}
Let $B=\prod_{i=1}^n[\ell_i,u_i]$ be a finite box. Fixed coordinates are
removed when necessary. A function is \emph{multiaffine} if it is affine in
each coordinate separately; equivalently, it is a multilinear polynomial
with no repeated variable in a monomial. For a continuous scalar function
$f$ on $B$, let $\vex_B f$ be its greatest convex underestimator and
$\cav_B f$ its least concave overestimator. The vertical section of the
convex hull of its graph is
\[
 \{y:(x,y)\in\conv\{(z,f(z)):z\in B\}\}
 =[\vex_B f(x),\cav_B f(x)].
\]
Define the \emph{scalar graph-hull gap} and, for a specified decomposition
$f=\sum_{e\in E}f_e$, the \emph{termwise gap} by
\begin{align}
 \hgap_B(f;x)&=\cav_B f(x)-\vex_B f(x),\label{eq:hgap}\\
 \tgap_B(f;x)&=\sum_{e\in E}
 [\cav_B f_e(x)-\vex_B f_e(x)].\label{eq:tgap}
\end{align}
The notation $\tgap_B$ always refers to the displayed decomposition.
It is the width obtained by giving each scalar factor its exact graph hull
and imposing $y=\sum_e y_e$. It can change when a factor is split into
smaller pieces, whereas $\hgap_B$ does not. Always,
\begin{equation}\label{eq:gap-order}
 0\le\hgap_B(f;x)\le\tgap_B(f;x),
\end{equation}
because the sum of convex underestimators is a convex underestimator of
the sum, and similarly for concave overestimators. Ratios are taken only
where $\hgap_B(f;x)>0$; gap inequalities apply at every point.

For a positive multilinear polynomial, the factors are its displayed
monomials $a_e\prod_{i\in e}x_i$, with $a_e\ge0$. For a signed bilinear
function, they are the signed products $a_{ij}x_ix_j$. A bilinear factor
has the familiar four-inequality McCormick hull. For a higher-degree factor,
its exact hull and a recursive McCormick relaxation are different objects.
A recursive product relaxation is exact for one monomial on a unit box,
but need not be exact on a general positive box
\citep[Sections 2--3]{Luedtke2012}. A bound for \eqref{eq:tgap} on that
box therefore does not automatically bound the gap of a recursive bilinear
factorization.

Adding an affine function to a factor, or adding a separate affine factor,
changes neither gap. An affine bijection between boxes preserves both
envelopes after composition and hence preserves the gaps for the transformed
factors. These facts follow from the graph-hull representation.

\subsection{Vertex distributions and compatible marginals}
Let $\mathcal V(B)$ denote the vertices of $B$, and let $\laws_B(x)$ be
the distributions on $\mathcal V(B)$ with mean $x$. On the unit cube these
are joint Bernoulli distributions with the prescribed singleton means;
coordinates need not be independent.

\begin{lemma}[Vertex representation]\label{lem:vertex-law}
For a multiaffine $f$ on $B$,
\begin{align}
 \vex_B f(x)&=\min_{P\in\laws_B(x)}\E_P f(X),\label{eq:vex-law}\\
 \cav_B f(x)&=\max_{P\in\laws_B(x)}\E_P f(X).\label{eq:cav-law}
\end{align}
Both extrema are attained. The envelopes are continuous piecewise-affine
functions on $B$.
\end{lemma}
\begin{proof}
Write $x_i=\ell_i+(u_i-\ell_i)p_i$. Independently choose $X_i=u_i$ with
probability $p_i$ and $X_i=\ell_i$ otherwise. Multiaffinity gives
$\E X=x$ and $\E f(X)=f(x)$. Every graph point therefore belongs to the
convex hull of vertex graph points, and the reverse inclusion in the full
graph hull is immediate. The vertical sections give the formulas. The
feasible distributions form a nonempty compact polytope, so both extrema
are attained. The lower and upper surfaces of the vertex graph polytope
are respectively maxima and minima of finitely many affine functions on
$B$, which also proves continuity at its boundary.
\end{proof}

These representations are classical \citep[Section 3]{Luedtke2012}.
They expose the compatibility issue: different factors may use different
minimizing distributions in \eqref{eq:tgap}, whereas \eqref{eq:vex-law}
requires one distribution for their sum.

For $m_e(x)=\prod_{i\in e}x_i$ with $e\ne\varnothing$, the unit-cube
formulas are
\begin{equation}\label{eq:monomial-envelopes}
 \cav m_e(x)=\min_{i\in e}x_i,\qquad
 \vex m_e(x)=\max\{0,\sum_{i\in e}x_i-|e|+1\}.
\end{equation}
To prove them, interpret $\E m_e(X)$ as the probability of joint success.
The upper bound is attained by $X_i=\ind\{U\le x_i\}$ for one uniform
$U\in[0,1]$. For the lower bound, place failure intervals of lengths
$1-x_i$ consecutively on the unit circle. If their total length is below
one they are disjoint, and if it is at least one they cover the circle.
Their union therefore has measure $\min\{1,\sum_{i\in e}(1-x_i)\}$.
Each interval has its prescribed measure, so the complement gives the
claimed minimum joint-success probability. Unused coordinates can be
sampled independently to complete the mean vector.

\begin{proposition}[Common upper attainment and deficiencies]
\label{prop:deficiency}
Let $f(x)=\sum_e a_e m_e(x)$ on $[0,1]^n$, with $a_e\ge0$. Discard
supports of size at most one. For each remaining support choose an anchor
$i_e$ of smallest mean and put
\[
 u_e=x_{i_e},\quad S_e=\sum_{j\in e\setminus\{i_e\}}(1-x_j),
 \quad t_e=\min\{u_e,S_e\},\quad
 D_e(X)=X_{i_e}-\prod_{i\in e}X_i.
\]
Then $D_e$ is the nonnegative indicator that the anchor succeeds and some
other support coordinate fails. Moreover,
\begin{align}
 \cav f(x)&=\sum_e a_eu_e,\label{eq:common-upper}\\
 \tgap(f;x)&=\sum_e a_et_e,\label{eq:term-deficiency}\\
 \hgap(f;x)&=\max_{P\in\laws_{[0,1]^n}(x)}
                     \sum_e a_e\E_P D_e(X).\label{eq:hull-deficiency}
\end{align}
\end{proposition}
\begin{proof}
The common-threshold law attains every single-product upper envelope
simultaneously. Positivity gives \eqref{eq:common-upper}, and subtracting
the two single-product envelopes gives \eqref{eq:term-deficiency}.
Finally, every admissible law has $\E X_{i_e}=u_e$, so minimizing the
expected polynomial value is equivalent to maximizing the weighted
sum of deficiencies.
\end{proof}

The common-upper statement is classical
\citep[Theorems 4--5 in the authors' manuscript]{Luedtke2012}. A distribution
used to prove a gap bound must work for all terms together; independently
optimizing the distribution for each term merely recovers the termwise
relaxation. Mixtures of admissible laws preserve the means, and nonnegative
deficiencies allow the guarantees of different laws to be added.

\begin{lemma}[Two classes handled by independence]\label{lem:easy-terms}
Put $c_0=1-e^{-1}$. Independent Bernoulli rounding gives
$\E D_e\ge c_0t_e/2$ if either every support mean exceeds $1/2$, or at
least two support means are at most $1/2$.
\end{lemma}
\begin{proof}
Independence gives $\E D_e=u_e(1-\prod_{j\ne i_e}x_j)$. In the second
case a nonanchor mean is at most $1/2$, giving at least $u_e/2\ge t_e/2$.
In the first case the expression is at least
$u_e(1-e^{-S_e})\ge c_0u_e\min(1,S_e)$. If $S_e\le1$, this is at least
$c_0S_e/2\ge c_0t_e/2$; otherwise it is at least $c_0u_e=c_0t_e$.
\end{proof}

\begin{proposition}[Nonnegative-box transfer]\label{prop:box-transfer}
For a positive multilinear polynomial on $B$ with $0\le\ell_i<u_i$,
substitute $x_i=\ell_i+(u_i-\ell_i)p_i$ and expand the original monomials.
The unit-cube polynomial $\widetilde f$ has nonnegative coefficients and
no larger degree. With its expanded monomials as factors,
\begin{equation}\label{eq:expanded-gap}
 \hgap_B(f;x)=\hgap(\widetilde f;p),\qquad
 \tgap_B(f;x)\le\tgap(\widetilde f;p).
\end{equation}
One common-threshold endpoint law attains the upper envelope of every
original monomial and their sum.
\end{proposition}
\begin{proof}
Expansion coefficients are products of nonnegative endpoints and interval
widths. The full-envelope equality follows from the affine bijection.
For an original factor expanded as $\sum_q g_q$, the sum of the lower
convex envelopes is a convex underestimator and the sum of the upper
concave envelopes is a concave overestimator. Subtract and sum over the
original factors to obtain \eqref{eq:expanded-gap}. Common-threshold
rounding attains all expanded upper envelopes together and therefore also
those of each original factor and their sum.
\end{proof}

Expansion is a proof device and need not be an efficient representation.
It changes the factor incidence structure. Degree bounds transfer through
\eqref{eq:expanded-gap}; bounds based on the original incidence graph
require separate arguments. For bilinear functions, arbitrary box endpoints
cause no such difficulty: an affine change introduces only scaled bilinear
terms and affine summands, regardless of signs.

\subsection{What a spatial certificate counts}\label{subsec:certificate-model}
Let $\mathcal F$ be a nonempty compact feasible set, let $f$ be continuous,
and put $f^*=\min_{x\in\mathcal F}f(x)$. A region oracle assigns a valid
lower bound $\operatorname{LB}(S)$ to feasible points in each region $S$.
A \emph{certificate at target $T_*$} is a finite cover $S_1,\ldots,S_N$
of $\mathcal F$ with $\operatorname{LB}(S_j)\ge T_*$ for all $j$.
Together with a feasible value $U$ and an absolute tolerance $\epsilon\ge0$,
the target $U-\epsilon$ certifies absolute gap at most $\epsilon$.
For $U>0$ and a relative tolerance $0\le\theta<1$, the target
$(1-\theta)U$ certifies $(U-f^*)/U\le\theta$. A size lower bound may
grant the best incumbent $U=f^*$, requiring $f^*>0$ in this relative
convention. Under these tolerance ranges, either target is nondecreasing
in $U$, so worse incumbents require weakly higher targets.

A complete spatial split tree supplies a cover by its final regions.
A lower bound for arbitrary covers is thus independent of branching order
and node selection within the specified region and oracle model. If a
procedure discards feasible regions by propagation or additional cuts,
those deductions must also be included in the certificate model. A lower
bound for one oracle does not automatically survive stronger coupled cuts,
different branch coordinates, or analytic solution of separate components.

\section{Signed bilinear gaps and interaction density}\label{sec:bilinear}
Let $G=(V,E)$ be a finite simple graph and
$b_a(x)=\sum_{ij\in E}a_{ij}x_ix_j$ on $[0,1]^V$, with arbitrary real
coefficients. Zero-coefficient edges may be removed. Let $c^*(a)$ be the
least $c\ge0$ with $\tgap(b_a;x)\le c\hgap(b_a;x)$ everywhere, and set
$c^*(0)=0$. Define
\begin{equation}\label{eq:density-definition}
 \dens(G)=\max_{\varnothing\ne W\subseteq V}\frac{|E(G[W])|}{|W|},
\end{equation}
with value zero for an empty graph. This is half the maximum subgraph
average degree, not fractional arboricity, whose denominator is $|W|-1$.
Let $d_G$ be degeneracy and $\Delta_G$ maximum degree. Then
\begin{equation}\label{eq:density-comparisons}
 d_G/2\le\dens(G)\le d_G,\qquad \dens(G)\le\Delta_G/2.
\end{equation}
The first lower bound uses a subgraph of minimum degree $d_G$; the upper
bounds use a degeneracy orientation and the degree sum identity.

\begin{theorem}[Bilinear upper bounds]\label{thm:bilinear-upper}
For every real coefficient vector on $G$,
\begin{equation}\label{eq:bilinear-upper}
 c^*(a)\le\min\{4\sqrt{\dens(G)},\ 4\sqrt{d_G},\ 2\sqrt{\Delta_G}\}.
\end{equation}
If $G$ is bipartite with parts $P,Q$ and has an edge, let $\Delta_P$ and
$\Delta_Q$ be the maximum degrees in the respective parts. Then
\begin{equation}\label{eq:bipartite-upper}
 c^*(a)\le\sqrt{2\min\{\Delta_P,\Delta_Q\}}
       \le\sqrt{2\min\{|P|,|Q|\}}.
\end{equation}
The factor $\sqrt2$ in the first bipartite bound is sharp. The conclusions
transfer to arbitrary finite boxes after removing fixed coordinates.
\end{theorem}
The proof refines the cut argument of \citet{Boland2017}. The density order
also follows from older operator theory, as shown in
Section~\ref{subsec:schur}; no new general graph-norm principle is claimed.

\subsection{An exact characterization by induced cuts}
For $W\subseteq V$, set
\[
 Q_{a,W}(s)=\sum_{ij\in E(G[W])}a_{ij}s_is_j,\qquad
 L_W(a)=\sum_{ij\in E(G[W])}|a_{ij}|,
\]
and define
\begin{equation}\label{eq:cut-range}
 R_W(a)=\tfrac12\left(\max_{s\in\{-1,1\}^W}Q_{a,W}(s)
                         -\min_{s\in\{-1,1\}^W}Q_{a,W}(s)\right).
\end{equation}
A cut encoded by $s$ has weight $\tfrac12\sum_{ij}a_{ij}-\tfrac12Q_{a,W}(s)$,
so $R_W$ is exactly the maximum cut weight minus the minimum cut weight.
If $L_W>0$ then $R_W>0$: distinct degree-two sign characters are orthogonal
under uniform signs, and their nonzero combination cannot be constant.
Set $L_W/R_W=0$ when both vanish.

\begin{lemma}[Induced-cut characterization]\label{lem:induced-cut}
For every coefficient vector,
\begin{equation}\label{eq:cut-characterization}
 c^*(a)=\max_{W\subseteq V}\frac{L_W(a)}{R_W(a)}.
\end{equation}
Each nonzero induced ratio is attained at $x_i=1/2$ on $W$ and $x_i=0$
outside $W$.
\end{lemma}
\begin{proof}
At that point any admissible vertex law fixes outside coordinates to zero.
Write $X_i=(1+s_i)/2$ on $W$. Since $\E s_i=0$,
\[
 \E b_a(X)=\tfrac14\left(\sum_{ij\in E(G[W])}a_{ij}+\E Q_{a,W}(s)\right).
\]
An extremizing sign vector and its negative, each of probability $1/2$,
attain either extreme of $Q_{a,W}$ with the prescribed means. Thus
$\hgap=R_W/2$, whereas each active product has width $1/2$ and
$\tgap=L_W/2$. These identities hold at every point of
$\{0,1/2,1\}^V$, with $W$ its set of half-valued coordinates: fixing
coordinates to one instead of zero adds only affine terms on $W$.

To extend the comparison, the width of the product $x_ix_j$ is
\[
 \min(x_i,x_j)-\max(0,x_i+x_j-1)
 =\min(x_i,x_j,1-x_i,1-x_j).
\]
Partition the cube by the hyperplanes $x_i=x_j$, $x_i+x_j=1$ ($i\ne j$),
and $x_i=1/2$. The termwise gap is affine on every closed cell.
Every cell vertex is in $\{0,1/2,1\}^V$. Indeed, active equations fix a
coordinate to $0,1/2$, or $1$, or relate two coordinates by equality or
complementation. In a connected component of these relations, one fixed
coordinate determines all values. If none is fixed, a cycle with an odd
number of complementations forces the value $1/2$. Otherwise a free
parameter remains and a small perturbation preserves all active equations,
contradicting extremality. This accounts for every cell vertex.

The graph-hull gap is concave, being $\cav b_a-\vex b_a$. Hence
$\tgap-c\hgap$ is convex on each cell for $c\ge0$. Its value at any
convex combination of cell vertices is at most the combination of its
values there. The grid-point bound therefore proves the upper direction
of \eqref{eq:cut-characterization}, and the attained ratios prove the other.
\end{proof}
The grid identities and upper implication are those of
\citet[Lemma 1 and Corollary 1]{Boland2017}, following
\citet{Luedtke2012}. The cell argument above supplies the boundary details.

\subsection{Cut ranges, row norms, and fractional orientations}
Let $A$ be the symmetric weighted adjacency matrix on $W$, with zero
diagonal and zero entries at nonedges. Write $a_{i,T}=(A_{ij})_{j\in T}$
for $T\subseteq W$ and $a_i=a_{i,W}$. For a rectangular matrix define
$\norm{M}_{\infty\to1}=\max_{u\text{ sign}}\norm{Mu}_1
=\max_{u,v\text{ sign}}v^{\mathsf T}Mu$; empty matrices have norm zero.

\begin{lemma}[Polarization and row estimates]\label{lem:cut-row}
Writing $R=R_W(a)$,
\begin{align}
 R&=\max_{S\subseteq W}\norm{A_{S,W\setminus S}}_{\infty\to1},
    \label{eq:cut-polarization}\\
 R&\ge\frac1{\sqrt2}\max_{S\subseteq W}
          \sum_{i\in S}\norm{a_{i,W\setminus S}}_2,
    \label{eq:cut-khinchin}\\
 R&\ge\frac14\sum_{i\in W}\norm{a_i}_2.\label{eq:cut-full-rows}
\end{align}
\end{lemma}
\begin{proof}
For signs $s,t$, put $u=(s+t)/2$, $v=(s-t)/2$. Their supports are
complementary, and symmetry gives $Q(s)-Q(t)=2v^{\mathsf T}Au$.
Conversely, arbitrary signs on complementary supports arise from
$s=u+v$, $t=u-v$. Maximizing half the difference proves
\eqref{eq:cut-polarization}.

The sharp real Khinchin inequality of \citet{Szarek1976} states that
independent uniform signs satisfy
\begin{equation}\label{eq:khinchin}
 \E\left|\sum_j c_j\varepsilon_j\right|
       \ge\frac1{\sqrt2}\left(\sum_j c_j^2\right)^{1/2}.
\end{equation}
Apply it to every row of $A_{S,T}$ for a fixed partition $S,T$, and sum.
The best sign vector performs at least as well as its average, proving
\eqref{eq:cut-khinchin}. Transposition gives the same estimate for side $T$.

Choose a partition locally maximizing the cut weight with nonnegative
weights $a_{ij}^2$. At every vertex at least half its squared incident
weight crosses, since otherwise moving the vertex improves the cut. Thus
\[
 X=\sum_{i\in S}\norm{a_{i,T}}_2,\qquad
 Y=\sum_{i\in T}\norm{a_{i,S}}_2
 \quad\text{satisfy}\quad
 X+Y\ge\frac1{\sqrt2}\sum_i\norm{a_i}_2.
\]
The two row estimates yield
$R\ge\max(X,Y)/\sqrt2\ge(X+Y)/(2\sqrt2)$, which proves
\eqref{eq:cut-full-rows}.
\end{proof}

\begin{lemma}[Fractional orientation]\label{lem:fractional-orientation}
Every finite graph $H$ admits nonnegative fractions
$\theta_{ij}+\theta_{ji}=1$ on each edge such that
$\sum_j\theta_{ij}\le\dens(H)$ for every vertex. This is the least
possible uniform bound on the vertex loads.
\end{lemma}
\begin{proof}
For $m=|E(H)|>0$, form a flow network with a source, edge nodes, vertex
nodes, and a sink. Source-to-edge capacities are one; arcs from an edge
node to its endpoints have capacity $m+1$; vertex-to-sink capacities are
$t$. A cut of capacity at most $m$ cannot cut an edge-to-endpoint arc.
For a fixed set $U$ of vertex nodes on the source side, placing all edge
nodes internal to $U$ on that side minimizes capacity, giving
$m-|E(H[U])|+t|U|$. At $t=\dens(H)$ every cut has capacity at least $m$.
Max-flow/min-cut supplies a flow saturating the source arcs; outgoing flows
at each edge define its fractions. Conversely, the loads over $U$ sum to
at least $|E(H[U])|$, proving minimality. For $m=0$ use zero loads.
\end{proof}

\begin{proof}[Proof of Theorem~\ref{thm:bilinear-upper}]
Fix an induced graph $H$ with nonempty weighted support and abbreviate
$L=\sum_{ij}|a_{ij}|$, $R=R_{V(H)}(a)$. Weighted Cauchy--Schwarz gives
\[
 \sum_j\theta_{ij}|a_{ij}|
 \le\left(\sum_j\theta_{ij}\right)^{1/2}
      \left(\sum_j\theta_{ij}a_{ij}^2\right)^{1/2}
 \le\sqrt{\dens(H)}\norm{a_i}_2.
\]
Summing counts each absolute edge weight once. By
\eqref{eq:cut-full-rows}, $L\le4\sqrt{\dens(H)}R$.
Also $\norm{a_i}_2\ge\norm{a_i}_1/\sqrt{\Delta_H}$ at every nonisolated
vertex, and $\sum_i\norm{a_i}_1=2L$. The same row bound therefore gives
$R\ge L/(2\sqrt{\Delta_H})$. The graph parameters do not increase in
induced subgraphs, so Lemma~\ref{lem:induced-cut} proves
\eqref{eq:bilinear-upper}.

For bipartite $H$ use its inherited parts $P',Q'$ in
\eqref{eq:cut-khinchin}. Every edge crosses, giving
\[
 R\ge\frac1{\sqrt2}\sum_{i\in P'}\norm{a_i}_2
       \ge\frac{L}{\sqrt{2\Delta_P}}.
\]
Transpose to obtain the bound with $\Delta_Q$ and apply the same induced
characterization. On a unit-weight four-cycle with one negative edge,
$L=4$, $R=2$, and $\Delta_P=\Delta_Q=2$. Its center ratio attains two,
so the stated bipartite constant is sharp. An affine box transformation
scales each bilinear coefficient by the two interval widths and introduces
only affine summands, proving the final assertion.
\end{proof}

The degeneracy bound has a distinct direct proof. Orient along a degeneracy
order with outdegree at most $d_G$ and color vertices in reverse order.
Place a vertex opposite the side receiving at least half the squared weight
of its outgoing edges. Its crossing row norm is at least its outgoing row
norm divided by $\sqrt2$. Summing and applying Cauchy--Schwarz to at most
$d_G$ outgoing entries gives $X+Y\ge L/\sqrt{2d_G}$. The final step of
Lemma~\ref{lem:cut-row} yields $R\ge L/(4\sqrt{d_G})$. This proof is
weaker than the density theorem but uses an ordinary acyclic orientation.

\subsection{Sharp order on every interaction graph}
Define
\[
 \Gamma(G)=\sup_{a\in\R^E\setminus\{0\}}c^*(a),\qquad
 \Gamma_\pm(G)=\max_{a\in\{-1,1\}^E}c^*(a),
\]
with both values zero for edgeless graphs.

\begin{theorem}[Worst coefficients on a fixed graph]\label{thm:graph-density}
For every graph with an edge,
\begin{equation}\label{eq:graph-density}
 \max\left\{1,\sqrt{\frac{\dens(G)}{2\log2}}\right\}
 \le\Gamma_\pm(G)\le\Gamma(G)\le4\sqrt{\dens(G)}.
\end{equation}
Requiring every coefficient to be nonzero does not change $\Gamma(G)$.
It also equals
\begin{equation}\label{eq:center-supremum}
 \sup_{a\ne0}L_V(a)/R_V(a).
\end{equation}
The lower bound has a full-support signing and a witness in $\{0,1/2\}^V$.
\end{theorem}
\begin{proof}
On an $h$-vertex, $m$-edge graph with $m>0$, choose independent edge signs
$\sigma_e$. For fixed vertex signs $s$ and $\lambda>0$,
$\E e^{\lambda Q_\sigma(s)}=\cosh(\lambda)^m\le e^{m\lambda^2/2}$,
and similarly for $-\lambda$. Put $Z=\max_s|Q_\sigma(s)|$. Global sign
reversal preserves $Q_\sigma$, so
\[
 e^{\lambda Z}\le\sum_{s:s_1=1}
       (e^{\lambda Q_\sigma(s)}+e^{-\lambda Q_\sigma(s)}).
\]
Taking expectations and using Jensen gives
$\E Z\le h\log2/\lambda+m\lambda/2$. At
$\lambda=\sqrt{2h\log2/m}$, at least one signing satisfies
\begin{equation}\label{eq:random-sign}
 \norm{Q_\sigma}_\infty\le\sqrt{2mh\log2}.
\end{equation}
The values at different sign vectors need not be independent.

Apply this to a densest induced graph $H=G[W]$. Since $L_W=m$ and
$R_W\le\norm{Q_\sigma}_\infty$,
$L_W/R_W\ge\sqrt{\dens(G)/(2\log2)}$. Give all remaining edges arbitrary
signs; they do not affect the witness face of Lemma~\ref{lem:induced-cut}.
The lower bound one follows from $R_W\le L_W$, and the upper bound is
Theorem~\ref{thm:bilinear-upper}.

Call the supremum in \eqref{eq:center-supremum} $C$. The full center shows
$\Gamma\ge C$. Conversely, every induced ratio is a whole-graph ratio
after zeroing all coefficients outside the induced subgraph. Thus
$\Gamma\le C$. On the compact set $L_V(a)=1$, $R_V$ is continuous and
strictly positive, so the ratio attains a maximum. Perturb zero entries
of a maximizer to nonzero values. Continuity of the full-graph ratio shows
that the full-support supremum is also $C$. This does not assume continuity
of $c^*$ as edges disappear or attainment under the full-support restriction.
\end{proof}

Uniform boundedness of all signed factors in a graph family is therefore
equivalent to uniform boundedness of its induced density. No closure
assumption on the family is needed: fixing outside coordinates exposes a
densest subgraph. Global average density is insufficient. Attaching $q^2$
leaves to one vertex of $K_{q,q}$ produces a connected graph with total
edge-to-vertex ratio less than two but induced density at least $q/2$.

The density is unweighted. Scaling all coefficients by a nonzero scalar
preserves the gap ratio, whereas a sum of coefficient magnitudes scales.
A single edge of arbitrarily small magnitude therefore rules out replacing
$\dens$ by $\max_W\sum_{e\in E(W)}|a_e|/|W|$. The row estimates remain
coefficient-sensitive certificates; a uniform gap bound must use them on
every induced subgraph, not only the full graph.

The order may equivalently be expressed using degeneracy or arboricity.
If $\alpha_G$ is the minimum number of forests partitioning the edges,
then $\dens(G)\le\alpha_G\le d_G\le2\dens(G)$. For the middle inequality,
label outgoing edges in a degeneracy orientation with distinct labels
from $1$ to $d_G$. Each label class is a forest: an undirected cycle would
have an earliest vertex with two outgoing edges of that label. The other
inequalities follow by counting forest edges and
\eqref{eq:density-comparisons}.

\subsection{Positive coefficients, exactness, and finite constants}
For positive coefficients a proper coloring with $\chi\ge2$ colors gives
the classical stronger bound
\begin{equation}\label{eq:positive-bilinear}
 c^*(a)\le
 \begin{cases}2-2/\chi,&\chi\text{ even},\\
 2-2/(\chi+1),&\chi\text{ odd}.
 \end{cases}
\end{equation}
Choose uniformly $\lfloor\chi/2\rfloor$ colors for one side of a cut.
Every edge crosses with probability
$p_\chi=2\lfloor\chi/2\rfloor\lceil\chi/2\rceil/[\chi(\chi-1)]$.
Thus the maximum cut is at least $p_\chi L$, and the minimum cut is zero.
The same coloring works on every induced subgraph, so
Lemma~\ref{lem:induced-cut} gives \eqref{eq:positive-bilinear}.
This is the bound of \citet[Theorem 8 in the authors' manuscript]{Luedtke2012}.
In particular, the positive bilinear factor is at most two and is one when
its nonzero support is nonempty and bipartite. For the all-positive $K_n$
with $n\ge2$, the center ratio is
$\binom n2/\lfloor n^2/4\rfloor$, tending to two. Hence the worst-over-signs
quantifier in Theorem~\ref{thm:graph-density} is essential.

Assume the effective support graph, obtained by deleting zero-coefficient
edges, has at least one edge. Then $c^*(a)=1$ if and only if every cycle
of this graph
has an even number of positive and an even number of negative edges
\citep[Theorem 4]{Boland2017}. Boland et al.\ identify this as a consequence
of the convex-envelope criterion in Theorem 3.10 of Misener, Smadbeck,
and Floudas, \emph{Dynamically generated cutting planes for mixed-integer
quadratically constrained quadratic programs and their incorporation into
GloMIQO 2}; our attribution to that predecessor is through Boland et al.,
who give the alternative cut proof used here.
Indeed, $R_V=L_V$ exactly when one cut
contains all positive edges and no negative edges, and another contains
all negative edges and no positive edges. An edge set is a cut exactly
when it intersects every cycle evenly. Necessity follows by traversing a
cycle. For sufficiency, label vertices by the parity of specified edges
on a path from a root in each component; cycle parity makes the label
path-independent and supplies the cut. The criterion passes to induced
subgraphs, proving sufficiency for $c^*=1$. Conversely, $c^*=1$ implies
$R_V=L_V$ by Lemma~\ref{lem:induced-cut}. In particular, every forest with
at least one edge has factor one for every signing. For empty effective
support the gaps coincide at zero and $c^*=0$ by convention.

The best universal constant in $\Gamma(G)\le K_\rho\sqrt{\dens(G)}$
lies in $[2,4]$: the lower bound is the four-cycle with one negative edge,
with density one and factor two. For degeneracy the corresponding interval
is $[\sqrt2,4]$. These are finite-parameter questions, distinct from the
best constants along sequences with density or degeneracy tending to
infinity; the four-cycle does not give a growing-density lower bound.

The bounds imply $c^*\le4\sqrt2$ for series-parallel graphs by degeneracy,
$c^*\le4\sqrt3$ for planar graphs by density, and
$c^*\le\sqrt{2\min(p,q)}$ for $K_{p,q}$. For $K_n$, the maximum-degree
bound is $2\sqrt{n-1}$, compared with $600\sqrt n$ in
\citet[Theorem 2]{Boland2017}. These are scalar vertical-gap comparisons,
not approximation guarantees for a constrained objective.

\subsection{Exact finite full-signing constants}\label{subsec:finite-signings}
For $n\ge2$, distinguish the largest full-center ratio among complete-graph
signings from the largest ratio over all points:
\[
 M_n=\max_{a\in\{-1,1\}^{E(K_n)}}\frac{\binom n2}{R_V(a)},
 \qquad F_n=\Gamma_\pm(K_n).
\]
All edges here have magnitude one. These maxima differ from the
arbitrary-real-coefficient supremum in \eqref{eq:center-supremum}.

\begin{proposition}[Finite exhaustive computation]\label{prop:finite-signings}
The exact values are
\[
\begin{array}{c|rrrrr}
 n&3&4&5&6&7\\\hline
 \min_a R_V(a)&2&4&4&5&8\\
 M_n&3/2&3/2&5/2&3&21/8\\
 F_n&3/2&3/2&5/2&3&3
\end{array}
\]
\end{proposition}
\begin{proof}
The first row is established by the exhaustive integer computation in
Appendix~\ref{app:finite-signings}, which proves completeness of the
enumeration, gives executable code, and supplies attaining signings.
Dividing $\binom n2$ by that row gives $M_n$.
Lemma~\ref{lem:induced-cut} gives
\[
 F_n=\max_{2\le k\le n} M_k.
\]
Indeed, each induced support of a full signing is a full signing of some
$K_k$, giving the upper bound. Conversely, any signing of $K_k$ extends
to a full signing of $K_n$ by assigning signs to all remaining edges,
and its ratio persists on the face fixing the new coordinates to zero.
Since $M_2=1$, this proves the last row.
\end{proof}

In particular $M_7<M_6$ while $F_7=F_6=3$. An explicit witness is the
$K_6$ signing in Table~\ref{tab:finite-witnesses}, extended by a new vertex
with all incident coefficients positive. At the point whose six old
coordinates are $1/2$ and whose new coordinate is zero, its ratio is three;
at the full center its ratio is only $21/10$. The finite computation
optimizes precisely these listed full-signing problems; it makes no claim
about larger $n$ or the optimal constants for arbitrary real coefficients.

\subsection{The relation to Schur-multiplier theory}\label{subsec:schur}
We record the older implication to separate the density order from the
explicit elementary constant. For a real matrix $M$, its Schur multiplier
acts by entrywise multiplication, with norm $\norm{M}_{\mathrm{Schur}}$
measured in the Euclidean operator norm. Define the real projective norm
\[
 \pi(M)=\inf\left\{\sum_k\norm{u_k}_\infty\norm{v_k}_\infty:
                         M=\sum_k u_kv_k^{\mathsf T}\right\}.
\]
The established results of
\citet[Theorems 1.2 and 2.4]{DavidsonDonsig2007} give
\begin{equation}\label{eq:schur-input}
 \norm{M}_{\mathrm{Schur}}\le2\sqrt{\beta(P)},\qquad
 \pi(M)\le K_G\norm{M}_{\mathrm{Schur}},
\end{equation}
for $|M_{ij}|\le1$ supported on a pattern $P$, where $K_G$ is the real
Grothendieck constant and
\[
 \beta(P)=\max_{R,C:\,|R|+|C|>0}
             \frac{|P\cap(R\times C)|}{|R|+|C|}.
\]
The first bound uses the weighted matrix theorem, which permits continuous
row and column bounds. Using the integer pattern theorem instead would
introduce rounding. These are established norm inequalities, not statements
about sizes of convex extended formulations.

For the symmetric adjacency pattern,
\[
 |P\cap(R\times C)|\le|E(G[R\cup C])|+|E(G[R\cap C])|
                   \le\dens(G)(|R|+|C|).
\]
An edge is counted at most once unless both endpoints lie in $R\cap C$,
when both matrix orientations are counted. Taking $R=C$ equal to a densest
set proves $\beta(P)=\dens(G)$. For weighted adjacency $A$, take
$M=\operatorname{sign}(A)$, with zeros where $A$ vanishes. The definition
of $\pi$ gives
\[
 2L_V=\ip{A}{M}\le\pi(M)\norm{A}_{\infty\to1}.
\]
For sign vectors $u,v$, put $p=(u+v)/2$, $q=(u-v)/2$. Symmetry yields
$u^{\mathsf T}Av=2[Q_{a,V}(p)-Q_{a,V}(q)]$. Multiaffinity puts these two
values between the vertex extrema on $[-1,1]^V$, so
$\norm{A}_{\infty\to1}\le4R_V$. Consequently
\begin{equation}\label{eq:old-density-transfer}
 L_V\le4K_G\sqrt{\dens(G)}R_V.
\end{equation}
Apply this on every induced graph to recover the density upper order.
The direct random-sign argument supplies the matching bilinear lower order.

Equivalently, let
$S(G)=\sup_{a\ne0}L_V(a)/\norm{Q_{a,V}}_\infty$ be the real Sidon
constant of the edge-supported degree-two sign characters, with $S(G)=0$
for edgeless graphs. Because $Q$
has mean zero,
$\norm{Q}_\infty/2\le R_V\le\norm{Q}_\infty$ and hence
$S(G)\le\Gamma(G)\le2S(G)$. For bipartite graphs, flipping all signs in
one part negates $Q$, making its range symmetric and giving
$S(G)=\Gamma(G)$. The constant four improves the particular transfer
\eqref{eq:old-density-transfer}; it is not asserted to be the best constant
available in all prior literature.

\section{Universal bounds for positive multilinear functions}
\label{sec:positive-universal}
Throughout this section the factors are the positive monomials, with equal
supports collected. Let $R_d$ be the supremum of $T/H$ over all dimensions,
positive multilinear polynomials of maximum degree at most $d$, finite
nonnegative boxes, and points with $H>0$. Let $C_n$ be the corresponding
supremum in $n$ variables, with no additional degree restriction. Fixed
coordinates and affine summands may be removed. Proposition~\ref{prop:box-transfer}
shows that both suprema can equivalently be taken on the unit cube: expansion
does not increase dimension or degree and can only increase the termwise gap.
All logarithms without a subscript are natural.

\begin{theorem}[Sharp leading growth]\label{thm:positive-growth}
As the indicated integer parameter tends to infinity,
\begin{equation}\label{eq:positive-growth}
 R_d\sim\frac{\log d}{\log\log d},\qquad
 C_n\sim\frac{\log n}{\log\log n}.
\end{equation}
For every $d\ge2$, put $\Lambda=\max\{16,1+\log d\}$ and $c_0=1-e^{-1}$.
The explicit finite bound
\begin{equation}\label{eq:original-harmonic-finite}
 T\le\frac{24\Lambda}{c_0\log\Lambda}\,H
\end{equation}
holds on every finite nonnegative box. On the unit cube the upper bounds
are witnessed by distributions defined from the complete marginal vector
and degree allowance, independently of the supports and coefficients.
\end{theorem}

We first exhibit the obstruction to a constant bound, then construct a
common distribution that matches its order. The lower family disproves
Conjecture~1 in the author manuscript of \citet[p.~22]{Luedtke2012}, which
asks for a uniform constant for positive multilinear functions on
nonnegative boxes. The signed bilinear examples of Section~\ref{subsec:schur}
concern a different coefficient class.

\subsection{A sparse multiscale obstruction}
\begin{theorem}[Exact dyadic hull gap]\label{thm:dyadic-exact}
For every integer $L\ge2$, there is a unit-coefficient polynomial in
$n_L=2^L+L$ variables, with $2^{L+1}-2$ monomials, maximum degree
$d_L=2^{L-1}+1$, and an interior point such that
\begin{equation}\label{eq:dyadic-exact}
 T=L,\qquad H=s+(L-s)2^{-s},
\end{equation}
where $s\in\{1,\ldots,L-1\}$ satisfies
\begin{equation}\label{eq:dyadic-cutoff}
 (L-s+1)2^{-(s+1)}\le1\le(L-s+2)2^{-s}.
\end{equation}
Either adjacent choice at a cutoff equality gives the same value.
The total number of variable occurrences is $L2^L+2^{L+1}-2$.
\end{theorem}
\begin{proof}
Put $m=2^L$. Use leaf variables $z_1,\ldots,z_m$ and anchors
$a_1,\ldots,a_L$. For each $j$, partition the leaves into $2^j$ blocks
of size $k_j=2^{L-j}$, and define
\begin{equation}\label{eq:dyadic-polynomial}
 f_L(a,z)=\sum_{j=1}^L\sum_{B\in\mathcal P_j}
                   a_j\prod_{i\in B}z_i,
 \qquad a_j=2^{-j},\quad z_i=1-1/m.
\end{equation}
The upper and lower envelopes of each monomial at this point are $2^{-j}$
and $\max(0,2^{-j}-k_j/m)=0$. Thus $T=\cav f_L=L$.
The support and occurrence counts follow by summing $2^j$ and
$2^j(k_j+1)$, respectively.

For an arbitrary binary law with these means, let $R$ count failed leaves
and $N_j$ count blocks hit by a failure. Then $\E R=1$ and
$N_j\le\min(2^j,R)$. The deficiency representation gives
\begin{equation}\label{eq:dyadic-surrogate}
 H=\max\E\sum_j A_jN_j
 \le\sup\E\sum_j A_j\min(2^j,R),\qquad \E A_j=2^{-j}.
\end{equation}
For a fixed law of $R$, an anchor of mass $u$ maximizes its expectation
against any increasing function of $R$ by selecting the upper $u$-tail,
splitting an atom if necessary. Indeed, moving selected mass from a smaller
value to a larger unselected value cannot reduce the objective. A common
nonincreasing quantile $r(t)$, $0<t<1$, therefore makes every anchor
$A_j=\ind\{t\le2^{-j}\}$ optimal simultaneously. Consequently the right
side of \eqref{eq:dyadic-surrogate} equals the supremum of
$\int_0^1 S_{l(t)}(r(t))\,dt$, where
\[
 S_l(r)=\sum_{j=1}^l\min(2^j,r),\qquad
 w_l=\Prob(l(t)=l)=
 \begin{cases}2^{-l-1},&0\le l<L,\\2^{-L},&l=L.\end{cases}
\]
Here $r$ is integer valued, bounded by $m$, and has integral one.

Define $B_q=(L-q+2)2^{-q}$. The sequence decreases strictly, with
$B_1\ge1\ge B_L$, so an $s$ as in \eqref{eq:dyadic-cutoff} exists.
For every $r\ge0$,
\begin{equation}\label{eq:dyadic-affine-certificate}
 S_l(r)\le sr+F_l(s),\qquad
 F_l(s)=\begin{cases}0,&l\le s,\\2^{l-s+1}-2,&l>s.\end{cases}
\end{equation}
For $l\le s$, bound each term by $r$. For $l>s$, cap the first $l-s$
terms at $2^j$ and the remaining $s$ terms at $r$. Integrating yields
\[
 H\le s+\sum_lw_lF_l(s)=s+(L-s)2^{-s}.
\]

To attain this bound, choose $l$ with probabilities $w_l$, let
$A_j=\ind\{j\le l\}$, and for $q\in\{s,s+1\}$ put
\[
 R_l^{(q)}=\begin{cases}0,&l<q,\\2^{l-q+1},&l\ge q.\end{cases}
\]
Each anchor has its required mean and $\E R^{(q)}=B_q$. Mix $q=s$ and
$q=s+1$ with respective probabilities
$(1-B_{s+1})/(B_s-B_{s+1})$ and $(B_s-1)/(B_s-B_{s+1})$.
This gives $\E R=1$. For $l>s$ the two failure counts are the endpoints
of $[2^{l-s},2^{l-s+1}]$, where \eqref{eq:dyadic-affine-certificate}
is an equality. For $l=s$ they are $0,2$, also equality points, and for
$l<s$ both are zero.

It remains to realize the hit counts, a step requiring nested partitions.
Label leaves by $L$-bit strings, with level-$j$ blocks defined by the first
$j$ bits. The first $R$ strings in bit-reversal order have exactly
$\min(2^j,R)$ different prefixes of length $j$: their prefixes are the
reversed residues $0,\ldots,R-1$ modulo $2^j$. Declare these leaves failed
and apply a uniform bitwise XOR shift to every label. The shift preserves
all hit counts and gives each leaf failure probability $R/m$, conditional
on $R$. Averaging gives failure mean $1/m$ for every leaf. The constructed
law therefore attains the bound in the original polynomial.
\end{proof}

The cutoff conditions imply $s=\log_2 L+O(1)$ and $(L-s)2^{-s}=O(1)$.
Thus the ratio in \eqref{eq:dyadic-exact} is asymptotic to
$L/\log_2 L$. The exact formula uses nested dyadic partitions.
For arbitrary equal-size partitions the same surrogate proves the weaker
but still divergent lower bound
\begin{equation}\label{eq:arbitrary-partitions}
 \frac{T}{H}\ge\frac{L}{5+\log_2(1+(L-1)\log2)};
\end{equation}
Appendix~\ref{app:dyadic-alternatives} gives its distinct integral proof
and the dense symmetric predecessor.

The construction may be made homogeneous of degree $d_L$: introduce
$d_L-2$ padding variables and multiply a degree-$k$ term by the first
$d_L-k$ of them. At padding means one both gaps agree exactly with those
above, since all representing laws fix the padding coordinates. Moving
them sufficiently close to one makes all means interior with ratio
arbitrarily close to \eqref{eq:dyadic-exact}. This adds $O(n_L)$ variables
and keeps $O(n_L)$ monomials, but the occurrence count after padding is
$O(n_L^2)$, not the original $O(n_L\log n_L)$.

\subsection{Harmonic failure distributions}
Call a coordinate low when $x_i\le1/2$ and high otherwise. Independence
handles the terms with zero or at least two low coordinates by
Lemma~\ref{lem:easy-terms}. For the remaining terms we require failures
at several probability scales simultaneously.

Fix a real $M\ge d-1$, and put
\begin{equation}\label{eq:harmonic-parameters}
 \Lambda=1+\log M,\qquad \eta=(d-1)/M\in(0,1].
\end{equation}
We use $\eta$ for this cutoff parameter, to distinguish it from graph
density and physical box aspect ratio. Both laws below draw one uniform
$U\in(0,1)$ and set every low coordinate to $\ind\{U\le x_i\}$.
In the \emph{threshold law}, a high coordinate with failure mean
$p=1-x_i$ fails exactly when $U\le p$. In the \emph{harmonic law}, a
coordinate with $p=0$ never fails; otherwise set
\begin{equation}\label{eq:harmonic-law}
 h_p=\min(Mp,1),\quad A_p=1+\log(h_p/p),\quad
 q_p(t)=\frac{\min(1,p/t)}{A_p}\ind\{t\le h_p\}.
\end{equation}
Conditional on $U=t$, high coordinates fail independently with
probabilities $q_p(t)$. Since $M\ge1$, $1\le A_p\le\Lambda$ and
$h_p\ge p$. In particular these probabilities are valid, and
\[
 \int_0^1q_p(t)\,dt
 =\frac{p+p\log(h_p/p)}{A_p}=p.
\]
Thus the two laws preserve the full marginal vector, including boundary
means, and require no information about a particular term.

\begin{lemma}[Complete harmonic guarantee]\label{lem:harmonic-curve}
Consider a term with one low anchor of mean $u$, and high failure means
$p_1,\ldots,p_r$, where $r\le d-1$. Put
$t_e=\min(u,\sum_jp_j)>0$ and $z=\min(\max_jp_j/t_e,1)$. The threshold
and harmonic deficiencies are at least $zt_e$ and $J_{\Lambda,\eta}(z)t_e$,
respectively, where
\begin{equation}\label{eq:harmonic-integral}
 F_{\Lambda,\eta}(v)=1-\exp\left(-\frac{1-\eta v}{\Lambda v}\right),
 \qquad J_{\Lambda,\eta}(z)=\int_z^1 F_{\Lambda,\eta}(v)\,dv,
 \quad F_{\Lambda,\eta}(0)=1.
\end{equation}
\end{lemma}
\begin{proof}
Threshold deficiency is $\min(u,\max_jp_j)\ge zt_e$. If
$\max_jp_j\ge t_e$, the harmonic assertion is nonnegativity. Otherwise
integrate over $\max_jp_j\le t\le t_e$. The anchor succeeds there and
all $p_j\le t$. An inactive coordinate has $p_j<t/M$, so the inactive
failure mass is at most $rt/M\le\eta t$. Active mass is consequently at
least $\sum_jp_j-\eta t\ge t_e-\eta t$. On active coordinates,
$q_{p_j}(t)\ge p_j/(\Lambda t)$. Conditional independence gives union
failure probability at least
$1-\exp(-(t_e-\eta t)/(\Lambda t))$. Substitution $v=t/t_e$ proves the
claim. Zero term gaps require no division and satisfy every asserted
nonnegative bound automatically.
\end{proof}

\begin{theorem}[Optimized finite harmonic bound]\label{thm:harmonic-fixed-point}
For the parameters in \eqref{eq:harmonic-parameters}, let $\zeta$ be the
unique solution of $J_{\Lambda,\eta}(\zeta)=\zeta$ in $(0,1)$. Then
\begin{equation}\label{eq:harmonic-fixed-point-bound}
 R_d\le\frac1\zeta+\frac2{c_0}.
\end{equation}
If $w>0$ solves $we^w=\Lambda e^{-\eta}$ and $w\ge1$, then
\begin{equation}\label{eq:lambert-bound}
 \zeta\ge\frac{w-1/(2w)}{\Lambda},\qquad
 R_d\le\frac{\Lambda}{w-1/(2w)}+\frac2{c_0}.
\end{equation}
The parameter $M$ can be optimized over $[d-1,\infty)$ in either bound.
For $N=1+\log(d-1)$, choosing $M=(d-1)N$ gives, as $d\to\infty$,
\begin{equation}\label{eq:second-order-upper}
 R_d\le\frac{N}{\log N-\log\log N+o(1)}+\frac2{c_0}.
\end{equation}
This is an upper certificate, with no matching second-order lower assertion.
\end{theorem}
\begin{proof}
$F$ is continuous and strictly decreasing, positive on $(0,1)$, and
nonnegative at one. Thus $J$ is strictly decreasing and strictly convex,
and $J(z)-z$ has opposite signs at zero and one. Its unique root is
$\zeta$. Set $a=F(\zeta)\in(0,1)$. The convex tangent inequality gives
\[
 J(z)+az\ge J(\zeta)+a\zeta=(1+a)\zeta.
\]
Mix the harmonic and threshold laws with weights $1/(1+a)$ and
$a/(1+a)$. Lemma~\ref{lem:harmonic-curve} gives deficiency at least
$\zeta t_e$ for every one-low term. This mixture is optimal for the two
\emph{scalar guarantees} $J(z)$ and $z$: at $z=\zeta$ every convex
combination equals $\zeta$. It does not prove optimality over all actual
couplings or imply simultaneous attainability of these lower curves by
an extremal polynomial. Mix this law and independence with weights
proportional to $1/\zeta$ and $2/c_0$. Each term receives at least
$t_e/(1/\zeta+2/c_0)$ from its appropriate component; the other
contributions are nonnegative. Proposition~\ref{prop:deficiency} and
then Proposition~\ref{prop:box-transfer} prove
\eqref{eq:harmonic-fixed-point-bound}.

For the finite certificate, use $x-x^2/2\le1-e^{-x}\le x$ for $x\ge0$
and, for $0<z\le1$,
\[
 \int_z^1(1/v-\eta)^2\,dv
 =1/z-1+2\eta\log z+\eta^2(1-z)\le1/z.
\]
The last inequality follows from $2\eta\log z\le0$ and
$-1+\eta^2(1-z)\le0$. For $0<A\le\Lambda$ this gives
\begin{equation}\label{eq:lambert-integral-certificate}
 \Lambda J(A/\Lambda)
 \ge\log\Lambda-\log A-\eta+\eta A/\Lambda-1/(2A)
 \ge\log\Lambda-\log A-\eta-1/(2A).
\end{equation}
Take $A=w-1/(2w)$. Since $w\ge1$, $A\ge1/2$ and $A<\Lambda$.
Using $w+\log w=\log\Lambda-\eta$,
\[
 A+\log A+\frac1{2A}-(w+\log w)
 =\log\left(1-\frac1{2w^2}\right)+\frac1{4Aw^2}\le0.
\]
The inequality follows from $\log(1-t)\le-t$ and $A\ge1/2$.
Hence $J(A/\Lambda)\ge A/\Lambda$; monotonicity of $J(z)-z$ proves
\eqref{eq:lambert-bound}.

For the chosen cutoff, $\Lambda=N+\log N$ and $\eta=1/N$. The defining
identity for $w$ first gives
$\log\Lambda-\eta-\log\log\Lambda\le w\le\log\Lambda$
for all sufficiently large $N$. Thus $w/\log N\to1$, and substitution
back gives $w=\log N-\log\log N+o(1)$. Since
$(\log N)^2/N\to0$,
\[
 \frac{N}{\Lambda}\left(w-\frac1{2w}\right)
 =\log N-\log\log N+o(1).
\]
This proves \eqref{eq:second-order-upper}.
\end{proof}

\begin{proof}[Proof of Theorem~\ref{thm:positive-growth}]
For its original explicit finite bound, take
$\Lambda=\max\{16,1+\log d\}$ and $M=e^{\Lambda-1}\ge d$.
For a one-low term, put $t_e=\min(u,\sum_jp_j)>0$.
If $\max_jp_j\ge t_e/\sqrt\Lambda$, threshold deficiency is at least
$t_e/\sqrt\Lambda\ge t_e\log\Lambda/\Lambda$; here
$\log\Lambda\le\sqrt\Lambda$ follows by maximizing
$(\log t)/\sqrt t$. Otherwise integrate the harmonic law over
$[t_e/\sqrt\Lambda,t_e/2]$. Inactive failure mass is at most $t$,
so active mass is at least $t_e/2$. The conditional union probability is
at least $c_0t_e/(2\Lambda t)$, using $1-e^{-z}\ge c_0\min(1,z)$.
The harmonic deficiency is at least
\[
 \frac{c_0t_e}{4\Lambda}(\log\Lambda-2\log2)
 \ge\frac{c_0t_e\log\Lambda}{8\Lambda}.
\]
An equal mixture of independence, threshold and harmonic laws now gives
\eqref{eq:original-harmonic-finite}, since independence supplies at least
$c_0t_e/2$ on the other two classes. Every distribution is global.

Theorem~\ref{thm:harmonic-fixed-point} implies
$R_d\le(1+o(1))\log d/\log\log d$.
For the reverse inequality choose the largest $L$ with
$2^{L-1}+1\le d$ in Theorem~\ref{thm:dyadic-exact}. Its degree differs
from $d$ by a factor at most two, asymptotically, and its ratio is
asymptotic to $L/\log_2 L$. For dimension, $C_n\le R_n$; choose the
largest $L$ with $2^L+L\le n$ and add unused coordinates. Again the
dimensions differ by at most a constant factor. These comparisons prove
both equivalences in \eqref{eq:positive-growth} for all integers tending
to infinity, not merely for a dyadic subsequence.
\end{proof}

For comparison, the earlier leading-constant mixture has a distinct explicit
parameterization. Set $\Lambda=\max\{e^6,1+\log d\}$,
$M=e^{\Lambda-1}$, $b=\log\Lambda$, $a=\Lambda/b^2$, $\beta=1/b$, and
\begin{equation}\label{eq:older-leading-mixture}
 h=\frac{(1-1/b)(1-1/(2b^2))(b-3\log b)}{\Lambda}.
\end{equation}
Here $h>0$ because $b\ge6$ and $b-3\log b>0$.
If $\max p_j\ge t_e/a$, threshold deficiency is at least $t_e/a$.
Otherwise, on $[t_e/a,\beta t_e]$, active mass is at least
$(1-\beta)t_e$ and the lower exponent
$(1-\beta)t_e/(\Lambda t)$ is at most $1/b^2$.
Integrating $1-e^{-z}\ge z(1-1/(2b^2))$ gives harmonic deficiency
at least $ht_e$. Weights proportional to $1/h,a,2/c_0$ on harmonic,
threshold, and independence give $R_d\le1/h+a+2/c_0$, also with leading
constant one. The full-curve theorem retains more of the same distributions'
information and improves this certificate's lower-order loss.

A useful interpretation is simultaneous interpolation between Fr\'echet
bounds. The global law with $Z=1/\zeta+2/c_0$ satisfies, for every
$e\subseteq[n]$ with $2\le|e|\le d$,
\begin{equation}\label{eq:simultaneous-frechet}
 \E\prod_{i\in e}X_i\le(1-Z^{-1})\min_{i\in e}x_i
       +Z^{-1}\max(0,\sum_{i\in e}x_i-|e|+1).
\end{equation}
The best fraction guaranteed for every marginal vector is asymptotic to
$\log\log d/\log d$: the theorem gives the lower guarantee, and summing
any proposed guarantee over the dyadic example gives the reverse order
and constant. No convex-envelope evaluation algorithm is inferred.

\subsection{Removing coefficients without changing the degree supremum}
\begin{theorem}[Unit coefficients and interior homogeneous examples]
\label{thm:coefficient-removal}
For each fixed $d\ge2$, the supremum $R_d$ is unchanged by requiring
unit coefficients, degree exactly $d$ for every term, and strictly interior
unit-cube evaluation points simultaneously. Dimension and support size
are unrestricted in this assertion.
\end{theorem}
\begin{proof}
By positive expansion it suffices to approximate any unit-cube example.
Delete affine terms and homogenize with $d-2$ padding variables as above.
At padding means one the gaps are unchanged. Continuity of the finite
polyhedral envelopes (Lemma~\ref{lem:vertex-law}) permits an interior
perturbation preserving any strict target-ratio inequality. Normalize by
the largest coefficient. We now have one fixed homogeneous polynomial
$f=\sum_{e\in E}a_e\prod_{i\in e}x_i$, $0<a_e\le1$, with $n$ variables,
$s$ distinct supports, interior means $x$, and $H(f;x)>0$.

Replace coordinate $i$ by $m$ clones $y_{i1},\ldots,y_{im}$, and set
\[
 \bar y_i=m^{-1}\sum_{r=1}^my_{ir},\qquad F_m(y)=m^df(\bar y),
 \qquad y_{ir}=x_i.
\]
There are $sm^d$ distinct candidate clone monomials: each uses one clone
from each of its $d$ distinct original groups. At the repeated point,
\begin{equation}\label{eq:clone-envelopes}
 \vex F_m=m^d\vex f,\quad \cav F_m=m^d\cav f,\quad
 T(F_m)=m^dT(f),\quad H(F_m)=m^dH(f).
\end{equation}
Indeed, any binary clone law induces random group averages $q\in[0,1]^n$
with $\E q=x$. Conditional independent rounding of the original variables
with means $q_i$ has expected polynomial $f(q)$, so every normalized clone
expectation is an original feasible expectation. Conversely, setting all
clones in a group equal to its original binary coordinate embeds every
original law. This proves equality of the entire feasible expectation
intervals, and hence of both envelopes. Each clone term has its original
list of means, proving the termwise identity.

Independently retain each candidate term arising from $e$ with probability
$a_e$ and give it coefficient one; call the result $G_m$. For $N_0$
independent random variables with range lengths at most one, the bounded-sum
inequality of \citet[Theorem~2]{Hoeffding1963} gives
\begin{equation}\label{eq:bounded-sum}
 \Prob\bigl(\abs{\textstyle\sum_j(Y_j-\E Y_j)}>t\bigr)
 \le2e^{-2t^2/N_0}.
\end{equation}
For completeness, if $Y\in[a,b]$, the second derivative of
$\log\E e^{\lambda Y}$ is its variance under exponential tilting, at most
$(b-a)^2/4$: its squared distance from the interval midpoint is at most
that quantity. Integrating twice gives
$\E e^{\lambda(Y-\E Y)}\le e^{\lambda^2(b-a)^2/8}$.
Multiplication over independent variables, Markov's inequality, and
optimization at $\lambda=4t/N_0$ prove each tail in
\eqref{eq:bounded-sum}. Thus no concentration result is needed without
its hypotheses or proof.

Apply this inequality with $N_0=sm^d$ at each of the $2^{nm}$ clone
vertices, and once to the termwise gap at the repeated point. The latter
is a sum of the same retention indicators with fixed weights in $[0,1]$.
A union bound gives failure probability at most
\begin{equation}\label{eq:clone-union-bound}
 2(2^{nm}+1)\exp(-2t^2/(sm^d))
\end{equation}
for the two simultaneous assertions
\[
 \max_{v\in\{0,1\}^{nm}}|G_m(v)-F_m(v)|\le t,
 \qquad |T(G_m)-T(F_m)|\le t.
\]
Take $t=K m^{(d+1)/2}$ with $K^2>sn\log2/2$. The failure probability
tends to zero as $m\to\infty$, with $f,n,s,d$ fixed. Uniform vertex
error bounds the error of each envelope at every marginal vector by $t$,
by comparing expectations under all feasible laws. Thus the hull-gap
error is at most $2t$. Since $t/m^d\to0$ for $d\ge2$, successful samples'
normalized gaps converge to those of $f$, and their ratios converge
because its hull gap is positive. This proves equality of suprema, not
preservation at a fixed dimension or existence of a small sampled support.
\end{proof}

\section{Cubic gaps and the role of unequal means}\label{sec:cubic}
The divergence in Theorem~\ref{thm:positive-growth} requires increasing
degree. Degree three already distinguishes positive multilinear functions
from the positive bilinear factor-two theorem, however. We prove
\begin{equation}\label{eq:cubic-interval}
 \frac{1610000}{743033}\le R_3\le\frac{31}{12}.
\end{equation}
The lower endpoint is a bound on a supremum supplied by a sequence of
certificates, not a claimed attained finite optimum. We also give exact
finite hull values and a homogeneous unit-coefficient example with only
52 variables. The exact value of $R_3$ remains open.

\subsection{Endpoint orientation and a cubic mixture}
One particularly simple global law applies in every degree. Draw
$U\in[0,1]$ uniformly and independently orient each coordinate left or
right with equal probability. A left coordinate is the indicator of
$[0,x_i]$ and a right coordinate of $[1-x_i,1]$. Call this law $O$.
Every coordinate has mean $x_i$.

\begin{proposition}[Endpoint-orientation bound]\label{prop:endpoint-orientation}
For maximum degree $d\ge2$, on every nonnegative box,
\begin{equation}\label{eq:endpoint-bound}
 T\le\frac{d-1}{1-2^{-(d-1)}}H.
\end{equation}
In particular the constant is $8/3$ for cubics and two for bilinear terms.
\end{proposition}
\begin{proof}
On the unit cube, fix a degree-$(k+1)$ term with minimum mean $u$ at its
anchor. Conditional on the anchor orientation, a same-oriented coordinate
imposes no further restriction on its success interval. An opposite
coordinate excludes an interval of length $a_j=\min(u,1-x_j)$ at the
opposite end. These excluded intervals are nested. If
$a_{(1)}\ge\cdots\ge a_{(k)}$, the first selected index has probability
$2^{-j}$ of being $j$, and hence
\[
 \E_O D_e=\sum_{j=1}^k2^{-j}a_{(j)}
 \ge\frac{1-2^{-k}}{k}\sum_ja_j
 \ge\frac{1-2^{-k}}{k}\min(u,\sum_j(1-x_j)).
\]
The first inequality follows by summing
$(2^{-i}-2^{-j})(a_{(i)}-a_{(j)})\ge0$ over pairs, and the second by
$\sum_j\min(u,p_j)\ge\min(u,\sum_jp_j)$. The fraction
$(1-2^{-k})/k$ decreases with $k$, since it is the average of the first
$k$ entries of $1/2,1/4,\ldots$. Sum the simultaneous guarantees and use
Proposition~\ref{prop:box-transfer}.
\end{proof}

\begin{theorem}[Three-law cubic bound]\label{thm:cubic-upper}
Every positive multilinear polynomial of degree at most three on a finite
nonnegative box satisfies $T\le(31/12)H$. On the unit cube, the global
law
\begin{equation}\label{eq:cubic-mixture}
 P=\frac{18}{31}O+\frac6{31}I+\frac7{31}B
\end{equation}
guarantees deficiency at least $12t_e/31$ for every quadratic or cubic
term. Here $I$ is independent rounding. For $B$, draw one uniform $U$,
set each low coordinate to $\ind\{U\le x_i\}$, and, independently
conditional on $U$, let each high coordinate fail with probability $1/2$
when $U\le2(1-x_i)$ and never otherwise. The weights are optimal among
fixed mixtures of these three laws for a uniform termwise guarantee.
\end{theorem}
\begin{proof}
The definition of $B$ gives each high coordinate failure mean $1-x_i$;
thus all three laws have the prescribed means. Consider sorted cubic
means $u\le v\le w$, and set $t=\min(u,2-v-w)$. Terms with $t=0$
need no estimate. Write $D_O,D_I,D_B$ for their expected deficiencies.

If $u\le v\le1/2$, then $t=u$. Opposite anchor and second-coordinate
orientations give $D_O\ge u/2$, while $D_I=u(1-vw)\ge u/2$.
These two contributions alone give $12t/31$.

If exactly one coordinate is low, set $a=1-v\ge b=1-w$. Directly
integrating the nested exclusion and failure intervals gives
\[
 D_O=\tfrac12\min(u,a)+\tfrac14\min(u,b),\qquad
 D_B=\tfrac12\min(u,2a)+\tfrac14\min(u,2b).
\]
We verify $18D_O+7D_B\ge12\min(u,a+b)$ in all cases. If $a\le u/2$,
the left side is $16a+8b\ge12(a+b)$. If $b\ge u/2$, it is at least
$18(3u/8)+7(3u/4)=12u$. If $b\le u/2\le a\le u$, it is
$9a+8b+7u/2$. When $a+b\le u$, use
$3a+4b\le4u-a\le7u/2$; when $a+b\ge u$, use
$9a+8b\ge8u+a\ge17u/2$. Finally, if $b\le u/2$ and $a\ge u$,
the expression is $25u/2+8b\ge12u$. These cases exhaust $a\ge b\ge0$,
including their boundaries. The $I$ contribution is nonnegative.

If all coordinates are high, put $c=1-u\ge a=1-v\ge b=1-w$.
Integrating the $B$ union probabilities $7/8,3/4,1/2$ over intervals of
lengths $2b,2(a-b),2(c-a)$ gives total failure probability
$c+a/2+b/4$. Subtracting the anchor failure probability $c$ yields
\[
 D_B=D_O=a/2+b/4\ge3(a+b)/8\ge3t/8.
\]
Independence gives $D_I=u(a+b-ab)\ge7t/16$. Indeed, set $s=a+b\le1$
and use $ab\le s^2/4$. For $s\le u$ the normalized expression is at
least $u(1-s/4)\ge u(1-u/4)\ge7/16$. For $s\ge u$ it is at least
$s-s^2/4\ge u-u^2/4\ge7/16$. Both comparisons use $u>1/2$ and
monotonicity on the unit interval. Thus the mixture fraction is at least
$(25/31)(3/8)+(6/31)(7/16)=12/31$.

For sorted quadratic means $u\le v$, set $t=\min(u,1-v)$; always
$D_O=t/2$. If both means are low, $D_I\ge t/2$. If exactly one is low,
$D_B=\min(u,2(1-v))/2\ge t/2$. If both are high, $D_B=t/2$ and
$D_I=ut\ge t/2$. The resulting mixture fractions are at least
$12/31$, $25/62$, and $1/2$. This proves the simultaneous guarantee.
Summation and positive expansion give the claimed box bound.

To check optimality in the stated family, let the arbitrary weights of
$(O,I,B)$ be $(w,r,v)$ with $w+r+v=1$. The following test means give the
listed exact or limiting normalized deficiencies:
\[
\begin{array}{c|ccc}
\text{means}&D_O/t&D_I/t&D_B/t\\\hline
(u,1/2),\ 0<u\le1/2&1/2&1/2&0\\
(\varepsilon,1-\varepsilon^2,1-\varepsilon^2),\ \varepsilon\downarrow0
 &3/8&0&3/4\\
(u,1-u/2,1-u/2),\ u\downarrow1/2\text{ from above}&3/8&7/16&3/8
\end{array}
\]
In the second row $t=2\varepsilon^2$ for small $\varepsilon$ and the
independent fraction is $\varepsilon-\varepsilon^3/2$. In the last row
$1/2<u<2/3$ ensures $t=u$ and the independent fraction is $u-u^2/4$.
Therefore a universal fraction $\alpha$ obeys
\[
 \alpha\le\tfrac12-\tfrac v2,\qquad
 \alpha\le\tfrac38+\tfrac{3v}8-\tfrac{3r}8,\qquad
 \alpha\le\tfrac38+\tfrac r{16}.
\]
Average these inequalities with weights $3/31,4/31,24/31$. The variable
coefficients cancel, giving $\alpha\le12/31$.
This conclusion concerns fixed-mixture termwise guarantees; it does not
identify $R_3$, or exclude a stronger coefficient-dependent analysis.
\end{proof}

\subsection{A scalar certificate and the strongest limiting lower bound}
For a set $Q$ of variables, write
$E_k(Q)=\sum_{S\subseteq Q,\,|S|=k}\prod_{i\in S}x_i$.
At a binary vertex this is $\binom{K}{k}$ when $K$ counts successes in
$Q$. Count-polynomial identities below are vertex identities, not
substitutions defining the multilinear functions at continuous points.

\begin{lemma}[Cubic scalar minorant with positive slack]
\label{lem:cubic-scalar}
On $[0,1]^3$, let
\begin{align*}
 F(a,b,c)&=6c^3+27bc^2+18b^2+30ac+20ab+9a^2,\\
 \ell(a,b,c)&=37a+\tfrac{79}{2}b+38c-\tfrac{103}{3}.
\end{align*}
Then $F\ge\ell+\delta_*$, where $\delta_*=901/120000$.
\end{lemma}
\begin{proof}
For fixed $c$, the Hessian of $h=F-\ell$ in $(a,b)$ is
$\left(\begin{smallmatrix}18&20\\20&36\end{smallmatrix}\right)$,
positive definite with determinant 248. For $0\le c\le3/10$, set
$a=1$ and $b_*(c)=13/24-3c^2/4\in(0,1)$. There $\partial_bh=0$ and
$\partial_ah=-49/6+30c-15c^2\le-31/60<0$.
The convex first-order inequality proves that this is the constrained
minimum on the $(a,b)$ square: $(a-1)\partial_ah\ge0$ for $a\le1$.
Its value is
\[
 p(c)=\frac{-972c^4+576c^3+1404c^2-768c+101}{96}.
\]
For $3/10\le c\le1$, minimizing over unrestricted real $a,b$ is a valid
lower bound. Completing the square gives
\[
 r(c)=\frac{-78732c^4+212256c^3-203796c^2+82032c-11275}{2976}.
\]
For example, with $v=30c-37$, $w=27c^2-79/2$ and
$q=6c^3-38c+103/3$, this value is
$q-(18v^2-20vw+9w^2)/248$, which verifies the elimination directly.

Table~\ref{tab:cubic-bernstein} supplies exact degree-four Bernstein
representations. A row with interval $[l,u]$, denominator $D$, and
numerators $v_0,\ldots,v_4$ means
\[
 p(c)\ \text{or}\ r(c)=D^{-1}\sum_{i=0}^4v_i\binom4i
     t^i(1-t)^{4-i},\qquad t=(c-l)/(u-l).
\]
These are rational polynomial identities, verified by expanding in $t$.
The five intervals cover the required ranges. The Bernstein basis is
nonnegative and sums to one, and the minimum of all 25 coefficients
$v_i/D$ is $1802/240000=\delta_*$. Thus both lower polynomials are at
least $\delta_*$ wherever used, proving the claim.
\end{proof}

\begin{table}[htbp]
\centering\small
\begin{tabular}{ccrl}
\toprule
Polynomial&Interval&$D$&$(v_0,v_1,v_2,v_3,v_4)$\\\midrule
$p$&$[0,1/5]$&60000&$(63125,39125,20975,9395,4133)$\\
$p$&$[1/5,3/10]$&240000&$(16532,6008,1802,3788,11597)$\\
$r$&$[3/10,1/2]$&7440000&$(215357,1285415,1434125,1250375,1008125)$\\
$r$&$[1/2,3/4]$&190464&$(25808,18056,7964,9230,15869)$\\
$r$&$[3/4,1]$&190464&$(15869,22508,34520,45920,31040)$\\\bottomrule
\end{tabular}
\caption{Exact scalar positivity certificate for Lemma~\ref{lem:cubic-scalar}.}
\label{tab:cubic-bernstein}
\end{table}

\begin{theorem}[Analytic cubic lower family]\label{thm:cubic-analytic}
For each multiple $m\ge9$ of nine, take three groups $U,V,W$ of $m$
variables and write $A,B,C$ for their sums. Define
\begin{equation}\label{eq:cubic-analytic-family}
 f_m=2E_3(W)+3B E_2(W)+2mE_2(V)
       +\frac{5m}{3}AC+\frac{10m}{9}AB+mE_2(U).
\end{equation}
At group means $1/4,1/2,3/4$, respectively,
\begin{equation}\label{eq:cubic-analytic-finite}
 \frac{T(f_m)}{H(f_m)}\ge
 \frac{161/4-135/(4m)+6/m^2}
 {223/12-\delta_*+135/(4m)+9/m^2}.
\end{equation}
Consequently $R_3\ge1610000/743033$.
\end{theorem}
\begin{proof}
All coefficients in \eqref{eq:cubic-analytic-family} are positive integers.
At binary counts put $a=A/m,b=B/m,c=C/m$. Expansion of binomial
coefficients gives the exact identity
\begin{equation}\label{eq:cubic-count-expansion}
 \frac{18f_m}{m^3}=F(a,b,c)
 -\frac{18c^2+27bc+18b+9a}{m}+\frac{12c}{m^2}
 \ge F(a,b,c)-72/m.
\end{equation}
Every feasible binary law has expected normalized counts
$(1/4,1/2,3/4)$. Lemma~\ref{lem:cubic-scalar} therefore gives
\[
 \frac{18}{m^3}\vex f_m\ge139/6+\delta_*-72/m.
\]
Counting supports and using \eqref{eq:monomial-envelopes} yields
\begin{align}
 \frac{18}{m^3}\cav f_m&=167/4-153/(4m)+9/m^2,
       \label{eq:cubic-analytic-cav}\\
 \frac{18}{m^3}T(f_m)&=161/4-135/(4m)+6/m^2.
       \label{eq:cubic-analytic-tgap}
\end{align}
For clarity, the six scaled upper contributions in the displayed family
order are
\[
 \frac92-\frac{27}{2m}+\frac9{m^2},\quad
 \frac{27}2-\frac{27}{2m},\quad
 9-\frac9m,\quad\frac{15}2,\quad5,\quad
 \frac94-\frac9{4m}.
\]
Only $2E_3(W)$ has a positive termwise lower envelope; its scaled value
is $3/2-9/(2m)+3/m^2$. This proves the two exact formulas. Subtraction
bounds the scaled hull gap by the denominator in
\eqref{eq:cubic-analytic-finite}. It is positive for $m\ge9$, and the
actual hull gap is positive as well: independence gives strictly smaller
expectations than common-threshold rounding for every positive nonlinear
term at interior means. Division is therefore valid. Taking arbitrarily
large multiples of nine yields
\[
 R_3\ge\frac{161/4}{223/12-901/120000}
       =\frac{4830000}{2229099}=\frac{1610000}{743033}.
\]
Dropping the positive slack recovers the weaker $483/223$; the same
certificate proves the stronger reduced fraction above. Neither minorant
is asserted to be attained in expectation, and neither limiting bound
asserts an attained finite gap ratio or convergence of the actual family
ratios to that lower endpoint.
At $m=36$, even dropping $\delta_*$ from
\eqref{eq:cubic-analytic-finite} gives $T/H\ge16985/8436>2$;
retaining the positive slack makes this finite bound strictly stronger.
\end{proof}

\subsection{Exact finite cubic witnesses}
Finite examples permit a stronger type of statement: matching primal and
dual certificates determine both envelopes exactly. The three-group
polynomials in Table~\ref{tab:cubic-finite} have group sizes $m$ and
means $(1/4,1/2,3/4)$, and take the form
\begin{equation}\label{eq:cubic-six-family}
 f=c_1E_3(W)+c_2B E_2(W)+c_3E_2(V)+c_4AC+c_5AB+c_6E_2(U).
\end{equation}

\begin{proposition}[Exact finite values]\label{prop:cubic-finite}
The envelope and ratio values in Table~\ref{tab:cubic-finite} are exact.
The 18- and 24-variable polynomials have 212 and 464 monomials,
respectively. Complete rational primal and dual certificates, and a
short executable integer/rational checker, appear in
Appendix~\ref{app:cubic-certificates}.
\end{proposition}
\begin{proof}
At binary counts $A,B,C\in\{0,\ldots,m\}$ the polynomial is
\[
 \Phi_m(A,B,C)=c_1\binom C3+c_2B\binom C2+c_3\binom B2
              +c_4AC+c_5AB+c_6\binom A2.
\]
The appendix certifies an affine minorant at every count state and a
probability distribution of count states with means $(m/4,m/2,3m/4)$
attaining it. Conditional uniform subsets in each group realize every
individual marginal, so these give matching lower-envelope certificates.
The upper and termwise lower values follow directly from the six orbit
sizes and monomial bounds in the appendix. Subtracting yields the ratios.
\end{proof}

\begin{table}[htbp]
\centering\small
\begin{tabular}{rllrrl}
\toprule
$n$&$(c_1,\ldots,c_6)$&$\cav f$&$\sum_e\vex f_e$&$\vex f$&$T/H$\\\midrule
18&$(2,3,9,10,7,7)$&$1647/4$&$10$&$2750/13$&$20891/10411$\\
24&$(2,3,13,12,8,7)$&$971$&$28$&$3572/7$&$6601/3225$\\
192&$(2,3,120,105,70,63)$&$587944$&$20832$&$34172072/105$&$7443345/3445256$\\\bottomrule
\end{tabular}
\caption{Exact cubic witnesses, all with ratio above two.}
\label{tab:cubic-finite}
\end{table}

The 24-variable example also gives a quantitative homogeneous interior
variant. Write it as $f=P+Q$, where $P$ is cubic and $Q$ quadratic,
and put $g(x,z)=P(x)+zQ(x)$ with $z=999/1000$.
All 25 coordinates are interior and all 464 monomials are cubic, with
integer coefficients at most 13. Their upper and lower termwise values
are unchanged, so $\cav g=971$ and $T(g)=943$.
The sum of quadratic coefficients is
\[
 13\binom82+12\cdot8^2+8\cdot8^2+7\binom82=1840.
\]
Thus for any feasible law,
$\E g\ge\vex f-1840/1000$, regardless of dependence involving $Z$.
Consequently
\begin{equation}\label{eq:homogeneous-25}
 H(g)\le3225/7+46/25=80947/175,\qquad
 \frac{T(g)}{H(g)}\ge\frac{165025}{80947}>2.
\end{equation}
This is a certified bound, rather than an exact new hull value.

For completeness, it yields a finite nonconstructive unit-coefficient
example through Theorem~\ref{thm:coefficient-removal}. Normalize by 13,
clone each variable $m=1000$ times, and take error allowance $t=m^3/10$.
The failure probability in \eqref{eq:clone-union-bound} is at most
\begin{equation}\label{eq:finite-25000-probability}
 2(2^{25m}+1)\exp(-m^3/23200)<1.
\end{equation}
Indeed its logarithm is at most $(25m+2)\log2-m^3/23200<0$ at $m=1000$,
even with $\log2<1$. Since $943-2(80947/175)=3131/175$, every
successful sample satisfies
\begin{equation}\label{eq:finite-25000-margin}
 T(G_m)-2H(G_m)\ge m^3\left(\frac{3131}{2275}-\frac12\right)>0.
\end{equation}
It is nonempty and has positive hull gap at its interior point. This
proves existence in 25,000 variables with $464m^3$ candidate monomials;
no successful sample is claimed to be listed. The explicit construction
below is much smaller, while coefficient removal remains useful for
approximating an arbitrary fixed-degree ratio.

\subsection{Two means and an explicit homogeneous unit-coefficient example}
\begin{theorem}[Two-level cubic family]\label{thm:cubic-two-level}
For each positive multiple $m$ of four, take two groups $U,W$ of $m$
variables with means $1/2,3/4$, and define
\begin{equation}\label{eq:two-level-family}
 f_m=A E_2(W)+\frac{5m}{4}E_2(U),\qquad A=\sum_{i\in U}x_i.
\end{equation}
Then
\begin{equation}\label{eq:two-level-ratio}
 \frac{243}{115}(1-1/m)\le\frac{T(f_m)}{H(f_m)}\le\frac{243}{115}.
\end{equation}
For $m=4,8,12,16$ its exact ratios are, respectively,
$27/16,21/11,99/50,135/67$. Among positive multiples of four,
this family's ratio exceeds two if and only if $m\ge16$.
There is an explicit polynomial
with 52 variables and 4,320 distinct homogeneous cubic monomials, every
coefficient one, at an interior point with ratio at least $2700/1343>2$.
\end{theorem}
\begin{proof}
For real $a,c\in[0,1]$, set $F_2(a,c)=ac^2+5a^2/4$ and
$\ell_2(a,c)=11a/6+20c/27-95/108$. The identity
\begin{equation}\label{eq:two-level-scalar}
 F_2-\ell_2=\frac54\left(a-\frac{11-6c^2}{15}\right)^2
       +\frac{(1-c)(3c-2)^2(3c+7)}{135}
\end{equation}
proves $F_2\ge\ell_2$ on the square. Equality holds at $(1/3,1)$ and
$(5/9,2/3)$; their mixture with weights $1/4,3/4$ has mean
$(1/2,3/4)$ and expected value $16/27$. Thus the scalar envelope value
at that mean is exactly $16/27$.

At binary counts $A=ma,C=mc$,
\[
 \frac{2f_m}{m^3}=F_2(a,c)-\frac{ac+(5/4)a}{m}.
\]
Every feasible law has $\E a=1/2$ and $\E c=3/4$, and $ac\le a$.
It follows that $2\vex f_m/m^3\ge16/27-9/(8m)$. Each monomial
has upper envelope $1/2$ and lower envelope zero; therefore
\begin{equation}\label{eq:two-level-cav}
 \frac{2}{m^3}\cav f_m=\frac{2}{m^3}T(f_m)=\frac98(1-1/m).
\end{equation}
Subtraction gives $2H(f_m)/m^3\le115/216$.
For the reverse bound, choose one of the two equality atoms above, then
sample all coordinates independently conditional on its two probabilities.
Distinct variables in each term have conditional product expectation,
so $2\E[f_m\mid a,c]/m^3=(1-1/m)F_2(a,c)$. This is a feasible law and
gives $2\vex f_m/m^3\le(1-1/m)16/27$. Hence
\[
 \frac{115}{216}(1-1/m)\le\frac{2H(f_m)}{m^3}\le\frac{115}{216}.
\]
Together with \eqref{eq:two-level-cav} this proves
\eqref{eq:two-level-ratio}, including positive denominators and convergence
of the actual ratios to $243/115$. For example, $m=20$ gives the lower
bound $4617/2300>2$ in 40 variables.

At $m=16$, the exact integer-grid certificate is
\begin{equation}\label{eq:two-level-16-dual}
 A\binom C2+20\binom A2\ge214A+\frac{177}{2}C-1686
 \qquad(A,C\in\{0,\ldots,16\}).
\end{equation}
Appendix~\ref{app:cubic-certificates} supplies its integer verification.
Uniform success subsets of sizes $8,12$ give equality, with value 1088.
The concave envelope and termwise gap are 2160, so $H=1072$ and
$T/H=135/67$.

For $m=4,8,12$, Table~\ref{tab:two-level-small} in
Appendix~\ref{app:cubic-certificates} gives matching affine minorants and
uniform fixed-count laws, proving the three exact ratios stated above.
All three are below two. For every $m\ge20$, the lower bound in
\eqref{eq:two-level-ratio} is at least $4617/2300>2$, since
$1-1/m$ is increasing. Together with $m=16$, these facts prove the
if-and-only-if assertion among positive multiples of four in
\eqref{eq:two-level-family}; they make no minimum-dimension claim over
other families.

Introduce 20 distinct coordinates $z_k$ and define explicitly
\begin{equation}\label{eq:unit-52}
 g=A E_2(W)+\left(\sum_{k=1}^{20}z_k\right)E_2(U),
 \qquad x_U=1/2,\quad x_W=3/4,\quad z_k=999/1000.
\end{equation}
There are $16\binom{16}{2}+20\binom{16}{2}=4320$ distinct cubic
supports, all with coefficient one. Every term still has upper envelope
$1/2$ and lower envelope zero, so $\cav g=T(g)=2160$. At binary vertices,
\[
 g=f_{16}-\left(\sum_k(1-z_k)\right)E_2(U),\qquad
 0\le E_2(U)\le120.
\]
Every feasible law projects to the original means, giving
$\vex g\ge1088-120(20/1000)=1088-12/5$.
Thus $0<H(g)\le5372/5$ and $T(g)/H(g)\ge2700/1343$.
Positivity of $H$ follows by comparing independent and common-threshold
laws at these interior means.
\end{proof}

The two-level examples show that means in the upper half of the unit
cube do not force factor two. Small upward perturbations of the $1/2$
means preserve a strict ratio above two by continuity, so all means can
be strictly above $1/2$. No quantitative perturbation radius or minimal
dimension claim is implied. The next result explains why at least two
distinct means are needed on the unit cube.

\section{Equal means and classical symmetric envelopes}\label{sec:equal-means}
The envelope of the complete elementary symmetric polynomial is classical.
Sherali's equation~(13) and Theorem~3 \citep[pp.~252--253]{Sherali1997}
give, on the entire unit cube,
\begin{equation}\label{eq:sherali-symmetric}
 \vex E_d(x)=\max\left(0,
  \max_{k=d-1,\ldots,n-1}
  \left\{\binom{k}{d-1}\sum_i x_i-(d-1)\binom{k+1}{d}\right\}\right).
\end{equation}
The following extremal gap formulation is a consequence of that
symmetric-envelope machinery. We give a direct proof at equal means to
make the common-law guarantee explicit.

\begin{theorem}[Exact equal-mean formula]\label{thm:equal-finite}
Fix $n\ge2$ and $0<u<1$. Put $b=\lfloor nu\rfloor$,
$\theta=nu-b$, and, for $2\le d\le n$,
\[
 t_d(u)=\min\{u,(d-1)(1-u)\},\qquad
 q_{n,d}(u)=\frac{(1-\theta)\binom bd+\theta\binom{b+1}d}{\binom nd},
\]
with $\binom{k}{d}=0$ for integers $0\le k<d$. The maximum of $T/H$
over positive multilinear polynomials with a nonzero nonlinear part at
$x=(u,\ldots,u)$ is
\begin{equation}\label{eq:equal-finite}
 R_n^{\mathrm{eq}}(u)=\max_{2\le d\le n}
                \frac{t_d(u)}{u-q_{n,d}(u)}.
\end{equation}
A maximum-degree restriction $2\le D\le n$ restricts this maximum to
$d\le D$. Each candidate ratio is attained by $E_d$.
\end{theorem}
\begin{proof}
Draw $K=b$ or $b+1$ with probabilities $1-\theta,\theta$, then choose
a uniformly random $K$-element success subset. Each mean is $u$ and
each degree-$d$ product expectation is $q_{n,d}(u)$. This single law gives
\[
 H(f;u\boldsymbol1)\ge\sum_ea_e(u-q_{n,|e|}(u)),\qquad
 T(f;u\boldsymbol1)=\sum_ea_et_{|e|}(u).
\]
The quotient is bounded by the largest term ratio, provided its
denominators are positive. To see positivity and attainment, the binary
value of $E_d$ is $\binom Kd$, whose first forward difference is
$\binom K{d-1}$ and is nondecreasing. Its linear interpolation on integer
counts is convex. Jensen's inequality shows that any integer-valued $K$
with mean $nu$ has expected value at least the adjacent-count
interpolation. Our uniform-subset law attains this value with all
individual means correct, establishing $\vex E_d=\binom nd q_{n,d}$.
Its upper envelope is $u\binom nd$. Comparing the minimum with the
independent law, whose product expectation is $u^d$, gives
$q_{n,d}\le u^d<u$. Thus the denominators are positive and every claimed
candidate is attained. The maximum is attained by a maximizing degree.
\end{proof}

\begin{corollary}[Dimension-free equal-mean supremum]\label{cor:equal-infinite}
For $0<u<1$ the exact supremum over dimensions, and the limit of
$R_n^{\mathrm{eq}}(u)$ as $n\to\infty$, is
\begin{equation}\label{eq:equal-infinite}
 R_\infty^{\mathrm{eq}}(u)
 =\max_{k\ge1}\frac{\min\{1,k(1-u)/u\}}{1-u^k}\le2.
\end{equation}
Writing $r_u=u/(1-u)$, a maximizing $k$ is one of
$\max(1,\lfloor r_u\rfloor)$ and $\max(1,\lceil r_u\rceil)$.
For $u\le1/2$ the value is $1/(1-u)$. The uniform constant two is sharp.
\end{corollary}
\begin{proof}
The comparison $q_{n,d}\le u^d$ gives the upper bound by the displayed
maximum after writing $k=d-1$. For fixed $d$,
$q_{n,d}(u)\to u^d$ as $n\to\infty$, since both adjacent counts divided
by $n$ tend to $u$. Thus any fixed-degree candidate is approached by
$E_d$ in increasing dimension.

For $k\ge r_u$ the expression is $1/(1-u^k)$, decreasing with $k$.
For $k\le r_u$ it equals $k/(u+u^2+\cdots+u^k)$, increasing with $k$
because the average of that decreasing sequence decreases. Hence a finite
maximizer occurs at one of the stated two integers. Its fixed-degree
limit proves both the supremum and convergence assertions.

Set $t=k(1-u)/u$. Bernoulli's inequality gives
$u^{-k}=(1+(1-u)/u)^k\ge1+t$, and therefore
\[
 \frac{\min(1,t)}{1-u^k}\le\min(1,t)\frac{1+t}{t}\le2.
\]
When $u\le1/2$ the maximizing integer is $k=1$, giving $1/(1-u)$.
At $u=1/2$, the complete bilinear polynomial for even $n$ has ratio
$2(n-1)/n$, tending to two. This is the classical positive bilinear
sharpness example, now viewed within the equal-mean formula.
\end{proof}

These exact formulas concern equal normalized means on the unit cube.
Positive expansion transfers the degree upper bounds to nonnegative
boxes, but does not identify an original positive-box factorization with
the elementary symmetric unit-cube model. In particular, the precise
finite equal-mean optimization above makes no such additional box claim.

\section{Incidence structure and marginal restrictions}\label{sec:incidence}
The \emph{incidence graph} of a specified factorization has variable vertices
$i$, factor vertices $e$, and one edge $ie$ for each $i\in e$. It is a
simple bipartite graph even when two factors have identical supports. This
is different from the primal variable graph, which joins two variables
occurring together in a factor. All widths below concern the incidence graph.
Its treewidth is the minimum, over tree decompositions, of the largest bag
size minus one; a tree decomposition covers every edge and has connected
bags for each vertex. Its degeneracy is the largest minimum degree of a
nonempty subgraph. Its orientation parameter is the minimum possible maximum
outdegree, with no acyclicity requirement. Its variable frequency $\nu$ is
the largest number of factors containing one variable. A feedback variable
set is a set of variable vertices whose deletion leaves a forest; write
$f_F$ for its size. Affine factors are discarded before these parameters
are measured, unless a factorization is explicitly retained.

\subsection{Ownership and asymmetric orientations}
Write $\kappa=2/(1-e^{-1})$. Let $B(d)$ be any valid universal degree-$d$
upper bound from Section~\ref{sec:positive-universal}, with $B(1)=0$.
In particular one can choose $B(d)\sim\log d/\log\log d$.

\begin{theorem}[Incidence orientation bound]\label{thm:incidence-orientation}
For a positive multilinear polynomial on a unit cube, suppose an incidence
orientation has variable outdegrees at most $r$ and factor outdegrees at
most $s$, where $r,s$ are nonnegative integers. Then
\[
 \tgap\le\bigl[r+\min\{s,B(s+1)\}+\kappa\bigr]\hgap.
\]
In particular frequency $\nu$ gives $\nu+\kappa$, and an orientation with
all outdegrees at most $k$ gives $k+\min\{k,B(k+1)\}+\kappa$.
These conclusions transfer to boxes with zero lower bounds by scaling.
\end{theorem}
\begin{proof}
Call a coordinate low if its mean is at most $1/2$. A term with exactly
one low coordinate will be called a one-low term. All other terms receive
at least $t_e/\kappa$ deficiency from independence by
Lemma~\ref{lem:easy-terms}. In a one-low term the low coordinate is its
anchor, of mean $u_e$. Write $q_i=1-x_i$ for the high-coordinate failure
means. Partition its high variables into $I_e$, whose incidence edges
point into the factor, and $O_e$, whose edges point out. Put
\[
 t_e^I=\min\{u_e,\sum_{i\in I_e}q_i\},\qquad
 t_e^O=\min\{u_e,\sum_{i\in O_e}q_i\}.
\]
Capped sums are subadditive, so $t_e\le t_e^I+t_e^O$.

For $r>0$, independently at each high variable label its incidences into
one-low factors by distinct labels in $\{1,\ldots,r\}$. In round $h$ a
variable is owned by its factor with label $h$, if any. Thus the owned
sets $J_{e,h}$ are disjoint over factors and partition $I_e$ over rounds.
Take one uniform $U$ and give every low coordinate its success interval
$[0,x_i]$. For the high variables owned by $e$, arrange failure sets so
that their union inside $[0,u_e]$ has length
$\min\{u_e,\sum_{i\in J_{e,h}}q_i\}$. This can be done by consecutive
arcs of lengths $\min(q_i,u_e)$ on the circle of circumference $u_e$;
when $q_i>u_e$, add measure $q_i-u_e$ outside that interval. For $u_e=0$
use arbitrary failure sets of the required measures. Unowned variables
also receive arbitrary marginal-correct failure sets. Disjoint ownership
makes these assignments compatible, including when several terms share
one low anchor. Their full deficiencies dominate the owned failure events.
A uniform mixture of rounds therefore gives
\[
 \E D_e\ge\frac1r\sum_h\min\{u_e,\sum_{i\in J_{e,h}}q_i\}
 \ge t_e^I/r.
\]
If $r=0$, the incoming gaps all vanish.

A different global law uses the same low success intervals and high
failure intervals $[0,q_i]$. Since $|O_e|\le s$, it gives
\[
 \E D_e\ge\min\{u_e,\max_{i\in O_e}q_i\}\ge t_e^O/s
\]
when $s>0$; empty outgoing sets have zero gap. Mixing these two laws and
independence with weights proportional to $r,s,\kappa$ proves the bound
$r+s+\kappa$, including zero parameters by omitting their components.

For the degree alternative, form the positive polynomial with terms
$a_eX_{i_e}\prod_{i\in O_e}X_i$ for one-low factors. Its degree is at most
$s+1$, its termwise gap is $\sum_e a_et_e^O$, and its deficiency under
\emph{every} common-marginal law is at most the original total deficiency.
Its hull gap is consequently at most the original hull gap. Thus the
outgoing sum is at most $B(s+1)\hgap$, as well as at most $s\hgap$.
The incoming sum is at most $r\hgap$ by the ownership law, and the other
terms' sum at most $\kappa\hgap$ by independence. Adding proves the
refinement. This argument needs only a polynomial-level degree bound,
not an unproved per-term guarantee from that bound. Orienting all edges
from variables to factors gives the frequency specialization.
\end{proof}

\begin{corollary}[Sharp ownership constant]\label{cor:ownership-sharp}
At points where every term has exactly one low coordinate, suppose each
high variable occurs in at most $r$ factors, with unrestricted low-variable
frequency. The sharp supremum of $\tgap/\hgap$ is $r$ for every integer
$r\ge1$.
\end{corollary}
The ownership proof gives the upper bound using only its $r$ rounds.
The lower bound for $r\ge2$ follows from the next construction with $L=r$;
for $r=1$ a single nonaffine product attains one. This parameter concerns
high variables at the stated point and is distinct from ordinary frequency.

\subsection{Variable radix and exact width}
\begin{theorem}[Variable-radix family]\label{thm:variable-radix}
Let $L\ge2$ and $b\ge L$ be integers, and put $m=b^L$. Index $m$ leaves
by length-$L$ base-$b$ words. Let $\mathcal P_j$ partition them by their
first $j$ digits. With separate anchors $a_1,\ldots,a_L$, define
\[
 f_{L,b}(a,z)=\sum_{j=1}^L\sum_{Q\in\mathcal P_j}
                       a_j\prod_{i\in Q}z_i,
 \qquad a_j=b^{-j},\quad z_i=1-m^{-1}.
\]
Then
\begin{equation}\label{eq:radix-gaps}
 \tgap=L,\qquad \hgap=1+(L-1)/b,
 \qquad \tw(G_{L,b})=L.
\end{equation}
\end{theorem}
\begin{proof}
A level-$j$ block has $m/b^j$ leaves, whose failure means sum to $b^{-j}$.
Its local lower value is zero and upper value $b^{-j}$; $b^j$ blocks give
one unit of termwise gap at each level. For an arbitrary binary law let
$R$ count failed leaves and $N_j$ count hit level-$j$ blocks. Then
$\E R=1$, $N_j\le\min(b^j,R)$, and total deficiency is
$\sum_j A_jN_j$. For a fixed distribution of $R$, a prescribed-mean
indicator has largest correlation with the nondecreasing function
$\min(b^j,R)$ by selecting its upper tail: exchanging selected probability
at a smaller count with unselected probability at a larger count proves
this assertion, with fractional selection at a tied count allowed. All
anchors can use these nested tails simultaneously. Represent the count by
its nonincreasing quantile $R(U)$ and select $A_j=\ind\{U\le b^{-j}\}$.
If exactly $l$ anchors are selected, then
\[
 \sum_{j=1}^l\min(b^j,R)\le R+\sum_{j=1}^{l-1}b^j.
\]
For $l=0$ use $0\le R$. Integrating gives
$\hgap\le1+\sum_{j=1}^{L-1}b^j b^{-j-1}=1+(L-1)/b$.

The probability of level $l$ is $w_l=(b-1)b^{-l-1}$ for $0\le l<L$
and $w_L=b^{-L}$. Two count profiles are
\[
 R_l^{(1)}=b^l\ind\{l\ge1\},\qquad
 R_l^{(2)}=b^{l-1}\ind\{l\ge2\}.
\]
Their means are $M_1=[(L-1)(b-1)+b]/b$ and
$M_2=[(L-2)(b-1)+b]/b^2$. Our parameter range ensures
$M_2\le1\le M_1$, so mixing the profiles with first-profile weight
$(1-M_2)/(M_1-M_2)$ gives mean count one. Each profile attains the
preceding pointwise inequality at every level, including $l=0,1$.
For an integer count $R$, fail the first $R$ leaves in digit-reversal
order. Their first $j$ digits are the reversals of the last $j$ digits of
$0,\ldots,R-1$, so they hit exactly $\min(b^j,R)$ level-$j$ blocks.
A uniform coordinatewise digit shift modulo $b$ preserves every hit count
and gives each leaf conditional failure probability $R/m$. The mixture
therefore has all the required marginals and attains the upper bound.
This argument works for composite radices as well as prime ones.

For treewidth, eliminate leaves first, filling their $L$ factor neighbors
into cliques. These neighbors form nested chains. Next eliminate factor
vertices from level $L$ down to level one. A level-$j$ factor's remaining
neighbors lie among its $j-1$ ancestor factors and anchors of levels
$j,\ldots,L$, at most $L$ vertices. Induction verifies this description:
eliminating a descendant joins its ancestor chain and deeper anchors,
never two unrelated factors. The final $L$ anchors also have elimination
width at most $L$. For the reverse bound choose a level-$(L-1)$ block
with $b$ leaves and contract each deepest bilinear factor into its leaf.
The resulting minor contains $K_{L,b}$, with the $L-1$ ancestor factors
and $a_L$ opposite those leaves. For $b\ge L$, this graph has a clique
minor of order $L+1$: pair $L-1$ vertices on opposite sides into branch
sets, then keep one unused vertex from each side. All branch sets are
connected and pairwise adjacent. Minor monotonicity proves the lower bound.
\end{proof}

The restriction $b\ge L$ is part of the exact formula just proved. In
particular it cannot be replaced silently by $b=2$ for arbitrarily many
levels; Theorem~\ref{thm:dyadic-exact} uses a different cutoff. More
generally, for integers $b\ge2,L\ge2$, the same radix proof gives
$\hgap=s+(L-s)b^{-s}$ whenever $1\le s<L$ and
\[
 M_{s+1}\le1\le M_s,\qquad
 M_q=\frac{(L-q)(b-1)+b}{b^q}.
\]
Indeed, cap the first $l-s$ summands and bound the last $s$ by $R$, giving
$\sum_{j\le l}\min(b^j,R)\le sR+\sum_{j\le l-s}b^j$.
Its integrated constant is $(L-s)b^{-s}$. The profiles
$R_l^{(q)}=b^{l-q+1}\ind\{l\ge q\}$ for $q=s,s+1$ have means $M_q$
and attain equality; mix them to mean one and use digit reversal as above.

\begin{corollary}[Sharp leading structural growth]\label{cor:structural-growth}
Let $W_k,D_k,P_k$ be the supremal unit-cube positive-monomial gap ratios
under, respectively, incidence treewidth, degeneracy, or orientation
parameter at most $k$. For every integer $k\ge2$,
\[
 W_k\ge k,\qquad D_k\ge k+1,\qquad P_k\ge k+1,
\]
whereas
\[
 W_k\le D_k\le P_k\le k+\min\{k,B(k+1)\}+\kappa.
\]
In particular $W_k\sim D_k\sim P_k\sim k$ as $k\to\infty$.
\end{corollary}
\begin{proof}
Treewidth at most $k$ implies degeneracy at most $k$: in a reduced tree
decomposition take a vertex exclusive to a leaf bag and delete it,
then continue on induced subgraphs. A degeneracy order oriented toward
later vertices has outdegrees at most $k$. This proves the upper chain.
Use $L=k$ in Theorem~\ref{thm:variable-radix} and let $b\to\infty$ for
$W_k\ge k$. For the stronger other bounds use $L=k+1$: orient each
deepest factor toward both its leaf and anchor, all other leaf incidences
toward their factor, and the remaining factor incidences toward their
anchor. Outdegrees are at most $k$, since deepest factors have outdegree
two. For degeneracy delete deepest factors of degree two, their isolated
anchor, then leaves of degree $L-1=k$, leaving stars. Letting $b$ grow
proves both lower bounds. This family still has treewidth $k+1$.
For example $L=3,b=5$ gives $15/7>2$ for degeneracy and orientation
parameter two, with treewidth three. Finally $B(k+1)=o(k)$.
\end{proof}

For ordinary frequency, the example in Proposition~\ref{prop:feedback-sharp}
below gives a lower bound $2-1/\nu$ for every $\nu\ge2$.
The dyadic family has maximum frequency $2^L$ (at its deepest anchor), so
choosing the largest power of two at most $\nu$ also gives the lower
asymptotic $(1-o(1))\log\nu/\log\log\nu$. Frequency two has the sharper
exact value $3/2$, proved in Section~\ref{sec:frequency}.

\subsection{A lower bound on the normalized means}
Let $C_\delta$ be the supremal positive-monomial gap ratio on unit cubes
at points whose coordinates are all at least $\delta$, with dimension
and degree unrestricted. This restricts the \emph{evaluation point}; the
envelopes remain those on the full unit cube.

\begin{theorem}[Marginal-floor law]\label{thm:marginal-floor}
For $0<\delta\le1/2$, set
\[
 M=\max\{2,\log(1/\delta)\},\quad \tau=\delta/M^2,\quad
 \Lambda=\log\frac{1+\tau}{\tau},\quad
 I_\delta=\int_0^1\left(1-e^{-1/[\Lambda(z+M^{-2})]}\right)\,dz.
\]
Then $I_\delta>0$ and $C_\delta\le I_\delta^{-1}+\kappa$.
For $1/2<\delta<1$, $C_\delta\le C_{1/2}$. Moreover
\[
 C_\delta\sim\frac{\log(1/\delta)}{\log\log(1/\delta)}
 \quad(\delta\downarrow0).
\]
The same leading asymptotic holds with every coordinate in
$[\delta,1-\delta]$. The upper bound transfers to nonnegative boxes when
all nonfixed normalized means $(x_i-\ell_i)/(u_i-\ell_i)$ are at least
$\delta$.
\end{theorem}
\begin{proof}
The density $h(t)=1/[\Lambda(t+\tau)]$ integrates to one on $(0,1)$.
Use low-coordinate success intervals $[0,x_i]$. For a high coordinate
with failure mean $q<1/2$ define
\[
 q_0(t)=\min\{1,qh(t)\},\quad m_q=\int_0^1q_0(t)\,dt,\quad
 q_1(t)=q_0(t)+\frac{q-m_q}{1-m_q}(1-q_0(t)).
\]
Here $0\le m_q\le q<1$, so $q_1\in[q_0,1]$ and its integral is exactly
$q$, including $q=0$. Conditional on the common uniform variable $U=t$,
sample high failures independently with probabilities $q_1(t)$.
For a set of high coordinates with failure sum $S$, the probability of
some failure is at least $1-e^{-Sh(t)}$. If any $qh(t)\ge1$, a failure
is certain; otherwise $q_1(t)\ge qh(t)$ and the usual product bound proves
it. The clipping case is needed for this conclusion.

A one-low term with anchor mean $u\ge\delta$ consequently has
\[
 \E D_e\ge\int_0^u(1-e^{-S/[\Lambda(t+\tau)]})\,dt.
\]
For $S>0$, substitute $t=uz$, $y=S/u$, and note
\[
 \frac{1-e^{-ay}}{\min(1,y)}\ge1-e^{-a}\quad(a\ge0,y>0).
\]
Concavity proves this for $y\le1$, monotonicity for $y\ge1$.
Since $\tau/u\le M^{-2}$, division by $t_e=u\min(1,y)$ now gives
$\E D_e/t_e\ge I_\delta$. Zero $S$ requires no division. Mix this law
and independence with weights proportional to $I_\delta^{-1},\kappa$;
Lemma~\ref{lem:easy-terms} handles all other factors.

As $\delta\downarrow0$, $\Lambda=M+2\log M+o(1)$ and
$\Lambda/M^2\to0$. The inequalities $1-e^{-a}\le\min(1,a)$ and
$1-e^{-a}\ge a-a^2/2$ give
\[
 I_\delta\le\frac{1+\log\Lambda}{\Lambda},\qquad
 I_\delta\ge\frac{\log\Lambda}{\Lambda(1+\Lambda/M^2)}
                   -\frac{\Lambda-1}{2\Lambda^2}.
\]
For the second bound integrate only over $[1/\Lambda,1]$, using
$z+M^{-2}\le z(1+\Lambda/M^2)$. Hence
$I_\delta\sim\log\Lambda/\Lambda$, proving the upper asymptotic.
The dyadic family with $L=\lfloor\log_2(1/\delta)\rfloor$ has all means
in $[\delta,1-\delta]$ and ratio asymptotic to $L/\log_2 L$, which gives
the matching lower bound. Monotonicity gives the finite large-$\delta$
claim. Finally positive affine expansion on a nonnegative box preserves
the lower restriction on normalized means, so
Proposition~\ref{prop:box-transfer} gives its stated extension.
\end{proof}

\begin{theorem}[Uniform joint asymptotic]\label{thm:joint-growth}
Let $C(n,d,k,\delta)$ be the supremal unit-cube positive-monomial ratio
in dimension at most $n$, degree at most $d$, incidence treewidth at most
$k$, and with every mean at least $\delta$. Assume integers $n,d\ge2$,
$k\ge1$, and $0<\delta\le1/2$. Put
$q=\min\{n,d,1/\delta\}\ge e^e$,
$\phi(q)=\log q/\log\log q$, and $h=\min\{k,\phi(q)\}$.
Along every joint parameter sequence with $h\to\infty$,
\[
 C(n,d,k,\delta)\sim h.
\]
The same statement holds for the two-sided strip and with degeneracy or
the orientation parameter replacing treewidth.
\end{theorem}
\begin{proof}
The separate dimension, degree, marginal-floor, and incidence estimates
give $C\le\min\{(1+o(1))\phi(q),k+o(k)\}$. Their relative errors vanish
uniformly along the stated sequences, since both $q$ and $k$ then diverge.
For a simultaneous lower construction take
\[
 L=\lfloor h\rfloor,\qquad b=\lfloor(q/2)^{1/L}\rfloor,\qquad m=b^L.
\]
Then $b\ge(1-o(1))\log q$, $L/b\le(1+o(1))/\log\log q\to0$,
and eventually $b\ge L\ge2$. The radix family has dimension
$m+L\le q/2+\phi(q)\le q\le n$, degree $m/b+1\le q\le d$, and
width $L\le k$. Its smallest anchor mean and leaf failure mean are
$1/m\ge2/q\ge\delta$, and all other means also lie in the two-sided
strip. Thus its ratio $L/[1+(L-1)/b]\sim h$. Its degeneracy and
orientation parameter are at most its treewidth, proving the alternatives.
This is not an asymptotic assertion for a fixed small value of $k$.
\end{proof}

\section{Feedback variables and frequency two}\label{sec:frequency}
\subsection{A universal repair law}
The next argument works on the original factor scopes, including on boxes
with positive lower endpoints.

\begin{lemma}[Feedback repair and gluing]\label{lem:feedback-repair}
Fix singleton Bernoulli means and a feedback variable set $F$ of size
$f_F$. For any nonnegative local functions $g_e$ and any prescribed-mean
local laws $P_e$ on their scopes, there is one global prescribed-mean law
$P'$ satisfying
\[
 \E_{P'}g_e\ge2^{-f_F}\E_{P_e}g_e\qquad\hbox{for every }e.
\]
\end{lemma}
\begin{proof}
Extend each local law to $F\cup e$ by independently sampling any missing
feedback coordinates. Put $C=2^{f_F}$. Construct a global reference law
$Q$ with a common uniform $U$: each feedback coordinate independently
chooses, fairly, its success interval $[0,x_j]$ or $[1-x_j,1]$;
every outside coordinate uses $[0,x_i]$. Write $Q_F$ for its feedback
marginal and $Q_i$ for its $(F,i)$ marginal. For a feedback state $s$
and outside state $b$, every prescribed-mean local law satisfies
\[
 P(s,b)\le m(s,b;i):=\min\bigl(\{x_j:s_j=1\},
 \{1-x_j:s_j=0\},\{x_i\text{ if }b=1;\ 1-x_i\text{ if }b=0\}\bigr).
\]
The minimum includes all feedback coordinates and the outside coordinate.
Choose the feedback orientations that place every desired event at the
left endpoint if $b=1$, and at the right endpoint if $b=0$. This single
orientation pattern has probability $1/C$, and its intersection has
length $m(s,b;i)$. Therefore $Q_i(s,b)\ge P(s,b)/C$.
The same endpoint argument gives $Q_F(s)\ge P_F(s)/C$, also when there
are no outside variables. For $F=\varnothing$ its unique state has mass
one.

For a fixed extended factor law $P_e$, define the nonnegative residuals
\[
 h_s=Q_F(s)-(P_e)_F(s)/C,\qquad
 r_{i,s,b}=Q_i(s,b)-(P_e)_{F,i}(s,b)/C\quad(i\in e\setminus F).
\]
For each $i$, $\sum_b r_{i,s,b}=h_s$. At state $s$ with $h_s>0$, place
residual mass $h_s$ and sample these outside coordinates independently
with probabilities $r_{i,s,b}/h_s$. For $h_s=0$ all these residuals
vanish. For an empty outside scope place just the mass $h_s$.
The resulting measure $R_e$ has total mass $1-1/C$, and
$P'_e=P_e/C+R_e$ has feedback marginal $Q_F$ and $(F,i)$ marginal $Q_i$,
while dominating $P_e/C$ entrywise.

Condition on any $s$ with $Q_F(s)>0$. The factor laws $P'_e(\cdot\mid s)$
have consistent outside singleton marginals $Q_i(s,b)/Q_F(s)$. Their
incidence graph is a forest. Root each nontrivial component at a factor,
sample its law, then traverse the tree. Each newly reached factor shares
exactly one already sampled variable; sample its new variables using its
local conditional law given that variable. Null conditioning states are
never encountered. Isolated variables use their stated conditional means;
factors with no outside variables need no sampling. This preserves every
conditional factor law. Mix over $s$ with weights $Q_F(s)$. The resulting
global law has marginals $P'_e$ on $F\cup e$, and nonnegativity gives
the claimed domination of expectations. There is one loss $C$, regardless
of the lengths of the forest paths.
\end{proof}

\begin{theorem}[Feedback-variable gap]\label{thm:feedback}
For positive monomials on any finite nonnegative box, deletion of $f_F$
variable vertices to an incidence forest implies $\tgap\le2^{f_F}\hgap$.
\end{theorem}
\begin{proof}
Substitute fixed coordinates first. In normalized endpoint variables $Y$,
an original physical product is $\varphi_e(Y)=\sum_{A\subseteq e}
c_A\prod_{i\in A}Y_i$ with $c_A\ge0$. This expansion is used only
inside the original scope. For each nonempty $A$ choose an index $k_A$
of smallest normalized mean, and let
\[
 \ell_e(Y)=c_\varnothing+\sum_{\varnothing\ne A\subseteq e}c_AY_{k_A}.
\]
On binary vertices $\ell_e\ge\varphi_e$, while every prescribed-mean
law has $\E\ell_e=\cav\varphi_e$, by common upper attainment.
Thus $g_e=\ell_e-\varphi_e$ is a nonnegative function on the original
scope whose local maximum expectation is the original factor's exact
gap. Apply Lemma~\ref{lem:feedback-repair} to maximizing local laws,
then sum with the positive coefficients. The full upper value is the
sum of local upper values, proving the result. Replacing the factors by
expanded hyperedges would change the incidence graph and would not
justify this proof.
\end{proof}

For $f_F=0$ this recovers classical incidence-forest exactness; exactness
of standard linearization on Berge-acyclic hypergraphs is established
in \citet{DelPiaKhajavirad2018}. The general domination lemma is sharp
for arbitrary nonnegative local payoffs: on $f_F$ feedback bits and one
outside bit, all fair, use the $2^{f_F+1}$ indicators of their states.
Every local maximum is $1/2$, while their sum is identically one.
Private variables fixed to one can make the scopes distinct. These
indicators have zero literals, so they do not prove sharpness for positive
monomials when $f_F>1$.

\begin{proposition}[One feedback variable is sharp]\label{prop:feedback-sharp}
For $n\ge2$, let $f(a,x)=a\sum_{i=1}^nx_i+\prod_{i=1}^nx_i$ and evaluate
at $a=1/n$, $x_i=1-1/n$ on the unit cube. Then
$\tgap=2-1/n$ and $\hgap=1$. Its incidence graph has treewidth two and
becomes a tree after deleting $a$.
\end{proposition}
\begin{proof}
The bilinear terms contribute one in total and the large product
contributes $(n-1)/n$. If $R$ counts failed leaves and $A$ is the
binary anchor, the full deficiency is
$\E[AR]+\Prob(R\ge1)-1/n$, with $\E R=1$ and $\E A=1/n$.
Pointwise $AR+\ind\{R\ge1\}\le R+A$, giving the upper bound one.
Exactly one uniformly chosen failed leaf and an independent anchor
attain equality. Write $e_0$ for the large factor and $e_i$ for $ax_i$.
Path bags $\{a,e_0,x_i\}$ with attached bags $\{a,x_i,e_i\}$ give width
two. A cycle rules out width one, and deleting $a$ leaves a tree.
\end{proof}

\subsection{Half-integral degree slabs}
For frequency at most two, form a loopless dual multigraph with factor
vertices. A shared variable is an edge between its two factors; a private
variable is an edge to its own new dummy vertex. Parallel edges are
allowed. Unused variables may be completed independently. Dummy factors
have zero weight or identically zero costs. Let $g$ be its odd girth,
with $g=\infty$ in the bipartite case.

\begin{lemma}[Degree-slab vertices]\label{lem:degree-slab}
For a loopless multigraph and $p\in[0,1]^E$, put
$s_v(p)=\sum_{i\in\delta(v)}p_i$. Every vertex of
\[
 \mathcal S(p)=\{z\in[0,1]^E:
       \lfloor s_v(p)\rfloor\le s_v(z)\le\lceil s_v(p)\rceil\ (v\in V)\}
\]
has fractional edges forming vertex-disjoint odd cycles, all valued $1/2$.
\end{lemma}
\begin{proof}
Let $J$ be the fractional edges at a vertex and $R$ the tight degree rows
incident to them, counting a repeated lower/upper equality once.
The restricted rows must have full column rank on $J$: otherwise a
small perturbation in a null direction preserves every tight constraint
and box bound, contradicting extremality. Hence $|J|\le|R|$.
An integer tight residual degree cannot contain exactly one fractional
edge, so every such row has at least two fractional incidences. Thus
$2|R|\le2|J|$. Equality forces each fractional edge to have both endpoints
in $R$, and each such vertex to have fractional degree two. These are
cycles; an even cycle admits an alternating perturbation and is excluded.
On an odd cycle the two incident fractional values sum to the integer
one, forcing all values to $1/2$. The argument excludes parallel
fractional two-cycles and degree-one dummy vertices automatically.
\end{proof}
This is the usual fractional matching rank argument, here stated for
integer degree slabs; compare the fixed-degree formulation in
\citet[Theorem 2.1]{Barrus2014}.

\begin{theorem}[Frequency two and odd girth]\label{thm:frequency-two}
For positive monomials of frequency at most two on a unit cube,
\[
 \tgap\le\frac{g}{g-1}\hgap\quad(g<\infty),
 \qquad \tgap=\hgap\quad(g=\infty).
\]
Every finite odd-girth constant is sharp, and the unrestricted
frequency-two supremum is $3/2$.
\end{theorem}
\begin{proof}
Use failure probabilities $p_i=1-x_i$, and define
$b_v(p)=\max_{i\in\delta(v)}p_i$ and $c_v(p)=\min\{1,s_v(p)\}$,
with empty values zero. For a failure edge set $J$, write $I_v(J)$ for
its coverage indicator. Exact single-product envelopes give
\begin{equation}\label{eq:coverage-baseline}
 \tgap=\sum_v a_v(c_v-b_v),\qquad
 \hgap=\max_{\Prob(i\in J)=p_i}\sum_v a_v(\E I_v(J)-b_v).
\end{equation}
In particular an approximation of unshifted coverage alone is insufficient.
Decompose $p=\E Z$ into vertices of $\mathcal S(p)$. On each allowed
slab $c_v$ is affine (linear below one, constant above), so
$\E c_v(Z)=c_v(p)$; convexity gives $\E b_v(Z)\ge b_v(p)$.

Given $Z=z$, keep integral edges. On each fractional odd cycle of length
$l$, choose one of its $l$ maximum matchings uniformly, then fairly
choose that matching or its complement in the cycle. Each edge has
marginal $1/2$. A matching leaves its uniformly selected vertex uncovered;
its complement covers every vertex. Thus coverage probability is
$1-1/(2l)$ at each cycle vertex. A vertex with another selected integral
edge is covered surely. Otherwise its $c_v(z)=1$ and $b_v(z)=1/2$.
All vertices off these cycles have equal integral $b_v(z)=c_v(z)$.
With $\alpha=1-1/g$, this proves
\[
 \E[I_v(J)\mid Z]\ge\alpha c_v(Z)+(1-\alpha)b_v(Z).
\]
Averaging retains the baseline:
$\E I_v(J)-b_v(p)\ge\alpha[c_v(p)-b_v(p)]$.
Sum this in \eqref{eq:coverage-baseline}. For a bipartite graph there
are no fractional cycles, so the slab decomposition itself attains
$c_v(p)$ simultaneously, proving exactness.

For sharpness take a unit-weight odd cycle of length $l$ with all edge
means $1/2$. Its termwise gap and coverage baseline both equal $l/2$.
For every selected edge set $J$,
$\sum_vI_v(J)\le|J|+(l-1)/2$: if $|J|\le(l-1)/2$ use $2|J|$,
otherwise use $l$. Expected coverage is therefore at most $l-1/2$,
attained by the matching/complement law. The hull gap is $(l-1)/2$.
\end{proof}
Bipartite equality is a statement about the scalar envelopes. It does
not assert compatibility of all separately feasible product-value
coordinates in the full lifted multilinear polytope. Zero-lower boxes
transfer by scaling; unequal positive aspect ratios need a separate analysis.

\subsection{Convex cardinality factors}
A discrete sequence $\varphi(0),\ldots,\varphi(d)$ is convex if
$\varphi(k+2)-2\varphi(k+1)+\varphi(k)\ge0$. Its cardinality factor
is the unique multiaffine interpolant of the values
$\varphi(\sum_iX_i)$ on binary vertices. It is generally not the literal
continuous composition $\varphi(\sum_ix_i)$. Sequences may be negative,
decreasing, or nonmonotone; only their discrete convexity is needed.

\begin{lemma}[Cardinality envelopes]\label{lem:cardinality-envelopes}
Let $f_e$ be such a factor and $\widehat\varphi_e$ its piecewise-affine
interpolation between integers. Then
\[
 \vex f_e(p)=\widehat\varphi_e\Bigl(\sum_{i\in e}p_i\Bigr).
\]
A prescribed-mean law supported on the two adjacent cardinalities attains
this minimum. One common-threshold law maximizes every such factor and
their sum. If $p_{(1)}\ge\cdots\ge p_{(d)}$ are the local means, its
upper value is
\[
 \varphi_e(0)+\sum_{j=1}^d[\varphi_e(j)-\varphi_e(j-1)]p_{(j)}.
\]
\end{lemma}
\begin{proof}
Jensen's inequality gives the lower bound for the convex function
$\widehat\varphi_e$. The cube slab between the two adjacent integer
sums is integral: two fractional coordinates admit an opposite
perturbation; one fractional coordinate cannot make an integer sum
bound tight and admits a small perturbation. Decompose $p$ into its
binary vertices to attain equality.
Discrete convexity implies supermodularity of the cardinality set
function. Among laws maximizing a factor, choose one maximizing also
$\E|X|^2$. If two incomparable support sets have positive mass, replace
equal mass on them by their intersection and union. Singleton means
are unchanged; supermodularity cannot decrease the primary value,
while the secondary value strictly increases by twice the product of
the two nonempty set differences' sizes. This is impossible. The
maximizer is chain-supported, and the chain law with prescribed means
is precisely the common-threshold law. The same law therefore maximizes
each factor individually, and summing establishes simultaneous upper
attainment. Integrating its cardinality gives the displayed formula.
\end{proof}
The single-factor principles are classical; symmetric envelope formulas
appear in \citet{Sherali1997}, and the common-aspect product formula in
\citet[Proposition 4.1]{AdamsGupteXu2019}.

\begin{theorem}[Cardinality odd-girth bound]\label{thm:cardinality-frequency}
For a sum of convex cardinality factors of frequency at most two,
with each entire factor kept in the relaxation, the sharp constants of
Theorem~\ref{thm:frequency-two} hold unchanged.
\end{theorem}
\begin{proof}
Use success probabilities in the degree slab, and decompose $p=\E Z$
as before. Each local lower envelope is affine on its allowed slab, so
\[
 \E\vex f_v(Z)=\vex f_v(p),\qquad \E\tgap(Z)\le\tgap(p),
\]
where the second assertion also uses concavity of each upper envelope.
At a fractional-cycle vertex $v$, let $m$ count its incident integral
ones. Its local lower value is $\varphi_v(m+1)$, and its upper value
is $[\varphi_v(m)+\varphi_v(m+2)]/2$, since precisely two coordinates
are half-valued. Matching/complement rounding on a length-$l$ cycle
selects zero or two of them with probability $1/(2l)$ each, and one
otherwise. Therefore the expected local cost is exactly
$\vex f_v(Z)+\tgap_v(Z)/l$. Off the cycles the local gap is zero.
Averaging gives
\[
 \vex f(p)\le\E f(X)\le\sum_v\vex f_v(p)+\tgap(p)/g.
\]
Subtract from the common upper value. Bipartite slabs are integral and
need no rounding. Positive bilinear odd cycles prove sharpness.
\end{proof}

\begin{corollary}[Common-aspect physical products]\label{cor:cardinality-box}
For original positive products on boxes $[\ell_i,\rho_e\ell_i]$,
$\ell_i>0$, assume each scope $e$ has a common coordinate aspect ratio
$\rho_e>1$. Then the frequency-two odd-girth theorem applies to the
original factors. Its constant is sharp for every fixed common ratio.
\end{corollary}
\begin{proof}
A normalized binary vertex gives the table
$a_e(\prod_{i\in e}\ell_i)\rho_e^K$, which is discrete convex.
Shared variables enforce ratio compatibility, so different connected
components may have different ratios. No new factors are introduced.
For a bilinear odd cycle on $[1,\rho]^n$, expansion gives affine terms
plus $(\rho-1)^2$ times the unit-cube bilinear polynomial. Both gaps
scale by the same positive number. Fixed coordinates are substituted
first, and ratio one is trivial.
\end{proof}

\subsection{Rational scalar-envelope algorithms}
\begin{proposition}[Frequency-two binary oracle]\label{prop:edge-cover-oracle}
For a rational positive polynomial of frequency at most two on a unit
cube and arbitrary rational $\lambda$, one can minimize
$f(X)-\lambda^{\mathsf T}X$ over binary $X$ in polynomial bit time.
\end{proposition}
\begin{proof}
With $Z=1-X$, the objective is
\[
 -\sum_i\lambda_i+\sum_i\lambda_iZ_i+
                     \sum_v a_v\ind\{v\text{ uncovered by }Z\}.
\]
Dummy vertices have zero penalties. Include every negative-cost edge:
adding it strictly lowers cost and cannot increase an uncovered penalty.
Add its cost to the constant, remove it from the remaining choices, and
set its endpoints' penalties to zero, retaining those vertices. Remaining
costs and penalties are nonnegative. Add a hub $h$, its mate $k$, their
zero-cost edge $hk$, and an edge $vh$ of cost equal to each vertex $v$'s
penalty. Any selected original edge set extends to an edge cover by
paying exactly its uncovered penalties and adding $hk$. Conversely,
retaining original edges of an augmented cover leaves uncovered only
vertices whose penalty edges were paid. These two comparisons prove
value equality with minimum-weight edge cover.
For completeness this reduces to matching: for a nonnegative-cost graph
without isolated vertices let $\mu(v)$ be its cheapest incident cost.
From a matching $M$, add a cheapest edge at each unmatched vertex. Its
cover cost is at most
$\sum_v\mu(v)-\sum_{uv\in M}(\mu(u)+\mu(v)-c_{uv})$.
Conversely an inclusion-minimal cover is a union of stars (a path of
three edges or a cycle would have a removable edge). Choose one edge
from each star as a matching. Every other edge can be charged to its
unmatched leaf, with $\mu(v)$ no larger than its cost. This shows the
minimum of the displayed expression is no larger than the minimum cover
cost. Hence equality holds and maximum-weight matching supplies an
optimal cover. Parallel edges can keep their cheapest representative.
Weighted matching is polynomial-time \citep{Edmonds1965}.
\end{proof}

\begin{proposition}[Convex-degree matching oracle]\label{prop:cardinality-matching}
For explicitly listed rational discrete-convex tables and arbitrary
rational edge costs $c_i$, the binary problem
\[
 \min_{J\subseteq E}\sum_v\varphi_v(\deg_J(v))+\sum_{i\in J}c_i
\]
is solvable in polynomial bit time.
\end{proposition}
\begin{proof}
At vertex $v$ create $d_v$ slots with nondecreasing charges
$\Delta_v(j)=\varphi_v(j)-\varphi_v(j-1)$. For each edge $i=uv$, create
mandatory vertices $i_u,i_v$, joined by an inactive zero-cost edge.
Connect $i_u$ to all slots of $u$ with cost $\Delta_u(j)+c_i$, and
$i_v$ to all slots of $v$ with cost $\Delta_v(j)$. A matching covering
all mandatory vertices either uses the inactive edge or uses two slots;
mixed activation cannot cover both endpoints. Thus it represents one
binary choice per original edge. Every chosen set is representable,
and the occupied slots at a vertex can be reassigned to its cheapest
$k$ slots because their available gadget neighbors are identical.
Their charges telescope to $\varphi_v(k)-\varphi_v(0)$.

To enforce mandatory coverage by ordinary unconstrained matching, let
$W$ be the sum of absolute original gadget edge costs and subtract
$M=2W+1$ per mandatory endpoint covered by each edge. The all-inactive
matching covers every mandatory vertex. Missing even one such vertex
loses bonus $M$, exceeding any possible original cost difference.
Thus an optimum covers them all, also with negative increments or
negative $c_i$. Adding $\sum_v\varphi_v(0)$ and undoing the fixed bonus
recovers the desired value. The gadget has $O(|E|)$ vertices and
$O(\sum_vd_v^2)$ edges with polynomial rational encoding length.
This is the incremental-slot matching mechanism of convex degree
optimization; compare \citet[Theorem 1.2 and Section 3]{DezaOnn2021}.
\end{proof}

\begin{corollary}[Exact scalar envelopes and supporting minorants]
\label{cor:scalar-envelope-algorithm}
Both preceding factor classes admit exact rational convex-envelope
evaluation, a supporting affine minorant, and scalar graph-hull
separation in polynomial bit time. This includes rational common-aspect
physical products. No polynomial-size explicit description of the full
polytope carrying all individual factor values is asserted.
\end{corollary}
\begin{proof}
The finite vertex-law LP has dual
\[
 \max_{\eta,\lambda}\{\eta+\lambda^{\mathsf T}p:
                 \eta+\lambda^{\mathsf T}z\le f(z)\ (z\in\{0,1\}^n)\}.
\]
Its separation oracle is precisely binary minimization of
$f(z)-\lambda^{\mathsf T}z$, with arbitrary signed linear costs.
Affine summands are absorbed into these costs and restored in the
value; unused coordinates optimize independently. Let
$A=\sum_v\max_j|\varphi_v(j)|$, or $A=\sum_va_v$ for positive monomials,
adding the absolute coefficient sum of separate affine terms if present.
Then $|f(z)|\le A$. The dual polyhedron contains no line, since the rows
$(1,z)$ span $\R^{n+1}$, and its attained optimal face contains a vertex.
Such a vertex solves $n+1$ independent tight equations with a zero-one
matrix. Cramer's rule bounds each coordinate by $(n+1)!A$; its nonzero
integer denominator determinant has magnitude at least one. Clearing
input denominators also gives polynomial bit length for each coordinate.
Thus the box of radius $\max\{1,(n+1)!A\}$ retains an optimum and has
polynomial encoding size. Rational separation-to-optimization
\citep[Theorem~6.4.9]{GLS1988} gives the exact optimum and minorant, also at boundary
means. Every vertex inequality extends to the whole cube by multiaffine
interpolation. Sorting gives the common-threshold upper value and a
supporting upper affine function; together with cube bounds these
separate the scalar graph hull. Rational common-ratio products generate
their convex tables by polynomially many rational multiplications.
\end{proof}

\section{The exact incidence-treewidth-two constant}\label{sec:width-two}
\begin{theorem}[One-sided cycle coloring]\label{thm:one-sided-coloring}
In a finite simple bipartite graph $G=(V,F;I)$ of treewidth at most two,
one can color the factor side $F$ with two colors so that every cycle
whose factor vertices all have one color contains an even number of
factor vertices. Consequently every zero-one matrix with this bipartite
support has a partition of its rows into two totally unimodular matrices.
\end{theorem}
The all-cycle assertion is stronger than excluding induced odd-factor
cycles (balancedness). The stronger statement will supply total
unimodularity for arbitrary integer cardinality slabs.

\begin{proof}
We prove a two-terminal series-parallel invariant. A network starts with
an edge; series composition identifies one terminal of each of two
otherwise disjoint networks, and parallel composition identifies both
terminal pairs. Intermediate networks are simple and bipartite subgraphs
of the eventual graph. A coloring is good if its monochromatic-factor
cycles have even factor count. A terminal path is $c$-monochromatic when
all its factor vertices, including factor terminals, have color $c$.
Its parity counts those factor vertices modulo two.

For each network there is a parity $\epsilon\in\{0,1\}$ with all the
following witnesses available. They are possibly different colorings.
For $V,V$ terminals and either color $c$, an \emph{active} good coloring
has at least one $c$-monochromatic terminal path, all such paths have
parity $\epsilon$, and no path is monochromatic in the other color.
If $\epsilon=0$, a \emph{blocking} good coloring with neither type of
monochromatic terminal path is also available. For mixed terminals,
with either prescribed color $c$ of the factor terminal, an active good
coloring has paths of color $c$, all of parity $\epsilon$; a blocking
witness with that prescribed terminal color is available whenever there
is no direct terminal edge. For $F,F$ terminals, with either prescribed
common color $c$, the same active witness exists; if $\epsilon=1$ a
blocking witness with those colors also exists. In addition, for either
assignment of distinct terminal colors, a good coloring exists. Such
colors automatically block monochromatic terminal paths. Finally a
network with a direct terminal edge is mixed and uses $\epsilon=1$.

The base edge satisfies the mixed invariant with parity one. For a
series composition let $\delta$ be one if the shared vertex is a factor,
and zero otherwise. The output parity is
$\epsilon=\epsilon_1\mathbin\oplus\epsilon_2\mathbin\oplus\delta$.
Every path concatenates component terminal paths; subtracting the
shared factor once accounts for $\delta$. Every cycle stays in a
component. Active witnesses concatenate by using the same desired color
and matching the shared factor's color when present. The required
additional witnesses have the following complete case analysis.
\begin{enumerate}
\item For outer types $V,V$ and a shared variable, use opposite active
colors in the two $V,V$ components. No monochromatic path can cross both.
\item For outer types $V,V$ and a shared factor, blocking is required
only if the component parities differ. The even mixed component has no
direct edge and can block with the prescribed shared-factor color;
use an active witness of that color in the other component.
\item For mixed outer types and a shared variable, use the active color
opposite to the outer factor's color in the $V,V$ component and use
that prescribed factor color in the mixed component. This blocks.
\item For mixed outer types and a shared factor, give the shared factor
the color opposite to the outer factor. The $F,F$ component uses its
distinct-terminal witness, and the mixed component uses a compatible
active witness. This again blocks.
\item For outer types $F,F$ of common color and a shared variable,
blocking is required only for output parity one, so the two mixed
component parities differ. Block the even component, which has no
direct edge, and activate the other, respecting both outer colors.
\item For outer types $F,F$ of common color and a shared factor, color
the shared factor oppositely. Both $F,F$ components have their required
distinct-terminal good witnesses and block.
\item For outer $F,F$ terminals of distinct colors and a shared
variable, activate each mixed component in its outer terminal's color.
If instead the shared vertex is a factor, give it either color and use
in each component the equal-color active or distinct-color witness
as appropriate.
\end{enumerate}
These cases include reversed mixed types and exhaust all triples of
terminal types. A nontrivial series composition has no direct outer edge,
so its mixed blocking requirement has also been checked.

In a parallel composition every new simple cycle consists of one
terminal path from each component. If the paths are monochromatic in
one color, its factor parity is their parity XOR plus $\tau$, the
number of factor terminals modulo two. Thus $\tau=0$ for equal terminal
types and $\tau=1$ for mixed types. We check the invariant as follows.
For $V,V$ terminals with equal component parities, activate both in the
same color. If both are even, both may also block. For different parities,
activate the odd component and block the even component. Output parity
is $\epsilon_1\lor\epsilon_2$; whenever both components contribute paths
their equal parities make new cycles even.
For $F,F$ terminals of common prescribed color, activate both when
parities agree; if both are odd, both may also block. If parities differ,
activate the even component and block the odd one. Output parity is
$\epsilon_1\land\epsilon_2$. With distinct prescribed terminal colors,
use their good witnesses in both components; a crossing monochromatic
cycle would contain both differently colored terminals and is impossible.
For mixed terminals with no direct edge, both components can block in
the prescribed factor color. Activate either one and block the other
for an active witness, choosing its parity as output, or block both
for the blocking witness. For mixed terminals with a direct edge,
simplicity implies exactly one component contains it. Activate its
parity-one witness and block the edge-free component. No blocking
output is required. These mixed choices create no crossing monochromatic
cycle. The induction is complete.

A treewidth-two graph has no $K_4$ minor, hence no subdivision of $K_4$.
The classical block characterization supplies a two-terminal
series-parallel representation for each nontrivial block
\citep[Theorem 3.1]{HassinTamir1986}. Apply the invariant to each block.
Root the block-cut forest. If a child block attaches at a factor,
interchange its two colors if needed to match the articulation color;
if it attaches at a variable no matching is needed. Bridges use the
base edge and isolated factors receive arbitrary colors. Every cycle
lies in one block, proving the graph assertion.

For the matrix assertion use Camion's criterion: a $0,\pm1$ matrix
is totally unimodular exactly when every Eulerian submatrix has entry
sum divisible by four \citep[Theorem 6.5]{Cornuejols2000}. In either
color class an Eulerian submatrix has even row and column degrees.
Its bipartite support decomposes into edge-disjoint simple cycles.
Each has an even number of factor vertices, hence length divisible
by four. Its total number of ones is consequently divisible by four,
as required. The argument covers rectangular Eulerian submatrices too.
\end{proof}

\begin{theorem}[Sharp width-two gap]\label{thm:width-two-gap}
A sum of discrete-convex cardinality factors with incidence treewidth
at most two satisfies $\tgap\le2\hgap$. This includes positive
monomials on unit cubes and zero-lower boxes, and original positive
products with a common aspect ratio within each scope. The sharp
supremum is two, including for every fixed common positive aspect ratio
$\rho>1$.
\end{theorem}
\begin{proof}
Partition factor rows into two TU classes by
Theorem~\ref{thm:one-sided-coloring}. For a class $C$ and prescribed
means $p$, the polytope
\[
 0\le z\le1,\qquad
 \Bigl\lfloor\sum_{i\in e}p_i\Bigr\rfloor
 \le\sum_{i\in e}z_i\le
 \Bigl\lceil\sum_{i\in e}p_i\Bigr\rceil\quad(e\in C)
\]
is integral and contains $p$. Indeed row sign changes, duplicate rows,
and unit rows preserve total unimodularity; each basic solution with
integer right-hand side is integral by its determinant $\pm1$.
A binary-vertex decomposition of $p$ attains all class lower envelopes
by Lemma~\ref{lem:cardinality-envelopes}. Under that law every other
factor has nonnegative expected deficiency from its own upper value.
The equal mixture of the two class laws therefore attains at least
half the sum of local gaps. Their common upper extremizer makes this
a lower bound for half the full termwise gap in the scalar hull gap.

For positive unit-cube monomials one can equivalently use failure rows
$\sum_{i\in e}z_i\le1$ when their mean sum is at most one and
$\sum_{i\in e}z_i\ge1$ otherwise. A classwise integral decomposition
attains the failure union bound in every row. The mixed unit-right-hand-side
integrality theorem for balanced matrices suffices for this particular
argument \citep[Theorem 6.13]{Cornuejols2000}; it would not justify
arbitrary integer cardinality slabs without the stronger TU conclusion.
Common-aspect products are cardinality factors by
Corollary~\ref{cor:cardinality-box}. Proposition~\ref{prop:feedback-sharp}
proves the unit-cube lower supremum two.

For sharpness at a fixed positive ratio, put $\epsilon=1/\rho$ and
$\alpha=1-\epsilon$. On $[\epsilon,1]^{n+1}$ use
\[
 f_n(a,x)=\alpha^{-1}a\sum_{i=1}^nx_i+\prod_{i=1}^nx_i,
 \qquad \E A=1/n,\quad\E X_i=1-1/n,
\]
where $a=\epsilon+\alpha A$ and $x_i=\epsilon+\alpha X_i$ are endpoint
variables. The normalized nonlinear bilinear terms are $\alpha AX_i$,
so their gaps sum to $\alpha$. The large product is $\epsilon^R$ for
failure count $R$, with $\E R=1$. Its lower value is $\epsilon$ by
adjacent-cardinality attainment, and its upper value
$1-(1-\epsilon^n)/n$. Hence
\[
 \tgap_n=2\alpha-\frac{1-\epsilon^n}{n},\qquad
 \hgap_n=\max\E[\alpha AR+1-\epsilon^R]-\frac{1-\epsilon^n}{n}.
\]
For any $M>0$, $1-\epsilon^r\le\alpha r$ on integer $r\ge0$ gives
\[
 \E[1-\epsilon^R]\le\alpha(1-\E[R;R>M])+1/M,
 \quad \alpha\E[AR]\le\alpha M/n+\alpha\E[R;R>M].
\]
Markov's inequality was used in the first bound. Thus
$\hgap_n\le\alpha+\alpha M/n+1/M$. Exactly one uniformly failed leaf
and an independent anchor give
$\hgap_n\ge\alpha+\alpha/n-(1-\epsilon^n)/n$.
Taking $M=\sqrt n$ yields $\hgap_n\to\alpha$ and
$\tgap_n/\hgap_n\to2$. The incidence graph is the same width-two graph
as before; all means are interior. Scaling physical coordinates by
$\rho$ maps the box to $[1,\rho]^{n+1}$ and preserves the gaps of the
correspondingly rescaled positive polynomial.
\end{proof}

\subsection{What fails for general nonnegative factors}
\begin{example}[An exact width ratio of three]\label{ex:parity-three}
Use six fair variables in pairs $AB,BC,CA$, and put
\[
 E(u,v)=1-u-v+2uv,\qquad N(u,v)=u+v-2uv.
\]
These are the multiaffine equality and inequality indicators, both in
$[0,1]$ on the square. Define
\[
 g_A=E(AB)E(CA),\qquad g_B=N(AB)E(BC),\qquad g_C=N(BC)N(CA).
\]
Each local upper envelope is one at the center: mix a satisfying
assignment and its complement. Each local lower envelope is zero,
by the analogous mixture of a falsifying assignment. No two factors
can equal one on a binary assignment, because adjacent factors demand
opposite parities on their common pair. Thus the sum's upper envelope
is at most one, attained by the all-zero/all-one mixture, which
satisfies $g_A$. The pair assignment $AB=01,BC=01,CA=00$ and its
complement make all factors zero, proving the sum's lower envelope
zero. Therefore $\tgap=3$ and $\hgap=1$ for these three entire factors.

A central bag $\{A,B,C\}$ and six attached bags
$\{A,B,X_{AB,j}\}$, $\{B,C,X_{BC,j}\}$, $\{C,A,X_{CA,j}\}$,
$j=1,2$, form a width-two decomposition. A four-cycle proves the width
is exactly two. Every variable has frequency two. The factors are
nonnegative on the full cube, but their monomial coefficients have mixed
signs, so neither the positive-monomial nor the convex-cardinality
factor theorem applies.
\end{example}

\begin{proposition}[Independent-set payoff family]\label{prop:independent-payoffs}
For a simple graph $G$ with nonnegative vertex weights $w_v$, place two
fair bits on each edge and assign its endpoints opposite demanded
parities. Give vertex $v$ payoff $w_v$ if all its incident demanded
parities hold, and zero otherwise. The sum of separate \emph{maximum
expected payoffs} is $\sum_vw_v$, while the global maximum is
$\alpha_w(G)$, the maximum independent-set weight. With all constructed
factor scopes retained, including zero-weight factors, a graph with
an edge has incidence treewidth $\max\{2,\tw(G)\}$.
\end{proposition}
\begin{proof}
Each local satisfaction pattern and its complement attain $w_v$ with
fair means. Satisfied vertices form an independent set, since the two
endpoints of an edge demand opposite parities. Conversely fix a maximum
independent set, give each incident edge its selected endpoint's demand,
and choose remaining parities arbitrarily. Complement-pair mixing
satisfies all selected factors with fair means, proving the maximum
$\alpha_w(G)$; any extra satisfied factors still form an independent
set and cannot improve that maximum.
To a decomposition of $G$ attach two size-three bags $\{u,v,X_{uv,j}\}$
for every edge at a bag containing $u,v$. This proves the treewidth
upper bound. Contracting one length-two path per edge and deleting
the other gives $G$ as a minor; each edge also creates a four-cycle,
proving both lower bounds.
\end{proof}
Thus $K_{k+1}$ with unit vertex weights gives maximum-payoff ratio $k+1$
at incidence width $k$ for $k\ge2$. A general payoff maximum ratio is not an envelope-width
ratio: the minima must also be computed. Example~\ref{ex:parity-three}
does compute them and gives the claimed exact width obstruction.

\subsection{Limits of the coloring method}\label{subsec:width-three-boundary}
A distinct open question at incidence treewidth two asks whether the
factors admit two classes whose incidence subgraphs are each chordal
bipartite, meaning they contain no induced cycle of length at least six.
The proved all-cycle parity property permits induced eight-cycles and
does not establish this target. The recorded experiment found such a
partition for all 5,000 generated supports; this is finite evidence only.

The all-cycle coloring invariant cannot extend with a bounded number
of colors to width three. In $K_{3,m}$, $m\ge3$, with its three-vertex
side interpreted as variables, every three factors lie on a six-cycle.
Each permitted color class therefore contains at most two factors,
requiring $\lceil m/2\rceil$ colors, although the graph has treewidth
three. Its all-ones matrix is already TU: determinants of order at
least two vanish. The six-cycles have chords and do not obstruct
balancedness. This separates the stronger coloring mechanism from a
TU-row-partition target.

Finite searches found three balanced classes in 3,000 generated
width-at-most-three supports (212 required more than two), and three
TU classes in 300 further supports. Their accepted widths were
certified by particular elimination orders; rejected greedy orders did
not certify larger width. The TU tests enumerated column signings by
the Ghouila-Houri criterion and used exact coloring backtracking.
These searches are finite evidence, not proofs of universal partitions.
The exact $W_k$ for $k\ge3$, $k$-class balanced or TU partitions, and
sharp planar constants remain open here. A planar bipartite incidence
graph does admit an orientation of maximum outdegree at most two:
Hall's theorem assigns every edge to one of two slots at an endpoint,
since every edge set uses at most twice as many edges as incident
vertices. Orient from the assigned endpoint. Thus the general
orientation theorem gives a finite planar bound, without settling its
sharp value. Auxiliary reductions and the simpler forest-coloring
obstruction are recorded in Appendix~\ref{app:structural-auxiliary}.

\section{Positive boxes and coefficient regularity}\label{sec:positive-boxes}
Let $C_{\rm box}(\rho_B)$ be the supremal gap ratio for positive monomial
factorizations on strictly positive boxes with every coordinate aspect
ratio $u_i/\ell_i\le\rho_B$, over arbitrary dimensions and evaluation
points with positive hull gap. Unlike $C_\delta$, this restricts the box,
not the normalized evaluation point. Fixed coordinates are removed.

\subsection{A coefficient inequality with one ambient orientation law}
Fix an ambient dimension $N\ge2$. Choose a uniform subset $J\subseteq[N]$
of size $\lfloor N/2\rfloor$ and an independent uniform $U\in[0,1]$.
Coordinate $i$ of mean $p_i$ succeeds on $[0,p_i]$ if $i\in J$ and on
$[1-p_i,1]$ otherwise. Call this balanced orientation law $O$. It has
all prescribed means, though for odd $N$ its individual orientation
choices are not fair. The probability that a fixed pair has opposite
orientations is
\[
 q_N=\frac{2\lfloor N/2\rfloor\lceil N/2\rceil}{N(N-1)},\qquad
 \beta_N=q_N^{-1}=
 \begin{cases}2-2/N,&N\text{ even},\\2-2/(N+1),&N\text{ odd}.
 \end{cases}
\]
All local supports and smaller-dimensional induction steps below use
restrictions of this one law. Deleting a deterministic coordinate
neither conditions nor resamples the orientation subset.

For a support $S$ of size $d\le N$, let $K=\sum_{i\in S}X_i$.
Write $C_j,P_j,O_j$ for $\E\binom Kj$ under common-threshold,
independent, and ambient orientation rounding. With
$s=\sum_{i\in S}p_i$, $k=\lfloor s\rfloor$, and $\theta=s-k$, put
\[
 V_j=(1-\theta)\binom kj+\theta\binom{k+1}j.
\]
Degree-zero moments are one; negative orders and orders above support
size are zero. By Lemma~\ref{lem:cardinality-envelopes}, $C_j$ and $V_j$
are the true upper and lower envelope values of the elementary
symmetric polynomial of degree $j$.

\begin{lemma}[Finite spreading inequality]\label{lem:box-coefficient}
For every support $S\subseteq[N]$, every mean vector, and every integer
$j\ge0$,
\begin{equation}\label{eq:box-coefficient}
 P_j-V_j\le\beta_N(C_j-O_j)+(C_{j-1}-P_{j-1}).
\end{equation}
The same statement holds with $\beta_N$ replaced by $2$ and independently
fair endpoint orientations in every coordinate.
\end{lemma}
\begin{proof}
Write $\beta$ for the indicated constant and set
$F_{S,j}=\beta C_j+V_j-P_j-\beta O_j+C_{j-1}-P_{j-1}$.
Orders zero and one vanish; order $d+1$ is $C_d-P_d\ge0$, and higher
orders vanish. For order two let
$L_2=\sum_{a<b}\max(0,p_a+p_b-1)$. Equal orientations give a pair's
upper Fr\'echet value; opposite orientations give its lower value.
Their opposite probability is $1/\beta$, including under restriction.
Thus $O_2=(1-1/\beta)C_2+L_2/\beta$, and
\[
 F_{S,2}=(C_2-P_2)+(V_2-L_2)\ge0.
\]
The first inequality follows from upper optimality, the second because
every law, including an adjacent-cardinality law, respects every pair
lower bound.

We induct on support size. A mean-zero coordinate is deterministic and
its deletion gives $F_{S,j}=F_{S\setminus\{i\},j}$.
A mean-one coordinate and Pascal's identity give
$F_{S,j}=F_{S\setminus\{i\},j}+F_{S\setminus\{i\},j-1}$.
The adjacent-cardinality values obey the same identity after shifting
the sum by one. These are restrictions of the ambient law, not newly
sampled balanced laws, so its correlated coins cause no exception.
All boundary points follow by induction; dimension one is immediate.

For an interior vector and $3\le j\le d$, select distinct coordinates
of global minimum $x$ and global maximum $y$. Move them to $x-h,y+h$
with $0\le h\le\min(x,1-y)$. The sum stays fixed, so $\Delta V_j=0$.
In increasing order of means,
\[
 C_j=\sum_{i=1}^d p_{(i)}\binom{d-i}{j-1}.
\]
Putting the selected minimum first and maximum last also handles ties.
Hence, with $B=\binom{d-1}{j-1}$ and $D=\binom{d-1}{j-2}$,
$\Delta C_j=-Bh$ and $\Delta C_{j-1}=-Dh$; the last coordinate has
zero coefficient since $j\ge3$. If $E_r$ is the elementary symmetric
polynomial in the other $d-2$ means, the unchanged sum and product
change $\Delta(xy)=-h(y-x+h)$ give exactly
\[
 \Delta(-P_j-P_{j-1})
 =h(y-x+h)(E_{j-2}+E_{j-3})\le Dh.
\]
Here $0\le y-x+h\le1$ and $E_r\le\binom{d-2}{r}$.
Finally couple both orientation vectors using the same $U$ and coins.
Increasing a single mean by $h$ changes only that bit, monotonically,
on an event of probability $h$. The increase in $\binom Kj$ is then
$\binom{K_{\rm rest}}{j-1}\le B$. Decreasing the minimum first and
increasing the maximum second therefore gives $\Delta O_j\ge-Bh$,
regardless of dependence among orientation coins. Combining the finite
changes yields $\Delta F_{S,j}\le-\beta Bh-Dh+Dh+\beta Bh=0$.
Take $h=\min(x,1-y)$ to reach a boundary point, whose coefficient is
nonnegative by induction. The original coefficient is at least that
boundary value. No differentiability assumption is used.
\end{proof}

The choice of global extrema in this proof matters. For independently
fair orientations, the stronger Schur-concavity assertion is false:
\begin{equation}\label{eq:schur-counterexample}
 \begin{split}
 F_{[4],3}(2/5,7/10,24/25,97/100)&=6201/12500,\\
 F_{[4],3}(2/5,7/10,959/1000,971/1000)&=4961031/10000000.
 \end{split}
\end{equation}
Spreading the last two coordinates by $1/1000$ increases the coefficient
by $231/10000000$. These rational values follow by averaging the sixteen
orientation patterns, integrating between their interval endpoints,
and evaluating the displayed symmetric moments. They refute arbitrary
pair spreading, without affecting the global-minimum/global-maximum
comparison.

\begin{theorem}[Positive-box upper bounds]\label{thm:positive-box-upper}
For $N\ge2$ nonfixed coordinates and $\rho_B>1$, every positive
multilinear polynomial on a strictly positive box of aspect ratios at
most $\rho_B$ satisfies
\[
 \tgap\le(\rho_B+\beta_N)\hgap\le(\rho_B+2)\hgap.
\]
The common mixture of independence and balanced orientation with weights
$\rho_B/(\rho_B+\beta_N)$ and $\beta_N/(\rho_B+\beta_N)$ captures this
fraction of each original monomial's exact gap simultaneously. For zero
or one nonfixed coordinate all gaps vanish.
\end{theorem}
\begin{proof}
First consider a support on $[1,\rho_B]^N$ and put $t=\rho_B-1$.
Its binary physical value is $(1+t)^K$, so its upper, independent,
orientation, and lower values are $C(t)=\sum_jC_jt^j$, $P(t)$,
$O(t)$, and $V(t)=\sum_jV_jt^j$. The last identity follows from a
single adjacent-cardinality law attaining every order together.
The coefficient of $t^j$ in
$(\beta_N+t)C+V-(1+t)P-\beta_NO$ is $F_{S,j}\ge0$.
Consequently the true local gap satisfies
\begin{equation}\label{eq:physical-local}
 C-V\le\rho_B(C-P)+\beta_N(C-O).
\end{equation}
The same mixture works on all supports in the ambient $N$ coordinates.
Summing with nonnegative coefficients and using common upper attainment
proves the common-aspect conclusion. Fair independent orientation coins
and $\beta=2$ prove $\rho_B+2$ directly in every dimension.

For unequal endpoints $0<\ell_i<u_i$ put
\[
 s_i=\frac{u_i/\ell_i-1}{\rho_B-1}\in(0,1],\qquad
 z_i=\ell_i[(1-s_i)+s_i y_i],\quad y_i\in[1,\rho_B].
\]
This is an affine bijection onto the original box. Each original
monomial becomes a positive sum of physical $y$-monomials. Its true
factor gap is at most the sum of their gaps, by
\eqref{eq:gap-order}, while the full hull gap is affine invariant.
Apply the common-aspect bound to the expansion, always using restrictions
of the same ambient orientation law. Moreover the original factor's
upper value is the sum of the expanded upper values by common-threshold
attainment. Its deficiency under either law is thus the sum of expanded
deficiencies, proving the claimed per-original-factor guarantee.
The certificate can be sampled without forming the expansion: the
endpoint laws pull back to the original endpoints. Box inclusion alone
would not prove any of these comparisons.
\end{proof}

\begin{theorem}[Coefficient-regularity cardinality bound]
\label{thm:coefficient-regularity}
In ambient dimension $N\ge2$, suppose each entire local factor is the
multiaffine interpolant of
\[
 \psi_e(k)=\sum_{j=0}^{d_e}a_{e,j}\binom kj,\qquad a_{e,j}\ge0,
 \qquad a_{e,j+1}\le L a_{e,j}\quad(j\ge2),
\]
where $L\ge0$ is common to all factors and coefficients beyond the
degree are zero. Then
\[
 \tgap\le(L+1+\beta_N)\hgap.
\]
Separate affine terms may be added freely. Fair independent endpoint
orientations give the dimension-free bound $L+3$. These are bounds
for the gaps of the entire cardinality factors, not the gaps obtained
by separately relaxing their expanded monomials.
\end{theorem}
\begin{proof}
The nonnegative binomial coefficients make each sequence discrete convex:
its second difference is $\sum_{j\ge2}a_{e,j}\binom{k}{j-2}\ge0$.
Lemma~\ref{lem:cardinality-envelopes} supplies a common upper extremizer,
and within each factor one adjacent-cardinality law attains every
$V_j$ together. Denote that factor's four values by $C_\psi,V_\psi,
P_\psi,O_\psi$. Multiply \eqref{eq:box-coefficient} by its $a_j$ and
sum. Since $C_j-P_j\ge0$ and the order-zero and order-one deficiencies
vanish, reindexing gives
\[
 \begin{split}
 P_\psi-V_\psi
 &\le\beta_N(C_\psi-O_\psi)
          +\sum_{j\ge2}a_{j+1}(C_j-P_j)\\
 &\le\beta_N(C_\psi-O_\psi)+L(C_\psi-P_\psi).
 \end{split}
\]
Add $C_\psi-P_\psi$ to obtain
\[
 C_\psi-V_\psi\le(L+1)(C_\psi-P_\psi)+\beta_N(C_\psi-O_\psi).
\]
One global mixture with weights proportional to $L+1$ and $\beta_N$
therefore captures the required fraction for all factors. Summing
and using common upper attainment proves the theorem. This also
handles $L=0$, deterministic means, and affine factors. With at most
one nonfixed coordinate the interpolants are affine and all gaps vanish.
The product sequence $\rho_B^k$ has $a_j=(\rho_B-1)^j$, so
$L=\rho_B-1$ recovers the common-aspect product inequality.
\end{proof}
This refinement is a direct consequence of the coefficient mechanism;
no optimality or independent priority is asserted. Distinct earlier
low/high and asymmetric proofs, and the precise limitation of fixed
independence/orientation mixtures, appear in
Appendix~\ref{app:positive-box-predecessors}.

\subsection{A lower family on each fixed positive box}
\begin{theorem}[Aspect-ratio lower bound]\label{thm:positive-box-lower}
For every $\rho_B>1$,
\[
 \max\{2,\rho_B\}\le C_{\rm box}(\rho_B)\le\rho_B+2.
\]
Thus $C_{\rm box}(\rho_B)=\rho_B+O(1)$ as $\rho_B\to\infty$.
The exact finite-aspect-ratio constant is not determined.
\end{theorem}
\begin{proof}
The lower bound two already holds at width two by
Theorem~\ref{thm:width-two-gap}. For the lower bound $\rho_B$, fix
$\epsilon=1/\rho_B\in(0,1)$ and put $\alpha=1-\epsilon$.
Use the radix supports with $L\ge2$, $b=L^2$, $m=b^L$, now on
$[\epsilon,1]^{m+L}$ with normalized means
$\E A_j=b^{-j}$ and $\E Z_i=1-m^{-1}$. A level-$j$ term has
$k=m/b^j$ leaves and total failure mean $(1-b^{-j})+k/m=1$.
Its vertex value is $\epsilon^{R_e}$, so its exact lower value is
$\epsilon$, attained by exactly one failure with the prescribed
categorical probabilities. Common-threshold rounding gives its upper
value
\[
 C_{j,Q}=\epsilon+\alpha b^{-j}
                      -\frac\epsilon m(1-\epsilon^k).
\]
Indeed its successive threshold states have product values
$1,\epsilon,\epsilon^{k+1}$ with interval lengths
$b^{-j},1-m^{-1}-b^{-j},m^{-1}$. Therefore
\begin{equation}\label{eq:positive-radix-term}
 \tgap_L=\alpha L-D_L,\qquad
 D_L=\epsilon\sum_{t=0}^{L-1}\frac{1-\epsilon^{b^t}}{b^t},
 \qquad0\le D_L\le\frac{\epsilon b}{b-1}.
\end{equation}
For any common endpoint law let $R_Q$ count failed leaves in block $Q$
and $R$ count them in total. Since $\E R=1$, subtracting expected
physical products from their common upper value gives exactly
\begin{equation}\label{eq:softened-coverage}
 \hgap_L=\max\E\left[
 \alpha\sum_jA_j\sum_{Q\in\mathcal P_j}(1-\epsilon^{R_Q})
 +\epsilon\sum_j\sum_{Q\in\mathcal P_j}(1-\epsilon^{R_Q})
 \right]-D_L.
\end{equation}
For its first term use $1-\epsilon^{R_Q}\le\ind\{R_Q>0\}$ and the
unit-cube radix hull upper bound. For its second use
$1-\epsilon^r\le\alpha r$ and $\sum_{Q\in\mathcal P_j}R_Q=R$.
This proves
\[
 \hgap_L\le\epsilon\alpha L+
                 \alpha[1+(L-1)/b]-D_L.
\]
Exactly one uniformly failed leaf, with independent anchors of the
required means, instead gives
\[
 \hgap_L\ge\epsilon\alpha L+
                    \alpha^2\sum_{j=1}^Lb^{-j}-D_L.
\]
The same geometric inequality gives $D_L\le\epsilon\alpha L$, so this
last bound is strictly positive for every finite $L$. The bounded
correction in \eqref{eq:positive-radix-term} now implies
$\tgap_L/L\to\alpha$, $\hgap_L/L\to\epsilon\alpha$, and their ratio
tends to $\rho_B$. Scale coordinates by $\rho_B$ to get positive
coefficients on $[1,\rho_B]^{m+L}$; both gaps are preserved. This is
a limit in dimension for each fixed aspect ratio, requiring no exchange
of those limits. The upper bound is Theorem~\ref{thm:positive-box-upper}.
\end{proof}

\subsection{Unequal aspect ratios at frequency two}
\begin{example}[A rational bipartite ratio $7/6$]\label{ex:unequal-seven-six}
For $f(x,y,z)=xy+xyz$ on $[1,2]\times[1,2]\times[1,3]$, evaluate at
$(5/4,5/4,5/2)$. Put $x=1+a,y=1+b,z=1+2c$, so the normalized mean is
$(1/4,1/4,3/4)$. The following valid affine bounds give exact values:
\[
 xy\ge1+a+b,\quad xyz\ge2a+2b+3c,\quad
 f\ge4a+4b+4c,\quad f\le2+8a+4b+2c.
\]
Validity follows by substitution at all eight binary vertices and
multiaffine interpolation. The respective attaining distributions are
\[
\begin{array}{c|l}
 xy\text{ lower}&000,001,011,101\text{ each of mass }1/4\\
 xyz\text{ lower}&001\text{ of mass }3/4;\ 110\text{ of mass }1/4\\
 f\text{ lower}&001\text{ of mass }1/2;\ 011,100\text{ each of mass }1/4\\
 f\text{ upper}&000\text{ of mass }1/4;\ 001\text{ of mass }1/2;
                         \ 111\text{ of mass }1/4.
\end{array}
\]
Each row has the required means. The separate upper bounds
$xy\le1+2a+b$ and $xyz\le1+6a+3b+2c$ are attained by the last law.
Thus $\tgap=13/2-19/4=7/4$, while $\hgap=13/2-5=3/2$.
The dual multigraph is bipartite: its two factor vertices share two
parallel edges, with a dummy leaf for $z$. This disproves unequal-box
bipartite scalar exactness, while leaving the common-aspect theorem intact.
\end{example}

\begin{proposition}[Bipartite positive-box boundary]\label{prop:unequal-bipartite}
For arbitrary nonnegative boxes with a bipartite dual graph, positive
monomial gaps satisfy $\tgap\le2\hgap$. Their worst ratio over strictly
positive boxes is at least $3/2$, already with two factors and three
variables. Its exact value between $3/2$ and $2$ remains open here,
as does a universal $3/2$ bound for unequal-box frequency-two instances.
\end{proposition}
\begin{proof}
More generally a proper coloring of the term conflict graph into $q$
classes gives $\tgap\le q\hgap$: terms in one class have disjoint
scopes, so join their minimizing laws independently and complete unused
means. They attain their entire class gap, and other factors have
nonnegative expected deficiency from their upper envelopes. Since those
upper values are commonly attained, $\hgap$ is at least each class's
total gap. Sum these inequalities. This also proves a bound two for
any two positive terms regardless of their boxes or frequency.

For the lower family let $0<\epsilon\le2$ and use
$f_\epsilon=(2/\epsilon)xy+xyz$, with
$x,y\in[1,1+\epsilon]$, $z\in[1,3]$, at normalized means
$(a,b,c)=(1/4,1/4,3/4)$. After removing affine parts the two factors are
\[
 A=2\epsilon ab,\qquad
 B=\epsilon^2ab+2\epsilon ac+2\epsilon bc+2\epsilon^2abc,
 \qquad C=A+B.
\]
Their lower bounds at the point are attained at the first three laws
of Example~\ref{ex:unequal-seven-six}, respectively. The affine
certificates are $A\ge0$, $B\ge\epsilon^2(a+b+c-1)$, and
$C\ge2\epsilon(a+b+c-1)$. In binary order
$000,001,010,011,100,101,110,111$, the last two residual lists are
\[
 \begin{split}
 &(\epsilon^2,0,0,\epsilon(2-\epsilon),0,
                  \epsilon(2-\epsilon),0,\epsilon(\epsilon+4)),\\
 &(2\epsilon,0,0,0,0,0,\epsilon^2,\epsilon(3\epsilon+2)),
 \end{split}
\]
which are nonnegative throughout the stated range. The local lower
values sum to $\epsilon^2/4$ and the full lower value is $\epsilon/2$.
The same common upper law as before gives $3\epsilon/2+3\epsilon^2/4$.
Hence
\[
 \tgap=\epsilon(3/2+\epsilon/2),\quad
 \hgap=\epsilon(1+3\epsilon/4),\quad
 \frac{\tgap}{\hgap}=\frac{2(\epsilon+3)}{3\epsilon+4}\longrightarrow3/2.
\]
To obtain unit coefficients divide the polynomial by $2/\epsilon$ and
put $z'=(\epsilon/2)z$. It becomes $xy+xyz'$ on
$[1,1+\epsilon]^2\times[\epsilon/2,3\epsilon/2]$, with the same ratio.
\end{proof}
For frequency-two products on unequal positive boxes, normalizing upper
endpoints to one gives a binary failure cost
$\prod_{i\in J\cap e}\alpha_i=\exp(-\sum_{i\in J\cap e}w_i)$,
where $\alpha_i=\ell_i/u_i$ and $w_i=-\log\alpha_i$. Each edge has
the \emph{same} weight $w_i$ at both incident factors. This is a weighted
convex degree potential, whereas the matching theorem above uses
unweighted cardinalities. Finite floating-point searches approaching
$3/2$ from below do not supply an upper bound. If the dual multigraph
is a forest, conditional gluing actually preserves all local minima
and maxima even for signed multiaffine factors on arbitrary boxes;
Appendix~\ref{app:structural-auxiliary} gives the elementary argument.

\section{Exact complexity of a single physical product}\label{sec:single-product}
A ratio bound can be small even when evaluating an individual factor's
exact envelope is difficult. For one factor the ratio is identically
one wherever its gap is positive. The next result concerns computing
that factor's envelope, not loss caused by separating interacting terms.
All complexity statements use rational binary encodings and bit time.

\begin{theorem}[Midpoint decision on narrow unequal boxes]
\label{thm:single-product-hardness}
Given rational $u_i>1$ and $q$, deciding whether
\[
 \vex_B m(\bar x)\le q,\qquad m(x)=\prod_{i=1}^nx_i,
 \quad B=\prod_i[1,u_i],\quad \bar x_i=(1+u_i)/2,
\]
is NP-complete. For every fixed rational $\eta>0$ it remains
NP-complete when $u_i\le1+\eta$ for all $i$.
\end{theorem}
\begin{proof}
Reduce from PARTITION, the classical NP-complete integer partition
problem \citep{Karp1972}. Let its positive integers be $a_1,\ldots,a_n$
with total $A$. Doubling them permits us to assume $A$ even and at least
two. Fix an integer $K\ge1$ and put
\[
 \epsilon=(16KA^3)^{-1},\qquad u_i=1+\epsilon a_i.
\]
These boxes satisfy $1<u_i\le1+1/(16KA^2)<2$.
In normalized endpoint variables $X_i=1+\epsilon a_iZ_i$, the midpoint
means are $\E Z_i=1/2$. Write $S(z)=\sum_i a_i z_i$. On a binary vertex,
\begin{equation}\label{eq:partition-expansion}
 \prod_i(1+\epsilon a_iz_i)
 =1+\epsilon S(z)+\frac{\epsilon^2}{2}
             \left(S(z)^2-\sum_i a_i^2z_i\right)+R(z).
\end{equation}
All remainder coefficients are nonnegative. Its order-$j$ elementary
symmetric sum is at most $A^j$, and $\epsilon A\le1/16$, so
\begin{equation}\label{eq:partition-remainder}
 0\le R(z)\le\frac{(\epsilon A)^3}{1-\epsilon A}
     \le2\epsilon^3A^3=\epsilon^2/(8K)\le\epsilon^2/8.
\end{equation}
Every admissible law has $\E S=A/2$, whence
\[
 \E m(X)=b+\frac{\epsilon^2}{2}\operatorname{Var}(S)+\E R(Z),
 \qquad b=1+\frac{\epsilon A}{2}
       +\epsilon^2\left(\frac{A^2}{8}-\frac14\sum_i a_i^2\right).
\]
A partition state and its complement, with equal probabilities, have
all means $1/2$ and $S=A/2$ deterministically. Thus a YES instance has
lower envelope at most $b+\epsilon^2/8$. In a NO instance every binary
state has integer $S\ne A/2$, so $\operatorname{Var}(S)\ge1$ under
every feasible law. Its lower envelope is at least $b+\epsilon^2/2$.
The rational threshold $q=b+\epsilon^2/4$ separates the cases.
All numbers have polynomial encoding length, since $\log A$ is
polynomial in the original input length. For a fixed $\eta>0$, choose
fixed $K$ with $1/(16K)\le\eta$ to enforce the narrower interval.

For NP membership, the vertex-law LP has equality rows
$1,z_1,\ldots,z_n$ and right-hand side $(1,1/2,\ldots,1/2)$.
An optimal basic solution has at most $n+1$ positive probabilities.
It provides that many binary states and rational weights of polynomial
bit length: a nonsingular subsystem has zero-one entries and
half-integer right-hand side, so determinant bounds apply. Each vertex
product has polynomial bit length as a product of at most $n$ input
rationals. Exact arithmetic verifies nonnegative weights, total mass,
all means, and expected product at most $q$ in polynomial time.
\end{proof}

The NO case also has an explicit affine certificate in normalized
coordinates:
\begin{equation}\label{eq:partition-dual}
 \ell(z)=1-\epsilon^2A^2/8+\epsilon^2/2
    +\sum_i\left[\epsilon a_i+\frac{\epsilon^2}{2}(Aa_i-a_i^2)\right]z_i.
\end{equation}
Subtracting it from \eqref{eq:partition-expansion} gives
$\epsilon^2[(S-A/2)^2-1]/2+R(z)\ge0$ at every binary vertex of a
NO instance, hence everywhere by multiaffine interpolation. Its midpoint
value is $b+\epsilon^2/2$.
The reduction is from weakly NP-complete PARTITION; it establishes
ordinary, not strong, NP-hardness. It concerns varying unequal boxes
contained in a narrow common box, not the envelope on that one common
box. The separating accuracy can be exponentially small in input length.

\begin{corollary}[Graph-hull membership]\label{cor:single-product-membership}
Deciding whether $(\bar x,q)$ belongs to the scalar graph hull of the
same single monomial is NP-complete under the same narrow-box restrictions.
\end{corollary}
\begin{proof}
At the midpoint its common-threshold upper value is
$C=(1+\prod_i u_i)/2$. Retaining the product's nonnegative quadratic
terms gives
\[
 C\ge1+\epsilon A/2+\epsilon^2(A^2-\sum_i a_i^2)/4,
 \qquad C-b\ge\epsilon^2A^2/8\ge\epsilon^2/2>q-b.
\]
Thus the upper-hull condition always holds on the reduced instances,
and membership is exactly the lower-envelope threshold question.
Membership lies in NP: a basic distribution for the normalization,
coordinate, and value equations has at most $n+2$ positive weights.
These equations have rational polynomial-bit entries and determinant
bounds give polynomial-bit weights. Exact verification proves membership.
\end{proof}

\begin{proposition}[A fixed number of ratio types]\label{prop:ratio-types}
For one product on a rational strictly positive box having at most $k$
distinct coordinate aspect ratios, its exact scalar envelope and a
supporting affine minorant can be computed in $n^{O(k)}$ times a
polynomial in the input bit length. In particular fixed $k$ is tractable.
\end{proposition}
\begin{proof}
After substituting fixed coordinates and normalizing endpoints, a
binary product with count $h_j$ in ratio group $j$ has value
$(\prod_i\ell_i)\prod_{j=1}^k r_j^{h_j}$. To minimize this value minus
$\sum_i\pi_i z_i$ for arbitrary dual coefficients, sort the coefficients
within each group. For fixed counts choose the $h_j$ largest coefficients,
and use prefix sums. Enumerate all $\prod_j(n_j+1)\le(n+1)^k$ count
tuples. Rational multiplication and sorting have polynomial bit cost;
this is an exact vertex oracle. The determinant-bounded dual argument
of Corollary~\ref{cor:scalar-envelope-algorithm} then computes the lower
envelope and its minorant. The common-threshold upper envelope is
explicit. Thus mere closeness of the ratios does not substitute for a
bound on their number of distinct values.
\end{proof}

\subsection{Rank-one transport and logarithmic accuracy dependence}
Multimarginal optimal transport with specified finite marginals minimizes
expected cost over their joint laws. With $n$ fair binary marginals and
factored cost tensor
\[
 C=\bigotimes_{i=1}^n(1,1+\epsilon a_i),\qquad
 C_z=\prod_i(1+\epsilon a_i z_i),
\]
it is exactly the midpoint envelope above. This tensor is strictly
positive and rank one, with no sparse correction, and
\[
 1\le C_z\le\prod_i(1+\epsilon a_i)
       \le\sum_{j\ge0}(\epsilon A)^j\le16/15.
\]
\begin{corollary}[Rank-one value and weighted dual accuracy]
\label{cor:rank-one-precision}
Exact threshold evaluation for these rank-one MOT instances is
NP-complete. Unless $P=NP$, neither their values nor their weighted
vertex minima admit a deterministic additive-$\delta$ approximation
algorithm polynomial in rational input length and
$\log(C_{\max}/\delta)$. For randomized bounded-error algorithms the
corresponding consequence requires $NP\not\subseteq BPP$.
\end{corollary}
\begin{proof}
Take $\delta=\epsilon^2/32$ in the YES and NO intervals of
Theorem~\ref{thm:single-product-hardness}. Their approximation intervals
remain disjoint and $\log(1/\delta)$ is polynomial in the PARTITION
input length. Such an algorithm would decide PARTITION.
The weighted vertex oracle is hard directly: set
\[
 \pi_i=\epsilon a_i+\frac{\epsilon^2}{2}(Aa_i-a_i^2),\qquad
 d_0=1-\epsilon^2A^2/8.
\]
Equation~\eqref{eq:partition-expansion} gives
\[
 C_z-\pi^{\mathsf T}z=d_0+\frac{\epsilon^2}{2}(S(z)-A/2)^2+R(z).
\]
Its minimum is at most $d_0+\epsilon^2/8$ in a YES instance and at
least $d_0+\epsilon^2/2$ in a NO instance. The same approximation margin
and polynomial encoding bounds apply. Exact MOT decision has the
NP certificates already given. If an input convention requires bit
lengths polynomial in the number of marginals, append a number of dummy
binary factors $(1,1)$ linear in the original input bit length. Give
them fair marginals and zero dual potentials. They change neither
rank, cost range, value, nor oracle minimum, and meet that convention.
\end{proof}

The published structured-MOT paper of \citet[Section 7.2, after
Theorem 7.4]{AltschulerBoixAdsera2023} asks whether constant-rank
approximate dual minimization can have logarithmic inverse-accuracy
dependence, leading to exact algorithms. The preceding reduction rules
out that dependence in the rational bit model under the stated
complexity assumption, already at rank one. Their Corollary 7.5 has
polynomial dependence on inverse accuracy and is consistent with it.
This comparison identifies the historical question; it is not a claim
that no subsequent work addressed it. Nor is it a lower bound in an
unrestricted unit-cost arithmetic model. The unweighted smallest entry
of this tensor is simply one; arbitrary dual weights are essential.
None of these arguments proves hardness at fixed additive error.

\subsection{A boundary of multilinearity}\label{subsec:rational-powers}
\begin{example}[Two positive rational-power terms]\label{ex:rational-powers}
On the fixed interval $[1,3]$, let
$f_\epsilon(x)=x^{1-\epsilon}+x^{1+\epsilon}$, $0<\epsilon<1$.
The first term is concave and the second convex, while
\[
 f_\epsilon''(x)=\epsilon x^{-1-\epsilon}
       [(1+\epsilon)x^{2\epsilon}-(1-\epsilon)]>0.
\]
At $x=2$ their separate widths and full width are consequently
\[
 \tgap_\epsilon=3\sinh(\epsilon\log3)-4\sinh(\epsilon\log2),
 \qquad
 \hgap_\epsilon=1+3\cosh(\epsilon\log3)-4\cosh(\epsilon\log2).
\]
These follow by subtracting each function from its endpoint secant,
with the opposite subtraction for the concave term. Taylor expansion
gives
\[
 \tgap_\epsilon=(3\log3-4\log2)\epsilon+O(\epsilon^3),\qquad
 \hgap_\epsilon=\tfrac12[3(\log3)^2-4(\log2)^2]\epsilon^2
                                      +O(\epsilon^4).
\]
The first leading coefficient is $\log(27/16)>0$; the second is
positive by strict convexity of $x(\log x)^2$ on $[1,3]$.
Thus their ratio diverges like a positive constant times $1/\epsilon$.
Taking $\epsilon=1/k$, $k\ge2$, keeps rational exponents in
$[1/2,3/2]$ with only two positive unit coefficients and one variable.
Under $x=e^y$, however, both terms become convex exponentials on
$[0,\log3]$. Their lower envelopes are exact and endpoint secants
add, so the termwise and full widths coincide in that representation.
The example concerns the original coordinate and does not preclude
useful logarithmic reformulation.
\end{example}

\section{Fractional cardinality and spatial certificates}
\label{sec:cardinality-spatial}
The following optimization problems have explicit solutions. Their role is
to distinguish the size of certificates produced by specified spatial
oracles from the difficulty of finding an optimizer. The obstruction uses a
concave quadratic penalty and an affine balance; it is a separate
construction from the positive multilinear gap examples above.

Fix integers $n\ge2$ and $0\le k\le n-1$, and write
\begin{equation}\label{eq:cardinality-problem}
 K=k+\tfrac12,\qquad
 \mathcal F_{n,k}=\{x\in[0,1]^n:\textstyle\sum_i x_i=K\},\qquad
 F(x)=\sum_{i=1}^n x_i(1-x_i).
\end{equation}
The dimension parameter $n$ counts original coordinates. A vertex of
$\mathcal F_{n,k}$ has exactly $k$ coordinates equal to one, one equal to
$1/2$, and the others zero: two strictly fractional coordinates would admit
a small opposite perturbation, whereas a point with at most one fractional
coordinate is determined by its active bounds and the balance. Concavity
and compactness therefore give
\begin{equation}\label{eq:cardinality-optimum}
 \min_{\mathcal F_{n,k}}F=\tfrac14.
\end{equation}
For each oracle below, a certified region is either infeasible or has
oracle bound at least $T_*$. A certificate is a finite cover of
$\mathcal F_{n,k}$ by such regions, as in
Section~\ref{subsec:certificate-model}. For absolute tolerance
$\epsilon\ge0$ we grant the best incumbent and set
$T_*=1/4-\epsilon$. Pruning at $U-\epsilon$, for any feasible incumbent
$U\ge1/4$, implies this condition. Infeasible regions contain no witness
and do not help cover the feasible set. All lower bounds also hold for a
weaker valid oracle: certifying a region with that oracle certifies it
with the stronger oracle used in the proof.

There are close precedents in other certificate models. Jarre's binary
knapsack construction retains an exponential branch-and-bound obstruction
with a max-cut semidefinite bound \citep[Sections 2--3]{Jarre2018}.
Coniglio's hyperplane-clustering analysis gives a related spatial
obstruction under midpoint branching on symmetric coordinate domains
\citep[Assumption 1, Propositions 2--4, Appendix C in the review version]{Coniglio2026}.
Our arguments concern arbitrary coordinate restrictions and a fixed
objective gap under the explicitly defined node oracles. We make no
priority claim for exponential spatial lower bounds in general. The
fractional-cardinality moments used later are classical; the additional
argument places them in spatial regions and counts how many witnesses a
certified region can contain.

\subsection{Chords and endpoint avoidance}
For a box $B=\prod_i[a_i,b_i]\subseteq[0,1]^n$, including degenerate
intervals, define
\begin{equation}\label{eq:cardinality-chord}
 \ell_i(x_i)=a_ib_i+(1-a_i-b_i)x_i,\qquad
 x_i(1-x_i)-\ell_i(x_i)=(x_i-a_i)(b_i-x_i).
\end{equation}
Its chord bound is
\[
 \operatorname{LB}_{\rm ch}(B)=
 \min_{x\in B\cap\mathcal F_{n,k}}\sum_i\ell_i(x_i),
\]
with $+\infty$ for an empty feasible intersection. If $g_i$ is any convex
underestimator of $x_i(1-x_i)$ on $[a_i,b_i]$, convexity between the
endpoints gives $g_i\le\ell_i$. This also holds on a singleton interval.
Thus the infimum for any separable convex underestimator $\sum_i g_i$
is at most $\operatorname{LB}_{\rm ch}(B)$; an infimum is used because
arbitrary convex underestimators need not attain their infima on a closed
interval without lower semicontinuity at its endpoints.

For clarity, two common formulations yield exactly these chords. For
$y_i=1-x_i\in[1-b_i,1-a_i]$, both McCormick lower planes for
$w_i=x_iy_i$ become $w_i\ge\ell_i(x_i)$. The upper planes remain
compatible: the exact product lies above the chord and below each upper
plane. The separable $\alpha$BB expression
$x_i(1-x_i)-\alpha(x_i-a_i)(b_i-x_i)$ has second derivative
$2(\alpha-1)$, so the smallest convexifying value $\alpha=1$ gives the
chord. Larger valid values give weaker underestimators. These statements
concern the specified separable formulations, without additional coupled
cuts or nonlinear subproblem solution rules.

We isolate the counting argument so that later oracle changes do not
obscure it. Let $n=k+m+z$, where $k,z\ge0$ and $m\ge1$ are integers.
For each ordered partition $(H,M,Z)$ of $[n]$ with these sizes, put
\begin{equation}\label{eq:cardinality-witness}
 p=\frac1{2m},\qquad w=\ind_H+p\ind_M.
\end{equation}
These $n!/(k!m!z!)$ distinct witnesses belong to $\mathcal F_{n,k}$.
For an arbitrary product $S=\prod_i S_i$, with nonempty
$S_i\subseteq[0,1]$, define
\begin{equation}\label{eq:endpoint-exclusions}
 A(S)=\{i:0\notin S_i\},\quad D(S)=\{i:1\notin S_i\},\quad
 R(S)=A(S)\cup D(S).
\end{equation}
For boxes these are $A=\{i:a_i>0\}$ and $D=\{i:b_i<1\}$.

\begin{lemma}[Endpoint-avoidance count]\label{lem:endpoint-count}
Suppose each certified product domain that contains a witness has
$|R(S)|\ge q$, where $q\ge0$ is real. Any certified cover of
$\mathcal F_{n,k}$ contains at least
\begin{equation}\label{eq:cardinality-cover}
 \mathcal N(q):=
 \min\left\{\left(\frac n{n-k}\right)^{q/2},
                 \left(\frac n{n-z}\right)^{q/2}\right\}
\end{equation}
regions. The same conclusion follows from the stronger hypothesis
$|R(S)|>q$. Both denominators are positive, and $\mathcal N(0)=1$.
\end{lemma}
\begin{proof}
Choose the ordered partition uniformly. A contained witness must have
$Z\cap A=\varnothing$ and $H\cap D=\varnothing$. For a fixed set of
$s$ excluded positions, a uniformly chosen $z$-subset avoids it with
probability
\[
 \frac{\binom{n-s}{z}}{\binom nz}
 =\prod_{j=0}^{s-1}\frac{n-z-j}{n-j}
 \le\left(\frac{n-z}{n}\right)^s
 \quad(s\le n-z).
\]
For $s>n-z$ the probability is zero, so the upper bound still holds;
for $s=0$ the product is one. The analogous inequality holds for $H$.
No independence between the two avoidance events is required. A fixed
region's witness fraction is at most
\[
 \min\left\{\left(\frac{n-z}{n}\right)^{|A|},
               \left(\frac{n-k}{n}\right)^{|D|}\right\}.
\]
Since $|A|+|D|\ge |R|\ge q$, at least one exponent is at least $q/2$.
The displayed fraction is consequently at most $1/\mathcal N(q)$.
Regions containing no witness have fraction zero. The union bound over a
cover proves the assertion, whether or not its regions overlap. Bases
equal to one when $k=0$ or $z=0$ cause no exception.
\end{proof}

\begin{theorem}[Separable spatial certificates]\label{thm:chord-cover}
Let $n=k+m+z$, with integers $k,z\ge0$, $m\ge1$, $n\ge2$, and
$0\le\epsilon<1/4$. For problem~\eqref{eq:cardinality-problem}, every
cover by boxes certified by the chord oracle at
$T_*=1/4-\epsilon$ has at least $\mathcal N(h)$ members, where
\begin{equation}\label{eq:chord-h}
 h=m(1/2-2\epsilon).
\end{equation}
This includes arbitrary real coordinate split points and all separable
convex underestimators dominated by the chord oracle. For
$k=m=z=t$, $n=3t$, the bound is
$(3/2)^{t(1/4-\epsilon)}$, hence $(3/2)^{n/24}$ at $\epsilon=1/8$.
\end{theorem}
\begin{proof}
In a box containing $w$, the chord values at its zero and unit witness
coordinates vanish. On $U=[n]\setminus R$ the interval is $[0,1]$,
whose chord vanishes identically. Consequently
\[
 T_*\le\operatorname{LB}_{\rm ch}(B)
 \le\sum_{i\in M\cap R}\ell_i(p)
 \le |M\cap R|p(1-p)<\frac{|M\cap R|}{2m}.
\]
The last strict inequality is applicable because $T_*>0$ forces
$M\cap R\ne\varnothing$. Thus $|R|>2mT_*=h$, and
Lemma~\ref{lem:endpoint-count} applies. A complete coordinate split tree
provides such a cover, regardless of split locations or branching order.
\end{proof}

\subsection{Upper certificates, tightening, and tolerances}
\begin{proposition}[Explicit upper certificates]\label{prop:chord-upper}
For $0<\epsilon<1/4$, a chord certificate for
\eqref{eq:cardinality-problem} has at most $2^{n+2}-3$ nodes and
$2^{n+1}-1$ leaves. For $\epsilon=0$, a midpoint certificate has exactly
$\binom{n+1}{k+1}$ leaves if only the stopping rules in the proof are used.
In particular its node count is
$O(n^{\min\{k+1,n-k\}})$ when that exponent is fixed.
\end{proposition}
\begin{proof}
Choose $0<\alpha<1/2$ with
$\alpha(1-\alpha)=1/4-\epsilon$. Process coordinates in order. At each
surviving node split the next coordinate at $\alpha$, then split its
upper child at $1-\alpha$. The middle interval has constant chord
$\alpha(1-\alpha)$, and all other chords are nonnegative, so the middle
child is certified. Continue on the low and high intervals. At a final
corner box the chord sum is $(1-\alpha)\|x-v\|_1$, where
$v\in\{0,1\}^n$ is its corner. On the demand slice,
$\|x-v\|_1\ge|K-\sum_i v_i|\ge1/2$. This gives a bound
$(1-\alpha)/2>1/4>T_*$. Processing coordinate $i$ adds four nodes per
surviving prefix, of which there are $2^{i-1}$. The stated counts follow.
For a strict pruning rule choose a slightly larger middle-chord target
still below $1/4$.

For exact certification split at $1/2$ instead. If a prefix has $c$ high
intervals and $d$ low intervals, its chord bound is at least each of
\[
 \frac{c-K}{2},\qquad \frac{K-(n-d)}2.
\]
Indeed, the sum on high coordinates is at most $K$, whereas the sum on
low coordinates is at least $K-(n-d)$; all other chord contributions
are nonnegative. Stop at $c=k+1$ or $d=n-k$. No path survives depth $n$.
With remaining high and low stopping counts $a,b$, the number of leaves
satisfies $L(a,b)=L(a-1,b)+L(a,b-1)$ and
$L(0,b)=L(a,0)=1$, hence $L(a,b)=\binom{a+b}{a}$. Start at
$(a,b)=(k+1,n-k)$. A full binary tree has twice as many leaves minus one
nodes. This proves the exact count and polynomial assertion.
\end{proof}

The balanced fixed-tolerance family therefore has certificate size
$2^{\Theta(n)}$ under the chord model. If both $k$ and $n-k$ are
proportional to $n$, choose $m,z$ proportional to $n-k$ to obtain the
same exponential order. Fixed demand level, or fixed complementary demand
level, admits the polynomial certificate just proved. At
$\epsilon\ge1/4$, the root chord bound zero already meets the target;
no positive exponent is asserted there.

The chord error is genuinely quadratic in interval width:
\begin{equation}\label{eq:second-order-chord}
 0\le F(x)-\sum_i\ell_i(x_i)
 \le\frac14\sum_i(b_i-a_i)^2
 \le\frac n4\max_i(b_i-a_i)^2.
\end{equation}
In particular, the difference between the exact node optimum and the
chord bound is at most this quantity, by evaluating $F$ at a chord
minimizer. Thus the fixed-tolerance lower bound persists despite this
second-order error estimate; it does not require a first-order local
approximation error.

\begin{proposition}[Charged reductions and equality tolerance]
\label{prop:charged-tolerance}
Suppose a run keeps feasible points covered by its final boxes together
with discarded boxes certified by the same oracle at $T_*$. If it processes
$N$ nodes and discards at most $a$ such boxes per node, a cover lower bound
$L$ implies $N\ge L/(a+1)$. Pure feasibility reductions need not be
charged when they discard no feasible point.

For a demand feasibility tolerance $0\le\delta<1/2$, replace
$\mathcal F_{n,k}$ by
$\mathcal F^\delta_{n,k}=\{x\in[0,1]^n:|\sum_i x_i-K|\le\delta\}$.
Its exact optimum is $1/4-\delta^2$. If
$\epsilon+\delta^2<1/4$, a chord certificate at absolute tolerance
$\epsilon$ for this enlarged set satisfies Theorem~\ref{thm:chord-cover}
with $\epsilon$ replaced by $\epsilon+\delta^2$.
\end{proposition}
\begin{proof}
The augmented cover has at most $(a+1)N$ members. A discarded open
coordinate slab may be charged by its closed box only when that box is
certified; this condition is part of the assertion, not an unproved
closure assumption. Subsequent product-domain results can charge the
removed coordinate set itself.

For a fixed sum $s\in[k+1/2-\delta,k+1/2+\delta]\subset(k,k+1)$,
the same vertex argument as above gives minimum penalty
$(s-k)(1-s+k)$. Its minimum on this interval is
$1/4-\delta^2$, attained by $k$ ones and a remaining coordinate
$1/2-\delta$ or $1/2+\delta$. A tolerance-feasible incumbent is therefore
at least this value. The exact-balance witnesses remain feasible for the
enlarged problem, and the chord bound over the enlarged feasible set is
at most its value at each witness. The proof of
Theorem~\ref{thm:chord-cover} applies at the positive target
$1/4-(\epsilon+\delta^2)$.
\end{proof}
The restriction $\delta<1/2$ is material: once the tolerance interval
contains an integer sum, the optimum is zero. An unrestricted formula
$1/4-\delta^2$ would then be false. Relative pruning in the convention
$\operatorname{LB}\ge(1-\theta)U$, $0\le\theta<1$, on the original
problem implies the absolute target with $\epsilon=\theta/4$.

\subsection{Full second moments and box RLT}
For $B=\prod_i[a_i,b_i]\subseteq[0,1]^n$, let
$\operatorname{LB}_{\rm SDP}(B)$ be the infimum of
$\sum_i(\mu_i-X_{ii})$ over symmetric $X$ and $\mu\in\R^n$ satisfying
\begin{align}
 a\le\mu\le b,\quad &\sum_i\mu_i=K,\quad
 \begin{pmatrix}1&\mu^T\\\mu&X\end{pmatrix}\succeq0,
 &\sum_jX_{ij}=K\mu_i\quad(i\in[n]),\label{eq:sdp-model}
\end{align}
and all pairwise box inequalities
\begin{align}
 X_{ij}-a_i\mu_j-a_j\mu_i+a_ia_j&\ge0,\notag\\
 b_j\mu_i-X_{ij}-a_ib_j+a_i\mu_j&\ge0,\notag\\
 b_i\mu_j-X_{ij}-b_ia_j+a_j\mu_i&\ge0,\notag\\
 b_ib_j-b_i\mu_j-b_j\mu_i+X_{ij}&\ge0
 \qquad(i,j\in[n]).\label{eq:box-rlt}
\end{align}
Indices may coincide. These are the linearizations of all products of
$x_i-a_i$ and $b_i-x_i$. Exact evaluation $(\mu,X)=(x,xx^T)$ shows
validity. As before the infimum of an empty relaxation is $+\infty$,
and a region is certified when its bound is at least
$T_*=1/4-\epsilon$ (or it is infeasible).

\begin{lemma}[Restricted fractional second moments]\label{lem:sdp-witness}
Let $k,m,z\ge1$, and let $B$ contain a witness~\eqref{eq:cardinality-witness}.
If $|R(B)|<\min\{k,z\}$, the model
\eqref{eq:sdp-model}--\eqref{eq:box-rlt} has a feasible point of value
$|M\cap R|p(1-p)$.
\end{lemma}
\begin{proof}
Fix $x_i=w_i$ deterministically on $R$ and put
\[
 U=[n]\setminus R,\quad s=|U|,\quad
 t=K-\sum_{i\in R}w_i=|H\cap U|+p|M\cap U|.
\]
At least one unit and one zero witness coordinate remain, so $s\ge2$
and $1\le t\le s-1$. On $U$ set
\[
 c=t/s,\quad d=t(t-1)/(s(s-1)),\quad
 \mu_i=c,\quad X_{ii}=c,\quad X_{ij}=d\ (i\ne j).
\]
On $R$ set $\mu_i=w_i$, and set $X_{ij}=\mu_i\mu_j$ whenever at least
one index lies in $R$. First moments meet the box bounds because the
intervals on $U$ are $[0,1]$ and $w\in B$. Their sum is $K$.
The covariance is zero on rows and columns in $R$, and on $U$ it is
\[
 X_{UU}-\mu_U\mu_U^T
 =(c-d)(I-J/s),\qquad c-d=\frac{t(s-t)}{s(s-1)}\ge0.
\]
It is positive semidefinite, proving augmented PSD by the Schur complement
of the entry one. For $i\in U$,
$c+(s-1)d=tc$ gives $\sum_jX_{ij}=Kc$; for $i\in R$, the equality
is $\sum_jX_{ij}=w_i\sum_j\mu_j=Kw_i$.

For distinct $i,j\in U$, the four slack-product expectations are
\[
 d,\quad c-d,\quad c-d,\quad
 1-2c+d=\frac{(s-t)(s-t-1)}{s(s-1)},
\]
all nonnegative. For $i=j\in U$ they are $c,0,0,1-c$.
If at least one index is in $R$, the expectation factors into a
nonnegative deterministic slack and a nonnegative expected slack; this
also covers repeated restricted indices and degenerate intervals.
Finally $\mu_i-X_{ii}=0$ on $U$, and on $R$ it is $w_i(1-w_i)$.
\end{proof}

\begin{theorem}[SDP--RLT cover bound]\label{thm:sdp-cover}
Let $n=k+m+z$, where $k,m,z\ge1$ are integers, and let
$0\le\epsilon<1/4$. Every cover for
\eqref{eq:cardinality-problem} certified by
\eqref{eq:sdp-model}--\eqref{eq:box-rlt} at $T_*=1/4-\epsilon$ has
at least $\mathcal N(q)$ members, where
$q=\min\{k,z,m(1/2-2\epsilon)\}$.
In the balanced case $k=m=z=t$, this is the chord lower bound in
Theorem~\ref{thm:chord-cover}.
\end{theorem}
\begin{proof}
If a containing box has $|R|<q$, Lemma~\ref{lem:sdp-witness} gives a
feasible objective at most $|R|p(1-p)<q/(2m)\le T_*$, where the
strict comparison holds also when $R=\varnothing$ because $q>0$.
The box cannot be certified. Therefore $|R|\ge q$ on every certified
box containing a witness, and Lemma~\ref{lem:endpoint-count} applies.
The diagonal mixed inequality gives
$X_{ii}\le(a_i+b_i)\mu_i-a_ib_i$, so the SDP bound dominates the
chord bound. The upper certificates of Proposition~\ref{prop:chord-upper}
remain valid.
\end{proof}
This resolves the behavior under spatial branching for this SDP--RLT
model; it is not merely a calculation of its root value.

One can also grant all products of two affine inequalities valid on
$B\cap\mathcal F_{n,k}$ without strengthening this oracle. Indeed,
linear-programming duality represents any such nonnegative affine
function, on a nonempty polytope, as a nonnegative constant plus a
nonnegative combination of box slacks plus a multiple of
$e(x)=\sum_i x_i-K$. The representation is an identity of affine
polynomials, including for degenerate boxes. Expanding the product of
two representations leaves only existing nonnegative slack moments and
terms $L[e v]=0$ with $v$ affine. This assertion concerns affine
inequalities in the original $x$ coordinates. It does not grant all
valid affine cuts in the lifted $(\mu,X)$ coordinates.

\section{Full preorderings and fractional-cardinality moments}
\label{sec:cardinality-preordering}
An order-$r$ oracle below acts on polynomials of total degree at most
$2r$, where $r\ge1$ is an integer. On a box
$B=\prod_i[a_i,b_i]\subseteq[0,1]^n$, it consists of linear functionals
$L:\R[x]_{\le2r}\to\R$ satisfying
\begin{align}
 L[1]&=1,\label{eq:preorder-normalization}\\
 L[e v]&=0\quad(\deg v\le2r-1),\qquad
 e=\sum_i x_i-K,\label{eq:preorder-equality}\\
 L[g P^2]&\ge0\quad
 \left(\sum_{j=1}^a\deg g_j+2\deg P\le2r,\quad
 g=\prod_{j=1}^a g_j\right),\label{eq:preorder-positivity}
\end{align}
where each $g_j$ is a coordinate-bound slack $x_i-a_i$ or $b_i-x_i$.
The empty product and repeated slack factors are included. The square
polynomial $P$ may couple every coordinate. Zero polynomials impose no
condition; constants have degree zero. The node bound
$\operatorname{LB}_r(B)$ is the infimum of $L[F]$ over these functionals.
This is the full truncated box preordering, including all allowed products
of slacks, not just the separate single-generator localizers. Evaluation
at a feasible point proves validity.

At $r=1$ this model is exactly
\eqref{eq:sdp-model}--\eqref{eq:box-rlt}: positivity on affine squares
is augmented PSD, products of at most two slacks give box RLT, single
slacks give first-moment bounds, and the balance times affine polynomials
gives its first and second moment equations. At every order $r\ge1$ the
model dominates that degree-two model and hence the chord oracle.
Certification again means an infeasible region or
$\operatorname{LB}_r(B)\ge T_*=1/4-\epsilon$ for
problem~\eqref{eq:cardinality-problem}. Original dimension remains $n$.

\subsection{A self-contained positivity proof}
The moments below are the classical fractional-cardinality moments
associated with Grigoriev's knapsack obstruction
\citep{Grigoriev2001}. Potechin gives the explicit construction and its
classical provenance \citep[Theorem 1, Example 18, Theorem 44]{Potechin2019}.
We prove the precise sufficient range required for the present full
preordering; no external positivity theorem is needed.

For real $v$ and integer $a\ge0$, write
$(v)_a=v(v-1)\cdots(v-a+1)$, with $(v)_0=1$.
In $s$ formal Boolean variables $u_1,\ldots,u_s$, let $\operatorname{red}$
replace every positive power of $u_i$ by $u_i$. This is the canonical
multilinear representative modulo $u_i^2-u_i$; it respects sums and
products modulo these identities and cannot increase degree. For
$u_S=\prod_{i\in S}u_i$ define
\begin{equation}\label{eq:fractional-moments}
 E_{s,t}[u_S]=\frac{(t)_{|S|}}{(s)_{|S|}},\qquad
 E_{s,t}[P]=E_{s,t}[\operatorname{red}P].
\end{equation}
All subset sizes are at most $s$, so denominators are positive. The
functional is defined on the Boolean quotient and is used only through
the degree stated in each assertion; no positivity on all polynomials is
assumed. In particular $E_{s,t}[1]=1$. For $s=0$, deterministic evaluation
of constants will be used instead of a fractional-cardinality functional.

\begin{lemma}[Cardinality identities and moment positivity]
\label{lem:fractional-positive}
Let $d\ge1$ be an integer, $s\ge2d$ an integer, and
$t,s-t\ge2d-1$. Then
\begin{equation}\label{eq:fractional-cardinality-identity}
 E_{s,t}\left[\left(\sum_i u_i-t\right)v\right]=0
 \quad(\deg v\le2d-1),
\end{equation}
and $E_{s,t}[P^2]\ge0$ for every $\deg P\le d$.
\end{lemma}
\begin{proof}
For squarefree $u_S$ of degree $a\le2d-1\le s-1$,
\[
 E_{s,t}\left[\left(\sum_i u_i\right)u_S\right]
 =a\frac{(t)_a}{(s)_a}
  +(s-a)\frac{(t)_{a+1}}{(s)_{a+1}}
 =t\frac{(t)_a}{(s)_a}.
\]
Boolean reduction and linearity prove
\eqref{eq:fractional-cardinality-identity} for every allowed multiplier.

To reduce positivity to homogeneous polynomials, for $|S|=a\le d$ put
\begin{equation}\label{eq:fractional-homogenization}
 H_S(u)=\frac{\sum_{T\supseteq S,\ |T|=d}u_T}
                  {\binom{t-a}{d-a}}.
\end{equation}
The denominator is positive under $t\ge2d-1$ (and is one for $a=d$).
In the Boolean quotient the numerator equals
$u_S\binom{\sum_i u_i-a}{d-a}$: at a Boolean vector it counts the
$d$-subsets containing $S$ when all coordinates of $S$ equal one, and
both sides vanish otherwise. The polynomial
$\binom{z-a}{d-a}-\binom{t-a}{d-a}$ is divisible by $z-t$ with quotient
degree at most $d-a-1$ when $a<d$. Consequently
\[
 u_S-H_S=\left(\sum_i u_i-t\right)Q_S
 \quad\hbox{in the Boolean quotient},\qquad \deg Q_S\le d-1.
\]
Replace each monomial of $\operatorname{red}P$ by its $H_S$, obtaining a
homogeneous squarefree polynomial $H$ of degree $d$. Then
$P-H=(\sum_i u_i-t)Q$ modulo Boolean identities with $\deg Q\le d-1$.
Since $\deg(Q(P+H))\le2d-1$, the cardinality identity gives
$E_{s,t}[P^2]=E_{s,t}[H^2]$. Every expression used here has degree at most
$2d$ before Boolean reduction.

Index the homogeneous moment matrix $N$ by $d$-subsets $I,J$ of $[s]$:
\[
 N(I,J)=\frac{(t)_{2d-|I\cap J|}}{(s)_{2d-|I\cap J|}}.
\]
Let $P_j(I,J)=\binom{|I\cap J|}{j}$ for $0\le j\le d$.
This is a Gram matrix: associate with $I$ the vector indexed by
$j$-subsets $A$, with entry $\ind\{A\subseteq I\}$.
The exact decomposition is
\begin{equation}\label{eq:fractional-gram}
 N=\sum_{j=0}^d
   \frac{(t)_{2d-j}(s-t)_j}{(s)_{2d}}P_j.
\end{equation}
Indeed, in an entry with $|I\cap J|=\ell$, the numerator on the right
is
\begin{align*}
 \sum_{j=0}^\ell\binom\ell j(t)_{2d-j}(s-t)_j
 &=(t)_{2d-\ell}
   \sum_{j=0}^\ell\binom\ell j
       (t-2d+\ell)_{\ell-j}(s-t)_j\\
 &=(t)_{2d-\ell}(s-2d+\ell)_\ell.
\end{align*}
The last equality is the falling-factorial Vandermonde identity, obtained
by counting ordered choices from two disjoint finite sets for nonnegative
integer arguments and then by polynomial identity for real arguments.
Division by $(s)_{2d}$ gives the claimed entry. No division by
$(t)_{2d-\ell}$ was used, so zero coefficients cause no exception.
All coefficients in \eqref{eq:fractional-gram} are nonnegative because
$t,s-t\ge2d-1$, and $(s)_{2d}>0$. Thus $N\succeq0$ and
$E_{s,t}[H^2]\ge0$.
\end{proof}

\begin{lemma}[Assignment indicators and every slack product]
\label{lem:assignment-localizer}
Suppose $r\ge1$, $s\ge2r$, and $t,s-t\ge2r-1$.
For every product $g$ of $u_i$ and $1-u_i$, including repetitions, and
every polynomial $P$ with
$\deg g+2\deg P\le2r$, one has $E_{s,t}[gP^2]\ge0$.
\end{lemma}
\begin{proof}
Boolean reduction makes $g$ zero if both $u_i$ and $1-u_i$ occur for
some $i$. Otherwise it becomes an assignment indicator
\[
 I_{A,B}=\prod_{i\in A}u_i\prod_{j\in B}(1-u_j),\qquad A\cap B=\varnothing.
\]
Put $a=|A|$, $b=|B|$, and $v=a+b\le\deg g$. For a squarefree monomial
$u_C$ on the remaining coordinates, direct expansion gives
\begin{equation}\label{eq:indicator-moment}
 E_{s,t}[I_{A,B}u_C]
 =\frac{(t)_{a+|C|}(s-t)_b}{(s)_{v+|C|}}
 \quad(v+|C|\le s).
\end{equation}
One verification is induction on $b$: inserting an additional factor
$1-u_j$ subtracts the expression with $j$ also fixed to one, and putting
the two terms over their common denominator supplies the next factor
$s-t-b$. The case $b=0$ is \eqref{eq:fractional-moments}.
In particular the indicator weight is
\[
 \pi=E_{s,t}[I_{A,B}]=\frac{(t)_a(s-t)_b}{(s)_v}\ge0.
\]
If $P$ is constant, the desired inequality is $\pi P^2\ge0$.
If $P$ is nonconstant, $v\le\deg g\le2r-2$, so all factors in this
weight are strictly positive. Write
$\widehat P=P(1_A,0_B,u_{[s]\setminus(A\cup B)})$.
For $v<s$, cancellation in \eqref{eq:indicator-moment} proves
\begin{equation}\label{eq:conditional-moments}
 E_{s,t}[I_{A,B}P^2]=\pi E_{s-v,t-a}[\widehat P^2].
\end{equation}
This identity is used only when its remaining moments are defined.
If $\widehat P$ is constant, its value is simply $\pi\widehat P^2$;
the same direct evaluation applies when $v=s$, without defining a
zero-variable conditional functional.
Otherwise let $d=\deg\widehat P\ge1$. Then
\[
 v+2d\le2r,\quad s-v\ge2d,\quad
 t-a\ge2r-1-a\ge2d-1,\quad s-t-b\ge2d-1.
\]
Lemma~\ref{lem:fractional-positive} applies to
$E_{s-v,t-a}$. This proves positivity, including for arbitrary squares
that couple all unassigned variables and for repeated original slacks.
\end{proof}

The sufficient degree range in Lemma~\ref{lem:fractional-positive} is
not the sharp classical range for the moment matrix alone. For the
\emph{full preordering} and this particular functional, the slack
products explain the two-sided restriction: if $0<t<s$ is noninteger,
$s\ge2r$, and $t<2r-1$, a product of $\lfloor t\rfloor+2\le2r$
distinct lower slacks has a negative falling-factorial moment. Applying
the same argument to upper slacks treats $s-t<2r-1$. Integer $t$ is
different: \eqref{eq:fractional-moments} is then expectation under the
uniform distribution of $t$-subsets. Sharper moment-matrix positivity
alone therefore does not extend this full-preordering construction past
the noninteger slack-product obstruction.

\subsection{The order versus region tradeoff}
\begin{theorem}[Full-preordering spatial cover bound]
\label{thm:preordering-cover}
For problem~\eqref{eq:cardinality-problem}, let
$n=k+m+z$ with positive integers $k,m,z$, let $r\ge1$ be an integer,
and let $0\le\epsilon<1/4$. Put
\begin{equation}\label{eq:qr}
 q_r=\min\{k-2r+2,\ z-2r+2,\ m(1/2-2\epsilon)\}.
\end{equation}
If $q_r>0$, every box cover certified by
\eqref{eq:preorder-normalization}--\eqref{eq:preorder-positivity} at
$T_*=1/4-\epsilon$ has at least $\mathcal N(q_r)$ members.
If $q_r\le0$, only the trivial bound of one region is asserted.
\end{theorem}
\begin{proof}
Take a box containing a witness~\eqref{eq:cardinality-witness} and
suppose $|R|<q_r$. Since $|R|$ and $k-2r+2$ are integers,
$|R|<k-2r+2$ implies $|H\cap U|\ge k-|R|\ge2r-1$;
likewise $|Z\cap U|\ge2r-1$. Set
\[
 s=|U|,\qquad t=|H\cap U|+p|M\cap U|.
\]
Then $t,s-t\ge2r-1$. Also $s\ge2r$: for $r=1$ a unit and a zero
coordinate remain, and for $r\ge2$ we have $s\ge4r-2\ge2r$.
Define $L$ by fixing $x_R=w_R$ and applying $E_{s,t}$ to $x_U=u$.
This degree-preserving substitution is normalized. It changes $e$ into
$\sum_{i\in U}u_i-t$, whose products with every substituted multiplier
of degree at most $2r-1$ vanish by
Lemma~\ref{lem:fractional-positive}. Bound slacks on $R$ are
nonnegative constants, and those on $U$ are Boolean slacks. Removing
constant factors cannot increase degrees, so
Lemma~\ref{lem:assignment-localizer} proves every required positivity
constraint.

The Boolean penalty on $U$ vanishes, leaving
\begin{equation}\label{eq:restricted-penalty}
 L[F]=|M\cap R|p(1-p)\le |R|p(1-p)
 <\frac{q_r}{2m}\le\tfrac14-\epsilon.
\end{equation}
The strict inequality follows from $|R|<q_r$ and $q_r>0$, including
$R=\varnothing$. This feasible functional precludes certification.
Thus every certified box containing a witness has $|R|\ge q_r$,
and Lemma~\ref{lem:endpoint-count} proves the bound. In particular
non-strict pruning at the target is covered: the constructed value is
strictly below it. No integer rounding of $q_r/2$ is necessary.
\end{proof}

For balanced parts $k=m=z=t$, positivity of $q_r$ is precisely
$r<(t+2)/2$ together with $\epsilon<1/4$. A uniform exponential
bound as $n=3t\to\infty$ follows whenever $\epsilon<1/4$ is fixed
and $r\le ct$ for a fixed $c<1/2$. At $\epsilon=1/8$, the exponent
$n/24$ survives whenever
\[
 t-2r+2\ge t/4,\qquad\hbox{equivalently}\qquad r\le3t/8+1.
\]
An order approaching the positivity boundary can instead give only a
constant $q_r$; positivity alone is not an exponential-size assertion.
When $\epsilon=1/4$, the root has an available certificate because the
mixed slacks give $L[F]\ge0$. We do not divide by $T_*$ or by a
vanishing $q_r$ in any of these cases. Increasing $r$ may itself make the
oracle costly; the theorem grants that oracle and counts regions only.

The affine-inequality closure following Theorem~\ref{thm:sdp-cover}
extends to this degree. Expand a product of valid affine inequalities
in their nonnegative box-slack and balance representations and multiply
by an allowed square. Terms with a balance factor vanish within degree
$2r$, and all remaining terms are existing preordering inequalities.
This still grants no arbitrary coupled nonlinear valid cut, no localizer
of an objective cutoff, and no analytic solution of the coupled problem.
The charged-cover convention of Proposition~\ref{prop:charged-tolerance}
applies to objective-based spatial reductions.

\subsection{Unique minimizers and increasing concave costs}
\begin{corollary}[Distinct perturbations]\label{cor:unique-cardinality}
Let $0<d_1<\cdots<d_n$ and $\sum_i d_i\le\eta$, where $\eta>0$.
On $\mathcal F_{n,k}$ replace $F$ by
$F_d=F+\sum_i d_ix_i$. Its unique optimizer has
\begin{equation}\label{eq:unique-cardinality}
 x_i^*=1\ (i\le k),\qquad x_{k+1}^*=1/2,\qquad
 x_i^*=0\ (i\ge k+2),
\end{equation}
with value $F_d^*=1/4+\sum_{i\le k}d_i+d_{k+1}/2$.
For the cover assertion take $n=k+m+z$ with positive integers $k,m,z$
and integer $r\ge1$, as in Theorem~\ref{thm:preordering-cover}.
Under the order-$r$ box oracle, a cover at
$T_*=F_d^*-\epsilon$ has at least $\mathcal N(q_{r,\eta})$ members
provided $0\le\epsilon$, $\epsilon+\eta<1/4$, and
\begin{equation}\label{eq:perturbed-qr}
 q_{r,\eta}=\min\{k-2r+2,\ z-2r+2,
                         m(1/2-2\epsilon-2\eta)\}>0.
\end{equation}
There is no nonidentity coordinate permutation preserving both feasible
set and objective, even when objective equality is required only on the
demand slice.
\end{corollary}
\begin{proof}
The Hessian of $F_d$ is $-2I$, so a nonvertex, lying strictly between
two distinct feasible points, cannot minimize this strictly concave
function. All vertices have base value $1/4$. Among their linear costs,
an exchange puts the largest coordinate value at the smallest coefficient:
if $i<j$ and $x_i<x_j$, exchanging them reduces cost by
$(d_j-d_i)(x_j-x_i)>0$. Hence the unique minimizing vertex is
\eqref{eq:unique-cardinality}. It also uniquely minimizes the linear
cost over the whole slice: if $i<j$, $x_i<1$, and $x_j>0$, transferring
$\min\{1-x_i,x_j\}>0$ from coordinate $j$ to coordinate $i$ strictly
decreases that cost. At a linear-cost minimizer no such pair remains,
so the balance forces exactly \eqref{eq:unique-cardinality}.

The node constraints themselves impose $0\le L[x_i]\le1$, since
$x_i$ and $1-x_i$ are nonnegative combinations of node slacks and
constants. The constructed functional therefore gives linear cost at
most $\eta$. If $|R|<q_{r,\eta}$, its total value is strictly below
\[
 \frac{m(1/2-2\epsilon-2\eta)}{2m}+\eta
 =1/4-\epsilon\le F_d^*-\epsilon,
\]
so the previous count applies.

The slice has an interior point relative to its balance hyperplane,
namely $x_i=K/n\in(0,1)$. A linear form constant on this slice has
coefficient vector proportional to the all-ones vector. If a coordinate
permutation preserves the objective there, the permuted coefficient
vector differs from $d$ by a constant vector. Comparing sums forces this
constant to be zero, and distinctness forces the identity permutation.
\end{proof}

An explicit choice is $d_i=2\eta i/[n(n+1)]$, whose sum is $\eta$.
For $n=3t$, $\epsilon=1/8$, and $\eta=1/32$, the lower bound is
$(3/2)^{n/32}$ whenever $r\le13t/32+1$.
If $\eta<\epsilon$, the chord certificate at tolerance
$\epsilon-\eta$ has base bound at least
$1/4+\eta-\epsilon\ge F_d^*-\epsilon$; the nonnegative linear term
preserves it. Thus this explicit perturbed family again has
$2^{\Theta(n)}$ certificate size under the specified oracles.

Adding the conserved flow gives increasing strictly concave coordinate
costs:
\begin{equation}\label{eq:increasing-allocation}
 C_d(x)=\sum_i[x_i(2-x_i)+d_ix_i]=F_d(x)+K
 \quad(x\in\mathcal F_{n,k}).
\end{equation}
Each derivative is $2-2x_i+d_i>0$ on $[0,1]$ and each second derivative
is $-2$. Every equality-preserving functional adds exactly $K$, so the
optimizer and all absolute-gap conclusions are unchanged. When $k$ is
proportional to $n$, this raises the optimum to order $n$, while the
unshifted obstruction has constant size. A fixed relative tolerance
cannot be inferred from that absolute bound. Section~\ref{sec:relative-blocks}
uses fixed-dimensional blocks to address precisely this distinction.

\section{Arbitrary coordinate domains and local graph lifts}
\label{sec:coordinate-domains}
The endpoint argument uses membership, not interval geometry. We first
strengthen the information available about each original coordinate, then
show that adding arbitrarily many auxiliaries depending on that coordinate
does not cost degree in this particular moment construction.

\subsection{All valid local polynomial information}
Let $S=\prod_i S_i$, where each $S_i$ is an arbitrary nonempty subset of
$[0,1]$. Closedness, connectedness, semialgebraicity, and a finite description
are not required. Define the order-$r$ product-domain oracle by normalized
functionals on $\R[x]_{\le2r}$, all balance products
\eqref{eq:preorder-equality}, and all inequalities
\begin{equation}\label{eq:product-domain-oracle}
 L\left[\left(\prod_{j=1}^a g_j\right)P^2\right]\ge0,
 \qquad\sum_j\deg g_j+2\deg P\le2r,
\end{equation}
where each $g_j$ is any univariate polynomial nonnegative on its
coordinate set $S_i$. The empty product is allowed, factors may repeat,
and $P$ may couple every coordinate. We also grant
\begin{equation}\label{eq:product-domain-equalities}
 L[h v]=0\qquad(\deg h+\deg v\le2r)
\end{equation}
for every univariate polynomial $h$ vanishing on its $S_i$ and every
polynomial multiplier $v$ in the original coordinates. In particular this
model grants all local information through degree, not merely a chosen
list of input generators. Its representation need not be finite or
computationally efficient. The node bound is the infimum of $L[F]$ or
$L[F_d]$, as appropriate. A certified cover has original feasible set
$\mathcal F_{n,k}$, original dimension $n$, and target
$1/4-\epsilon$ or $F_d^*-\epsilon$ respectively.

\begin{theorem}[Product-domain certificates]\label{thm:product-domain-cover}
Under the respective parameter assumptions of
Theorem~\ref{thm:preordering-cover} and
Corollary~\ref{cor:unique-cardinality}, the cover bounds
$\mathcal N(q_r)$ and $\mathcal N(q_{r,\eta})$ hold for the
product-domain oracle~\eqref{eq:product-domain-oracle}--\eqref{eq:product-domain-equalities}.
\end{theorem}
\begin{proof}
For a domain containing a witness, use the exclusions
\eqref{eq:endpoint-exclusions}. If $|R|<q_r$, or $|R|<q_{r,\eta}$ in
the perturbed case, fix $x_R=w_R$ and use $E_{s,t}$ on $U$ exactly as
in Theorem~\ref{thm:preordering-cover}. Its parameter and balance
calculations are unchanged.

A local inequality on a restricted coordinate becomes a nonnegative
constant because $w_i\in S_i$. For $i\in U$, both endpoints belong to
$S_i$, and its Boolean reduction is
\begin{equation}\label{eq:local-endpoint-interpolation}
 \operatorname{red}g(u_i)=g(0)(1-u_i)+g(1)u_i,
 \qquad g(0),g(1)\ge0.
\end{equation}
Group factors belonging to the same coordinate and expand their product
in endpoint indicators. After removing constant factors, let $J$ be the
coordinates with nonconstant reduced factors. The expansion has the form
\[
 \sum_{\sigma\in\{0,1\}^J} c_\sigma
       \prod_{i\in J:\sigma_i=1}u_i
       \prod_{i\in J:\sigma_i=0}(1-u_i),\qquad c_\sigma\ge0.
\]
Every such coordinate requires at least one original degree, so
$|J|\le\sum_j\deg g_j$. If $\widehat P$ is the substituted square
polynomial, $\deg\widehat P\le\deg P$ and
$|J|+2\deg\widehat P\le2r$.
Lemma~\ref{lem:assignment-localizer} gives nonnegativity of each
indicator times $\widehat P^2$, hence every inequality in
\eqref{eq:product-domain-oracle}. This includes repeated factors and
squares coupling restricted and unrestricted coordinates.

A locally vanishing $h$ becomes zero at $w_i$ on $R$, or has zero
values at both endpoints on $U$. Its Boolean reduction is zero. This
remains true after multiplication by every allowed, possibly globally
coupled, $v$, proving \eqref{eq:product-domain-equalities}.
The objective value and perturbation estimate are exactly
\eqref{eq:restricted-penalty} and Corollary~\ref{cor:unique-cardinality}.
Thus certification forces the same restriction count, and
Lemma~\ref{lem:endpoint-count} finishes the proof.
\end{proof}
The constructed first moment need not itself belong to a disconnected
$S_i$; it is a moment satisfying the stipulated oracle. The proof uses only
endpoint and witness membership. It invokes no compact-domain hull theorem
and takes no closure of $S_i$.
A branch on $\psi_i(x_i)\le\beta$ versus $\psi_i(x_i)\ge\beta$ changes
only $S_i$, for any real-valued univariate function $\psi_i$. In particular,
branching on $x_i(1-x_i)$, with its disconnected endpoint-side region,
is covered. For an objective-based reduction $S_i\to T_i\subset S_i$,
the removed product region has coordinate set $S_i\setminus T_i$.
If nonempty and certified at the target, it can be charged directly under
Proposition~\ref{prop:charged-tolerance}, without a closure convention.

\subsection{Endpoint interpolation in lifted coordinates}
For each original coordinate $i$, introduce finitely many auxiliaries
\begin{equation}\label{eq:local-graph-map}
 y_{i,j}=\psi_{i,j}(x_i),\qquad j=1,\ldots,a_i,
 \qquad \psi_{i,j}:[0,1]\to\R.
\end{equation}
Every function is defined and finite on the entire interval. Neither
continuity nor polynomiality is required. Write
$\Gamma_i(v)=(v,\psi_{i,1}(v),\ldots,\psi_{i,a_i}(v))$.
A node restricts each local graph independently to
$\Gamma_i(S_i)$ for arbitrary nonempty $S_i\subseteq[0,1]$.
Thus restrictions in the lifted variables are measured by their exact
preimages. Original coordinates $x_i$ are retained and the original
balances remain affine. The number of auxiliaries is finite for each
instance but otherwise unrestricted.

The lifted order-$r$ oracle acts on polynomials in all retained and
auxiliary coordinates of total \emph{lifted} degree at most $2r$. It
imposes normalization; all multiples of the original affine balance
through that degree; every local polynomial equality vanishing on
$\Gamma_i(S_i)$ times every allowed global polynomial multiplier; and
\begin{equation}\label{eq:lifted-local-oracle}
 L\left[\left(\prod_j g_j\right)P^2\right]\ge0,
 \qquad\sum_j\deg g_j+2\deg P\le2r,
\end{equation}
where each $g_j$ is a polynomial in one local graph block, nonnegative
on that block's restricted graph. Degrees in this display are measured
before substitution. Squares can couple every lifted block, and factors
may repeat. Graph equations are included whenever their lifted degrees
allow them. Arbitrary coupled valid inequalities are not supplied.

Let $\widetilde F(x,y)$ be a polynomial of lifted degree at most $2r$
that agrees exactly with $F(x)$, or with $F_d(x)$, on the \emph{full}
product graph $\prod_i\Gamma_i([0,1])$. Agreement merely on the
balance-feasible graph is not the stated assumption. The node bound is
$\inf L[\widetilde F]$. A graph-region cover means that its preimages
cover the original compact set $\mathcal F_{n,k}$. Full-graph objective
agreement preserves its optimum and an attaining optimizer, even if the
lifted graph itself is noncompact because some $\psi_{i,j}$ is
noncontinuous or unbounded near a point. Targets are still
$T_*=1/4-\epsilon$ or $T_*=F_d^*-\epsilon$.

\begin{theorem}[No order loss for coordinatewise graph lifts]
\label{thm:local-graph-lift}
For the lifted model just defined, under the respective assumptions of
Theorem~\ref{thm:preordering-cover} and
Corollary~\ref{cor:unique-cardinality}, every certified cover has at
least $\mathcal N(q_r)$ or $\mathcal N(q_{r,\eta})$ regions.
The parameter is the lifted order $r$ itself, independent of the number
or polynomial degrees of the univariate auxiliaries. The dimension in
the bound is the number $n$ of original coordinates.
\end{theorem}
\begin{proof}
Fix a node containing a graph witness $\Gamma(w)$ and form $A,D,R,U$
from its preimage sets. Suppose $|R|$ is below the appropriate threshold.
The cardinality parameters $s,t$ satisfy the same hypotheses as before.
Define an algebraic substitution $\Phi$ by
\begin{align}
 \Phi(x_i)&=w_i,&
 \Phi(y_{i,j})&=\psi_{i,j}(w_i) &&(i\in R),\notag\\
 \Phi(x_i)&=u_i,&
 \Phi(y_{i,j})&=\psi_{i,j}(0)
       +(\psi_{i,j}(1)-\psi_{i,j}(0))u_i &&(i\in U).
 \label{eq:graph-interpolation-map}
\end{align}
Every image is affine or constant, so
$\deg\Phi(P)\le\deg P$ for every lifted polynomial. Define
\[
 L[P]=E_{s,t}[\Phi(P)]\qquad(\deg P\le2r).
\]
This is well-defined and normalized. The original affine balance maps
to $\sum_{i\in U}u_i-t$, and any allowed multiplier maps to degree at
most $2r-1$. Lemma~\ref{lem:fractional-positive} proves all balance
products.

A local valid inequality $g$ maps to a nonnegative constant on $R$.
For $i\in U$, the two values of $\Phi(g)$ are
$g(\Gamma_i(0))$ and $g(\Gamma_i(1))$. Both are nonnegative because
both endpoint graph tuples belong to the restricted local graph.
Thus its Boolean reduction has the nonnegative endpoint form
\eqref{eq:local-endpoint-interpolation}. Grouping and expanding all
local factors gives nonnegative combinations of assignment indicators
on at most $\sum_j\deg g_j$ distinct Boolean coordinates. Moreover
$\deg\Phi(P)\le\deg P$, so each indicator times $\Phi(P)^2$ lies
in the degree range of Lemma~\ref{lem:assignment-localizer}.
This proves every inequality~\eqref{eq:lifted-local-oracle}, including
squares coupling all original and auxiliary variables.

For a local equality $h$, $\Phi(h)$ vanishes at the deterministic
witness tuple on $R$ or at both endpoint graph tuples on $U$. Its
Boolean reduction is zero. For any allowed multiplier $v$,
$\operatorname{red}(\Phi(h)\Phi(v))=0$, and the degree bound ensures
$L[hv]$ is defined. Thus every granted local equality product holds.
The same reasoning would also respect any polynomial identity valid on
the full product graph, even a coupled identity, and its allowed
multiples: at every Boolean assignment the image is an actual graph
point. This observation does not grant coupled inequalities that use
the balances or objective cutoff.

At each Boolean assignment of $u_U$, \eqref{eq:graph-interpolation-map}
is an actual full-graph tuple with $x_R=w_R$. Therefore
$\Phi(\widetilde F)$ agrees at every Boolean point with the polynomial
obtained by substituting $x_R=w_R,x_U=u$ in $F$ or $F_d$. Two
polynomials agreeing on the Boolean cube have the same multilinear
reduction: successively evaluating at zero and one in each coordinate
proves uniqueness of that representation. Both polynomials have degree
at most $2r$. Their functional values are consequently identical.
The penalty is $|M\cap R|p(1-p)$ and any perturbation contributes at
most $\eta$. The strict below-target estimate and the endpoint count
from the preceding proofs finish the theorem.
\end{proof}

The substitution $\Phi$ constructs a moment functional. It is not an
identity $\psi_i(x_i)=\psi_i(0)+(\psi_i(1)-\psi_i(0))x_i$ for continuous
$x_i$. The proof checks true graph identities only at deterministic or
Boolean endpoint tuples. For example, it covers a square-root auxiliary
$y_i=\sqrt{x_i}$ with its polynomial identity $y_i^2=x_i$, and a
penalty auxiliary $y_i=x_i(1-x_i)$ with linear objective $\sum_i y_i$.
Both maps are finite on all of $[0,1]$ and their objective agreement, when
used to encode the intended objective, is checked on the full graph.
It also covers finite-valued nonpolynomial maps that have few polynomial
identities; the oracle still grants every local polynomial inequality
valid on the restricted graph.

The theorem requires independent local restrictions and retained original
coordinates. Eliminating $x_i$ through a nonlinear inverse, branching on
functions of several original coordinates, granting arbitrary coupled
valid cuts, or replacing the oracle by the exact hull of the coupled
feasible graph changes the model. Arbitrarily many auxiliaries also mean
that an exponential bound in original dimension need not be exponential
in lifted dimension.

\begin{proposition}[Literal polynomial substitution as a coarser bound]
\label{prop:literal-lift}
If all $\psi_{i,j}$ are polynomials of degree at most an integer $D\ge1$,
with the originals counted as degree one, literal graph substitution
also proves the lifted cover bounds with effective order $rD$ in place
of $r$. This is a valid weaker conclusion of the same type.
\end{proposition}
\begin{proof}
Let $\Psi$ replace each auxiliary by its defining polynomial and fix
no coordinates yet. Then $\deg\Psi(P)\le D\deg P$.
Take the product-domain functional of order $rD$ from
Theorem~\ref{thm:product-domain-cover} and set $\widetilde L[P]=L[\Psi(P)]$.
A lifted local inequality pulls back to a univariate polynomial valid on
$S_i$, and
\[
 \sum_j\deg\Psi(g_j)+2\deg\Psi(P)
 \le D\left(\sum_j\deg g_j+2\deg P\right)\le2rD.
\]
Thus every localizer holds; local equalities pull back to equalities
vanishing on $S_i$, and graph identities vanish identically.
The original balance has degree one after substitution, so for a lifted
multiplier $v$ of degree at most $2r-1$,
\[
 \deg(e\Psi(v))\le1+D(2r-1)\le2rD.
\]
All equality products are therefore available. Objective agreement on
the full graph means the substituted objective equals $F$ or $F_d$
on $[0,1]^n$. Since both are polynomials and the cube has interior, this
is a polynomial identity. Hence its evaluated objective is the required
one. The product-domain count at order $rD$ applies.
\end{proof}
For quadratic penalty auxiliaries the literal argument uses
$\min\{k-4r+2,z-4r+2,m(1/2-2\epsilon)\}$.
Theorem~\ref{thm:local-graph-lift} improves this to $q_r$ by using
endpoint interpolation instead of literal substitution. It is this
specific construction, not invariance of polynomial hierarchies under
arbitrary nonlinear reformulations, that removes the degree loss.

\section{Relative tolerance, tensor blocks, and escape cuts}
\label{sec:relative-blocks}
A fixed relative-gap obstruction for increasing concave costs follows by
using many independent blocks of fixed size. The node oracle is still
global: its square multipliers can couple variables from every block.
The independence is a feature of the optimization problem, and also gives
a short alternative certificate outside the single-cover model.

Fix integers $r\ge1$, $t\ge2r-1$, and $G\ge1$. There are $G$ blocks,
each with $\ell=3t$ original coordinates and its own balance:
\begin{equation}\label{eq:relative-problem}
 \begin{split}
 n&=3tG,\qquad K=t+\tfrac12,\qquad
 \sum_{i=1}^{3t}x_{b,i}=K\quad(b=1,\ldots,G),\qquad 0\le x_{b,i}\le1,\\
 C(x)&=\sum_{b=1}^G\sum_{i=1}^{3t}
       \bigl[x_{b,i}(2-x_{b,i})+d_{b,i}x_{b,i}\bigr].
 \end{split}
\end{equation}
Assume $d_{b,1}<\cdots<d_{b,3t}$ are positive and
$\sum_i d_{b,i}\le\eta$ for every block, with $\eta>0$.
Every coordinate cost is strictly increasing and strictly concave, by
its derivatives $2-2x+d_{b,i}>0$ and $-2$.
Corollary~\ref{cor:unique-cardinality} applied blockwise gives the exact
optimum
\begin{equation}\label{eq:relative-optimum}
 C^*=G(K+\tfrac14)+
       \sum_{b=1}^G\left(\sum_{i\le t}d_{b,i}+\tfrac12d_{b,t+1}\right)
 \ge G(K+\tfrac14).
\end{equation}

On an arbitrary coordinate product domain, grant the order-$r$ oracle
\eqref{eq:product-domain-oracle}--\eqref{eq:product-domain-equalities}
with all $G$ balance equalities and every polynomial multiplier through
\emph{total} degree $2r$. Squares may involve all $n$ variables. This
includes the box preordering and, for boxes at $r=1$, full global
SDP--RLT with every balance product. Let $\operatorname{LB}_r(S)$ be
the infimum of $L[C]$. A relative certificate at tolerance
$0<\theta<1$ is a finite cover of the feasible set by regions that are
infeasible or satisfy
\begin{equation}\label{eq:relative-target}
 \operatorname{LB}_r(S)\ge T_*=(1-\theta)C^*.
\end{equation}
Pruning at $(1-\theta)U$ for any feasible incumbent $U\ge C^*$ implies
this condition. The dimension in all ensuing bounds is original
$n=3tG$.

\begin{lemma}[Tensor positivity with a total-degree truncation]
\label{lem:tensor-preordering}
For each block $b$, suppose $L_b$ is a normalized functional on block
polynomials through degree $2r$, satisfying the block balance products,
local equalities, and every block preordering inequality. Define $L$ on
global monomials of degree at most $2r$ by
\begin{equation}\label{eq:tensor-moment}
 L\left[\prod_b x_b^{\alpha_b}\right]=
             \prod_b L_b[x_b^{\alpha_b}],
\end{equation}
and extend linearly. Then $L$ satisfies the corresponding global oracle,
including every globally coupled square multiplier.
The same assertion holds when a block contains coordinatewise graph
auxiliaries and degree is measured in the lifted variables.
\end{lemma}
\begin{proof}
Each block degree is at most the global degree, so the definition is
well-defined and normalized. To test a balance equality from block $b$
times a global monomial, factor the latter across blocks. Its expectation
is the block-$b$ balance product, of degree at most $2r$, times moments
of the other blocks, so it vanishes. Linearity treats every multiplier.
The argument for local equalities is identical and also covers products
involving several equalities within the total degree bound.

For a preordering expression let $g=\prod_b g_b$ group all local
inequality factors according to their block, and let $d=\deg P$ with
$\sum_b\deg g_b+2d\le2r$. For every $b$, form the matrix
\[
 M_b(\alpha,\beta)=L_b[g_b x_b^\alpha x_b^\beta],
 \qquad |\alpha|,|\beta|\le d.
\]
It is defined and positive semidefinite: the degree bound
$\deg g_b+2d\le2r$ allows the block preordering inequality for every
polynomial in this index span. Their tensor product is positive
semidefinite, as follows for example by tensoring Gram representations.
Keep only tuples of indices whose total degree across blocks is at most
$d$. The corresponding principal submatrix has entries exactly
$L[gx^\alpha x^\beta]$ by \eqref{eq:tensor-moment}; each such entry has
available global degree at most $2r$. Its quadratic form on the
coefficient vector of $P$ is $L[gP^2]\ge0$.
The unused entries of the larger tensor matrix need not represent
globally available moments: it is only a Gram construction whose required
principal submatrix has the stated total-degree interpretation. Nothing
in this proof uses that the coordinates within a block were unlifted.
\end{proof}

\begin{theorem}[A fixed relative-gap cover bound]
\label{thm:relative-cover}
For \eqref{eq:relative-problem} and the global product-domain oracle
above, set
\begin{equation}\label{eq:relative-parameters}
 q_0=t-2r+2\ge1,\qquad
 \tau=\tfrac14-\theta(K+\tfrac14)-\eta.
\end{equation}
If $\tau>0$, every certificate at target
\eqref{eq:relative-target} has at least
\begin{equation}\label{eq:relative-cover-bound}
 (3/2)^{q_0G\tau}=(3/2)^{q_0\tau n/(3t)}
\end{equation}
regions. For fixed $r,t,\theta,\eta$ satisfying these conditions, this
is exponential in $n$. If $\tau\le0$, only the trivial bound one is
asserted. The same conclusions hold for coordinatewise graph lifts
satisfying the local oracle and full-graph objective-agreement assumptions
of Theorem~\ref{thm:local-graph-lift}, retaining all original balances;
there is no auxiliary-degree loss in $r$.
\end{theorem}
\begin{proof}
Choose independently in each block a uniform ordered partition
$(H_b,M_b,Z_b)$ into three sets of size $t$, and set
$w_b=\ind_{H_b}+p\ind_{M_b}$ with $p=1/(2t)$.
Every combined witness is feasible. Fix a product domain containing one
and define $A_b,D_b,R_b$ by endpoint exclusion within block $b$.
If $|R_b|<q_0$, the construction of
Theorem~\ref{thm:product-domain-cover} gives a feasible block functional
whose penalty is at most
\[
 |R_b|p(1-p)\le\frac{|R_b|}{2t}
                    \le\frac{|R_b|}{2q_0},
\]
since $1\le q_0\le t$. No absolute pruning target is needed to construct
this functional: only the endpoint restriction count is used.
If $|R_b|\ge q_0$, instead use exact evaluation at the actual block
witness. Its penalty is
$tp(1-p)=1/2-1/(4t)\le1/2\le |R_b|/(2q_0)$.
Thus in either case the block functional is feasible through degree
$2r$, its first moments lie in $[0,1]$, and its penalty has the same
upper bound. Lemma~\ref{lem:tensor-preordering} gives a feasible global
functional of objective at most
\begin{equation}\label{eq:relative-witness-bound}
 GK+\frac1{2q_0}\sum_b|R_b|+G\eta.
\end{equation}
A certified region containing this witness must therefore satisfy
\[
 GK+\frac1{2q_0}\sum_b|R_b|+G\eta
 \ge\operatorname{LB}_r(S)\ge(1-\theta)C^*
 \ge(1-\theta)G(K+\tfrac14).
\]
It follows that $\sum_b|R_b|\ge2q_0G\tau$.

For a fixed domain, witness membership factors across blocks. In block
$b$, endpoint avoidance bounds its witness fraction by
\[
 \min\{(2/3)^{|A_b|},(2/3)^{|D_b|}\}
 \le(2/3)^{|R_b|/2}.
\]
Multiplying over independent blocks bounds the full witness fraction in
a certified region by $(2/3)^{q_0G\tau}$. The union bound over a cover
proves \eqref{eq:relative-cover-bound}. The restriction counts belong
to the region, not to the selected witness, so this is a uniform bound
for every certified region with at least one witness.

For lifted regions work with the exact original-coordinate preimages.
In a lightly restricted block use the affine endpoint substitution
\eqref{eq:graph-interpolation-map}; in a heavily restricted block use
evaluation at its true graph witness. These block functionals satisfy all
lifted localizers and equality products by the proof of
Theorem~\ref{thm:local-graph-lift}, with the original order $r$.
Lemma~\ref{lem:tensor-preordering} proves every global lifted localizer,
including squares coupling auxiliaries from all blocks. It also proves
all balance and local equality products.

For completeness, a lifted polynomial objective agreeing with $C$ on
the full product graph need not be written as a sum of block polynomials.
After the block substitutions, it agrees with the substituted $C$ at
every Boolean assignment of all remaining formal variables; the heavily
restricted blocks are actual fixed graph tuples. Affine substitution
preserves its total degree at most $2r$, and uniqueness of Boolean
reduction shows the two polynomials have the same value under the tensor
functional. That functional annihilates Boolean identities in every
remaining coordinate because each factor does. Hence
\eqref{eq:relative-witness-bound} still holds, and the preimage witness
count is identical. This proves the lifted assertion without assuming
blockwise separability of the written lifted objective.
\end{proof}

For the explicit parameters
\begin{equation}\label{eq:relative-example}
 r=1,\quad t=2,\quad\theta=1/32,\quad\eta=1/32,
 \qquad K=5/2,\quad q_0=2,\quad\tau=17/128,
\end{equation}
the lower bound is $(3/2)^{17n/384}$, even for full global SDP--RLT.
For any other fixed order choose fixed $t\ge2r-1$, $0<\eta<1/4$,
and $0<\theta<(1/4-\eta)/(K+1/4)$.
There is no assertion of a common positive relative tolerance while
$r$ grows without bound in this fixed-block construction.

\begin{proposition}[Uniqueness and symmetry modulo the balances]
\label{prop:relative-asymmetry}
Problem~\eqref{eq:relative-problem} has a unique optimizer: in every
block its first $t$ coordinates equal one, coordinate $t+1$ equals
$1/2$, and the rest equal zero. For the explicit choice
\begin{equation}\label{eq:squared-perturbation}
 d_{b,i}=\frac{\eta\bigl((b-1)3t+i\bigr)^2}{3t\,n^2},
 \qquad b\in[G],\ i\in[3t],
\end{equation}
no nonidentity permutation of the original coordinates preserves both
the feasible set and the objective restricted to that set.
\end{proposition}
\begin{proof}
Each independent block has the unique optimizer from
Corollary~\ref{cor:unique-cardinality}; their concatenation is therefore
the unique global optimizer. The coefficients in
\eqref{eq:squared-perturbation} are positive, globally distinct, and
increasing in each block. Each is at most $\eta/(3t)$, so its block
sum is at most $\eta$ as required.

The point with every coordinate $K/(3t)$ is strictly between zero and
one. Thus the affine hull of the feasible set is exactly the displayed
balance space. Its normal row space consists of vectors constant on each
block. A feasible-set-preserving coordinate permutation preserves this
row space. The image of a block indicator is a zero-one vector in it
with precisely $3t$ ones, so it is the indicator of one entire block.
Hence such a permutation maps whole blocks to whole blocks.

The unperturbed quadratic and linear parts are invariant under every
coordinate permutation. If the perturbed objective is preserved on the
feasible set, its coefficient difference is constant within each block;
this necessary condition follows by varying coordinates within their
balance hyperplane. Consequently the coefficient multiset of each block
must equal the coefficient multiset of its image block up to translation.
Put $c=\eta/(3t\,n^2)$ and $\ell=3t$. The sorted successive differences
in block $b$ are
\[
 c\bigl(2((b-1)\ell+i)+1\bigr),\qquad i=1,\ldots,\ell-1.
\]
Their first entries distinguish all blocks. Translation preserves these
differences, so every block is fixed. Within a fixed block, comparison
of the minimum shows that a finite coefficient set cannot equal a
nonzero translate of itself. The translation is zero, and its distinct
entries force the identity permutation.
\end{proof}
Distinctness alone would not prove the statement modulo several balances:
two blocks with coefficient lists differing by additive offsets could be
exchanged without changing the objective on their equal-demand slices.
The squared-index choice explicitly excludes this possibility. The
symmetry assertion is about the original optimization problem; redundant
auxiliary variables can of course introduce symmetries of their own.

\subsection{A short alternative certificate and the scope of the obstruction}
The blocks have fixed dimension. Solving each separately and summing its
exact certified value gives a certificate of size proportional to $G$,
if reuse and addition of component certificates is allowed. Even spatial
certification of each separate block by
Proposition~\ref{prop:chord-upper} has fixed cost for fixed $t$.
Theorem~\ref{thm:relative-cover} counts regions in one global cover with
the stated node oracle. It does not constrain a method that retains
separate component certificates and adds their bounds. In this family the
explicit sorted optimizer is an even simpler way to solve each component.

There is also a classical coupled cut that closes every root under study.
Its low written degree does not mean that its nonnegativity follows from
a low-degree local preordering.

\begin{proposition}[A Boolean-quadric clique cut closes the root]
\label{prop:clique-escape}
For integers $1\le k\le n-1$, the inequality
\begin{equation}\label{eq:clique-cut}
 2\sum_{i<j}X_{ij}-2k\sum_i\mu_i+k(k+1)\ge0
\end{equation}
is valid on the continuous quadratic graph
$(\mu,X)=(x,xx^T)$ for $x\in[0,1]^n$.
Adding it to the root equality-product relaxation yields the exact lower
bound $1/4$ for $F$, and the exact root bound for $F_d$ and $C_d$.
One such inequality per block makes the root bound for
\eqref{eq:relative-problem} exact.
\end{proposition}
\begin{proof}
This is Padberg's clique inequality
\citep[Lemma 2, equation (17)]{Padberg1989}, with the whole coordinate
set and parameter $k$, after multiplying by two and changing sides.
Its continuous validity also has a direct proof. After substituting
$X_{ij}=x_ix_j$, the left side is multiaffine. Successively choosing
endpoints can only decrease its value, so its minimum on the cube is
attained at a Boolean vertex. At a vertex with $s$ ones its value is
$(s-k)(s-k-1)\ge0$, since $s-k$ is an integer.

The balance and its linearized products imply
$\sum_{i,j}X_{ij}=K\sum_i\mu_i=K^2$.
Equation~\eqref{eq:clique-cut} therefore gives
\[
 \sum_iX_{ii}\le K^2-2kK+k(k+1)=k+\tfrac14,
 \qquad \sum_i(\mu_i-X_{ii})\ge\tfrac14.
\]
For the perturbed objective, first moments meet the unit bounds and
balance, and their linear cost is at least
$\sum_{i\le k}d_i+d_{k+1}/2$, the linear-program minimum proved in
Corollary~\ref{cor:unique-cardinality}. Adding the two bounds gives
$F_d^*$. The actual optimizer supplies an attaining exact moment matrix,
so the root bound is exact. Adding the conserved $K$ treats $C_d$,
and applying the argument separately to all block balances treats $C$.
The deduction does not use PSD. For the optional boundary case $k=0$,
the same polynomial is $2\sum_{i<j}x_ix_j\ge0$ and the algebra still
applies directly.
\end{proof}

Thus these easy-to-optimize families exhibit limits of particular
spatial certificate systems, including full node preorderings at the
specified orders, while known global cuts avoid the obstruction. In a
penalty lift the coupled affine inequality $\sum_i y_i\ge1/4$ would
also close the unperturbed root; its validity is precisely the global
conclusion, so it is not supplied as a local graph inequality.
These escape certificates motivate seeking different families when
stronger coupled cuts are granted. In particular, a later signed cubic
XOR construction serves a different spatial purpose and must not be
identified with the positive multilinear gap examples of the earlier
sections.

\section{Signed cubic certificates with exact quadratic moment information}
\label{sec:xor-quadratic}
The clique inequality in Proposition~\ref{prop:clique-escape} closes the
fractional-cardinality example. A signed cubic objective permits a different
obstruction: all its first and second moments can be genuine moments while
its three-variable correlations still give a false objective value. The
Boolean moment construction used below is classical
\citep{Schoenebeck2008}. We prove its transfer to continuous coordinate
regions, including the exact quadratic moment hull at each region.

Let $E$ be a multiset of $m>0$ clauses on $[n]$, each containing three
\emph{distinct} coordinates, with signs $b_e\in\{-1,1\}$. Occurrences count
multiplicity, and $\Delta=\max_i|\{e:i\in e\}|\ge1$. Set
\begin{equation}\label{eq:xor-objective}
 c_e(x)=\frac{1-b_e\prod_{i\in e}x_i}{2},\qquad
 F(x)=\frac1m\sum_{e\in E}c_e(x),\qquad x\in[-1,1]^n.
\end{equation}
Every $c_e$ lies in $[0,1]$ on the cube. Successively minimizing in each
coordinate moves to a Boolean vertex without increasing $F$, so
\begin{equation}\label{eq:xor-vertex-min}
 \OPT=\min_{[-1,1]^n}F=\min_{w\in\pmcube}F(w).
\end{equation}
Here clauses define the objective; the continuous problem does not impose
perfect clause satisfaction as a feasibility constraint.

\subsection{The node oracle and the general transfer}
A Boolean pseudoexpectation through degree $d$ is a linear functional
$\mathcal E$ on $\R[x]_{\le d}$ with $\mathcal E[1]=1$, respecting
$x_i^2=1$ through that degree and positive on $P^2$ whenever
$2\deg P\le d$. Boolean reduction replaces a monomial by the character
$\chi_A(x)=\prod_{i\in A}x_i$, where $A$ consists of its odd exponents;
$\chi_\varnothing=1$. Respecting the Boolean identities means that a
polynomial and its reduction have the same functional value.

For a nonempty coordinate product $B=\prod_i S_i\subseteq[-1,1]^n$, define
\begin{equation}\label{eq:xor-moment-hull}
 \mathcal Q(B)=\conv\{(x,xx^T):x\in B\}.
\end{equation}
Our order-$r$ node functional $L$ is defined through degree $2r$, normalized,
and satisfies
\begin{align}
 L\left[\left(\prod_{j=1}^a g_j\right)P^2\right]&\ge0,
 &\sum_j\deg g_j+2\deg P&\le2r,\label{eq:xor-node-preorder}\\
 L[qv]&=0,&\deg q+\deg v&\le2r,\label{eq:xor-node-local-equality}\\
 \bigl((L[x_i])_i,(L[x_ix_j])_{ij}\bigr)&\in\mathcal Q(B).
 &&\label{eq:xor-node-quadratic}
\end{align}
In \eqref{eq:xor-node-preorder}, each $g_j$ is any univariate polynomial
nonnegative on its coordinate set. Products may repeat factors, and the
empty product and globally coupled $P$ are allowed. In
\eqref{eq:xor-node-local-equality}, $q$ is any univariate polynomial
vanishing on its coordinate set, and $v$ is arbitrary through degree.
For boxes, the weaker oracle using just the two bound slacks per coordinate,
with arbitrary repeated products, is also covered by all ensuing lower
bounds. Let $\operatorname{LB}_r(B)=\inf L[F]$, where $r\ge2$ makes the
cubic objective available. An empty feasible set of functionals has bound
$+\infty$.

For compact boxes, \eqref{eq:xor-node-quadratic} is equivalent to adding
\emph{linearly on moments} every quadratic polynomial inequality valid
on the box: separation applies to the compact convex hull in
\eqref{eq:xor-moment-hull}. Products or square localizers of arbitrary
coupled quadratic inequalities are not part of this grant. For nonclosed
coordinate sets we impose actual hull membership, which can be stronger
than all closed halfspace inequalities. No closure is taken.

\begin{theorem}[Quadratic-hull cover transfer]\label{thm:xor-transfer}
Let $r\ge2$ be an integer. Suppose a Boolean pseudoexpectation $\mathcal E$
through degree $4r$ satisfies
\begin{equation}\label{eq:xor-source-assumptions}
 \mathcal E[\chi_e]=b_e\quad(e\in E),\qquad
 \bigl(\mathcal E[x],\mathcal E[xx^T]\bigr)
 \in\conv\{(w,ww^T):w\in\pmcube\}.
\end{equation}
Every finite cover of $[-1,1]^n$ by coordinate products certified by
\eqref{eq:xor-node-preorder}--\eqref{eq:xor-node-quadratic} at a target
$T_*>0$ has at least
\begin{equation}\label{eq:xor-transfer-count}
 2^{mT_*/\Delta}
\end{equation}
regions. The coordinate sets may be disconnected, nonclosed, or have no
finite polynomial description.
\end{theorem}
\begin{proof}
Use all $2^n$ Boolean vertices as equally weighted witnesses. For a domain
$B$ containing a witness $w$, define
\[
 R=\{i:\{-1,1\}\not\subseteq S_i\},\qquad U=[n]\setminus R,
 \qquad L[p]=\mathcal E[p(w_R,x_U)].
\]
This fixes coordinates deterministically and then takes a marginal of
$\mathcal E$. It is not conditioning on an event under $\mathcal E$ and
requires no positive mass for the chosen signs. Substitution does not
increase degree, and $L[1]=1$.

First consider a box-slack product $gP^2$. Slacks in $R$ become
nonnegative constants. In $U$, the interval contains both endpoints and
therefore equals $[-1,1]$. Each remaining slack is $1+x_i$ or $1-x_i$.
Modulo Boolean identities, their product is zero if opposite signs occur
on a coordinate, and otherwise is a nonnegative multiple of an assignment
indicator $I$. Repeated factors change only the nonnegative multiplier.
If $a=\deg g$ and $d=\deg P$, the indicator has degree at most $a$.
Writing $\widehat P=P(w_R,x_U)$, Boolean idempotence gives
\begin{equation}\label{eq:xor-indicator-square}
 I\widehat P^2=(I\widehat P)^2\pmod{x_i^2-1},\qquad
 \deg(I\widehat P)\le a+d\le2r.
\end{equation}
Both sides are within degree $4r$, so Boolean reduction and square
positivity of $\mathcal E$ prove the inequality.

For the stronger local oracle, a factor on $R$ is again a nonnegative
constant. On $U$ its Boolean reduction is
\begin{equation}\label{eq:xor-local-endpoints}
 g(x_i)=g(1)\frac{1+x_i}{2}+g(-1)\frac{1-x_i}{2}
 \pmod{x_i^2-1},\qquad g(1),g(-1)\ge0.
\end{equation}
Remove constant factors and expand the product. Every summand is a
nonnegative multiple of a product of endpoint indicators. The number of
nonconstant factors is at most the original sum of generator degrees;
hence its indicator degree is at most that sum. Applying
\eqref{eq:xor-indicator-square} proves every allowed localizer. This
argument includes arbitrary repeated factors. A local equality is zero
at the fixed witness on $R$, or has both endpoint values zero on $U$;
its Boolean reduction, also after any allowed multiplier, vanishes.

Let $\mu$ be a Boolean distribution realizing the original moments through
degree two. Marginalize $\mu$ on $U$ and set its other coordinates to
$w_R$. The resulting finite distribution is supported in $B$. Its means,
diagonal moments, and cross moments are exactly those of $L$, proving
\eqref{eq:xor-node-quadratic} even for nonclosed sets.

An untouched clause keeps cost zero. At most $\Delta|R|$ clauses touch
$R$, including multiplicity. After substitution each cost has the form
$(1-\sigma\chi_A)/2$ with $|A|\le3$. Since $\chi_A^2=1$,
\[
 0\le\mathcal E[(1\pm\chi_A)^2]
       =2(1\pm\mathcal E[\chi_A])
\]
proves its cost lies in $[0,1]$; these squares have degree at most six,
within $4r$. Thus
\begin{equation}\label{eq:xor-node-cost}
 \operatorname{LB}_r(B)\le L[F]\le\frac{\Delta|R|}{m}.
\end{equation}
Certification forces $|R|\ge mT_*/\Delta$. Every coordinate in $R$
permits exactly one Boolean endpoint when the domain contains a witness,
so the witness fraction in $B$ is $2^{-|R|}$ and is at most
$2^{-mT_*/\Delta}$. Domains containing no witnesses have fraction zero.
The union bound for any cover proves \eqref{eq:xor-transfer-count}.
\end{proof}

This result counts certified regions in the sense of
Section~\ref{subsec:certificate-model}. It permits arbitrary real coordinate
split points and variable choices. A complete split tree has at least this
many leaves; regions discarded by objective propagation must also be
certified and charged. The theorem does not bound leaves remaining after
free deletion of feasible regions.

\subsection{Classical signed moments and an occurrence-controlled family}
\begin{lemma}[Realizing signed quadratic moments]\label{lem:signed-moment-law}
Let an augmented moment matrix indexed by $1,z_1,\ldots,z_N$ be positive
semidefinite, with all diagonal entries one and all entries in
$\{0,-1,1\}$. There is a probability distribution on $\{-1,1\}^N$
with exactly its prescribed first and second moments.
\end{lemma}
\begin{proof}
Represent the matrix as a Gram matrix of unit vectors $v_0,v_1,\ldots,v_N$.
Inner product $1$ or $-1$ forces equality or opposition. Partition the
vectors into classes up to sign and choose a representative for each.
Distinct representatives are orthogonal because their inner products
belong to $\{0,-1,1\}$ and cannot have absolute value one. Choose $v_0$
as representative of its own class. Give this class the deterministic
random value one; give every other class an independent unbiased sign.
If $v_j=\sigma_j v_{C(j)}$, set $Z_j=\sigma_j\xi_{C(j)}$.
A coordinate in the constant class has mean $\sigma_j$, and other
coordinates have mean zero. Two coordinates in one class have product
mean $\sigma_i\sigma_j$; different classes give product mean zero,
including when one is the constant class. These are exactly all Gram
entries, and the distribution has finite support.
\end{proof}

We spell out the width convention of the source because two degree losses
will occur in the spatial argument. The following is the signed-character
construction in the proof of Schoenebeck's Lemma~13
\citep[Theorems 11--12 and Lemma 13]{Schoenebeck2008}, written as a
polynomial functional.

\begin{lemma}[Width and character degree]\label{lem:xor-width-moments}
Suppose a signed XOR formula has no width-$w$ XOR-resolution contradiction,
where $w\ge3$ is an integer and all initial clause widths are at most $w$.
There is a normalized Boolean functional through degree $w$ with every
character moment in $\{0,-1,1\}$, satisfying every clause in expectation
and positive on squares of degree at most $\lfloor w/2\rfloor$.
\end{lemma}
\begin{proof}
In multiplicative notation, the resolution rule derives
$\chi_{A\mathbin\triangle B}=st$ from $\chi_A=s$ and $\chi_B=t$
when $|A\mathbin\triangle B|\le w$. Let $\mathcal C_w$ be its closure,
including $\chi_\varnothing=1$. Absence of contradiction means that a
support can have at most one derived sign: opposite signs would derive
$\chi_\varnothing=-1$. Define $y_A$ for $|A|\le w$ to be that sign
when derived, and zero otherwise. In particular $y_\varnothing=1$.
Define $\mathcal E$ by Boolean reduction and $\mathcal E[\chi_A]=y_A$.
This proves normalization, signed moments, and the clause means.

Index characters by supports of size at most $\ell=\lfloor w/2\rfloor$.
Declare $I\sim J$ if a sign for $I\mathbin\triangle J$ is derived.
This is reflexive and symmetric. For transitivity, multiply the relations
for $I\mathbin\triangle J$ and $J\mathbin\triangle K$; the resulting
support $I\mathbin\triangle K$ has size at most $2\ell\le w$.
Thus the result is in the closure, with the product sign. Choose a
representative $I_0$ per class and orient $I$ by the derived sign
$s_I$ of $I\mathbin\triangle I_0$. Multiplication of two such relations
shows that $s_Is_J=y_{I\mathbin\triangle J}$ within a class. Give
distinct classes orthogonal unit vectors $e_C$ and put $v_I=s_Ie_{C(I)}$.
Then for all $|I|,|J|\le\ell$,
\begin{equation}\label{eq:xor-character-gram}
 \ip{v_I}{v_J}=y_{I\mathbin\triangle J}.
\end{equation}
For any polynomial $P$ of degree at most $\ell$, write its Boolean
reduction as $\sum_I p_I\chi_I$. Equation~\eqref{eq:xor-character-gram}
gives
\[
 \mathcal E[P^2]=\sum_{I,J}p_Ip_Jy_{I\mathbin\triangle J}
              =\left\|\sum_Ip_Iv_I\right\|^2\ge0.
\]
This proves positivity for arbitrary polynomial squares, including sums
whose terms jointly involve more than $\ell$ variables. The constant
character has a unit Gram vector, as normalization requires.
\end{proof}

At density eight, Theorem~11 of \citet{Schoenebeck2008} gives a constant
$a>0$ such that width at least $an$ is available with probability $1-o(1)$,
after reducing $a$ to absorb integer rounding. Take the source parameters
$k=3$, $d=8$, $\delta=\gamma=1/4$, and its density exponent
$\varepsilon_{\rm src}=0$. Theorem~12's density condition is
$8\ge1+8\log2$, and the denominator $k/2-\gamma-1=1/4$ is positive.
Theorem~11 supplies the positive width constant; Lemma~13 converts width
$w$ to level $\lfloor w/2\rfloor$. Lemma~\ref{lem:xor-width-moments}
verifies directly that this means moments through degree $w$ and square
positivity through degree $\lfloor w/2\rfloor$, not moments only through
half the width. We need $4r\le w$ below and $4rD\le w$ for the lifts.
This choice avoids the boundary $\gamma=1/2$ in the source's illustrative
three-variable specialization.

\begin{theorem}[A bounded-occurrence cubic family]\label{thm:xor-family}
There are absolute constants $a>0$ and $n_0$ such that for every integer
$n\ge n_0$ there is an objective \eqref{eq:xor-objective} with
\begin{equation}\label{eq:xor-family-parameters}
 7n\le m\le8n,\qquad \Delta\le64,\qquad \OPT\ge1/8,
\end{equation}
and a Boolean functional of available degree at least $an$ as in
Lemma~\ref{lem:xor-width-moments}. For every integer $r\ge2$ with
$4r\le an$, any certificate at absolute tolerance $1/16$, or relative
tolerance $1/2$, under the quadratic-hull node oracle has at least
\begin{equation}\label{eq:xor-family-count}
 2^{7n/1024}
\end{equation}
regions. The same instance works simultaneously for all these orders.
\end{theorem}
\begin{proof}
Sample $m_0=8n$ independent clauses, each with a uniform three-element
support and an independent uniform sign. This is the random-XOR model in
Schoenebeck's definition: ordered distinct variables and independent literal
negations induce the same support and parity-sign distribution.
For a fixed Boolean assignment, its violation count $V$ is
$\operatorname{Bin}(8n,1/2)$, regardless of overlap among supports. For
$t>0$, the generating function and $\cosh u\le e^{u^2/2}$ give
\[
 \Prob\{V\le2n\}
 \le e^{-2nt}\E e^{-t(V-4n)}
 =e^{-2nt}\cosh(t/2)^{8n}
 \le e^{-2nt+nt^2}.
\]
The inequality for $\cosh$ follows by integrating
$(\log\cosh u)'=\tanh u\le u$ for $u\ge0$ and using symmetry.
Taking $t=1$ and a union bound over all assignments bounds the probability
of any assignment with at most $2n$ violations by
$2^ne^{-n}=e^{-(1-\log2)n}=o(1)$. Independently of needing this elementary
bound, the source gives the signed functional with linear available degree
with probability $1-o(1)$.

Let $D_i$ be the original number of clauses containing coordinate $i$;
$D_i\sim\operatorname{Bin}(8n,3/n)$. Delete every clause incident to a
coordinate with $D_i>64$. If $Z$ is the number deleted, then
$Z\le\sum_iD_i\ind_{D_i>64}$. Differentiating the binomial generating
polynomial yields
\begin{align}
 \E[D_i\ind_{D_i>64}]
 &\le2^{-64}\E[D_i2^{D_i}]
   =\frac{48(1+3/n)^{8n-1}}{2^{64}}\nonumber\\
 &\le\frac{48e^{24}}{2^{64}}<\frac18.
 \label{eq:xor-deletion}
\end{align}
For example $e<3$ makes the last inequality follow from
$384\cdot3^{24}<2^{64}$, an integer inequality. Hence
$\Prob\{Z>n\}\le\E Z/n<1/8$ by Markov's inequality.
No independence between different $D_i$ is used. For every sufficiently
large $n$, the sum of this failure bound and the two preceding $o(1)$
failure probabilities is less than one. Fix an instance in their
intersection.

At least $7n$ clauses remain, all remaining occurrences are at most $64$,
and every assignment still violates at least $n$ clauses. Equation
\eqref{eq:xor-vertex-min} gives $\OPT\ge n/m\ge1/8$.
Deleting clauses preserves the original functional, its positivity, all
character moments, and the means of retained clauses. Its degree-two
moments have an actual Boolean realization by
Lemma~\ref{lem:signed-moment-law}. Apply Theorem~\ref{thm:xor-transfer}
with any $4r\le an$. At the best incumbent the absolute target is
$\OPT-1/16\ge1/16$ and the relative target is $\OPT/2\ge1/16$.
Thus $mT_*/\Delta\ge7n/(16\cdot64)$, proving the count. Restricting
one source functional to lower degrees proves simultaneous validity.
\end{proof}

\begin{proposition}[Sensitivity and curvature]\label{prop:xor-sensitivity}
Throughout $[-1,1]^n$, the objective \eqref{eq:xor-objective} satisfies
\begin{equation}\label{eq:xor-sensitivity}
 |\partial_iF|\le\frac{\Delta}{2m},\qquad
 \sum_{j\ne i}|\partial_i\partial_jF|\le\frac{\Delta}{m},\qquad
 \|\nabla^2F\|_2\le\frac{\Delta}{m}.
\end{equation}
In the family of Theorem~\ref{thm:xor-family}, the last bound is at most
$64/(7n)$.
\end{proposition}
\begin{proof}
An incident clause contributes at most $1/(2m)$ to a first derivative
and at most two mixed second derivatives of that magnitude. The diagonal
second derivatives vanish. Summing absolute contributions gives the first
two bounds, and the symmetric Hessian's spectral norm is at most its
maximum absolute row sum.
\end{proof}
These estimates record the normalization of this family. They do not
assert a general implication from small curvature to large certificates.

\section{Certificates after bounded monomial reformulation}
\label{sec:monomial-reformulation}
The preceding transfer survives coordinate branching on bounded products.
The number of auxiliary coordinates can be arbitrary; the count depends
on the original variables represented by independent parity restrictions.

Consider a map $h:[-1,1]^n\to[-1,1]^N$ of the form
\begin{equation}\label{eq:signed-monomial-map}
 h_j(x)=\sigma_j\chi_{A_j}(x),\qquad
 \sigma_j\in\{-1,1\},\qquad 1\le|A_j|\le D,
 \qquad j=1,\ldots,N,
\end{equation}
where $D\ge1$ is an integer. Include every original coordinate as an
identity coordinate. The nonempty supports are squarefree; repeated or
overlapping supports and either sign are allowed, and $N$ is finite but
otherwise unrestricted. The continuous feasible set is the full graph
$\Gamma_h=\{h(x):x\in[-1,1]^n\}$. Retaining originals makes its Boolean
witnesses $h(w)$ distinct. Let $\Phi\in\R[z]_{\le2r}$ satisfy the exact
polynomial identity
\begin{equation}\label{eq:lift-objective-pullback}
 \Phi(h(x))=F(x).
\end{equation}
Agreement merely on Boolean points is not the hypothesis here.

For $B=\prod_{j=1}^N S'_j\subseteq[-1,1]^N$, each $S'_j$ nonempty,
use the local oracle \eqref{eq:xor-node-preorder}--\eqref{eq:xor-node-quadratic}
in the $z$ coordinates, with objective $\Phi$. In addition grant
\begin{equation}\label{eq:lift-all-identities}
 L[qv]=0\quad\hbox{if }q(h(x))\equiv0,
 \qquad\deg q+\deg v\le2r.
\end{equation}
Thus every polynomial graph identity through available degree, with every
allowed multiplier, can be imposed. One may use just a particular generating
set instead; the stronger version makes the lower bound more informative.
A polynomial vanishing on the entire continuous graph has identically
zero pullback because the original cube has nonempty interior.
The quadratic requirement is
\begin{equation}\label{eq:lift-box-hull}
 (L[z],L[zz^T])\in\conv\{(z,zz^T):z\in B\},
\end{equation}
the hull of the entire coordinate domain. It is not the hull of
$\Gamma_h\cap B$. The latter would already optimize a linear or quadratic
objective exactly. These quadratic inequalities enter linearly on moments;
the oracle does not close them under arbitrary products or localizers.

\begin{theorem}[Bounded monomial transfer, including order one]
\label{thm:monomial-transfer}
Let $r,D\ge1$ be integers with $rD\ge2$. Suppose a Boolean
pseudoexpectation $\mathcal E$ is available through degree $4rD$, is
positive on squares through degree $2rD$, satisfies
$\mathcal E[\chi_e]=b_e$ for every clause, and has every character moment
through its available degree in $\{0,-1,1\}$. For any map
\eqref{eq:signed-monomial-map} and objective
\eqref{eq:lift-objective-pullback}, every finite cover of $\Gamma_h$ by
coordinate products certified by the preceding order-$r$ oracle at
$T_*>0$ contains at least
\begin{equation}\label{eq:monomial-transfer-count}
 2^{mT_*/(\Delta D)}
\end{equation}
regions. Arbitrary nonclosed coordinate sets and arbitrary numbers of
repeated or overlapping supports are allowed.
\end{theorem}
\begin{proof}
Let $B$ contain $h(w)$ for $w\in\pmcube$. Put
\begin{equation}\label{eq:lift-restricted-supports}
 R=\{j:\{-1,1\}\not\subseteq S'_j\},\qquad
 C=\bigcup_{j\in R}A_j,\qquad U=[n]\setminus C,
 \qquad \widetilde h(x_U)=h(w_C,x_U).
\end{equation}
Each coordinate in $R$ is now fixed to its witnessed sign. An unrestricted
coordinate becomes a signed character on $U$, possibly constant, and both
of its endpoints belong to its coordinate domain. Set
\begin{equation}\label{eq:lift-node-functional}
 L[p(z)]=\mathcal E[p(\widetilde h(x_U))].
\end{equation}
Pullback increases degree by at most $D$, so $L$ is defined through
lifted degree $2r$ and is normalized. This again uses deterministic
substitution followed by marginalization, with no conditioning step.

We verify all localizers, keeping the degree budget explicit. Suppose
$\sum_j\deg g_j+2\deg P\le2r$, where every $g_j$ is a locally valid
univariate polynomial in a lifted coordinate. A factor on a restricted
coordinate becomes a nonnegative constant. For any other coordinate,
Boolean reduction of $g_j$ after pullback equals
\[
 g_j(1)\frac{1+\widetilde h_{k_j}}{2}
 +g_j(-1)\frac{1-\widetilde h_{k_j}}{2},
 \qquad g_j(1),g_j(-1)\ge0,
\]
where $k_j$ is the coordinate used by factor $g_j$; these indices may
repeat. A constant factor needs no indicator and is removed first. Expanding the
nonconstant factors expresses the product as a sum of nonnegative
coefficients times products $I$ of parity indicators. Overlapping supports
cause no difficulty: all these factors commute and are idempotent in the
Boolean quotient, so $I^2=I$, even when $I=0$. With
$a=\sum_j\deg g_j$ and $d=\deg P$,
\begin{equation}\label{eq:lift-degree-accounting}
 \deg I\le Da,\qquad \deg\widehat P\le Dd,
 \qquad\widehat P=P\circ\widetilde h,
 \qquad\deg(I\widehat P)\le D(a+d)\le2rD.
\end{equation}
Consequently $I\widehat P^2=(I\widehat P)^2$ modulo Boolean equations,
within degree $4rD$. Original square positivity proves the localizer.
In particular, this checks every repeated bound-slack product, not just
individual bound multipliers. If a univariate equality vanishes on its
coordinate set, it vanishes at the fixed value or at both permitted
endpoints; its pulled-back Boolean reduction times any allowed multiplier
is zero. All graph identities in \eqref{eq:lift-all-identities} vanish
as literal polynomials after pullback and substitution. Their multiplied
pullbacks have degree at most $2rD$, so every stated equality is respected.

To verify \eqref{eq:lift-box-hull}, form the augmented first/second moment
matrix of $L$. A lifted linear polynomial pulls back to degree at most
$D\le2rD$, so this matrix is positive semidefinite. Every diagonal entry
is one because $\widetilde h_j^2=1$ after Boolean reduction. Every other
entry is an original character moment, multiplied by a fixed sign, and
therefore belongs to $\{0,-1,1\}$. Lemma~\ref{lem:signed-moment-law}
provides a finite Boolean distribution on the lifted coordinates with
exactly this matrix. For $j\in R$, its coordinate mean equals the fixed
sign $h_j(w)$. A Boolean variable with mean $1$ or $-1$ equals that sign
almost surely. All other coordinates have both endpoints in their sets.
Thus the realizing distribution is supported in $B$, including when the
sets are nonclosed. It may leave $\Gamma_h$: the graph moment identities
have already been proved for $L$, and no common graph-supported
distribution realizing its higher moments is asserted.

Finally, \eqref{eq:lift-objective-pullback} gives
$L[\Phi]=\mathcal E[F(w_C,x_U)]$. Only clauses meeting $C$ change their
costs. Their number is at most $\Delta|C|$. As in the unlifted proof,
each substituted clause is $(1-\sigma\chi_A)/2$ with $|A|\le3$.
The squares $(1\pm\chi_A)^2$ have degree at most six; they are available
because $2rD\ge4\ge3$. Therefore all such pseudo-costs lie in $[0,1]$,
and
\begin{equation}\label{eq:lift-cost-restriction}
 \operatorname{LB}_r(B)\le L[\Phi]\le\frac{\Delta|C|}{m}.
\end{equation}
Certification forces $|C|\ge mT_*/\Delta$.

For witness counting, write each original Boolean sign as $w_i=(-1)^{t_i}$
with $t_i\in\mathbb F_2$. Each restricted coordinate imposes one equation
with row $\ind_{A_j}$ and a right-hand side determined by $\sigma_j$ and
its permitted endpoint. Let $A_R$ be this binary row matrix and
$q=\rank_{\mathbb F_2}A_R$. The system is consistent, since it contains
$w$, and has exactly $2^{n-q}$ solutions. Every such solution meets all
endpoint restrictions and all unrestricted sets contain both endpoints,
so the witness fraction of $B$ is exactly $2^{-q}$.

Select $q$ basis rows from the original rows of $A_R$. If a column is zero
in each basis row, it is zero in their entire span, hence zero in every
row. The union of all supports is therefore the union of the basis-row
supports. Each basis support has at most $D$ elements, giving
\begin{equation}\label{eq:parity-rank-support}
 |C|\le Dq,\qquad
 2^{-q}\le2^{-|C|/D}\le2^{-mT_*/(\Delta D)}.
\end{equation}
The union bound over the graph witnesses proves the theorem. Counting
restricted lifted coordinates alone would not prove this bound: repeated
parities can create arbitrarily many restrictions with rank one.
\end{proof}

The proof needs $r\ge1$ and $rD\ge2$, together with availability of
$\Phi$ through degree $2r$. It uses no separate lower bound $r\ge2$.
The unlifted cubic objective still requires $r\ge2$ because it must be
evaluated by the node functional. The all-character assumption is stronger
than quadratic realizability in Theorem~\ref{thm:xor-transfer}; a lifted
pair of coordinates can represent a character on as many as $2D$ original
variables. Lemma~\ref{lem:xor-width-moments} supplies precisely the required
signed moments whenever $4rD\le w$.

\begin{corollary}[A quadratic formulation with a linear objective]
\label{cor:xor-quadratic-formulation}
For every sufficiently large original dimension $n$, the family of
Theorem~\ref{thm:xor-family} has an equivalent continuous formulation
with $N=n+2m\le17n$ coordinates, $2m$ quadratic equalities, and a linear
objective. At every integer order $r\ge1$ with $12r\le an$, certificates
at absolute tolerance $1/16$ or relative tolerance $1/2$ require at least
\begin{equation}\label{eq:quadratic-formulation-count}
 2^{7n/3072}\ge2^{7N/52224}
\end{equation}
coordinate regions under the oracle of Theorem~\ref{thm:monomial-transfer}.
This includes branching on all pair and clause coordinates at order one.
\end{corollary}
\begin{proof}
Order the three distinct variables of each clause as $(i,j,k)$ and introduce
\begin{equation}\label{eq:xor-factorable-formulation}
 u_e=x_ix_j,\qquad v_e=u_ex_k,\qquad
 \Phi(x,u,v)=\frac1m\sum_e\frac{1-b_ev_e}{2}.
\end{equation}
All coordinates have bounds $[-1,1]$. For every original $x$, these
equalities determine exactly $u_e=x_ix_j$ and $v_e=x_ix_jx_k$, which lie
in the stated bounds. Conversely every feasible tuple has those values.
Thus the feasible set is the continuous monomial graph, the objectives
agree as polynomials after substitution, and the degree parameter is
$D=3$. Repeated clauses can have separate auxiliaries; neither support
repetitions nor auxiliary occurrence counts alter the theorem. The source
condition is $4rD=12r\le an$, which allows $r=1$ for all sufficiently
large $n$. Using $T_*\ge1/16$, $m\ge7n$, and $\Delta\le64$ in
\eqref{eq:monomial-transfer-count} gives the first bound. The identity
$N=n+2m\le17n$ gives the second.
\end{proof}

\begin{corollary}[Support degree versus region count]\label{cor:degree-regions}
For the same family, whenever $r\ge1$, $rD\ge2$, and $4rD\le an$,
the fixed-tolerance lower bound for degree-$D$ monomial coordinates is
\begin{equation}\label{eq:degree-region-tradeoff}
 2^{7n/(1024D)}.
\end{equation}
Within this degree regime, a number of regions polynomial in original
dimension $n$ requires $D=\Omega(n/\log n)$. For fixed $D$, the lower
bound is exponential in $n$ even with arbitrarily many auxiliaries and
node order up to a sufficiently small constant multiple of $n$.
\end{corollary}
\begin{proof}
Apply Theorem~\ref{thm:monomial-transfer} with the family parameters. If
the number of regions is at most $Cn^k$ for constants $C,k$, taking
base-two logarithms in \eqref{eq:degree-region-tradeoff} gives
$7n/(1024D)\le\log_2 C+k\log_2 n$. This is the asserted necessary
condition. It is not a sufficient construction and makes no assertion
outside $4rD\le an$.
\end{proof}
For an arbitrarily inflated $N$, the exponent is in the original dimension
$n$, not automatically in $N$; likewise ``polynomial'' in this corollary
means polynomial in $n$. The endpoint interpolation of
Theorem~\ref{thm:local-graph-lift} concerned functions of one original
coordinate at a time. Its absence of auxiliary-degree loss does not
extend by that proof to these multivariable monomials. Here literal
pullback and the parity-rank argument supply a different, explicit
$4rD$ degree budget.

\section{Finite upper certificates and the boundary of affine branching}
\label{sec:xor-upper-affine}
The fixed-tolerance lower bounds have finite upper certificates of the
same exponential order in dimension. We give two certificates because
they use different information. The Bernstein proof uses the order-two
preordering and polynomial graph identities. A separate order-one proof
for the quadratic formulation uses the exact quadratic hull of the box.

\begin{proposition}[Order-two Bernstein certificates]
\label{prop:xor-bernstein-upper}
For every objective \eqref{eq:xor-objective} and $\epsilon>0$, at most
$\lceil3/\epsilon\rceil^n$ boxes certify the target $\OPT-\epsilon$
at order two. This holds both in the original variables and in the
quadratic formulation \eqref{eq:xor-factorable-formulation}. In the latter,
only original coordinates need subdivision. The quadratic moment hull
is unnecessary for this upper bound.
\end{proposition}
\begin{proof}
Set $M=\lceil3/\epsilon\rceil$ and partition each original interval
into $M$ equal intervals, of width $h=2/M\le2\epsilon/3$. Fix a resulting
box $B=\prod_i[a_i,a_i+h]$. Every clause cost has derivatives of
magnitude at most $1/2$ in each of its three coordinates. Moving between
points of its three-coordinate box one coordinate at a time changes its
value by at most $3h/2\le\epsilon$. Let $\beta_e$ be its minimum on
that box. In particular
\begin{equation}\label{eq:xor-grid-corner-bound}
 \frac1m\sum_e\beta_e\ge F(a)-3h/2\ge\OPT-\epsilon.
\end{equation}
For $i\in e$, write $\lambda_i^1=(x_i-a_i)/h$ and
$\lambda_i^0=(a_i+h-x_i)/h$. Multiaffine interpolation is the exact
polynomial identity
\begin{equation}\label{eq:xor-bernstein-identity}
 c_e(x)-\beta_e=
 \sum_{s\in\{0,1\}^e}\bigl(c_e(a_e+hs)-\beta_e\bigr)
       \prod_{i\in e}\lambda_i^{s_i}(x_i).
\end{equation}
Each coefficient is nonnegative, and each basis product is a positive
multiple of three coordinate-bound slacks. The minimum $\beta_e$ is
attained at a corner by multiaffinity. Thus the node preordering through
degree four certifies $L[c_e]\ge\beta_e$. Averaging and
\eqref{eq:xor-grid-corner-bound} give the original upper certificate.

For the lifted formulation, keep every auxiliary interval $[-1,1]$.
The exact identity
\begin{equation}\label{eq:bernstein-graph-transfer}
 v_e-x_ix_jx_k=(v_e-u_ex_k)+x_k(u_e-x_ix_j)
\end{equation}
has graph multipliers of total degree at most three. Order two therefore
implies $L[v_e]=L[x_ix_jx_k]$. The same Bernstein certificate in the
original bound slacks proves $L[\Phi]\ge\OPT-\epsilon$. There are
$M^n$ lifted boxes, and together they cover the full feasible graph.
\end{proof}

\begin{proposition}[Order-one certificates from the quadratic box hull]
\label{prop:xor-order-one-upper}
The quadratic formulation \eqref{eq:xor-factorable-formulation} admits
an order-one certificate at target $\OPT-\epsilon$ with at most
$\lceil3/\epsilon\rceil^n$ lifted boxes under the exact quadratic-box-hull
oracle, for every $\epsilon>0$. This certificate uses only the original
coordinate subdivisions and the two graph equations per clause imposed
on moments.
\end{proposition}
\begin{proof}
Use the same grid, with original lower corner $a$ and width $h=2/M$,
and leave auxiliaries in $[-1,1]$. Any feasible order-one functional
has its first and second moments realized by an actual distribution $\mu$
on this entire lifted box. The graph moment equations give
\begin{equation}\label{eq:order-one-graph-moments}
 L[u_e]=L[x_ix_j],\qquad L[v_e]=L[u_ex_k].
\end{equation}
We use these equalities only in expectation; $\mu$ need not satisfy the
graph equations pointwise. On the full box,
\[
 |x_ix_j-a_ia_j|
 \le |x_i-a_i|\,|x_j|+|a_i|\,|x_j-a_j|\le2h,
 \qquad |u_e(x_k-a_k)|\le h.
\]
Taking expectations, all of which are degree at most two, yields
\begin{align}
 |L[v_e]-a_ia_ja_k|
 &\le |L[u_ex_k]-a_kL[u_e]|
          +|a_k|\,|L[u_e]-a_ia_j|\nonumber\\
 &\le h+2h=3h.\label{eq:order-one-corner-estimate}
\end{align}
For either clause sign this gives
$L[(1-b_ev_e)/2]\ge c_e(a)-3h/2$.
Averaging proves $L[\Phi]\ge F(a)-3h/2\ge\OPT-\epsilon$.
The inequality holds for every feasible node functional and hence for
its infimum. The $M^n$ boxes cover the graph.

There is also an explicit quadratic inequality certifying each step. Set
\begin{equation}\label{eq:order-one-quadratic-cut}
 q_e=\frac{3h}{2}-\frac{b_e}{2}
       \bigl[u_e(x_k-a_k)+a_k(x_ix_j-a_ia_j)\bigr].
\end{equation}
The preceding pointwise bounds give $q_e\ge0$ on the \emph{full} lifted
box. Moreover,
\begin{equation}\label{eq:order-one-cut-identity}
 \frac{1-b_ev_e}{2}-c_e(a)+\frac{3h}{2}
 =q_e-\frac{b_e}{2}
       \bigl[(v_e-u_ex_k)+a_k(u_e-x_ix_j)\bigr].
\end{equation}
All polynomials in this identity have lifted degree at most two. The
quadratic hull gives $L[q_e]\ge0$, and the graph moment equations kill
the two remaining terms. Thus the proof is a certificate using a valid
quadratic box inequality linearly, without any higher-degree localizer
or graph-supported realizing distribution.
\end{proof}

At $\epsilon=1/16$, both propositions give $48^n$ regions. For the
family with $\OPT\ge1/8$, these also certify relative tolerance $1/2$
with the best incumbent: $\OPT-1/16\ge\OPT/2$. Consequently the
quadratic formulation has both an order-one lower bound
$2^{7n/3072}$ and an order-one upper bound $48^n$, for all sufficiently
large $n$. These counts measure regions, with no claim that the exact
quadratic-box-hull oracle is computationally tractable. For the original
cubic objective the distinct order-two Bernstein certificate remains the
applicable statement.

\subsection{Two obstructions to the same proof under affine branching}
The parity-rank count depends on restricted coordinates being signed
products with bounded support. A balanced affine halfspace behaves
differently even under an actual probability distribution.

\begin{proposition}[A halfspace rejects the preserved quadratic moments]
\label{prop:affine-quadratic-obstruction}
For every $n\ge1$, let $\mathcal E$ be uniform expectation on
$\pmcube$ and put $S(x)=\sum_i x_i$. The first and second moments of
$\mathcal E$ have no actual realization on
\[
 H=\{x\in[-1,1]^n:S(x)\ge0\},
\]
although $H$ contains at least half the Boolean witnesses. Equivalently,
adding the affine coordinate $z=S/n$ and restricting it to $[0,1]$
rejects the preserved moments by the valid quadratic inequality $z^2\le z$.
\end{proposition}
\begin{proof}
The uniform moments give $\mathcal E[S]=0$ and
$\mathcal E[S^2]=n$. A distribution supported where $S\ge0$ with
mean $S=0$ must have $S=0$ almost surely, contradicting that second
moment. In the affine coordinate the same contradiction is
$\mathcal E[z]=0$ and $\mathcal E[z^2]=1/n>0$.
The involution $w\mapsto-w$ proves $\Prob\{S(w)\ge0\}\ge1/2$;
for even $n$ the boundary can add mass.
\end{proof}

\begin{proposition}[Substitution cost for a halfspace localizer]
\label{prop:affine-substitution-obstruction}
Start from the same uniform Boolean expectation, fix an arbitrary
coordinate set $C$ to arbitrary signs $w_C$, and leave the others
independently uniform. If the resulting functional $L$ satisfies
\[
 L[SP^2]\ge0\qquad(1+2\deg P\le2r)
\]
for an integer $r\ge1$, then necessarily
\begin{equation}\label{eq:affine-substitution-count}
 |C|\ge\min\{r-1,\lceil n/2\rceil\}.
\end{equation}
\end{proposition}
\begin{proof}
Write $t=|C|$, $a=\sum_{i\in C}w_i$, and $s=n-t$. Suppose both
$t<r-1$ and $t<\lceil n/2\rceil$. Then $t\le r-2$ and $s\ge t+1$.
If $a<0$, $P=1$ already violates the localizer. Otherwise $a$ is an
integer in $[0,t]$. Select $a+1$ unrestricted coordinates, which is
possible because $a+1\le t+1\le s$, and set
\[
 P(x)=\prod_{i\text{ selected}}\frac{1-x_i}{2}.
\]
Its degree is $a+1\le r-1$, so the localizer has permitted total degree.
On its indicator event the selected variables sum to $-(a+1)$, all
other unrestricted variables have conditional mean zero, and the fixed
ones sum to $a$. Boolean idempotence therefore gives the exact value
\[
 L[SP^2]=(a-(a+1))\,2^{-(a+1)}=-2^{-(a+1)}<0.
\]
This is the contradiction. For $r=1$ the conclusion is the stated trivial
bound $t\ge0$.
\end{proof}

At order proportional to $n$, this halfspace requires linearly many
substitutions in the displayed construction while retaining at least half
the witnesses. Thus large substitution cost cannot itself imply
exponentially small witness mass for general affine inequalities.
These propositions obstruct this proof method. They do not prove that a
different node functional is impossible, or that the constant-gap XOR
objective has a short affine-branch certificate.

\subsection{Comparison with other disjunctive certificate models}
\label{subsec:xor-literature-boundary}
Disjunctive sum-of-squares certificates can keep polynomial degree fixed
while increasing the number of regions. Ahmadi, Dash, Hua, and Stellato
\citep[Sections 2--4 and 6]{AhmadiEtAl2026} develop such certificates
and a spatial framework on simplicial cones intersected with the sphere,
with a corresponding simplex method for matrix copositivity. Their
Section~6 states termination of the algorithms at positive prescribed
tolerance. A convergence result of that kind is compatible with
exponentially many regions. The coordinate-product bounds here do not
apply to all simplicial or algebraic disjunctions in that framework.

Stabbing Planes instead branches on an integer linear inequality
$ax\ge b$ and its integer negation $ax\le b-1$. It may discard the
intervening slab because the slab has no integer points
\citep[Section 1]{BeameEtAl2023}. Continuous subdivision must account
for those points and therefore has a different covering requirement.
Beame et al.\ give quasipolynomial-size Stabbing Planes refutations of
Tseitin formulas \citep[Theorem 1]{BeameEtAl2023}. Fleming et al.\ give
quasipolynomial-size Cutting Planes refutations for CNF encodings of
inconsistent finite-field linear systems, and size lower bounds for a
coefficient-restricted Stabbing Planes subsystem
\citep[Theorems 1.3--1.4]{FlemingEtAl2021}. These source statements do
not transfer a parity SOS obstruction directly to unrestricted affine
branching, nor do they give a constant-gap certificate for our objective.
A refutation of perfect Boolean satisfiability proves only that some
clause fails. Its direct normalized objective consequence is $F\ge1/m$
on Boolean vertices, and hence on the cube by multiaffinity; it does not
establish $F\ge1/16$ when $m$ grows. Proving or disproving comparable
constant-gap region bounds for a stronger affine-branching model requires
an additional argument.

\section{Representation and other computational resources}\label{sec:comparisons}
The preceding lower bounds specify what a node can certify. Other forms
of difficulty require different statements. Global lift size counts cone
factors or matrix order; integer precision counts discrete coordinates;
bound propagation studies a limit or a sequence of local updates. None
of these counts can be substituted for certified spatial regions.

Several short geometric comparisons also show why the representation must
be fixed. Appendix~\ref{app:point-packing} derives the four established
point-packing RLT/SDP values, including symmetry-adjusted secants and
attaining covariances. This benchmark illustrates the gain from moment
consistency and the care needed when the underlying box changes.
Appendix~\ref{app:scaling} gives exact scale and multiplicity hulls by
classical disaggregation, with a retained common intensive variable.
It separates this hull identity from cost separation, which requires
additional structure. Appendix~\ref{app:p-split} gives a retained-domain
counterexample to the universal nonexactness assertion in the published
P-split theorem, a sufficient directional repair, and exact translated-ball
projection and distance formulas. A rational change of coordinates can
turn a growing gap into a bounded one; convexifying the exact auxiliary
image can then be exact in aligned coordinates. These claims concern the
stated epigraph model and retained domain.

Appendix~\ref{app:rank-one} exhibits a stable correlation-polytope face
inside the unit-capacity rank-one hull. Established correlation lower
bounds transfer to exact conic lifts, while the explicit stability estimate
transfers fine entrywise-$\ell_1$ approximations at order $m^{-2}$.
The fixed-block exponential bound and the unrestricted PSD
superpolynomial bound are separate conclusions. Appendix~\ref{app:fbbt}
transfers least monotone fixed points to fair forward FBBT, proves
PosSLP-hard constant-error approximation of the limiting lower bound,
and gives a doubly exponential primitive-update example valid under every
fair ordering. Its primitive contractor restriction is essential.
Finally, Appendix~\ref{app:integer-precision} contrasts scalar and
simultaneous-bilinear approximation accuracy with spatial certification.
It records short midpoint and cover proofs, without extending the present
paper to the broader theory of mixed-integer representability.

\section{Consequences and open questions}\label{sec:synthesis}
The central distinction is whether locally admissible representations can
be made simultaneous. For envelope comparisons this is a probability law
with all the prescribed means. Positive coefficients permit common upper
attainment and reduce the comparison to lower-envelope rounding; signs
instead lead to cut ranges. Structural restrictions improve the available
couplings, while positive physical lower bounds change the coefficients
and may change the relevant incidence structure. In spatial certification,
a pseudoexpectation can satisfy a node's prescribed moments and polynomial
constraints without supplying a feasible graph distribution. Counting
regions that preserve it requires both a degree calculation and a witness
cover argument.

Several bounded refinements sharpen this picture. The fixed-ambient
balanced orientation improves the positive-box upper bound by replacing
two with $\beta_N<2$, and the coefficient-regularity cardinality argument
uses the same law under its stated coefficient ratios. The general-radix
cutoff construction gives exact attainment beyond the simpler large-radix
regime. Coordinatewise graph lifts admit endpoint interpolation without
polynomial-degree loss, including everywhere-finite nonpolynomial
coordinate functions. Bounded signed-monomial lifts admit the lower-bound
transfer at order $r\ge1$ when $rD\ge2$, and the quadratic formulation has
a separate order-one finite upper certificate. The unlifted cubic still
requires $r\ge2$. These completions retain their domain, oracle and
full-graph agreement hypotheses.

The remaining questions are quantitative or concern stronger methods.
The exact cubic constant lies between the certified lower value
$1610000/743033$ and $31/12$; optimality of the three-law mixture does not
determine it. The optimized harmonic upper asymptotic has no matching
second-order lower term. For a fixed physical aspect ratio the interval
$\max\{2,\rho_B\}\le C_{\rm box}(\rho_B)\le\rho_B+2$ is not closed by
the large-aspect leading constant. Exact structural constants for incidence
treewidth at least three remain open; the width-two coloring proof does
not extend merely by increasing the number of colors. On unequal-aspect
positive boxes even bipartite frequency-two exactness fails, and the
bipartite supremum remains between $3/2$ and two.

For spatial methods, arbitrary affine branching, additional coupled-cut
localizers, and convex hulls of feasible graphs exceed the proved oracle
scope. The balanced-halfspace calculation identifies an obstruction to the
preserved-moment proof under affine branching; it provides neither a fast
affine tree nor a lower bound for all such trees. The examples therefore
suggest precise questions about stronger local information, rather than
a conclusion about every global optimization algorithm.

\appendix
\section{Exhaustive verification of the finite signing constants}
\label{app:finite-signings}
This appendix gives the finite computational proof used in
Proposition~\ref{prop:finite-signings}. Label vertices $0,\ldots,n-1$.
For vertex signs $t_i$, switching replaces $a_{ij}$ by
$a'_{ij}=a_{ij}t_it_j$. The identity $Q_{a'}(s)=Q_a(t\cdot s)$ shows
that switching preserves both extrema of $Q$, and hence $R_V$.
Taking $t_0=1$ and $t_j=a_{0j}$ makes every coefficient incident to zero
positive. This normalized representative is unique within each switching
class: a switching that keeps those coefficients positive has $t_j=t_0$
for all $j$, so changes no edge. Thus it suffices to enumerate the
$2^{\binom{n-1}{2}}$ choices of negative edges among vertices $1,\ldots,n-1$.
There are respectively $2,8,64,1024,32768$ representatives for $n=3,\ldots,7$.

Encode a cut by the subset on side one, always leaving vertex zero on
side zero. Complementation changes no cut, so the $2^{n-1}$ such subsets
include every cut. If a cut has $c$ crossing edges, of which $h$ are
negative, its weight is exactly $c-2h$. The following complete Python
3.10+ program computes the first row of Proposition~\ref{prop:finite-signings}
using unbounded integers. The operation \texttt{bit\_count()} counts set
bits in an integer; each free-edge mask records which edges cross the cut.
No floating-point arithmetic or solver is used.

\begin{small}
\begin{verbatim}
from itertools import combinations
def min_range(n):
    edges = list(combinations(range(n), 2))
    free = [(i, j) for i, j in edges if i != 0]
    cuts = []
    for side in range(1 << (n - 1)):
        side <<= 1  # Vertex zero stays on side zero.
        cross = lambda i, j: ((side >> i) ^ (side >> j)) & 1
        size = sum(cross(i, j) for i, j in edges)
        mask = sum(cross(i, j) << k
                   for k, (i, j) in enumerate(free))
        cuts.append((size, mask))
    best = len(edges) + 1
    for negative in range(1 << len(free)):
        weights = [size - 2 * (negative & mask).bit_count()
                   for size, mask in cuts]
        best = min(best, max(weights) - min(weights))
    return best
result = [min_range(n) for n in range(2, 8)]
assert result == [1, 2, 4, 4, 5, 8]
print(result)
\end{verbatim}
\end{small}

Every signed cut weight lies between $-\binom n2$ and $\binom n2$;
its range is at most $\binom n2$, because two cut weights differ by a
sum of at most that many coefficients in $\{-1,1\}$. The initial value
\texttt{len(edges) + 1} therefore exceeds every candidate range. The two
loops cover all normalized signings and all cuts, so their minimum is
exactly $\min_a R_V(a)$, rather than merely an upper bound supplied by
the best witness found. This is an exhaustive finite proof, not an
analytic formula for general $n$.

\begin{table}[htbp]
\centering
\begin{tabular}{clrr}
\toprule
$n$ & Negative edges; all other edges have coefficient $+1$
    & $\min Q$ & $\max Q$\\
\midrule
3 & $\varnothing$ & $-1$ & $3$\\
4 & $\varnothing$ & $-2$ & $6$\\
5 & $12,14,23$ & $-4$ & $4$\\
6 & $14,15,23,25,34$ & $-5$ & $5$\\
7 & $12,13,16,23,25,34$ & $-7$ & $9$\\
\bottomrule
\end{tabular}
\caption{Attaining full-center signings, with vertices numbered from zero.
An entry $ij$ denotes the unordered pair $\{i,j\}$.}
\label{tab:finite-witnesses}
\end{table}

The extrema in Table~\ref{tab:finite-witnesses} can be checked by direct
evaluation of $Q$ on the $2^{n-1}$ vertex-sign vectors with $s_0=1$.
They give $R_V=(\max Q-\min Q)/2$ equal to the computed minimum in each
case. Extending the listed $K_6$ witness by vertex six and positive new
edges gives $\min Q=-9$, $\max Q=11$ at the full center, hence ratio
$21/10$; its original six-vertex face retains ratio three.

The accompanying script \path{verification/check_complete_signings.py}
implements the same integer enumeration, records a histogram of all
candidate ranges, and checks the displayed witnesses by direct evaluation
of $Q$ on every induced face. Running it with Python 3.10+ writes
\path{verification/complete_signings.json}, including the witnesses,
counts, exact rational ratios, and the extended six-vertex example.

\section{Distinct multiscale arguments and finite refinements}
\label{app:dyadic-alternatives}
\subsection{Arbitrary partitions and the dense predecessor}
The exact dyadic formula in Theorem~\ref{thm:dyadic-exact} relies on a
simultaneous prefix-hitting construction. A separate logarithmic estimate
works for arbitrary partitions in \eqref{eq:dyadic-polynomial}.

\begin{proposition}[Partition-independent estimate]
\label{prop:arbitrary-dyadic}
For the polynomial and means in \eqref{eq:dyadic-polynomial}, with any
choice of equal-size level partitions,
\[
 0<H\le5+\log_2(1+(L-1)\log2),\qquad T=L.
\]
\end{proposition}
\begin{proof}
Use the quantile surrogate from \eqref{eq:dyadic-surrogate}, writing
$r(t)$ for the nonincreasing failure-count quantile and $\delta=2^{-L}$.
The interval $0<t\le\delta$ contributes at most
$\delta\sum_{j=1}^L2^j<2$, and $t>1/2$ contributes zero. For
$\delta<t\le1/2$, let $l$ be the largest integer with $t\le2^{-l}$.
Write $\log_2^+v=\max\{0,\log_2v\}$ for $v>0$.
Then $2^l\le1/t$, and for $r>0$,
\begin{equation}\label{eq:dyadic-log-sum}
 \sum_{j=1}^l\min(2^j,r)
 \le r\left(3+\log_2^+\frac1{tr}\right).
\end{equation}
If $r\ge2^l$, the sum is less than $2r$. If $r<2$, it is $lr$ and
$l\le1+\log_2(2^l/r)$. Otherwise, with $h=\lfloor\log_2r\rfloor$,
the first $h$ capped terms sum to less than $2r$ and the others sum to
$(l-h)r$; use $h>\log_2r-1$. This proves
\eqref{eq:dyadic-log-sum}, interpreted as zero at $r=0$.

Put $D=(\delta,1/2)$ and $M_0=\int_Dr(t)\,dt\le1$. As
$\log_2^+v\le\log(1+v)/\log2$, the bound becomes
\[
 H\le2+3M_0+\frac1{\log2}
       \int_Dr(t)\log\left(1+\frac1{tr(t)}\right)\,dt.
\]
For $M_0>0$, Jensen's inequality under the probability measure
$r(t)\,dt/M_0$ gives
\[
 \int_Dr(t)\log\left(1+\frac1{tr(t)}\right)\,dt
 \le M_0\log\left(1+\frac{(L-1)\log2}{M_0}\right)
 \le\log(1+(L-1)\log2).
\]
The final step uses that $M\log(1+B/M)$ increases for $M>0$, since
$\log(1+s)\ge s/(1+s)$. At $M_0=0$ the integral is zero. This proves
the upper bound. Independent rounding and common-threshold rounding
have distinct expected polynomial values at these interior means, so $H>0$.
\end{proof}

A dense symmetric construction gives a different exact reduction.
With $m=2^L$, $u_j=2^{-j}$, $k_j=mu_j$, define
\begin{equation}\label{eq:dense-predecessor}
 \widetilde f_L(a,z)=\sum_{j=1}^L\frac1{u_j\binom m{k_j}}
             \sum_{|K|=k_j}a_j\prod_{i\in K}z_i,
 \qquad a_j=u_j,\quad z_i=1-1/m.
\end{equation}
Again $T=\cav\widetilde f_L=L$. Conditional on $r$ failed leaves, the
fraction of size-$k_j$ subsets hit is
\[
 g_j(r)=1-\frac{\binom{m-r}{k_j}}{\binom m{k_j}}
       \le\min(1,u_jr).
\]
The inequality is the union bound: a uniform size-$k_j$ subset contains
each specified failed leaf with probability $k_j/m=u_j$.
Thus the deficiency is $\sum_j A_jg_j(R)/u_j$, bounded by the same
surrogate $\sum_j A_j\min(2^j,R)$ and consequently by
Proposition~\ref{prop:arbitrary-dyadic}.

In this dense model the hull gap itself is exactly the optimum of the
finite rational program
\begin{equation}\label{eq:dense-count-lp}
 \begin{aligned}
 \max\quad&\sum_{j=1}^L\sum_{r=0}^m z_{jr}g_j(r)/u_j\\
 \text{subject to}\quad& p_r\ge0,\quad 0\le z_{jr}\le p_r,\\
 &\sum_{r=0}^mp_r=1,\quad \sum_{r=0}^mrp_r=1,
       \quad\sum_{r=0}^mz_{jr}=u_j\quad(j=1,\ldots,L).
 \end{aligned}
\end{equation}
Every feasible binary law supplies $p_r=\Prob(R=r)$ and
$z_{jr}=\Prob(A_j=1,R=r)$. Conversely draw $R$ with probabilities $p_r$,
choose a uniform $r$-subset of failed leaves, and, conditional on $R=r$,
choose anchors independently with probabilities $z_{jr}/p_r$ when
$p_r>0$. These choices realize all means and the objective; zero-mass
states can be ignored. Thus \eqref{eq:dense-count-lp} is an exact reduction,
not a formula for the sparse model. No floating-point values from this
program are needed for the divergent bound.

\subsection{Dyadic conditional-failure rounding and its tail lemma}
\label{app:dyadic-rounding}
The following proof uses finitely many dyadic scales rather than a
harmonic kernel. Though its $O(\log d)$ bound is weaker, its tail
argument is distinct.

\begin{lemma}[Tail integral]\label{lem:dyadic-tail}
For a finite collection $p_j\ge0$, let $N(s)=|\{j:p_j\ge s\}|$ and
$S=\sum_jp_j$. For $u>0$,
\begin{equation}\label{eq:dyadic-tail}
 \int_0^u\min\{N(s),u/s\}\,ds\ge\min(u,S).
\end{equation}
\end{lemma}
\begin{proof}
If $S\le u$, then every $p_j\le u$ and $\int_0^uN(s)\,ds=S$.
Monotonicity gives $sN(s)\le\int_0^sN(t)\,dt\le S\le u$, so the
integrand equals $N(s)$ throughout. If $S>u$, then
$\int_0^uN(s)\,ds\ge u$: either a $p_j\ge u$, or that integral equals
$S$. By continuity choose $v\le u$ with $\int_0^vN(s)\,ds=u$.
For $s\le v$ the same monotonicity bound gives $sN(s)\le u$.
Integrating on $(0,v)$ therefore gives $u$.
\end{proof}

\begin{proposition}[Dyadic degree bound]\label{prop:dyadic-degree}
Let $K=1+\lfloor\log_2(d-1)\rfloor$ for $d\ge2$. Every positive
multilinear polynomial of maximum degree at most $d$ on a nonnegative
box satisfies $T\le2(K+1)H/c_0$.
\end{proposition}
\begin{proof}
For each scale $B\in\{1,2,\ldots,2^{K-1}\}$ restricted to powers of
two, draw a common uniform $U$. Low coordinates succeed when $U\le x_i$.
A high coordinate with failure mean $p>0$ has $h=\min(Bp,1)$ and,
conditional on $U$, fails independently with probability $p/h$ when
$U\le h$, and never otherwise. Coordinates with $p=0$ never fail.
The marginal failure is $h(p/h)=p$ and the probabilities are valid
because $B\ge1$.

For a one-low term with anchor mean $u>0$, let $p_1,\ldots,p_r$ be
its high failure means. At $U=t\le u$, each $p_j\ge t/B$ is active
and has conditional failure probability at least $1/B$. Conditional
independence and $1-e^{-z}\ge c_0\min(1,z)$ give deficiency
\[
 G_B\ge c_0\int_0^u\min(1,N(t/B)/B)\,dt
       =c_0\int_0^{u/B}\min(B,N(s))\,ds.
\]
For $0<s\le u$ with $N(s)>0$, choose the largest power of two
$B\le\min(N(s),u/s)$. It is in the allowed set, since
$N(s)\le r\le d-1$, and exceeds half this minimum. Therefore
\[
 \sum_B\ind\{s\le u/B\}\min(B,N(s))
       \ge\tfrac12\min(N(s),u/s).
\]
Integrating and invoking Lemma~\ref{lem:dyadic-tail} yields
$\sum_BG_B\ge c_0\min(u,S)/2=c_0t_e/2$.
Terms with $u=0$ have zero gap. Independence supplies the same guarantee
for terms with zero or at least two low coordinates by
Lemma~\ref{lem:easy-terms}. The uniform mixture of the $K$ scale laws
and independence thus guarantees $c_0t_e/(2(K+1))$ for every term.
Sum and apply Proposition~\ref{prop:box-transfer}.
\end{proof}

\subsection{The simplest harmonic cutoff and its reciprocal correction}
\label{app:reciprocal-cutoff}
The general finite theorem permits $M=d-1$, so $\eta=1$ and
$\Lambda=1+\log(d-1)$. Let $z_\Lambda$ solve
$J_{\Lambda,1}(z_\Lambda)=z_\Lambda$ and put
$\alpha=\Lambda z_\Lambda$, $\ell=\log\Lambda$. For $\ell\ge4$,
\begin{equation}\label{eq:rho-one-finite}
 z_\Lambda\ge\frac{\ell-\log\ell-2}{\Lambda}.
\end{equation}
Indeed $A=\ell-\log\ell-2\ge1/2$ and $A\le\ell<\Lambda$.
Equation~\eqref{eq:lambert-integral-certificate} with $\eta=1$ gives
$\Lambda J(A/\Lambda)\ge\ell-\log\ell-1-1/(2A)\ge A$.
Monotonicity of $J(z)-z$ proves the finite bound.

The integrated linear upper bound and quadratic lower bound imply
\[
 \ell-\log\alpha-1-\frac1{2\alpha}
 \le\alpha\le\ell-\log\alpha-1+\frac\alpha\Lambda.
\]
Together with \eqref{eq:rho-one-finite}, these show $\alpha/\ell\to1$:
once $\alpha\ge1$, the upper bound gives
$\alpha\le\ell/(1-1/\Lambda)$, while the finite lower bound is
$\ell-\log\ell-2$. Substitution yields
\begin{equation}\label{eq:rho-one-asymptotic}
 \Lambda z_\Lambda=\log\Lambda-\log\log\Lambda-1+o(1).
\end{equation}
The constant $-1$ belongs to this fixed cutoff, whereas the tuned
cutoff in \eqref{eq:second-order-upper} removes that constant after
expressing the bound with numerator $N$.

There is also a reciprocal refinement. If $w e^w=\Lambda/e$, then
\begin{equation}\label{eq:rho-one-reciprocal}
 \alpha=w-\frac1{2w}+O(w^{-2})\qquad(\Lambda\to\infty).
\end{equation}
To prove it, integrate the cubic Taylor bound for $1-e^{-x}$ with
$x=(1-v)/(\Lambda v)$. After multiplication by $\Lambda$, the error
beyond the quadratic term is at most
\[
 \Lambda\int_{\alpha/\Lambda}^1\frac{x^3}{6}\,dv
 \le\frac1{12\alpha^2}.
\]
The exact quadratic integral is
$\int_z^1(1/v-1)^2\,dv=1/z+2\log z-z$.
Its omitted terms, and the linear endpoint term, are
$O(\log\Lambda/\Lambda)=o(\alpha^{-2})$ because $\alpha\sim\ell$.
Consequently
\[
 \alpha+\log\alpha=\log\Lambda-1-\frac1{2\alpha}
                       +O(\alpha^{-2}).
\]
Compare this with $w+\log w=\log\Lambda-1$. Since $\alpha\sim w$,
the mean value theorem for $t+\log t$ first gives
$\alpha-w=O(1/w)$ and, after substitution, gives
\eqref{eq:rho-one-reciprocal}. This refines the scalar certificate only;
it is not a second-order lower bound on $R_d$.

\subsection{A second rational two-level parameter choice}
\label{app:two-level-variant}
A related finite variant illustrates that the two-level mechanism is not
tied to the particular coefficients in Theorem~\ref{thm:cubic-two-level}.
For real $a,c\in[0,1]$, direct expansion gives
\begin{equation}\label{eq:two-level-variant-scalar}
\begin{aligned}
 ac^2+\frac{24}{25}a^2
 -\left(\frac{41}{25}a+\frac45c-\frac{68}{75}\right)
 &=\frac{24}{25}\left(a-\frac{41-25c^2}{48}\right)^2\\
 &\quad+\frac{(1-c)(5c-3)^2(5c+11)}{480}.
\end{aligned}
\end{equation}
Both terms on the right are nonnegative. The equality atoms
$(1/3,1)$ and $(2/3,3/5)$, each with probability $1/2$, have mean
$(1/2,4/5)$ and scalar envelope value $83/150$.
For $25\mid m$, use groups of $m$ variables with these two means and
$f_m=A E_2(W)+(24m/25)E_2(U)$. At binary normalized counts,
\[
 \frac{2f_m}{m^3}=ac^2+\frac{24}{25}a^2
                   -\frac{ac+(24/25)a}{m}.
\]
The expected correction is at most $49/(50m)$ because $ac\le a$.
The scaled concave envelope and termwise gap are, respectively,
$(49/50)(1-1/m)$ and $(22/25)(1-1/m)$; cubic terms have lower envelope
$1/10$, while quadratic terms have lower envelope zero. The affine
minorant gives scaled hull gap at most $49/50-83/150=32/75$.
Conditional independent sampling from the equality atoms, as in
Theorem~\ref{thm:cubic-two-level}, gives scaled hull gap at least
$(32/75)(1-1/m)$. Thus
\[
 \frac{33}{16}(1-1/m)\le\frac{T(f_m)}{H(f_m)}\le\frac{33}{16}.
\]
At $m=50$ the lower certificate is $1617/800>2$. This variant has a
smaller limiting ratio than the two-level family in the main text.

\section{Exact finite cubic certificates}\label{app:cubic-certificates}
The following certificates use success counts, so their verification does
not enumerate all $2^{192}$ binary vectors. Within each group, every
binary vector with the same count has the same polynomial value. The
finite enumeration therefore covers every vertex, including every
possible dependence pattern relevant to the envelope.

For the six-term form \eqref{eq:cubic-six-family}, let
$\Phi_m(A,B,C)$ be its binary count polynomial as defined in the proof of
Proposition~\ref{prop:cubic-finite}. Table~\ref{tab:cubic-duals} gives
integers $D,b_A,b_B,b_C,b_0$ such that
\begin{equation}\label{eq:cubic-integer-dual}
 D\Phi_m(A,B,C)\ge b_AA+b_BB+b_CC+b_0
       \quad(0\le A,B,C\le m\text{ integers}).
\end{equation}
Table~\ref{tab:cubic-primals} gives attaining laws. Each listed weight is
nonnegative, the weights sum to one, and their mean count vector is
$(m/4,m/2,3m/4)$. Uniform success subsets conditional on these counts
therefore give the required singleton marginals. The checker below
verifies the inequality at every integer count triple and equality of
primal and dual expectations in rational arithmetic.

\begin{table}[htbp]
\centering
\begin{tabular}{rrrrrr}
\toprule
$m$&$D$&$b_A$&$b_B$&$b_C$&$b_0$\\\midrule
6&13&962&778&842&$-4816$\\
8&7&793&770&798&$-5882$\\
64&105&871710&900446&899046&$-51743768$\\\bottomrule
\end{tabular}
\caption{Integer affine minorants for the three-group witnesses.}
\label{tab:cubic-duals}
\end{table}

\begin{table}[htbp]
\centering
\begin{tabular}{rll}
\toprule
$m$&Count states $(A,B,C)$&Corresponding probabilities\\\midrule
6&$(1,4,4),(2,1,6),(6,2,1)$&$17/26,\ 4/13,\ 1/26$\\
8&$(1,2,8),(1,5,6),(8,4,2)$&$2/7,\ 4/7,\ 1/7$\\
64&$(4,19,64),(5,19,64),(16,40,43),(64,31,17)$
 &$241/735,\ 12/735,\ 419/735,\ 63/735$\\\bottomrule
\end{tabular}
\caption{Primal laws attaining the affine bounds.}
\label{tab:cubic-primals}
\end{table}

For $m=8$, minimizing the residual in
\eqref{eq:cubic-integer-dual} over $A,B$ for each $C=0,\ldots,8$ gives
\[
 168,\ 42,\ 0,\ 30,\ 22,\ 0,\ 0,\ 16,\ 0.
\]
The three primal states have values $405,507,734$ respectively, with
expected value $3572/7$. In all three examples the orbit sizes, upper
values per unweighted monomial, and lower values per unweighted monomial
are respectively
\begin{align*}
 s&=\left(\binom m3,m\binom m2,\binom m2,m^2,m^2,\binom m2\right),\\
 u&=(3/4,1/2,1/2,1/4,1/4,1/4),\qquad
 l=(1/4,0,0,0,0,0).
\end{align*}
Thus $\cav f=\sum_ic_is_iu_i$ and $\sum_e\vex f_e=\sum_ic_is_il_i$.
These formulas prove the remaining columns of
Table~\ref{tab:cubic-finite}. The exact hull gaps for dimensions
18, 24, 192 are $10411/52$, $3225/7$, $27562048/105$, respectively.

For the two-level $m=16$ certificate
\eqref{eq:two-level-16-dual}, the integer residual is
\[
 2\left[A\binom C2+20\binom A2\right]-428A-177C+3372.
\]
Its minimum over $C=0,\ldots,16$ for successive $A=0,\ldots,16$ is
\[
\begin{gathered}
 540,\ 352,\ 204,\ 96,\ 28,\ 0,\ 9,\ 7,\ 0,\\
 0,\ 19,\ 63,\ 132,\ 235,\ 369,\ 540,\ 742.
\end{gathered}
\]
At counts $(8,12)$ the value and its affine bound both equal 1088.

For the smaller members of \eqref{eq:two-level-family}, write
$\Phi_m(A,C)=A\binom C2+(5m/4)\binom A2$.
Table~\ref{tab:two-level-small} gives affine minorants valid at every
integer pair $0\le A,C\le m$. At the listed counts each minorant equals
$\Phi_m$. Independently choosing uniform success subsets of those fixed
sizes in the two groups gives every singleton mean $1/2,3/4$ and attains
the bound. Taking expectations of the minorant under any feasible law
therefore proves the stated exact convex-envelope value. Every term has
lower envelope zero and upper envelope $1/2$, so
$\cav f_m=T(f_m)=(9m/8)\binom m2$; subtraction gives the exact ratios.
The checker verifies all $5^2+9^2+13^2$ minorant inequalities and these
attainment and envelope equalities in rational arithmetic.

\begin{table}[htbp]
\centering
\begin{tabular}{rllrrl}
\toprule
$m$&Affine minorant&Counts $(A,C)$&$\vex f_m$&$\cav f_m=T$&$T/H$\\\midrule
4&$9A+4C-19$&$(2,3)$&11&27&$27/16$\\
8&$47A+20C-188$&$(4,6)$&120&252&$21/11$\\
12&$(231/2)A+48C-684$&$(6,9)$&441&891&$99/50$\\\bottomrule
\end{tabular}
\caption{Exact smaller two-level certificates and attaining counts.}
\label{tab:two-level-small}
\end{table}

\newpage
Here is a complete Python 3 checker using only the standard library.
It enumerates $7^3+9^3+65^3$ three-group states and
$5^2+9^2+13^2+17^2$ two-group
states. The data are exactly those in the manuscript tables; no numerical
optimization or proprietary solver is involved. The same executable text
is supplied as \texttt{verification/check\_stage02\_finite.py}.
Concatenate the two code blocks in order to obtain that file.
\par\medskip\noindent\begin{minipage}{\textwidth}\small
\begin{verbatim}
from fractions import Fraction as Q
from itertools import product
from math import comb

# m, coefficients, denominator, integer affine coefficients,
# probability denominator, (state, probability numerator), ratio
cases = [
 (6, (2,3,9,10,7,7), 13, (962,778,842,-4816), 26,
  [((1,4,4),17), ((2,1,6),8), ((6,2,1),1)],
  Q(20891,10411)),
 (8, (2,3,13,12,8,7), 7, (793,770,798,-5882), 7,
  [((1,2,8),2), ((1,5,6),4), ((8,4,2),1)],
  Q(6601,3225)),
 (64, (2,3,120,105,70,63), 105,
  (871710,900446,899046,-51743768), 735,
  [((4,19,64),241), ((5,19,64),12),
   ((16,40,43),419), ((64,31,17),63)],
  Q(7443345,3445256))]

def phi(a, b, c, coef):
    values = (comb(c,3), b*comb(c,2), comb(b,2),
              a*c, a*b, comb(a,2))
    return sum(v*w for v,w in zip(values,coef))

\end{verbatim}
\end{minipage}
\newpage
\noindent\begin{minipage}{\textwidth}\small
\begin{verbatim}
for m,coef,D,dual,P,atoms,ratio in cases:
    for a,b,c in product(range(m+1), repeat=3):
        rhs = sum(v*w for v,w in zip((a,b,c,1),dual))
        assert D*phi(a,b,c,coef) >= rhs
    assert all(p >= 0 for _,p in atoms)
    assert sum(p for _,p in atoms) == P
    means = [sum(Q(p,P)*state[j] for state,p in atoms)
             for j in range(3)]
    assert means == [Q(m,4), Q(m,2), Q(3*m,4)]
    primal = sum(Q(p,P)*phi(*state,coef) for state,p in atoms)
    lower = sum(Q(v,D)*w for v,w in zip(dual,means+[1]))
    assert primal == lower
    sizes = (comb(m,3), m*comb(m,2), comb(m,2),
             m*m, m*m, comb(m,2))
    upper = (Q(3,4), Q(1,2), Q(1,2), Q(1,4), Q(1,4), Q(1,4))
    cav = sum(c*s*u for c,s,u in zip(coef,sizes,upper))
    term_lower = Q(coef[0]*comb(m,3),4)
    assert cav > primal
    assert (cav-term_lower)/(cav-primal) == ratio
    assert ratio > 2
    print(3*m, cav, term_lower, primal, ratio)

# m, affine coefficients (A,C,1), attaining value, cav=T, ratio
small = [
 (4, (9,4,-19), 11, 27, Q(27,16)),
 (8, (47,20,-188), 120, 252, Q(21,11)),
 (12, (Q(231,2),48,-684), 441, 891, Q(99,50))]
for m,dual,vex,cav,ratio in small:
    def value(a,c):
        return a*comb(c,2)+Q(5*m,4)*comb(a,2)
    slacks = [value(a,c)-dual[0]*a-dual[1]*c-dual[2]
              for a,c in product(range(m+1), repeat=2)]
    assert min(slacks) == 0
    a,c = m//2,3*m//4
    assert Q(a,m) == Q(1,2) and Q(c,m) == Q(3,4)
    assert value(a,c) == dual[0]*a+dual[1]*c+dual[2] == vex
    assert (m+Q(5*m,4))*comb(m,2)/2 == cav
    assert cav > vex and Q(cav,cav-vex) == ratio < 2
    print(2*m, cav, vex, ratio)

mins = []
for a in range(17):
    residuals = [2*(a*comb(c,2)+20*comb(a,2))
                 -428*a-177*c+3372 for c in range(17)]
    mins.append(min(residuals))
assert mins == [540,352,204,96,28,0,9,7,0,0,19,63,
                132,235,369,540,742]
assert 8*comb(12,2)+20*comb(8,2) == 1088
assert 214*8+Q(177,2)*12-1686 == 1088
assert Q(2160,2160-1088) == Q(135,67)
assert Q(135,67) > 2
assert Q(243,115)*(1-Q(1,20)) == Q(4617,2300) > 2
print("All finite cubic certificates passed.")
\end{verbatim}
\end{minipage}

\section{Auxiliary structural reductions}\label{app:structural-auxiliary}
The following arguments are useful independently of the width-two proof.
They concern the stated local laws or reductions and make no claim about
arbitrary larger-width extensions.

\begin{lemma}[Canonical pair laws on a forest]\label{lem:canonical-pairs}
On a forest of binary variables of prescribed means $p_i$, let $Q_{ij}$
be the midpoint of the two Fr\'echet endpoint pair laws on each edge:
\[
 Q_{ij}(1,1)=\frac{\min(p_i,p_j)+\max(0,p_i+p_j-1)}2.
\]
For any target set $S$ there is a global law with these pair marginals
and target union probability at least
$[\min(1,\sum_{i\in S}p_i)+\max_{i\in S}p_i]/2$, with both expressions
zero for $S=\varnothing$. For any nonnegative target payoff, such a law
also attains at least half its unconstrained singleton-feasible maximum.
The resulting law may depend on the target.
\end{lemma}
\begin{proof}
Every pair cell is affine in its $(1,1)$ probability and nonnegative
on the Fr\'echet interval. Its midpoint value is therefore at least
half every feasible cell value. Choose a singleton-feasible law $P$
maximizing the target union, using consecutive circular event intervals
to attain $c=\min(1,\sum_Sp_i)$. The edge measures
$R_{ij}=2Q_{ij}-P_{ij}$ are nonnegative probability measures with the
prescribed singleton means. They glue on the forest by rooting each
component and sampling each child conditionally on its parent. Null
conditioning states are irrelevant. Thus a law $R$ with these edge
marginals exists. The mixture $(P+R)/2$ has exactly the desired $Q$
marginals, and its target union is at least $(c+b)/2$ since every law
has union at least $b=\max_Sp_i$. For a nonnegative payoff use its
maximizing law $P$ and $\E_Rg\ge0$ in the same construction.
\end{proof}
This residual construction cannot simply be applied on a triangle:
for three fair bits, take $P$ equally on the two constant states.
Each canonical pair is uniform, so $2Q_{ij}-P_{ij}$ requires perfect
anticorrelation. Three bits cannot be pairwise anticorrelated. This
is a failure of that residual-gluing step, not a counterexample to a
positive-monomial width-two bound.

\begin{lemma}[Identical-neighborhood compression]\label{lem:twin-compression}
In a positive unit-cube monomial factorization, suppose a nonempty
variable block $B$ occurs in every incident factor in its entirety.
Replace it by one bit $W$ of mean
$w=\max(0,\sum_{i\in B}x_i-|B|+1)$ in those factors.
The full lower envelope and sum of local lower envelopes equal those
of the reduced instance. For some $\Delta\ge0$,
\[
 \tgap_{\rm original}=\tgap_{\rm reduced}+\Delta,\qquad
 \hgap_{\rm original}=\hgap_{\rm reduced}+\Delta.
\]
Thus every reduced gap inequality with constant at least one transfers.
\end{lemma}
\begin{proof}
For any original binary law, the product $V=\prod_{i\in B}X_i$ has
mean $q\ge w$. If $q>0$, independently retain its ones with probability
$w/q$, producing $W\le V$ of mean $w$; if $q=0$ use $W=0$.
The reduced objective is nondecreasing in $W$, so this operation cannot
increase it. Conversely choose a local law on $B$ attaining product
mean $w$ with all the original means. Attach it to a reduced law by
conditioning on its product value $W$. This lifts the law with the
same objective and original means. These two operations prove full
lower equality, including the endpoint cases.
For a factor with scope $B\cup C$, its original local lower value is
\[
 \max(0,\sum_Bx_i+\sum_Cx_i-|B|-|C|+1)
 =\max(0,w+\sum_Cx_i-|C|),
\]
where if $\sum_Bx_i-|B|+1<0$ both sides vanish. The upper value can
only decrease under compression because $w\le\min_Bx_i$; this also
holds when $C$ is empty and the reduced factor is affine.
Common upper attainment permits summing these decreases into one
$\Delta\ge0$. The two gap identities follow, and
$\tgap_{\rm reduced}+\Delta\le c(\hgap_{\rm reduced}+\Delta)$ for
$c\ge1$ proves the transfer without dividing by a gap.
\end{proof}

\begin{proposition}[Signed forest gluing on arbitrary boxes]
\label{prop:signed-forest}
For any multiaffine local factors on finite boxes whose incidence graph
is a forest, both upper and lower envelopes add. In particular this
holds for frequency-two factors whose dual multigraph is a forest.
\end{proposition}
\begin{proof}
Choose an optimizing vertex law for each factor and normalize its
coordinate endpoints. Root each incidence-tree component at a factor
and traverse it, sampling each new factor conditionally on its single
already sampled variable. Consistent singleton means ensure consistency
at that separator; every other variable is new, or there would be an
incidence cycle. Null separator states have zero probability. This
simultaneously preserves all chosen optimizing laws. Use minimizing
laws for the lower identity and maximizing laws for the upper identity.
Different components and unused variables can be sampled independently.
Fixed coordinates are substituted first. Signs do not enter the proof.
\end{proof}

A factor partition into $q$ incidence forests would give a $q$ bound
by mixing their exact laws for positive monomials, but treewidth does
not bound such a partition number. In $K_{2,m}$, $m\ge2$, with the two-vertex
side interpreted as variables, any two factors create a four-cycle;
all $m$ factors need distinct forest colors, although treewidth is two.
Adding one private variable to each factor makes their monomial scopes
distinct and preserves the obstruction. This explains why the stronger
matrix argument is useful.

For the radix construction, the general cutoff following
Theorem~\ref{thm:variable-radix} also records a fixed-radix limitation.
For integers $b\ge2,L\ge2$ its profile capacities satisfy
\[
 M_s=\frac{(L-s)(b-1)+b}{b^s},\qquad
 \sum_{q=s+1}^L M_q=\frac{L-s}{b^s}.
\]
The sum identity follows by telescoping
$M_q=(L-q+1)b^{1-q}-(L-q)b^{-q}$.
If $s$ brackets one between $M_s$ and $M_{s+1}$, the exact hull is
$s+(L-s)b^{-s}$, as proved by its two attaining profiles.
For fixed $b$, $s=\log_b L+O_b(1)$: the inequality $M_s\ge1$ gives
$s\le\log_b((b-1)L+b)$, hence $L-s\sim L$; the reverse bracket then
gives $b^{s+1}\ge(b-1)(L-s-1)+b$. The remainder term is $O_b(1)$,
so the hull is $\log_b L+O_b(1)$. This is a limitation of this
fixed-radix construction, not a sharp second-order lower theorem for
all positive polynomials.

\section{Distinct positive-box comparisons}\label{app:positive-box-predecessors}
These proofs are superseded as upper bounds by
Theorem~\ref{thm:positive-box-upper}, but use different estimates.
Throughout this appendix $\rho>1$, $t=\rho-1$, $\alpha=1/\rho$, and
$\eta=1-\alpha$. Work first on $[1,\rho]^n$. The law $I$ is independent
endpoint rounding and $O$ is independently fair endpoint orientation.
For one product write $C,V,P,O$ for its upper, lower, independent,
and orientation values, and $D_I=C-P,D_O=C-O$.

\subsection{The symmetric low/high argument}
\begin{proposition}[A coarse finite constant]\label{prop:coarse-box}
For positive monomials on a box whose aspect ratios are at most $\rho$,
\[
 \tgap\le K(\rho)\hgap,\qquad
 K(\rho)=2(1+1/\rho)+\max\{4(1+\rho),8\eta^{-3}\}.
\]
In particular, with $R=\max\{4,\max_i u_i/\ell_i\}$,
$\tgap\le(4R+6+2/R)\hgap$, which is $45/2$ for aspect ratios at most
four. Direct use of $\rho=2$ instead gives the older constant $67$.
\end{proposition}
\begin{proof}
Classify a normalized mean as low when it is at most $1/2$. Put $q_i=p_i$
for lows and $q_i=1-p_i$ for highs, so $q_i\le1/2$. Divide the product's
values by $\rho^h$, where $h$ is the number of high coordinates.
Let $L(s),H(s)$ count the respective deviations at least $s$, and let
$Q_L,Q_H$ be their sums. Integrals in this paragraph are over $[0,1/2]$.
Define
\[
 \begin{array}{ll}
 C_L=1+\int(\rho^{L(s)}-1)\,ds,&P_L=\prod_{\rm low}(1+tq_i),\\
 C_H=1+\int(\alpha^{H(s)}-1)\,ds,&P_H=\prod_{\rm high}(1-\eta q_i).
 \end{array}
\]
Common-threshold low successes and high failures occupy opposite ends
of the interval. Thus $C=C_L+C_H-1$, $P=P_LP_H$, and
\begin{equation}\label{eq:coarse-components}
 D_I=(C_L-P_L)+(C_H-P_H)+J_I,
 \quad J_I=(P_L-1)(1-P_H)\ge0.
\end{equation}
All three terms are nonnegative. Set
$A=C_L-1-tQ_L$, $B=C_H-1+\eta Q_H$, and let $q_L,q_H$ be the respective
largest deviations, zero for empty groups. Integer-count inequalities give
\begin{equation}\label{eq:dispersion-coarse}
 A\ge t^2(Q_L-q_L),\qquad B\ge\eta^2(Q_H-q_H).
\end{equation}
For the first, use $(1+t)^l-1-tl\ge t^2(l-1)_+$; for the second,
$\alpha^l-1+\eta l$ has increments $\eta(1-\alpha^l)\ge\eta^2$
for $l\ge1$. Integrating proves both.
Conditional orientation expectations at either end of the interval are
$b_L=(1+t/2)^{L(s)}$ and $b_H=(1-\eta/2)^{H(s)}$, so
\[
 O=1+2\int(b_Lb_H-1),\quad
 D_O=D_{OL}+D_{OH}+J_O,\quad J_O=2\int(b_L-1)(1-b_H).
\]
Here $D_{OL},D_{OH}\ge0$ are the separate-group deficiencies.
Expansion in powers of $t$ gives
$(1+t)^l-2(1+t/2)^l+1\ge[(1+t)^l-1-tl]/2$, since each order-$j\ge2$
coefficient is multiplied by $1-2^{1-j}\ge1/2$. Therefore
\begin{equation}\label{eq:orientation-coarse}
 D_O\ge A/2+J_O,\qquad J_O\ge(t\eta/2)\min(q_L,q_H).
\end{equation}
The high separate-group deficiency is also nonnegative directly from
convexity of the integer power at the midpoint of $\alpha$ and one.

Write $D_{IH}=C_H-P_H$. If $Q_H\le1$, put $q=q_H$ and $R'=Q_H-q$.
Bonferroni's bound gives
\[
 D_{IH}\ge B-\eta^2\sum_{i<j}q_iq_j
 \ge B-\eta^2(qR'+R'^2/2)
 \ge(1-q-R'/2)B\ge B/4.
\]
Here $B\ge\eta^2R'$, $q\le1/2$, and $q+R'\le1$; no division by
$R'$ is needed. If $Q_H\ge1$, order deviations decreasingly.
At a fixed sum, the linear common-threshold high value is minimized
by filling the largest-weight coordinates to their cap $1/2$ and at
most one remainder. Convex interpolation then gives
\[
 C_H\ge\tfrac12+\tfrac12\alpha^{2Q_H},\qquad P_H\le e^{-\eta Q_H}.
\]
The difference $g(Q)=1/2+\alpha^{2Q}/2-e^{-\eta Q}$ is nondecreasing
for $Q\ge1$. Indeed $-\log\alpha\le\eta/\alpha$ and
$\alpha^{2Q-1}\le e^{-\eta Q}$ imply
$g'(Q)\ge\eta[e^{-\eta Q}-\alpha^{2Q-1}]\ge0$.
The fourth-order exponential remainder yields
\[
 D_{IH}\ge g(1)=1-\eta+\eta^2/2-e^{-\eta}
 \ge\eta^3/6-\eta^4/24\ge\eta^3/8.
\]

Let $\phi_\rho$ interpolate $\rho^k$ between every consecutive integer,
including negative integers. The exact normalized lower envelope is
$V=\phi_\rho(Q_L-Q_H)$, because the original count sum is shifted by
integer $h$. Supporting slopes at zero are $\eta$ and $t$, so
\begin{equation}\label{eq:coarse-target}
 C-V\le A+B+t\eta\min(Q_L,Q_H).
\end{equation}
For $Q_H\le1$, use
$\min(Q_L,Q_H)\le\min(q_L,q_H)+A/t^2+B/\eta^2$ to obtain
\[
 T\le(1+1/\rho)A+(1+\rho)B+2J_O
 \le2(1+1/\rho)D_O+4(1+\rho)D_I.
\]
For $Q_H\ge1$, $T\le C_L=1+A+tQ_L$ and
$J_I\ge tQ_L(1-e^{-\eta})\ge t\eta Q_L/2$.
The preceding bounds therefore imply
\[
 T\le2D_O+8\eta^{-3}D_{IH}+2\eta^{-1}J_I
      \le2D_O+8\eta^{-3}D_I.
\]
Combining the cases proves the per-product inequality with the two
coefficients in $K(\rho)$. Mix the global laws proportionally, undo
the product's normalization, and sum to prove the common-aspect bound.
The positive affine substitution of
Theorem~\ref{thm:positive-box-upper} transfers it to unequal aspects.
For $R\ge4$, $8(1-1/R)^{-3}\le512/27<20\le4(1+R)$, giving
$K(R)=4R+6+2/R$. For $\rho=2$ the two coefficients are $3$ and $64$.
\end{proof}
The high-mass argument in fact permits replacing $8\eta^{-3}$ by
$1/g(1)$. Indeed $g(1)\le1-e^{-\eta}$, since their difference is
$\eta-\eta^2/2\ge0$; thus the same coefficient controls both the
constant and the $J_I$ term in $C_L$. This gives the distinct sharper
finite expression
$2(1+1/\rho)+\max\{4(1+\rho),1/g(1)\}$.

A tempting all-high simplification $T\le2D_I$ is false. On $[1,2]^3$
with failure probabilities $(1/2,49/100,1/100)$, divide values by
$2^3$. Exact values are
\[
 C=501/800,\quad V=1/2,\quad P=90147/160000,
 \quad T=20200/160000>2D_I=20106/160000.
\]
The boundary mean $1/2$ is allowed in this all-high calculation.
It does not invalidate the mixed-law inequality.

\subsection{The asymmetric split}
\begin{proposition}[An earlier leading-order bound]\label{prop:asymmetric-box}
For $\rho\ge64$, positive products on boxes of aspect ratios at most
$\rho$ satisfy
\[
 \tgap\le\left(2+\frac{\rho+1}{1-3/\sqrt\rho}\right)\hgap.
\]
\end{proposition}
\begin{proof}
Put $\delta=\rho^{-1/2}\le1/8$ and $b=(\rho+1)/(1-3\delta)$.
Now classify normalized means as low if $p_i\le1-\delta$, and high
otherwise. High deviations are $q_i=1-p_i<\delta$. Normalize by
$\rho^h$ again. Reset the counts to
$L(s)=\#\{i\text{ low}:p_i\ge s\}$ and
$H(s)=\#\{i\text{ high}:q_i\ge s\}$, with
$Q_L=\sum_{\rm low}p_i$ and $Q_H=\sum_{\rm high}q_i$.
The separate groups' common-threshold upper values and independent
values are now
\[
 \begin{array}{ll}
 C_L=1+\int_0^{1-\delta}(\rho^{L(s)}-1)\,ds,
 &P_L=\prod_{\rm low}(1+tp_i),\\
 C_H=1+\int_0^\delta(\alpha^{H(s)}-1)\,ds,
 &P_H=\prod_{\rm high}(1-\eta q_i).
 \end{array}
\]
Low successes and high failures under common-threshold rounding are
disjoint, so \eqref{eq:coarse-components}, the definitions
$A=C_L-1-tQ_L$, $B=C_H-1+\eta Q_H$, and
\eqref{eq:coarse-target} remain valid. The dispersion bounds
\eqref{eq:dispersion-coarse} hold with maxima $a_L,a_H$ of low means
and high failures; they did not require low means at most $1/2$.
A different estimate now replaces $D_{OL}\ge A/2$:
\[
 D_{IL}=C_L-P_L\ge\delta A.
\]
Expand the low physical product only for this calculation. Every
subset of degree at least two has independent product at most
$(1-\delta)$ times its minimum mean, and the sum of its common upper
values is exactly $A$. Constant and affine terms cancel.

For $s\in[0,\delta]$ and $U=s$ or $1-s$, conditional low expectations
are $1+(t/2)\ind\{s\le p_i\}$, while high normalized expectations
are $1-(\eta/2)\ind\{s\le q_i\}$. Outside these two end strips
the conditional high product is one. Conditional independence of the
orientation coins gives $D_O=D_{OL}+D_{OH}+J_O\ge J_O$, where
\[
 J_O=2\int_0^\delta
 \left(\prod_{\rm low}[1+(t/2)\ind\{s\le p_i\}]-1\right)
 \left(1-\prod_{\rm high}[1-(\eta/2)\ind\{s\le q_i\}]\right)ds
 \ge(t\eta/2)\min(a_L,a_H).
\]
Combining this with dispersion and the target inequality yields
$T\le(1+1/\rho)A+(1+\rho)B+2J_O$.
If $Q_H\le4\delta$, the same Bonferroni argument with $q=a_H\le\delta$
and $R'=Q_H-q\le4\delta$ gives
$D_{IH}\ge(1-q-R'/2)B\ge(1-3\delta)B$.
Since $(1+1/\rho)/\delta\le b$, it follows that
$T\le bD_I+2D_O$.

If $Q_H\ge4\delta$, the common high failure set has measure at most
$\delta$, so $C_H\ge1-\delta$. Also
$P_H\le e^{-4\eta\delta}\le1-2\delta$, by
$e^{-z}\le1-z+z^2/2$ and $4\eta-8\eta^2\delta\ge3-1=2$.
Hence $D_{IH}\ge\delta$ and $J_I\ge\delta tQ_L$, giving
$D_I\ge\delta(A+1+tQ_L)=\delta C_L$.
As $T\le C_L$, again $T\le bD_I+2D_O$.
Mix the two global laws with weights proportional to $b,2$, sum over
products, and use positive affine transfer for unequal aspects.
\end{proof}
Its expansion is $\rho+3\sqrt\rho+O(1)$; the coefficient theorem
improves the additive error to two.

\subsection{Optimality within a fixed two-law mixture}
\begin{proposition}[A limitation of fixed independent/fair-orientation mixtures]
\label{prop:fixed-mixture-optimal}
Fix $\rho>1$ and an orientation weight $w\in[0,1]$. For mixtures of
$I$ and independently fair $O$, the best uniform per-physical-monomial
captured fraction over all dimensions is at most
$\min\{w/2,(1-w)/\rho\}$. Optimizing the weight gives exactly
$1/(\rho+2)$, attained at $w=2/(\rho+2)$.
\end{proposition}
\begin{proof}
For a bilinear product at normalized means $(\epsilon,1-\epsilon)$,
$0<\epsilon<1/2$, deleting affine parts gives
$D_I/T=\epsilon$ and $D_O/T=1/2$. Letting $\epsilon\downarrow0$
bounds the mixture fraction by $w/2$.
For the other obstruction use means $\epsilon$ and $k$ copies of
$1-\epsilon/k$, where $k\ge1$. Divide physical values by $\rho^{k+1}$,
and write $\alpha=1/\rho$, $\eta=1-\alpha$, $\gamma=(1+\alpha)/2$.
The sum of success means is $k$, giving $V=\alpha$.
Common-threshold and fair-orientation integration, split at
$\epsilon/k$ and $\epsilon$, and direct independent multiplication give
\[
 \begin{split}
 C&=\alpha+\epsilon[\eta-\alpha(1-\alpha^k)/k],\\
 P&=\alpha[1+(\rho-1)\epsilon][1-\eta\epsilon/k]^k,\\
 O&=\alpha+\epsilon[\eta-2\gamma(1-\gamma^k)/k].
 \end{split}
\]
For clarity, in the orientation law the two end intervals of length
$\epsilon/k$ have normalized expected product $\gamma^{k+1}$,
the two next intervals of total length $2(\epsilon-\epsilon/k)$
have value $\gamma$, and the middle interval of length $1-2\epsilon$
has value $\alpha$; these yield the stated $O$.
For each fixed $k$, divide deficiencies by $\epsilon$ and let
$\epsilon\downarrow0$. Then
\[
 \begin{split}
 T/\epsilon&=\eta-\alpha(1-\alpha^k)/k,\\
 D_I/\epsilon&\longrightarrow\alpha\eta-\alpha(1-\alpha^k)/k,\\
 D_O/\epsilon&=[2\gamma(1-\gamma^k)-\alpha(1-\alpha^k)]/k.
 \end{split}
\]
The denominator coefficient is at least $\eta^2>0$ because
$1-\alpha^k\le k\eta$. Now let $k\to\infty$ to obtain
$D_I/T\to1/\rho$ and $D_O/T\to0$, bounding the mixture by
$(1-w)/\rho$. The maximum of the minimum of these two bounds is
$1/(\rho+2)$ at their intersection. Lemma~\ref{lem:box-coefficient}
with fair orientations supplies the matching guarantee.
\end{proof}
This is only an optimality statement for a fixed two-law mixture's
per-monomial guarantee over unrestricted dimensions. Other laws,
coefficient-dependent mixtures, and the true optimum joint law for a
sum of monomials are outside it. In particular it does not prove
$C_{\rm box}(\rho)=\rho+2$, nor conflict with the finite-ambient
balanced-orientation refinement.

\section{A classical point-packing benchmark}\label{app:point-packing}
The following four values were conjectured by \citet[Section 4]{Anstreicher2009}
and proved by \citet[Propositions 1(i--ii) and 3]{Khajavirad2024}.
We include a variance proof and explicit feasible covariances to illustrate
how an exact relaxation value differs from a spatial certificate bound.
For $n\ge2$ points $(x_i,y_i)\in[0,1]^2$, maximize $\gamma$ subject to
$(x_i-x_j)^2+(y_i-y_j)^2\ge\gamma$ for $i<j$. Replace the two sets of
products by symmetric matrices $X,Y$ and write
\[
 D_{ij}=X_{ii}-2X_{ij}+X_{jj}+Y_{ii}-2Y_{ij}+Y_{jj}\ge\gamma.
\]
RLT imposes all four products of coordinate-bound slacks, including
repeated indices, as in \eqref{eq:box-rlt}. SDP instead imposes
$X\succeq xx^T$, $Y\succeq yy^T$ and the diagonal secants
$X_{ii}\le(\ell_i+u_i)x_i-\ell_i u_i$, and likewise for $Y$.
The symmetry bounds, denoted SYM, restrict $x_i\in[1/2,1]$ for
$i\le n_x=\lceil n/2\rceil$ and $y_i\in[1/2,1]$ for
$i\le k=\lceil n/4\rceil$. Reflection gives the lower-half convention
in Khajavirad's paper. Both RLT and SDP secants are recomputed on these
intervals; adding only first-moment bounds would define different models.

\begin{proposition}[Four exact values]\label{prop:packing-values}
The optimal values are
\[
\begin{array}{c|c|c}
\text{model}&\text{value}&\text{range}\\\hline
\mathrm{RLT}&2&n\ge2\\
\mathrm{SDP}\text{ (with or without RLT)}&n/(n-1)&n\ge2\\
\mathrm{RLT+SYM}&1/2&n\ge5\\
\mathrm{SDP+SYM}\text{ (with or without RLT)}&k/[4(k-1)]&n\ge5.
\end{array}
\]
Here $k-1=\lfloor(n-1)/4\rfloor$.
\end{proposition}
\begin{proof}
For two coordinates in a common interval $[\ell,\ell+L]$, the RLT
maximum of $X_{ii}-2X_{ij}+X_{jj}$ at fixed means is
\[
 L\min\{x_i+x_j-2\ell,\ 2\ell+2L-x_i-x_j\}\le L^2.
\]
Take both diagonals at their secants and the off-diagonal at the larger
of its two lower McCormick planes. This is feasible for all pairs at once:
the diagonal secant dominates both diagonal lower planes, and each pair's
lower planes lie below its upper planes. Thus the bound two holds without
symmetry and is attained at $x=y=\boldsymbol1/2$, with diagonal $1/2$
and off-diagonal zero in each matrix.
With symmetry, two of the first $k\ge2$ points give the bound $1/2$.
Put the three groups $Q=[k]$, $H=[n_x]\setminus Q$, and
$F=[n]\setminus[n_x]$ at $(3/4,3/4)$, $(3/4,1/2)$, and $(1/2,1/2)$.
Set every diagonal to its interval secant and each off-diagonal to its
lower RLT bound. In group-pair order $QQ,QH,QF,HH,HF,FF$, the resulting
squared distances are
\[
 1/2,\quad7/8,\quad5/4,\quad5/4,\quad13/8,\quad2.
\]
Pairs from groups with fewer than two members impose no extra requirement.
This proves both RLT values.

For SDP, restrict to any $k\ge2$ points sharing $[\ell,\ell+L]$ in one
coordinate, and let $e=\boldsymbol1\in\R^k$. With
$s=\sum_i(x_i-\ell)$, PSD and the secants give
\begin{align*}
 \sum_{i<j}(X_{ii}-2X_{ij}+X_{jj})
 &=k\operatorname{tr}X-e^TXe\\
 &\le k[(2\ell+L)\sum_i x_i-k\ell(\ell+L)]-(\sum_i x_i)^2\\
 &=kLs-s^2\le k^2L^2/4.
\end{align*}
The pair average per coordinate is at most $kL^2/[2(k-1)]$.
Add the two coordinates, first using $k=n,L=1$ and then the symmetry
subset with $L=1/2$, to obtain the two asserted SDP upper bounds.

Without symmetry, set $x=y=e/2$ and
\begin{equation}\label{eq:packing-covariance}
 X=Y=ee^T/4+Z,\qquad
 Z=\frac{n}{4(n-1)}(I-ee^T/n)\succeq0.
\end{equation}
Its diagonals equal $1/2$ and its off-diagonals equal
$1/4-1/[4(n-1)]$. Every pair has total distance $n/(n-1)$.
These entries also satisfy all RLT bounds.

For symmetry use the same three group means as above. In each coordinate
set the moment matrix to mean outer product plus a covariance that is
block diagonal in $Q,H,F$. On $Q$ use in both coordinates
\[
 Z_Q=\frac{k}{16(k-1)}(I-ee^T/k).
\]
On $H$ use $I/16$ in $x$ and $I/4$ in $y$; on $F$ use $I/4$ in both.
These matrices are PSD because $I-ee^T/k$ is an orthogonal projection.
Every diagonal equals its tightened secant. The group-pair distances are
\[
 \frac{k}{4(k-1)},\quad\frac12,\quad\frac34,\quad
 \frac58,\quad\frac78,\quad1.
\]
They are at least the first value for every $k\ge2$.
For RLT feasibility, cross-group and distinct $H,F$ entries are products
of their means and hence obey the four planes. On $Q$ an off-diagonal is
$9/16-1/[16(k-1)]\in[1/2,5/8]$, exactly within the tightened RLT range
at mean $3/4$. Diagonal secants satisfy their diagonal lower planes as
above. Thus adding RLT cannot change either SDP optimum.
\end{proof}

In dimension $d$, the same one-coordinate average and the same covariance
in each coordinate give RLT value $d$ and SDP value
$dn/[2(n-1)]$ on $[0,1]^d$, for $n\ge2$. A subset of $k\ge2$ points
in a common side-$L$ cube forces an SDP bound at most
$dL^2k/[2(k-1)]$. This includes $L=0$. These are dimension-scaled
averaging statements, not claims about all symmetry or ordering models.
Numerical symmetry values at $n<5$ and ordering formulations are not
needed here. The benchmark explains weakness of these particular root
relaxations; it establishes no lower bound on every possible strengthening
or spatial algorithm.

\section{Scale disjunctions and shared operating variables}\label{app:scaling}
Scaling gives an example where a representation can remove an artificial
gap. The arguments are affine interpolation and classical disjunctive
convexification \citep[Section 2, Theorem 2.1]{Balas1998}; we give them
explicitly to identify the role of shared intensive variables.
Let $E\subset\R^d\times\R^p$ be nonempty compact, with per-unit extensive
coordinates $\xi$ and intensive coordinates $T$. Define
\[
 \lambda\odot E=\{(\lambda\xi,T):(\xi,T)\in E\},\qquad
 F(\Lambda,E)=\{(\lambda,\lambda\xi,T):\lambda\in\Lambda,(\xi,T)\in E\}.
\]
The nonnegative scale set $\Lambda$ is nonempty compact, with extremes
$\lambda_-,\lambda_+$. At zero scale, $X=0$ and $T\in\operatorname{proj}_T E$;
a different off-state domain must be modeled explicitly.

\begin{proposition}[Endpoint hull and multiplicities]\label{prop:scale-hull}
With $C=\conv E$, the hull of $F(\Lambda,E)$ is the hull of just its two
extreme-scale slices, with $E$ replaced by $C$. It is exactly the projection
of
\begin{equation}\label{eq:scale-disaggregation}
 \begin{gathered}
 0\le\theta\le1,\quad\lambda=(1-\theta)\lambda_-+\theta\lambda_+,
 \quad X=X^-+X^+,\quad T=T^-+T^+,\\
 (X^-,T^-)\in(1-\theta)(\lambda_-\odot C),\qquad
 (X^+,T^+)\in\theta(\lambda_+\odot C),
 \end{gathered}
\end{equation}
where $0K=\{0\}$ for nonempty $K$. For $N\ge1$ identical units sharing
$T$, allow integer $0\le n\le N$, $X=\sum_{i=1}^n\xi_i$, and
$(\xi_i,T)\in E$. Their hull is the same as $\conv F(\{0,\ldots,N\},E)$.
Writing $P_T=\operatorname{proj}_T C$, it is
\begin{equation}\label{eq:multiplicity-hull}
 0\le n\le N,\qquad
 \exists T^+:\ (X,T^+)\in\frac nN(N\odot C),\quad
 T-T^+\in(1-n/N)P_T.
\end{equation}
\end{proposition}
\begin{proof}
For fixed $(\xi,T)$, the point $(\lambda,\lambda\xi,T)$ depends affinely
on $\lambda$. If the extremes differ, interpolate with
$\theta=(\lambda-\lambda_-)/(\lambda_+-\lambda_-)$; if they coincide,
there is only one slice. Every original point is therefore a combination
of endpoint points. A fixed-scale map is linear in $(\xi,T)$, so its
hull is that slice of $C$. Combining two compact convex slices gives
\eqref{eq:scale-disaggregation}, including zero weights.
For an integer $n>0$ and fixed $T$, a sum of $n$ points of $E_T$ lies in
$n\conv E_T$, and identical copies include $n E_T$. Equivalently, a
multiplicity point is the average of its $n$ same-scale points
$(n,n\xi_i,T)$, while every same-scale point is a feasible multiplicity
point. Their hulls therefore agree at fixed $n$ after convexification.
Taking the hull over $n$ and applying the endpoint result proves
\eqref{eq:multiplicity-hull}. At $n=0$ both models use the stated
projection off-state; $N=0$ consists only of that slice.
\end{proof}
If $C=\{(\xi,T):A\xi+BT\le b\}$ and $P_T=\{T:DT\le e\}$ are
polytopes, \eqref{eq:multiplicity-hull} becomes
\[
 AX+NB T^+\le bn,\qquad D(T-T^+)\le e(1-n/N),\qquad0\le n\le N.
\]
Compact conic-representable slices can also be homogenized at zero weight.
Indeed, if $K_0=\{z:\exists y,Az+By+b\in\mathcal K\}$ is nonempty compact,
then at weight $t>0$ division in $Az+Bw+tb\in\mathcal K$ gives $tK_0$.
At $t=0$ a feasible $(z,w)$ is a recession direction of the lifted set:
add any nonnegative multiple to a feasible lift. Its projected direction
must be zero by compactness of $K_0$. Conversely $z=w=0$ is feasible.
This checks zero weights even when the auxiliary lift is unbounded.
The aggregation mechanism also appears in \citet[Section 4.1, Lemma 3
and Theorem 3, under Assumption 1]{WuEtAl2026}.

For example, let per-unit load $\xi\in[\ell,u]$, $0\le\ell<u$, common
$T\in[a,b]$, $a<b$, and per-unit extensive output $\omega=\xi T$.
Then $C$ is the McCormick hull of this bilinear graph. Formula
\eqref{eq:multiplicity-hull} supplies its exact local aggregate hull with
one disaggregated intensive coordinate. Its endpoint vertices are the two
off-state points and the four corners at scale $N$. Besides
$n\ell\le X\le nu$, $0\le n\le N$, $a\le T\le b$ and the global
McCormick inequalities for $W=XT$ on $[0,Nu]\times[a,b]$, this hull has
\begin{equation}\label{eq:scale-bilinear-cuts}
 W\le bX+\ell[NT-Na-n(b-a)],\qquad
 W\ge aX+\ell[NT-Nb+n(b-a)].
\end{equation}
To check sufficiency without a facet calculation, introduce $S=NT^+$.
The scaled per-unit McCormick system is equivalent to
\[
 an\le S\le bn,\quad N T-b(N-n)\le S\le NT-a(N-n),
\]
\[
 \ell S+aX-a\ell n\le W\le\ell S+bX-b\ell n,
 \qquad uS+bX-bun\le W\le uS+aX-aun.
\]
Eliminating $S$ by comparing each lower and upper bound gives precisely
the displayed system (redundant comparisons follow from the stated
bounds). When $\ell=0$, use these inequalities directly without dividing
by $\ell$. Alternatively \eqref{eq:multiplicity-hull} itself is a complete
formulation at that boundary.

\subsection{When scale costs do and do not separate}
\begin{example}[A scale-cost gap on a convex per-unit set]
\label{ex:scale-cost}
Take $E=\{(\xi,T):\xi=T\in[0,1]\}$, $\Lambda=\{0,1,2\}$, and
$f(0)=f(2)=1$, $f(1)=0$. The point $(\lambda,X,T)=(1,1,1/2)$ is the
average of $(0,0,0)$ and $(2,2,1)$. Its true minimum convexified cost
is one although the separate lower scale-cost envelope has $\check f(1)=0$.
In any representation let $p_i$ be the mass at scale $i$. Mean scale one
gives $p_0=p_2$. On scales zero, one, and two, respectively, $X-T$ is
at most zero, zero, and one. Thus the target $X-T=1/2$ forces
$p_2\ge1/2$, and the cost is at least $p_0+p_2\ge1$. The endpoint
representation attains it. Scaling $f$ scales this error arbitrarily.
\end{example}
An operating cost must also be lifted before convexifying a nonconvex
per-unit set. With $E=\{-1,1\}$, scale one and cost $|\xi|$, its
convexified cost at $X=0$ is one, not $|0|$. Further,
$\lambda c(X/\lambda,T)$ need not be convex with an unscaled intensive
coordinate: the convex cost $c(\xi,T)=T^2$ gives $\lambda T^2$.

\begin{proposition}[Extensive-only perspective]\label{prop:scale-perspective}
Let $E\subset\R^d$ be nonempty compact, $c:E\to\R$ continuous, and
$\Lambda\subset[0,\infty)$ finite and nonempty. Let $h$ be the lower
boundary of $\conv\{(\xi,c(\xi)):\xi\in E\}$ and $\check f$ the lower
convex envelope of the finite scale-cost data. For $\lambda>0$, the exact
hull of $Z\ge\lambda c(\xi)+f(\lambda)$, $X=\lambda\xi$, is
\[
 \lambda\in\conv\Lambda,\quad X\in\lambda\conv E,\qquad
 Z\ge\lambda h(X/\lambda)+\check f(\lambda).
\]
If $0\in\Lambda$, its zero-scale slice is $X=0$, $Z\ge f(0)$.
\end{proposition}
\begin{proof}
For a representation with weights $\mu_j$ and positive mean scale
$\lambda$, the size-biased weights $\mu_j\lambda_j/\lambda$ sum to one.
They give per-unit cost at least $h(X/\lambda)$, while the original
weights give scale cost at least $\check f(\lambda)$. Conversely choose
attaining per-unit atoms with weights $\nu_i$ and attaining scale weights
$\mu_j$. Their product weights $\mu_j\nu_i$ attain both bounds and the
required $X$. Compactness gives finite attaining representations; adding
vertical slack gives every larger $Z$. At zero mean scale, nonnegativity
forces all positively weighted scales to be zero.
\end{proof}
This is the usual perspective argument after the cost has been included
in the per-unit graph. The multiplicity statement follows by treating
that cost as another additive extensive coordinate.

\begin{proposition}[Rectangularity criterion]\label{prop:scale-rectangularity}
Let $E\subset\R^d\times\R^p$ be nonempty compact and let the finite
catalogue $\Lambda\subset[0,M]$ satisfy $\{0,s,M\}\subseteq\Lambda$
and $0<s<M=\max\Lambda$. Put $C=\conv E$,
$Q=\operatorname{proj}_\xi C$, and $P_T=\operatorname{proj}_T C$.
Every scale-only cost separates, in the precise sense
\begin{equation}\label{eq:scale-cost-separation}
 \begin{split}
 &\conv\{(\lambda,X,T,Z):(\lambda,X,T)\in F(\Lambda,E),Z\ge f(\lambda)\}\\
 &\qquad=\{(\lambda,X,T,Z):(\lambda,X,T)\in\conv F(\Lambda,E),
                         Z\ge\check f(\lambda)\},
 \end{split}
\end{equation}
if and only if $C=Q\times P_T$. It suffices to test the single cost
$f(s)=0$, $f(\lambda)=1$ at every other catalogue scale.
\end{proposition}
\begin{proof}
The epigraph hull equals the compact hull of the cost graph plus the
nonnegative vertical ray. Hence all required minimum-cost representations
are attained; no closure operation is hidden. If the diagnostic cost
separates, every point of the scale-$s$ base-hull slice has a zero-cost
representation. All its constituents must have scale $s$, so this slice
is exactly $s\odot C$. Put $\alpha=s/M\in(0,1)$. For
$(\xi,T_1)\in C$ and $T_0\in P_T$, mixing the extreme slices gives
$(s,s\xi,\alpha T_1+(1-\alpha)T_0)$ in that slice. Thus the fibre of
$C$ at $\xi$ contains $\alpha T_1+(1-\alpha)T_0$. Iteration gives
$\alpha^jT_1+(1-\alpha^j)T_0$ in the fibre for every $j$. Closedness
puts $T_0$ in it, proving rectangularity.
Conversely, under rectangularity the base hull is
$\{(\lambda,X,T):0\le\lambda\le M,X\in\lambda Q,T\in P_T\}$.
For $\lambda>0$ choose $\xi=X/\lambda$ and use any attaining scale-cost
weights with the same $(\xi,T)\in C$ in every scale slice. Decompose
that point into $E$ within each slice; the resulting representation
attains $\check f(\lambda)$. At zero choose any $\xi\in Q$ and use
the zero slice. Every representation has cost at least $\check f$ by
convexity, proving equality. Universal separation implies the diagnostic
case, completing the equivalence.
\end{proof}
Without an interior scale, two endpoint weights are already fixed by
$\lambda$, so all scale-only costs separate even for nonrectangular $C$.
Affine scale costs always separate. The criterion concerns all scale-only
costs with the stated off-state domain, not arbitrary operating costs.

\subsection{Integral counts and aggregate rounding}
For nonempty compact extensive-only unit types $E_j\subset\R^d$ and an integral polytope
$P\subset\R_+^k$ of allowable counts, the exact aggregate hull is
\begin{equation}\label{eq:integral-count-hull}
 \{(n,X):n\in P,\ X\in\textstyle\sum_j n_j\conv E_j\}.
\end{equation}
Indeed it is convex and contains every integer multiplicity point.
For the reverse inclusion write $X=\sum_j n_j\xi_j$ with
$\xi_j\in\conv E_j$, decompose $n$ into integral vertices of $P$, and
keep these $\xi_j$ fixed. At each integral vertex the sum of scaled
convex hulls is the hull of the corresponding finite Minkowski sum.
This proves \eqref{eq:integral-count-hull}; no integer-decomposition
property beyond integrality of $P$ is needed.

For $n$ identical compact extensive sets in $\R^d$, the classical
Shapley--Folkman lemma \citep[Appendix 2, Lemma 2 and corollary]{Starr1969}
represents any $X\in n\conv E$ as a sum with at most $\min(d,n)$
constituents outside $E$. Replacing each exceptional constituent by any
point of $E$ changes the sum by at most $\min(d,n)\operatorname{diam}(E)$
in a chosen norm. Thus the same bound holds for the Hausdorff distance
between the two aggregate sets. An $L$-Lipschitz objective has gap at most
$L\min(d,n)\operatorname{diam}(E)$ when minimized over those sets alone.
A relative $O(d/n)$ bound additionally requires a positive normalization
of order $n$ and uniformly bounded $L$ and diameter.

All these identities concern the stated local or aggregate sets. Added
balance constraints can exclude the rounded point, and
$\conv(F\cap H)$ can be strictly smaller than $\conv F\cap H$.
Endpoint hull equality also does not license removing intermediate
catalogue choices from the feasible model. If the per-unit set depends
on scale, even the hull equality can fail: $\xi=\lambda$ at scales
$1,2,3$ gives $(\lambda,X)=(2,4)$ outside the segment joining $(1,1)$
and $(3,9)$.

\section{Retained domains and coordinate dependence in P-split}
\label{app:p-split}
Kronqvist, Misener, and Tsay's P-split construction splits convex additive
constraints by groups of variables, introduces epigraph auxiliaries for
the groups, and convexifies the resulting bounded auxiliary disjunction
\citep[Sections 2--4]{KronqvistEtAl2026}. The original domain $x\in X$
remains a global constraint. The bounds are computed over this retained
compact convex $X$, and identical or scaled group functions share an
auxiliary under the paper's minimal-sharing Assumption 3 and Remark 1.
In the following comparisons, ``basic'' means these epigraph links,
exact global interval bounds, and the exact auxiliary disjunction hull.
Extra original--auxiliary links change this model.

\subsection{A domain counterexample and a local nonexactness criterion}
The published Definition 4 calls disjuncts fully disjoint when they are
pairwise disjoint. Theorem 6 states universal nonexactness when the
bounds are additive and the constraint functions strictly convex. Its
universal statement fails with the retained-domain convention; this
concerns that theorem's scope, not validity of the P-split formulations.

\begin{example}[An exact retained-box relaxation]\label{ex:psplit-exact-box}
On $X=[0,3]\times[-1,1]$, use
\[
 D_0=\{(t,w)\in X:t^2+w^2\le1\},\qquad
 D_3=\{(t,w)\in X:(t-3)^2+w^2\le1\}.
\]
The disjuncts are nonempty and disjoint, and are defined by strictly convex
quadratics, with additive component bounds on the box. Together they contain
all four vertices of $X$. Hence $\conv(D_0\cup D_3)=X$, and every convex
valid relaxation retaining $X$, including every P-split partition, is exact.
For a direct full-split lift use shared $w^2$ and set
\[
 (\alpha_0,\alpha_3,\alpha_w)=(3t,9-3t,1).
\]
This lies in the auxiliary hull as the combination of $(0,9,1)$ and
$(9,0,1)$ with weights $1-t/3,t/3$. The links follow from
$t^2\le3t$, $(t-3)^2\le9-3t$, and $w^2\le1$. In dimension $n\ge2$
replace $w^2$ by $(n-1)^{-1}\sum_{i=2}^n x_i^2$ on
$[0,3]\times[-1,1]^{n-1}$. The same box-vertex argument applies,
with rational positive definite quadratics and one constraint per disjunct.
\end{example}
The source explicitly calls its ``few constraints'' Assumption 4
technically unnecessary, and the last construction also satisfies its
intended large-dimension setting. In the proof of its Theorem 6, strict
Jensen slack is claimed in every component. Strict convexity supplies
that slack only when the two arguments differ. Along the top hull
segment in the example, $w$ is fixed at one. Further, a neighborhood
allowed by epigraph slack need not remain in $X$: an outward step
through its top facet is forbidden.

\begin{proposition}[A sufficient directional repair]\label{prop:psplit-repair}
Let $X$ and $Q$ be convex, $A$ finite, and $F_a$ continuous convex. Define
$R=\{x\in X:\exists\alpha\in Q,F_a(x)\le\alpha_a\ (a\in A)\}$.
Assume $R$ contains the entire true feasible set. Suppose true feasible points $x^0,x^1$ have lifts $F(x^0),F(x^1)\in Q$.
At $x^\theta=(1-\theta)x^0+\theta x^1$, $0<\theta<1$, let
$\Delta_a=(1-\theta)F_a(x^0)+\theta F_a(x^1)-F_a(x^\theta)$.
If a hull-supporting inequality $c^Tx\le\beta$ is tight at $x^\theta$
and a direction $v$ has $c^Tv>0$, $x^\theta+\epsilon v\in X$ for all
small $\epsilon\ge0$, and
$F_a(x^\theta+\epsilon v)\le F_a(x^\theta)$ for every zero-slack link
and all small $\epsilon\ge0$, then $R$ strictly contains the true hull.
\end{proposition}
\begin{proof}
Keep $\alpha=(1-\theta)F(x^0)+\theta F(x^1)\in Q$. Every positive
slack persists for a sufficiently small step by continuity. There are
finitely many links, so one common positive radius suffices. The assumptions
preserve the remaining links and domain, while the support violation is
$\epsilon c^Tv>0$. If varying links have directional Lipschitz constants
$L_a$ on a common radius, imposing $\epsilon L_a\le\Delta_a$ for every
positive-slack link gives an explicit allowable step within that radius.
\end{proof}

\subsection{Exact basic relaxation of two translated balls}
\begin{proposition}[Projection and Hausdorff error]\label{prop:psplit-balls}
For $n\ge2$, radius $r>0$ and separation $d>2r$, let $x=(t,w)$ and
\[
 D_0=\{t^2+\|w\|_2^2\le r^2\},\qquad
 D_d=\{(t-d)^2+\|w\|_2^2\le r^2\},
 \quad X=[-r,d+r]\times[-r,r]^{n-1}.
\]
Their hull is $C=[0,d]e_1+r\mathbb B_2^n$. The basic full coordinate
split, sharing the transverse square auxiliaries, projects exactly to
\begin{equation}\label{eq:psplit-ball-projection}
 R=\{(t,w):\|w\|_2\le r,\quad
                2(t-d/2)^2+\|w\|_2^2\le2(d/2+r)^2\}.
\end{equation}
The two-group partition $\{t\},\{w_1,\ldots,w_{n-1}\}$ gives the same
set. With $e_*=[(d/2+r)^2-r^2/2]^{1/2}-d/2>0$,
\begin{equation}\label{eq:psplit-ball-distance}
 d_H(R,C)=\sqrt{r^2+e_*^2}-r,
 \qquad d_H(R,C)/r\longrightarrow\sqrt2-1\quad(d/r\to\infty).
\end{equation}
\end{proposition}
\begin{proof}
Put $U=(d+r)^2$. The auxiliaries are $a\ge t^2$, $b\ge(t-d)^2$,
$c_i\ge w_i^2$, with $0\le a,b\le U$, $0\le c_i\le r^2$.
The auxiliary disjunction is $a+\sum_i c_i\le r^2$ or
$b+\sum_i c_i\le r^2$. Its hull is exactly
\begin{equation}\label{eq:psplit-aux-hull}
 0\le a,b\le U,\quad c_i\ge0,\quad\sum_i c_i\le r^2,
 \qquad a+b+\sum_i c_i\le U+r^2.
\end{equation}
Necessity holds in each disjunct. For sufficiency fix $c$ and put
$h=r^2-\sum_i c_i\in[0,r^2]$. At this fixed $c$, the union in $(a,b)$
is $[0,h]\times[0,U]\cup[0,U]\times[0,h]$. Its hull is the square
cut by $a+b\le U+h$: every vertex of that truncated square lies in
the union. This proves sufficiency with the same $c$ in each constituent.
The set \eqref{eq:psplit-aux-hull} is downward closed, so its epigraph
projection follows by substituting the actual squares. Completing the
square gives \eqref{eq:psplit-ball-projection}; its cylinder and ellipsoid
already imply all the retained box bounds. In the two-group partition
replace $\sum_i c_i$ by $c\ge\|w\|^2$. Its global upper bound
$(n-1)r^2$ is reduced by either disjunct to $r^2$, giving the same proof.
The additive-bound refinement property
\citep[Theorem 4 and Corollary 3]{KronqvistEtAl2026} puts this full-split
set inside every basic coarser coordinate partition for this box.

Between the centers every point of the cylinder is in $C$. By reflection,
consider only $t\le0$ and put $\rho=\|w\|\in[0,r]$. The leftmost allowed
coordinate is $t=-e(\rho)$ with
$e(\rho)=\sqrt{(d/2+r)^2-\rho^2/2}-d/2>0$.
The squared distance to the nearer center is
\[
 e(\rho)^2+\rho^2=(d/2+r)^2+d^2/4+\rho^2/2
                         -d\sqrt{(d/2+r)^2-\rho^2/2}.
\]
Its derivative with respect to $\rho^2$ is
$1/2+d/[4\sqrt{(d/2+r)^2-\rho^2/2}]>0$. The maximum occurs at
$\rho=r$, where the distance outside the capsule is
$\sqrt{r^2+e_*^2}-r$. At $\rho=0$ that distance is zero, so no positive
part was lost. Since $C\subset R$, directed and ordinary Hausdorff
distances coincide. Finally $e_*/r\to1$, proving the limit.
\end{proof}

\subsection{An unbounded coordinate effect with rational data}
First take unit balls centered at $0$ and $(D,D)$, $D\ge5$, retaining
$[-1,D+1]^2$. The basic split introduces four separate squares
$a_i\ge x_i^2$, $b_i\ge(x_i-D)^2$ with upper bound $U=(D+1)^2$,
and convexifies $a_1+a_2\le1$ or $b_1+b_2\le1$.
At $p=(2D/3,D/3)$ its true square vectors $a,b$ have
$2a_i,2b_i\le8D^2/9\le U$. Their average representation
$(a,b)=\tfrac12(0,2b)+\tfrac12(2a,0)$ proves feasibility in the
auxiliary hull. Projection onto the center segment is $(D/2,D/2)$,
so the Hausdorff error is at least $D/(3\sqrt2)-1$.
Under the orthogonal change $t=(x_1+x_2)/\sqrt2$,
$w=(x_1-x_2)/\sqrt2$, the retained domain is exactly
\[
 -\sqrt2\le t+w,t-w\le\sqrt2(D+1),
\]
a diamond, and the centers are $0,\sqrt2D e_t$. The shared $w^2$
auxiliary is at most one throughout either disjunct, hence $|w|\le1$
in the relaxation. For $t\le0$, the domain implies
$-t+|w|\le\sqrt2$, giving
$t^2+w^2\le(\sqrt2-|w|)^2+w^2\le2$; the convex quadratic in
$|w|\in[0,1]$ has its maximum at an endpoint. The right end is
identical. Between centers the cylinder lies in the capsule. The new
error is therefore at most $\sqrt2-1$. No additive-bound hierarchy is
asserted on the transformed diamond.

\begin{proposition}[Rational coordinate comparison]\label{prop:psplit-rational}
For rational $D\ge2$, unit balls with centers $0,v=(3D,4D)$ in
$X=[-1,3D+1]\times[-1,4D+1]$ have basic original-coordinate Hausdorff
error at least $4D/5-1$. Under the fixed rational orthogonal map
\begin{equation}\label{eq:psplit-rational-map}
 t=(3x_1+4x_2)/5,\qquad w=(4x_1-3x_2)/5,
\end{equation}
using the exact transformed domain and tight recomputed bounds, the error
is at most $\sqrt2-1$.
\end{proposition}
\begin{proof}
The witness $p=(2D,4D/3)$ has each coordinate at one or two thirds of
$v_i$. Twice each original square is at most $8v_i^2/9\le(v_i+1)^2$,
so the preceding half-weight lift is feasible. Its transformed coordinates
are $t(p)=34D/15\in[0,5D]$ and $w(p)=4D/5$, proving the lower bound.
In aligned coordinates the common transverse auxiliary again forces
$|w|\le1$. For $t\le0$, the inverse map gives $x_1\le4/5$ and
$x_2\le3/5$, while the retained domain gives $x_1,x_2\ge-1$.
Thus $t^2+w^2=x_1^2+x_2^2\le2$. For $t\ge5D$, apply the same argument
to $x-v$: its two components lie in $[-4/5,1]$ and $[-3/5,1]$.
The error outside either end is at most $\sqrt2-1$, and between centers
it is zero. The map has determinant $-1$; changing the sign of $w$
gives a proper rotation with identical conclusions. All bounds and domain
inequalities remain rational for rational $D$.
\end{proof}

\subsection{The exact auxiliary image is still different from the graph}
\begin{proposition}[Strongest auxiliary-only comparison]
\label{prop:psplit-exact-image}
In either original-coordinate example, let
$F(x)=(x_i^2,(x_i-v_i)^2)_{i=1}^2$. Every convex auxiliary set $Q$
containing $F(0)$ and $F(v)$ preserves the original witness in
\[
 R(Q)=\{x\in X:\exists\alpha\in Q,F(x)\le\alpha\}.
\]
In particular the lower error survives $Q=\conv F(D_0\cup D_v)$.
In aligned coordinates, however, the exact feasible auxiliary image hull
for $F'(t,w)=(t^2,(t-d)^2,w^2)$ gives exactly the true capsule, for $d>2$.
This holds with either exactly transformed retained domain above.
\end{proposition}
\begin{proof}
Convexity forces $[F(0)+F(v)]/2\in Q$, whose entries are $v_i^2/2$.
At the witness, both coordinate squares are at most $4v_i^2/9$ and
hence satisfy all epigraph links to that point. This argument uses only
the two feasible center images, so every valid convex auxiliary-only
strengthening of the same epigraph map retains the witness.

For alignment write $(a,b,c)=F'(t,w)$ and set
\[
 f(c)=d^2+1-c+2d\sqrt{1-c},\qquad0\le c\le1.
\]
At an actual feasible point, $t\in[-q,q]\cup[d-q,d+q]$ with
$q=\sqrt{1-c}$, so $a,b\le(d+q)^2=f(c)$. The function $f$ is
concave and decreasing. Its two hypographs therefore contain the whole
convex auxiliary image hull. If the epigraph links are feasible there,
then $|w|\le1$ and, by monotonicity,
\[
 |t|\le d+\sqrt{1-w^2},\qquad |t-d|\le d+\sqrt{1-w^2}.
\]
Their intersection is precisely
$-\sqrt{1-w^2}\le t\le d+\sqrt{1-w^2}$, the capsule.
Conversely the epigraph relaxation is convex and contains both balls,
so it contains their hull. The retained domain contains that hull,
which proves equality. In fact, just $0\le c\le1$ and $a,b\le f(c)$
suffice: each upper inequality has a conic representation with
$s\ge0$, $s^2+c\le1$ and, respectively,
$a+c-d^2-1\le2ds$ or $b+c-d^2-1\le2ds$.
\end{proof}
The exact auxiliary image hull retains only values of the chosen functions;
it is not the hull of their full feasible graph in $(x,\alpha)$.
Original-space cuts, constraints coupling original and auxiliary variables,
or identities replacing the epigraph links can remove the witness.
These calculations establish explicit coordinate effects in this family,
without an algorithm for choosing good coordinates in general and without
an assertion about every strengthened P-split formulation.

\section{A stable correlation face and global lift size}\label{app:rank-one}
This appendix compares global formulation size with the local gap and
spatial resources in the main text. All matrix $\ell_1$ norms here are
\emph{entrywise} sums, and $\|H\|_\infty=\max_{ij}|H_{ij}|$ for a cost
matrix. Scalar inequalities count as nonnegative cone factors.

\subsection{The face and inherited exact lower bounds}
For $n\ge1$ define the compact sets
\[
 \mathcal R_n=\{W\ge0:\rank W\le1,\ We\le e,\ W^Te\le e\},
 \qquad K_n=\conv\mathcal R_n.
\]
For a nonzero generator, its row and column marginals $r,c$ and total
$S$ satisfy $W=rc^T/S$, $0\le r,c\le e$, and $e^Tr=e^Tc=S\le n$.
Indeed, writing $W=uv^T$ and computing its marginals proves this identity.
Then $\operatorname{tr}W=r^Tc/S\le1$. Equality forces
$r_i(1-c_i)=c_i(1-r_i)=0$ for every $i$, so $r=c=1_A$ for a nonempty
set $A$. Thus the trace-one face is the hull of $1_A1_A^T/|A|$.

\begin{proposition}[Explicit correlation face]\label{prop:correlation-face}
In $K_{2m}$ put
\[
 T(W)=e^TWe,\quad D(W)=1-\operatorname{tr}W,\quad
 B(W)=\sum_{i=1}^m W_{i,m+i},
\]
\[
 F=\{W\in K_{2m}:D=B=0,\ T=m\}.
\]
Then $F$ is a face affinely isomorphic to
$\operatorname{COR}(m)=\conv\{xx^T:x\in\{0,1\}^m\}$.
The forward map is $X=mW_{[m],[m]}$, and, with $x=\diag X$, its inverse is
\begin{equation}\label{eq:correlation-inverse}
 W=\frac1m\begin{pmatrix}
 X&xe^T-X\\ ex^T-X&ee^T-xe^T-ex^T+X
 \end{pmatrix}.
\end{equation}
Moreover, the affine functional
\[
 g(W)=m-T(W)+(2m+1)B(W)+(4m+1)D(W)
\]
satisfies $g\ge B+D\ge0$ on $K_{2m}$ and exposes exactly $F$.
\end{proposition}
\begin{proof}
Within the trace face, $B=0$ selects sets $A$ with at most one element
from each pair $\{i,m+i\}$. Their totals are $|A|\le m$; maximizing the
total selects exactly one element per pair. These successive faces give
atoms $(x,e-x)(x,e-x)^T/m$. Expanding their blocks gives
\eqref{eq:correlation-inverse}, which extends affinely to their hull.
For exposure, on a nonzero generator the inequality
$|a-b|\le a(1-b)+b(1-a)$ gives $\|r-c\|_1\le2SD$.
Also $r_i+r_{m+i}-r_ir_{m+i}\le1$, so
\[
 S\le m+\sum_i r_ir_{m+i}\le m+SB+\|r-c\|_1
       \le m+2mB+4mD.
\]
The zero generator satisfies the same affine inequality, hence so does
$K_{2m}$. Substitution proves $g\ge B+D$ and the asserted zero set.
\end{proof}

Adding affine equations to a lift and composing its output map leaves its
cone unchanged. Thus every lift of $K_{2m}$ gives one of
$\operatorname{COR}(m)$ over the same cone, without a strict-feasibility
assumption. The established correlation-polytope bounds therefore imply:
\begin{itemize}
\item For fixed $d$ and $m\ge d$, a lift over $(\mathbb S_+^d)^q$ has
\[
 q\ge(3^d-1)^{-(1-1/d)}(1-3^{-d})^{-m/d}.
\]
For $d=2$ the stronger bound is $q\ge7^{-1/2}(9/7)^{m/2}$
\citep[Theorem 1]{FawziParrilo2013}.
\item The total dimension of an exact second-order-cone lift is at least
$7^{-1/2}(9/7)^{m/2}$. A Lorentz cone
$\{(t,z):\|z\|_2\le t,\ z\in\R^k\}$ decomposes into $k-1$ three-dimensional
Lorentz cones for $k\ge2$, each isomorphic to $\mathbb S_+^2$;
one-dimensional nonnegative factors and the two-dimensional Lorentz cone
require at most their dimension in such blocks. Apply the preceding bound.
\item An unrestricted PSD lift has matrix order at least
$2^{\alpha m^{2/13}}$ for an absolute $\alpha>0$, for sufficiently large
$m$ \citep[Theorem 5.4]{LRS2015}. Several PSD blocks count by total order.
\end{itemize}
These are inherited conic lower bounds through the displayed face. They
assert neither nonrepresentability by finitely many second-order cones nor
an exponential number of Lorentz cones of unrestricted dimension.
For comparison, $K_2$ is already nonpolyhedral: the face where its first
row and column sums equal one is the hull of
$W(t)=(1,t)^T(1,t)/(1+t)$, $0\le t\le1$. With $q=W_{11}$ its lower-right
entry is $q+1/q-2$, a strictly convex function on $[1/2,1]$. Its infinitely
many graph points are extreme. Positive-error LP bounds below consequently
carry information that exact finite LP nonrepresentability alone does not.

\subsection{A linear error bound and a quantitative section transfer}
\begin{proposition}[Stability]\label{prop:correlation-stability}
For every $W\in K_{2m}$,
\begin{equation}\label{eq:face-error}
 \operatorname{dist}_1(W,F)\le(136m+10)g(W).
\end{equation}
If a convex set $Q$ satisfies $K_{2m}\subset Q\subset K_{2m}+\epsilon U_1$,
and $H_F=\{D=B=0,T=m\}$, then each $Y\in Q\cap H_F$ is within
$(184m+6)\epsilon$ of $F$. In particular, with $A_m=m(184m+6)$,
\begin{equation}\label{eq:face-outer-transfer}
 \operatorname{COR}(m)\subset\{mY_{[m],[m]}:Y\in Q\cap H_F\}
 \subset\operatorname{COR}(m)+A_m\epsilon U_1.
\end{equation}
\end{proposition}
\begin{proof}
For a nonzero generator $W=rc^T/S$, round $r$ to a binary vector $a$.
Since $\min(t,1-t)\le2t(1-t)$ and $\|r-c\|_1\le2SD$,
\[
 u:=\|r-a\|_1\le6SD,\qquad v:=\|c-a\|_1\le8SD.
\]
Indeed $\sum r_i(1-r_i)\le\sum r_i(1-c_i)+\|r-c\|_1\le3SD$.
Write $k=\sum a_i$, and let $h,z$ count pairs with two and zero ones.
Then $k=m+h-z$ and the product difference inequality on $[0,1]^2$ gives
$h\le SB+\|r-c\|_1+u\le SB+8SD$.
Change one coordinate in each full or empty pair to obtain $b$ with exactly
one one per pair. Thus
\[
 \|a-b\|_1=h+z=2h+m-k\le2SB+22SD+|m-S|.
\]
Set $V=bb^T/m\in F$ and $\widetilde b=(S/m)b$. Outer-product norm
factorization, using that $r,c,\widetilde b$ have total $S$, gives
\begin{align*}
 \|W-V\|_1
 &\le\|r-b\|_1+\|c-b\|_1+3|S-m|\\
 &\le4SB+58SD+5|S-m|.
\end{align*}
Here the first two terms bound the comparison with
$\widetilde b\widetilde b^T/S$, and its distance from $V$ is $|S-m|$.
Since $S\le2m$ and
$|S-m|\le g+(2m+1)B+(4m+1)D$, the last bound is at most
$(18m+5)B+(136m+5)D+5g\le(136m+10)g$.
For $W=0$ the distance to any paired atom is $m$, while $g(0)=5m+1$.
Taking finite convex combinations proves \eqref{eq:face-error}.

For the stronger transfer retain the bound
$8mB+116mD+5|S-m|$ on each generator, including zero. Exposure gives
$(S-m)_+\le2mB+4mD$. In a convex decomposition of $W$, therefore,
\[
 \sum_\ell\lambda_\ell|S_\ell-m|
 \le4mB(W)+8mD(W)+m-T(W).
\]
The corresponding combination $V\in F$ consequently obeys
\[
 \|W-V\|_1\le28mB(W)+156mD(W)+5[m-T(W)].
\]
For $Y\in Q\cap H_F$ choose $W\in K_{2m}$ with
$\|Y-W\|_1\le\epsilon$. The coefficient norms of $B,D,T$ give
$0\le B(W),D(W)\le\epsilon$ and $m-T(W)\le\epsilon$.
The triangle inequality yields $(184m+6)\epsilon$.
The block projection has norm at most $m$, proving the last assertion.
\end{proof}
An immediate exact-penalty consequence is that for
$\lambda>(136m+10)\|H\|_\infty$, minimizing
$\langle H,W\rangle+\lambda g(W)$ over $K_{2m}$ has all minimizers in $F$
and value $\min_F\langle H,V\rangle$. Compare $W$ with the point $V$
provided by \eqref{eq:face-error}; the excess is at least
$[\lambda-(136m+10)\|H\|_\infty]g(W)$.

\subsection{Fine LP and SDP approximations}
\begin{proposition}[Explicit accuracy scales]\label{prop:rank-one-approximation}
Let $K_{2m}\subset Q\subset K_{2m}+\epsilon U_1$ with $Q$ convex.
If $\epsilon\le1/A_m$, any LP lift of $Q$ has $2^{\Omega(m)}$
inequalities. If $\epsilon\le1/(2A_m)$, any PSD lift of $Q$ has matrix
order at least $\exp(c(m/\log m)^{2/13})$ for an absolute $c>0$ and all
sufficiently large $m$.
\end{proposition}
\begin{proof}
For the LP assertion apply \eqref{eq:face-outer-transfer}. The inequalities
$\langle2\diag(a)-aa^T,X\rangle\le1$, $a\in\{0,1\}^m$, have coefficient
norm at most one and slack $(1-a^Tx)^2$ at $xx^T$. The projected set is
therefore between $\operatorname{COR}(m)$ and the polyhedron defined by
these inequalities with right side two. The fixed-dilation lower bound of
\citet[Theorem 6(i), author version p.\ 18]{BraunEtAl2015} gives the result.
This particular sandwich has a small SDP relaxation and does not imply
the next assertion.

For that assertion use the sign correlation polytope
$\operatorname{SCOR}(m)=\conv\{yy^T:y\in\{-1,1\}^m\}$.
We first establish its fixed-error consequence of the pseudo-density
bounds in \citet[Theorem 3.8, equation (3.11), and Theorem 5.3]{LRS2015}.
For odd $k\ge3$ their construction supplies a degree-$k$ pseudo-density
$D$ on $\{0,1\}^k$ with
\[
 \E D=1,\quad\|D\|_\infty\le k^{3/2},\quad
 \E Df_k=-1/(4k^2),\qquad
 f_k(x)=\frac{(\sum x_i-k/2)^2-1/4}{k^2}.
\]
Here degree-$k$ means $\E Dp^2\ge0$ for $\deg p\le k/2$.
The cited quantitative theorem states that if $f:\{0,1\}^k\to[0,1]$
and $\E Df<-\delta$, then the matrix $M(S,x)=f(x_S)$, indexed by
$k$-subsets of $[m]$ and binary $x$, has, for $m\ge2k$, PSD rank at least
\begin{equation}\label{eq:lrs-input}
 \left[\frac{c_0\delta m}{k^3\|D\|_\infty\log m}\right]^{k/4}
 \left[\frac{\delta}{\|D\|_\infty}\right]^{3/2}\sqrt{\E f}.
\end{equation}
These pseudo-density and rank statements are established inputs.

Suppose $\operatorname{SCOR}(m)\subset R\subset
\operatorname{SCOR}(m)+\eta U_1$ with $\eta\le1/2$.
For every $k$-subset $S$ the affine function
\[
 L_S(Y)=\frac{\sum_{i,j\in S}Y_{ij}-1}{4k^2}
\]
is nonnegative on the sign polytope. Its coefficient norm is $1/(4k^2)$,
so $L_S+\theta\ge0$ on $R$, where $\theta=\eta/(4k^2)$.
At $y=2x-e$ its slack is $f_k(x_S)+\theta$. This shifted function is
in $[0,1]$, has mean at least $(k-1)/(4k^2)\ge1/(6k)$, and has
pseudo-expectation at most $-1/(8k^2)$. Taking $\delta=1/(16k^2)$
in \eqref{eq:lrs-input} gives
\begin{equation}\label{eq:shifted-psd-rank}
 \rank_{\rm psd}M\ge c_1k^{-23/4}
 \left[\frac{c_2m}{k^{13/2}\log m}\right]^{k/4}.
\end{equation}
If $R$ has a PSD lift of order $q$, these valid affine slacks have a PSD
factorization of order at most $q+1$. To justify this without a proper-lift
assumption, restrict the lift to the smallest PSD face containing its
feasible set. A finite combination of feasible matrices spanning all
feasible ranges is positive definite on this face. Conic duality for the
minimization of each nonnegative affine slack then writes that slack on
the affine slice as $\langle B,Y\rangle+\tau$, $B\succeq0$, $\tau\ge0$.
Appending the scalar factors $\tau$ and $1$ supplies the asserted
factorization. Choose odd $k$ comparable to
$a(m/\log m)^{2/13}$ with $a>0$ small enough that the bracket in
\eqref{eq:shifted-psd-rank} is at least 16. This proves the claimed
superpolynomial order for $R$.

Finally partition $W$ into four $m$-by-$m$ blocks and set
$J(W)=m(W_{11}-W_{12}-W_{21}+W_{22})$.
Its entrywise norm is at most $m$, and $J(F)=\operatorname{SCOR}(m)$.
The full-matrix transfer in Proposition~\ref{prop:correlation-stability}
shows that $R=J(Q\cap H_F)$ satisfies the preceding sandwich with
$\eta=A_m\epsilon\le1/2$. Affine sections and images preserve the lift
order, completing the proof.
\end{proof}
The scale in both assertions is explicit order $m^{-2}$. For compact
convex outer sets this norm guarantee is equivalent, by separation, to
uniform additive objective error $\epsilon$ over all costs with
$\|H\|_\infty\le1$. It does not cover only one objective, a restricted
cost family, or a larger accuracy tolerance. No spatial region-count
lower bound follows solely from these global lift-size bounds.

At the same SDP accuracy scale, a second-order-cone lift also needs
superpolynomial total cone dimension: decompose its Lorentz cones into
three-dimensional cones as above and embed them block-diagonally into
a PSD cone of order at most twice the original total dimension. This
transfers only the displayed superpolynomial approximate bound, whereas
the exact second-order-cone bound is exponential in $m$.

\section{Limiting bounds and primitive FBBT updates}\label{app:fbbt}
Feasibility-based bound tightening (FBBT) contracts a box by propagating
individual constraints. Its limit, the number of primitive updates, and
the finite tolerance used by an implementation are distinct objects.
When the limiting box is nonempty, continuous linear FBBT limits can be
obtained by an LP \citep[Theorem~4.1]{BelottiEtAl2012}, although primitive
iteration need not terminate finitely.
Here all variables start in $[0,1]$. A \emph{primitive update} computes the
interval hull for one affine equality or one nonnegative product equality
and intersects it with the current box. For the limiting-bound theorem
we allow any additional sound contracting updates, provided every defining equation
receives its forward natural interval-arithmetic update infinitely often.
This is fairness, with no bounded-delay assumption. Exact real arithmetic
specifies the mathematical limit; an approximation algorithm is measured
in the ordinary binary Turing model.

\begin{lemma}[Least fixed point transfer]\label{lem:fbbt-fixed-point}
Let $P$ have nonnegative polynomial coefficients, and suppose $x=P(x)$
has a solution in the unit cube. Under sound contracting updates and fair forward
propagation, its limiting lower vector is the least nonnegative fixed
point $q$ of $P$. This includes lifted arithmetic-node equations.
\end{lemma}
\begin{proof}
The sequence $s^0=0$, $s^{j+1}=P(s^j)$ increases and is bounded by every
nonnegative fixed point. Its limit $q$ is a fixed point by continuity and
is the least one. Soundness preserves $q$, so every lower vector $l$
satisfies $l\le q$. The natural interval lower endpoint is exactly $P_i(l)$.
Inductively, once $l\ge s^j$, fairness supplies a finite later time at
which all coordinates have received a forward update and $l\ge s^{j+1}$.
Hence the lower limit is at least every $s^j$ and equals $q$. Apply the
same argument to the augmented polynomial system for an expression graph.
\end{proof}

\begin{proposition}[Restricted limiting-bound hardness]\label{prop:fbbt-hardness}
Computing a designated limiting lower bound to additive error $1/4$ is
PosSLP-hard, even for feasible affine and bilinear defining systems in
the unit cube with at most three variable names per equation,
nonnegative coefficients in $\{1/2,1\}$, constants in $\{0,1/2,1\}$,
and distinct names in every product. Every strongly connected component
of the directed defining-equation dependency graph has at most four
vertices. The feasible set has at most two points, all rational.
\end{proposition}
\begin{proof}
PosSLP asks whether an integer arithmetic circuit using $0,1,+,-,\times$
has positive output. Following the circuit normalization and amplification
mechanism of \citet[Theorem 5.2]{EtessamiYannakakis2009}, replace every
gate by positive and negative parts. Multiplication uses
$(uv)^+=u^+v^++u^-v^-$ and $(uv)^-=u^+v^-+u^-v^+$; addition and
subtraction use the corresponding sums. It suffices to compare two
nonnegative integer outputs $U,V$.

Pad these fan-in-two circuits to common depth $L\ge1$ with alternating
addition and multiplication layers, using identities $v+0,v\cdot1$.
This costs polynomial size. Define proof-only normalizers $M_0=1$,
$M_j=2M_{j-1}$ on addition layers and $M_j=M_{j-1}^2$ on multiplication
layers. Encode each gate by its normalized value $p$ and complement $r$.
At inputs they are zero or one; the equations for subsequent gates are
\[
 p=(p_j+p_k)/2,\quad r=(r_j+r_k)/2
\]
for addition and
\[
 p=p_jp_k,\quad v=p_jr_k,\quad r=r_j+v
\]
for multiplication. These identities give $p=\operatorname{value}/M_j$
and $r=1-p$, with unique values in $[0,1]$. Introduce a copy variable
whenever a product would repeat a name. Each equation has at most three
names and all these dependencies are acyclic. With $M=M_L\le2^{2^L}$,
put
\[
 c=(p_U+r_V)/2=\tfrac12+(U-V)/(2M),\qquad
 d=(r_U+p_V)/2=1-c.
\]
The four-equation component
\[
 a'=a,\quad t=aa',\quad v=ct,\quad a=d+v
\]
implements $a=d+ca^2$. Its least nonnegative root $a_*$ equals one if
$c\le1/2$ and $d/c$ if $c>1/2$. In the latter case, writing
$\Delta=U-V\ge1$,
$1-a_*=2\Delta/(M+\Delta)\ge1/M$.

Construct $b=2^{-2^{L+2}}$ and its complement by $L+2$ paired squarings
starting from $1/2$. These numbers require polynomially many equations,
not their full rational expansions as input constants. They satisfy
$b\le1/(8M)$. Add $e=(1-b)a$, represented using the complement variable,
and the two-equation component $w=ez$, $z=b+w$.
At the least detector root it has the unique solution
\[
 z_*=\frac{b}{1-(1-b)a_*}
 =1\quad(U\le V),\qquad z_*\le bM\le1/8\quad(U>V).
\]
All displayed values lie in the cube and form a fixed point. The least
fixed point of the full system has the unique acyclic values, the least
detector root (bounded above by this displayed fixed point), and then
this unique final value. Lemma~\ref{lem:fbbt-fixed-point} identifies $z_*$
with the FBBT limit. An additive-$1/4$ answer is at least $3/4$ in the
first case and at most $3/8$ in the second; threshold $1/2$ decides PosSLP.
The largest directed component has the four detector variables, and the
last has two. The detector has at most two rational roots, and each
admissible root determines all other variables uniquely and rationally.
\end{proof}
The directed graph has an edge from a defined variable to each argument
on its right side; this is not an undirected treewidth restriction.
PosSLP is not known to be NP-hard. The proposition does not assert
hardness of executing a prescribed finite number of sweeps or of obtaining
an approximately stationary box. A unique rational feasible point is
possible if one further affine inequality is allowed: introduce $y=ca$
and impose $y\le1/2$. It retains the least detector root and excludes
$a=1$ when $c>1/2$; soundness and the fair original forward updates retain
the same limiting-bound proof.

\begin{proposition}[Every primitive schedule can be slow]\label{prop:fbbt-iterations}
For each $n\ge1$ a constant-coefficient affine--bilinear system with
$O(n)$ variables and equations requires at least $2^{2^n-1}$ applications
of one affine constraint before a designated lower bound reaches $1/2$,
although its limit is one. This holds for every fair ordering of the
bidirectional primitive interval-hull updates defined above. The only
feedback component has two variables and is linear after its upstream
values are fixed.
\end{proposition}
\begin{proof}
Start with $b_0=c_0=1/2$ and, for $i=1,\ldots,n$, impose
\[
 h_i=b_{i-1},\quad b_i=b_{i-1}h_i,\quad
 v_i=b_{i-1}c_{i-1},\quad c_i=c_{i-1}+v_i.
\]
These acyclic equations give $b=b_n=2^{-2^n}$, $c=c_n=1-b$.
Add $w=cz$, $z=b+w$. The unique feasible values are $z=1,w=c$.
For a stronger initialization fix all acyclic variables exactly, leaving
$z,w\in[0,1]$. Primitive interval-hull contractors are monotone under box
inclusion, so the original run with the same schedule has no larger lower
bounds.

In the stronger run the invariant is $0\le l_w\le cl_z$, $l_z\le1$.
A product update sets $l_w$ to $cl_z$ at most; its inverse bound
$l_z\ge l_w/c$ cannot improve $l_z$. An affine update sets
\[
 l_z^{\rm new}=\max(l_z,b+l_w)\le b+cl_z,\qquad
 l_w^{\rm new}=\max(l_w,l_z-b).
\]
Since $l_z-b-cl_z=b(l_z-1)\le0$, both terms defining the new lower bound
on $w$ are at most $cl_z^{\rm new}$. This proves the invariant, including
simultaneous hull updates. Upper endpoints cannot improve these formulas:
feasibility forces $u_z=1$, while $u_w$ is either one or $c$.
After $K$ affine updates, induction and Bernoulli's inequality give
\[
 l_z\le1-c^K=1-(1-b)^K\le Kb.
\]
Thus $l_z\ge1/2$ requires $K\ge1/(2b)=2^{2^n-1}$.
\end{proof}
The precise bound is in the displayed parameter $n$; encoding indexed
variable names as binary strings gives ordinary input length $O(n\log n)$.
At $l_z=l_w=0$ with upstream variables fixed, every primitive update changes
any lower endpoint by at most $b$, although $l_z$ is distance one from its
limit. Small local residual and small distance to the limiting box are
therefore different guarantees.

Slow and nonfinite linear FBBT iteration is classical, as are repeated
squaring and error amplification. The slow-Newton examples of
\citet[Section 7, Theorem 7.1]{EsparzaEtAl2010} and
\citet[Section 4.1, equation (18)]{StewartEtAl2015} also imply a doubly
exponential first-bit bound for ordinary Kleene iteration. This is an
inference from their displayed systems, not their stated Newton theorem.
For Stewart et al.'s chain, $x_0=(x_0^2+1)/2$ and
$x_i=(x_i^2+x_{i-1})/2$, the errors from the fixed point one satisfy
$e_0(k+1)=e_0(k)-e_0(k)^2/2$ and $e_i(k)^2\ge e_{i-1}(k)$, the latter
because each iterate is below its next update. The first recurrence gives
$e_0(k)\ge1/(k+1)$ by induction (the map is increasing on $[0,1]$).
Hence $e_n(k)\ge(k+1)^{-1/2^n}>1/2$ when $k+1<2^{2^n}$. The
specific scope proved here is every ordering of the stated bidirectional
primitive contractors, with one feedback component that becomes linear.
Symbolic substitution, equation aggregation, fixed-point acceleration,
and global contractors are outside this iteration bound.

\section{A short comparison with integer precision}\label{app:integer-precision}
Here the resource is the number $p$ of integer coordinates in a convex
lift, or binaries in a linear lift. A relaxation contains the graph and
has vertical error at most $\epsilon>0$ at every projected point over its
box. There is no bound on its continuous auxiliary dimension or on the
ranges of general integer coordinates. The published parity midpoint
argument \citep[Lemma 4.1]{LubinEtAl2022} groups graph points by the parity
of a feasible integer lift. The midpoint of two lifts in one class has
integer coordinates and remains feasible by convexity. Taking closures
of these at most $2^p$ classes in the compact input box preserves every
continuous midpoint-error inequality and avoids any measurability
assumption on the original lift.

For $x^2$ on $[0,1]$, the midpoint error is $(a-b)^2/4$. Each closed class
therefore has diameter at most $2\sqrt\epsilon$. Covering the interval
gives $1\le2^p2\sqrt\epsilon$, or
\begin{equation}\label{eq:integer-square}
 \epsilon\ge2^{-2p-2}.
\end{equation}
This is attained by a binary linear relaxation: divide the interval into
$2^p$ equal intervals of length $h$. On each interval use its secant as
upper bound and the maximum of the endpoint tangents as lower bound.
Both vertical errors are at most $h^2/4$, with equality in the upper
bound at the midpoint. The finite union of these bounded polyhedra has
a linear encoding with $p$ binary labels (unused auxiliary dimension and
row count are unrestricted). Thus the minimum integer dimension and the
minimum binary count both equal
$\max\{0,\lceil(\log_2(1/\epsilon)-2)/2\rceil\}$.
The more compact sawtooth construction has the same exact worst error
\citep[combined 2022 preprint, Section 5.1.1, PDF p.\ 21]{BeachEtAl2024};
that inspected preprint has different numbering from the published parts.

For $xy$, the corresponding midpoint inequality is
$|(a_1-b_1)(a_2-b_2)|\le4\epsilon$.
A compact class $S$ has area at most $16\epsilon\log2$. Indeed choose
points $a,b$ attaining the extreme first coordinates, with width $X>0$.
At first coordinate $a_1+t$ its section is contained in intervals centered
at $a_2,b_2$ of lengths $8\epsilon/t,8\epsilon/(X-t)$, respectively.
Integrating their minimum gives
$2\int_0^{X/2}8\epsilon/(X-t)\,dt=16\epsilon\log2$.
The zero-width case has zero area. Hence
\begin{equation}\label{eq:integer-product}
 \frac{1}{16(\log2)2^p}\le\epsilon_{\rm best}(p)
 \le2^{-p-2}.
\end{equation}
For the upper bound partition the first coordinate into $2^p$ strips
and use the four McCormick inequalities in each. On a rectangle of side
lengths $h,k$ their maximum vertical error is $hk/4$, as follows by
scaling the unit-square formulas. Thus both general integer dimension
and binary count are $\log_2(1/\epsilon)+O(1)$ for this scalar product.

The simultaneous graph $w_{ij}=x_ix_j$, $ij\in E(G)$, has a different
leading coefficient: the fractional vertex-cover number
\[
 \tau^*(G)=\min\{\textstyle\sum_i a_i:a_i\ge0,\ a_i+a_j\ge1\ (ij\in E)\}.
\]
To see the quantitative comparison, for positive edge errors define
\[
 L_C=\min\left\{\sum_i s_i:s_i\ge0,\quad
 s_i+s_j\ge\max\{0,\log_2(1/(C\epsilon_{ij}))\}\ (ij\in E)\right\}.
\]
Then the least general integer dimension and binary linear count obey
\begin{equation}\label{eq:integer-graph}
 \lceil L_{20}\rceil\le p_{\min}\le p_{\rm bin}\le L_4+|V|.
\end{equation}
For completeness, if a compact planar set satisfies
$|\Delta x\Delta y|\le4\epsilon$, its coordinate widths satisfy
$d_xd_y\le20\epsilon$: two extreme-$x$ points differ in $y$ by at most
$4\epsilon/d_x$, and every other point is within $8\epsilon/d_x$ in
$y$ of one of them. Apply this to each edge projection of a closed parity
class. If all its widths $d_i$ are positive, $s_i=-\log_2d_i$ is feasible
for $L_{20}$, so its volume is at most $\prod_i d_i\le2^{-L_{20}}$;
a zero width gives zero volume. The cover proves the lower bound.

For the upper bound take an optimum $s$ of $L_4$, put $p_i=\lceil s_i\rceil$,
and share the binary expansion
$x_i=\sum_{k=1}^{p_i}2^{-k}\beta_{ik}+r_i$, $0\le r_i\le2^{-p_i}$,
across every incident product. For each edge, expand first in $x_i$ and
then expand the residual term in $x_j$. Binary--continuous products have
exact four-row linear descriptions; only $r_ir_j$ needs a McCormick
relaxation. Its error is at most $2^{-p_i-p_j}/4\le\epsilon_{ij}$.
This proves \eqref{eq:integer-graph} with at most
$O(|V|+|E|+\sum_i\deg(i)p_i)$ continuous variables and rows.
For fixed nonempty $G$ and common error tending to zero,
$L_C=\tau^*(G)\log_2(1/(C\epsilon))$. Consequently
\[
 p_{\min}(\epsilon)=p_{\rm bin}(\epsilon)+O_G(1)
 =\tau^*(G)\log_2(1/\epsilon)+O_G(1).
\]
These arguments count binary assignments or integer parity classes.
Their feasible midpoint belongs to the mixed-integer projection; it
need not be a region certified by any spatial bounding oracle.

\bibliographystyle{plainnat}
\bibliography{references}
\end{document}
